%% Thinking Forth
%% Copyright (C) 2004 Leo Brodie
%% Initial transcription by Ed Beroset
%% 
%% test harness
%% pdflatex harness

\newif\ifpdf
\ifx\pdfoutput\undefined
  \pdffalse % we are not running PDFLaTeX
\else
  \pdfoutput=1 % we are running PDFLaTeX
  \pdftrue
\fi

\ifpdf
\documentclass[10pt,american,pdftex]{book}
\else
\documentclass[10pt,american,dvips]{book}
\fi
\usepackage{pst-all}
\usepackage{pst-blur}
\usepackage{pst-grad}
\usepackage[T1]{fontenc}
\usepackage[latin1]{inputenc}
\usepackage{graphics}
\usepackage{framed}
\usepackage{babel}
\usepackage{fancyvrb}
%\usepackage[0.21]{listings}
\usepackage{longtable}
\usepackage[colorlinks=true,linkcolor=blue,citecolor=blue,urlcolor=blue]{hyperref}
\usepackage{geometry}
\usepackage{babel}
\usepackage{mdwlist}
\geometry{verbose,paperwidth=6.8125in,paperheight=9.125in,tmargin=0.4in,%
bmargin=1.3in,lmargin=0.7in,rmargin=1.4in,headheight=0cm,headsep=0cm,%
footskip=0.5in,marginparwidth=0.6in,marginparsep=0.1in}
\usepackage{makeidx}
\usepackage[2004,tip,tipno]{tf}

% Fonts

%\usepackage{times}
%\renewcommand{\poorbf}[1]{\textbf{#1}}

%\lstloadlanguages{Forth}
%\input{epsf}
\usepackage{array}

% for pdflatex
\ifpdf
\DeclareGraphicsExtensions{.pdf,.png}
\else
\DeclareGraphicsExtensions{.eps}
\fi

\makeindex

% stuff left to Bernd:

% chapter title page
% page footer
% section and subsection header style
% caption style

\begin{document}
\title{Thinking Forth\\
\large A Language and Philosophy for Solving Problems}
\author{Leo Brodie%
\thanks{Includes interviews with \Forth{}'s inventor, Charles H. Moore,
and other \Forth{} thinkers}
\\ \small author of Starting \Forth{}}
\date{1984}

%% external title page is here

\pagestyle{empty}
{\makeatletter
\HUGE
    \fontfamily{phv}\selectfont
  \vbox{\advance\textwidth1em\hrule \@height 3pt
    \vspace*{0.2ex}\hbox to \textwidth{\bfseries Leo Brodie\hss}\vspace*{0.2ex}
    \hrule \@height 3pt}
  \vspace*{100pt}
  \vbox{%\hsize\linewidth
    \parindent\z@\hyphenpenalty 10000\raggedright
  {\THUGE\bfseries\slshape THINKING FORTH\\[1em]}
   \Huge A Language \\ and Philosophy \\ for Solving Problems}
   \vspace*{1em plus 1fill}
   {\Large\noindent Includes interviews with \Forth{}'s inventor, \person{Charles H. Moore},
and other \Forth{} thinkers}
%  \vspace*{1em plus 1fill}
  \pagebreak
}
%\maketitle

%%

\cleardoublepage
\setcounter{page}{1}

\pagenumbering{roman}
%% about the author

\cleardoublepage
\vspace*{100pt}
\noindent \person{Leo Brodie} is a writer, programmer, consultant, teacher, and
world-renowned authority on \Forth{}. He was a technical writer for
\Forth{}, Inc., and has been an independent consultant (since 1981) for
such clients as IBM, NCR, and Lockheed. He is also the author of
\emph{Starting \Forth{}} (Prentice-Hall, 1982).
\vfill

%% end about the author

%% internal title page is here

\cleardoublepage
{\makeatletter
\HUGE
    \fontfamily{phv}\selectfont
  \vbox{\advance\textwidth1em\hrule \@height 3pt
    \vspace*{0.2ex}\hbox to \textwidth{\bfseries Leo Brodie\hss}\vspace*{0.2ex}
    \hrule \@height 3pt}
  \vspace*{100pt}
  \vbox{%\hsize\linewidth
    \parindent\z@\hyphenpenalty 10000\raggedright
  {\THUGE\bfseries\slshape THINKING FORTH\\[1em]}
   \Huge A Language \\ and Philosophy \\ for Solving Problems}
  \vspace*{1em plus 1fill}
  \pagebreak
}


%% copyright page

\pagebreak

\vspace*{100pt}

\begin{center}\copyright 1984, 1994, and 2004 by \person{Leo Brodie}

\bigskip

\includegraphics{img-cc}

\bigskip

\textbf{Attribution-NonCommercial-ShareAlike 2.0}
\end{center}
\textbf{You are free:}

\begin{itemize}
\item to copy, distribute, display, and perform the work
 
\item to make derivative works
\end{itemize}
\textbf{Under the following conditions:}

\begin{quote}
\begin{description}
\bigskip
\item[\smash{\lower12pt\hbox to 0pt{\hss\includegraphics{img-by} }}Attribution.]
You must give the original author credit.
\bigskip
\bigskip
\bigskip

\item[\smash{\lower12pt\hbox to 0pt{\hss\includegraphics{img-nc} }}Noncommercial.]
You may not use this work for commercial purposes.
\bigskip
\bigskip

\item[\smash{\lower12pt\hbox to 0pt{\hss\includegraphics{img-sa} }}Share Alike.]
If you alter, transform, or build upon this work, you may distribute
the resulting work only under a license identical to this one.
\bigskip
\end{description}
\end{quote}

\begin{itemize}
\item For any reuse or distribution, you must make clear to others the
license terms of this work.
 
\item Any of these conditions can be waived if you get permission from the copyright holder.
\end{itemize}

\begin{center}
\textbf{Your fair use and other rights are in no way affected by the above.}

\medskip
This is a human-readable summary of the Legal Code
{\urlstyle{rm}\def\UrlLeft#1\UrlRight{(the\ full\ license).}
\url{http://creativecommons.org/licenses/by-nc-sa/2.0/legalcode}}

\bigskip
\vfill
You can download the electronic version from
\url{http://thinking-forth.sourceforge.net/}


\end{center}

%% Antoine Lavoisier quote

\cleardoublepage

\vspace*{100pt}
\begin{quotation}
\noindent\emph{%
It is impossible to disassociate language from science or science from
language, because every natural science always involves three things:
the sequence of phenomena on which the science is based the abstract
concepts which call these phenomena to mind, and the words in which
the concepts are expressed. To call forth a concept, a word is needed;
to portray a phenomenon, a concept is needed. All three mirror one and
the same reality.}

\begin{flushright}
---\person{Antoine Lavoisier}, 1789
\end{flushright}
\end{quotation}

%% end quote

\cleardoublepage
\pagestyle{headings}\markboth{Contents}{Contents}
\addcontentsline{toc}{section}{\contentsname}
\tableofcontents
\makeatletter
\renewcommand\numberline[1]{\hb@xt@\@tempdima{#1\hfil}}
\makeatother
\setcounter{tocdepth}{4}
%\listoffigures\vfill\pagebreak
\sectionlabelfalse
\chapter*{Preface to the 1994 Edition}\label{preface94}
\pagestyle{headings}

\initial I'm honored that the \Forth{} Interest Group is reprinting
\emph{Thinking \Forth{}.} It's gratifying to know that the book may
have value to appreciators and users of \Forth{}.

This edition is a reproduction of the original edition, with only
minor typographical corrections. A lot has happened in the ten years
since the book's original publication, rendering some of the opinions
obsolete, or at best archaic. A ``revised, updated edition'' would
have entrailed a rewrite of many sections, a larger effort than I'm
able to make at the time.

Off all the opinions in the book, the one that I most regret seeing in
print is my criticism of object-oriented programming. Since penning
this book, I've had the pleasure of writing an application in a
version of \Forth{} with support for object-oriented programming,
developed by Digalog Corp. of Ventura, California. I'm no expert, but
it's clear that the methodology has much to offer.

With all this, I believe that many of the ideas in \emph{Thinking
\Forth{}} are as valid today as they were back then. Certainly \person{Charles
Moore}'s comments remain a telling insight on the philosophy that
triggered the development of \Forth{}.

I with to thank \person{Marlin Ouverson} for his excellent job, patiently
struggling against incompatible file formats and OCR errors, to bring
this reprint to life.

\part*{Preface}

Programming computers can be crazy-making. Other professions give
you the luxury of seeing tangible proof of your efforts. A watchmaker
can watch the cogs and wheels; a seamstress can watch the seams come
together with each stitch. But programmers design, build, and repair
the stuff of imagination, ghostly mechanisms that escape the senses.
Our work takes place not in RAM, not in an editor, but within our
own minds.

Building models in the mind is both the challenge and the joy of programming.
How should we prepare for it? Arm ourselves with better debuggers,
decompilers, and disassemblers? They help, but our most essential
tools and techniques are mental. We need a consistent and practical
methodology for \emph{Thinking about} software problems. That is what
I have tried to capture in this book. \emph{Thinking Forth} is meant
for anyone interested in writing software to solve problems. It focuses
on design and implementation; deciding what you want to accomplish,
designing the components of the system, and finally building the program.

The book stresses the importance of writing programs that not only
work, but that are also readable, logical, and that express the best
solution in the simplest terms.

Although most of the principles described here can be applied to any
language, I've presented them in the context of FORTH. FORTH is a
language, an operating system, a set of tools, and a philosophy. It
is an ideal means for thinking because it corresponds to the way our
minds work. Thinking FORTH is thinking simple, thinking elegant, thinking
flexible. It is not restrictive, not complicated, not over-general.
You don't have to know FORTH to benefit from this book. \emph{Thinking
Forth} synthesizes the FORTH approach with many principles taught
by modern computer science. The marriage of FORTH's simplicity with
the traditional disciplines of analysis and style will give you a
new and better way to look at software problems and will be helpful
in all areas of computer application.

If you want to learn more about FORTH, another book of mine, \emph{Starting
FORTH}, covers the language aspects of FORTH. Otherwise, Appendix
A of this book introduces FORTH fundamentals.

A few words about the layout of the book; After devoting the first
chapter to fundamental concepts, I've patterned the book after the
software development cycle; from initial specification up through
implementation. The appendixes in back include an overview of FORTH
for those new to the language, code for several of the utilities described,
answers to problems, and a summary of style conventions. Many of the
ideas in this book are unscientific. They are based on subjective
experience and observations of our own humanity. For this reason,
I've included interviews with a variety of FORTH professionals, not
all of whom completely agree with one another, or with me. All these
opinions are subject to change without notice. The book also offers
suggestions called ''tips.'' They are meant to be taken only as they
apply to your situation. FORTH thinking accepts no inviolable rules.
To ensure the widest possible conformity to available FORTH systems,
all coded examples in this book are consistent with the FORTH-83 Standard.
One individual who greatly influenced this book is the man who invented
FORTH, Charles Moore. In addition to spending several days interviewing
him for this book, I've been privileged to watch him at work. He is
a master craftsman, moving with speed and deftness, as though he were
physically altering the conceptual models inside the machine---building,
tinkering, playing. He accomplishes this with a minimum of tools (the
result of an ongoing battle against insidious complexity) and few
restrictions other than those imposed by his own techniques. I hope
this book captures some of his wisdom. Enjoy!


\section{Acknowledgments}

Many thanks to all the good people who gave their time and ideas to
this book, including: Charles Moore, Dr. Mark Bernstein, Dave Johnson,
John Teleska, Dr. Michael Starling, Dr. Peter Kogge, Tom Dowling,
Donald Burgess, Cary Campbell, Dr. Raymond Dessy, Michael Ham, and
Kim Harris. Another of the interviewees, Michael LaManna, passed away
while this book was in production. He is deeply missed by those of
us who loved him. 


\vfill\pagebreak
\sectionlabeltrue
\setcounter{sectionlr}{4}
\pagenumbering{arabic}
\setcounter{page}{1}
\chapter{The~Philosophy of~FORTH}\Chapmark{1}


FORTH is a language and an operating system. But that's not all; It's
also the embodiment of a philosophy. The philosophy is not generally
described as something apart from FORTH. It did not precede FORTH,
nor is it described anywhere apart from discussions of FORTH, nor
does it even have a name other than ``FORTH.''

What is this philosophy? How can you apply it to solve your software
problems? Before we can answer these questions, let's take 100 steps
backwards and examine some of the major philosophies advanced by computer
scientists over the years. After tracing the trajectory of these advances,
we'll compare---and contrast---FORTH with these state-of-the-art programming
principles.


\section{An Armchair History of Software Elegance}
In the prehistoric days of programming, when computers were dinosaurs,
the mere fact that some genius could make a program run correctly
provided great cause for wonderment. As computers became more civilized,
the wonder waned. Management wanted more from programmers and from
their programs. As the cost of hardware steadily dropped, the cost
of software soared. It was no longer good enough for a program to
run correctly. It also had to be developed quickly and maintained
easily. A new demand began to share the spotlight with correctness.
The missing quality was called ``elegance.'' In this section we'll
outline a history of the tools and techniques for writing more elegant
programs.


\subsection{Memorability}
The first computer programs looked something like this:

\begin{Code}
00110101
11010011
11011001
\end{Code}
Programmers entered these programs by setting rows of switches---``on''
if the digit was {}``1,'' ``off'' if the digit was ``0.'' These
values were the {}``machine instructions'' for the computer, and
each one caused the computer to perform some mundane operation like
``Move the contents or Register B to Register A,'' or {}``Add the
contents of Register C into the contents of Register A.'' This proved
a bit tedious. Tedium being the stepmother of invention, some clever
programmers realized that the computer itself could be used to help.
so they wrote a program that translated easy-to-remember abbreviations
into the hard-to-remember bit patterns. The new language looked something
like this:

\begin{Code}
MOV B,A
ADD C,A
JMC REC1
\end{Code}
The translator program was called an \emph{assembler}, the new language
\emph{assembly language}. Each instruction {}``assembled'' the appropriate
bit pattern for that instruction, with a one-to-one correspondence
between assembly instruction and machine instruction. But names are
easier for programmers to remember. For this reason the new instructions
were called \emph{mnemonics}.


\subsection{Power}
Assembly-language programming is characterized by a one-for-one correspondence
between each command that the programmer types and each command that
the processor performs. In practice, programmers found themselves
often repeating the same \emph{sequence} of instructions over and
again to accomplish the same thing in different parts of the program.
How nice it would be to have a name which would represent each of
these common sequences.

This need was met by the {}``macro assembler,'' a more complicated
assembler that could recognize not only normal instructions, but also
special names ({}``macros''). For each name, the macro assembler
assembles the five or ten machine instructions represented by the
name, just as though the programmer had written them out in full.


\subsection{Abstraction}
\index{A!Abstraction}
A major advance was the invention of the {}``high-level language.''{}
Again this was a translator program, but a more powerful one. High-level
languages make it possible for programmers to write code like this:

\begin{Code}
X = Y (456/A) - 2
\end{Code}

\wepsfigp{img1-004}{So then I typed GOTO 500---and here I am!}

which looks a lot like algebra. Thanks to high-level languages, engineers,
not just bizarre bit-jockeys, could start writing programs. BASIC
and FORTRAN are examples of high-level languages. High-level languages
are clearly more {}``powerful'' than assembly languages in the sense
that each instruction might compile dozens of machine instructions.
But more significantly, high-level languages eliminate the linear
correspondence between source code and the resulting machine instructions.
The actual instructions depend on each entire {}``statement'' of
source code taken as a whole. Operators such as + and = have no meaning
by themselves. They are merely part of a complex symbology that depends
upon syntax and the operator's location in the statement. This nonlinear,
syntax-dependent correspondence between source and object code is
widely considered to be an invaluable step in the progress of programming
methodology. But as we'll see, the approach ultimately offers more
restriction than freedom.


\subsection{Manageability}
Most computer programs involve much more than lists of instructions
to work down from start to finish. They also involve testing for various
conditions and then {}``branching'' to the appropriate parts of the
code depending upon the outcome. They also involve {}``looping''
over the same sections of code repeatedly, usually testing for the
moment to branch out of the loop. Both assembler and high-level languages
provide branching and looping capabilities. In assembly languages
you use {}``jump instructions''; in some high-level languages you
use {}``GO TO'' commands\index{G!GOTO commands}. When these capabilities are used in the
most brute-force way, programs tend to look like the jumble you see
in \Fig{fig1-1}. This approach, still widely used in languages like
FORTRAN and BASIC, suffers from being difficult to write and difficult
to change if corrections need to be made. In this {}``bowl-of-spaghetti''
school of programming, it's impossible to test a single part of the
code or to figure out how something is getting executed that isn't
supposed to be getting executed. Difficulties with spaghetti programs
led to the discovery of {}``flow charts.'' These were pen-and-ink
drawings representing the {}``flow'' of execution used by the programmer
as an aid to understanding the code being written. Unfortunately the
programmer had to make the translation from code to flow chart and
back by hand. Many programmers found old-fashioned flow charts less
than useful.

\wtexfigt{fig1-1}{Unstructured code\index{G!GOTO commands},
using jumps or {}``GOTOs.''}

\subsection{Modularity}
A significant advance arose with the invention of {}``Structured
Programming,'' a methodology based on the observation that large problems
are more easily solved if treated as collections of smaller problems
\cite{dahl72}. Each piece is called a module. Programs consist of modules
within modules.

Structured programming eliminates spaghetti coding by insisting that
control flow can be diverted only within a module. You can't jump
out from the middle of one module into the middle of another module.
For example, \Fig{fig1-2} shows a structured diagram of a module to
{}``Make Breakfast,'' which consists of four submodules. Within each
submodule you'll find a whole new level of complexity which needn't
be shown at this level.

A branching decision occurs in this module to choose between the
\wepsfiga{fig1-2}{Design for a structured program}
{}``cold cereal'' module and the {}``eggs'' module, but control f1ow stays
within the outer module. Structured programming has three premises:

\begin{enumerate}
\item Every program is described as a linear sequence of self-contained
functions, called \emph{modules}. Each module has exactly one entry
point and one exit point.
\item Each module consists of one or more functions, each of which has exactly
one entry point and one exit point and can itself be described as
a module.
\item A module can contain:

\begin{enumerate}
\item operations or other modules
\item decision structures (IF THEN statements)
\item looping structures
\end{enumerate}
\end{enumerate}
The idea of modules having {}``one-entry, one-exit'' is that you
can unplug them, change their innards, and plug them back in, without
screwing up the connections with the rest of the program. This means
you can test each piece by itself. That's only possible if you know
exactly where you stand when you start the module, and where you stand
when you leave it.

In {}``Make Breakfast'' you'll either fix cereal or make eggs, not
both. And you'll always clean up. (Some programmers I know circumvent
this last module by renting a new apartment every three months.)

%
\wtexfigt{fig1-3}{Structured programming with a non-structured language}

Structured programming was originally conceived as a design approach.
Modules were imaginary entities that existed in the mind of the programmer
or designer, not actual units of source code. When structured programming
design techniques are applied to non-structured languages like BASIC,
the result looks something like \Fig{fig1-3}.


\subsection{Writeability}
Yet another breakthrough encouraged the use of structured programs:
structured programming languages. These languages include control
structures in their command sets, so you can write programs that have
a more modular appearance. Pascal is such a language, invented by
Niklaus Wirth to teach the principles of structured programming to
his students. \Fig{fig1-4} shows how this type of language would allow
{}``Make Breakfast'' to be written.

\wtexfigb{fig1-4}{Using a structured language.}

Structured programming languages include control structure operators
such as IF and THEN to ensure a modularity of control flow. As you
can see, indentation is important for readability, since all the instructions
within each module are still written out rather than being referred
to by name (e.g., ``i''). The finished program might take ten pages,
with the ELSE on page five.


\subsection{Designing from the Top}
How does one go about designing these modules? A methodology called
``top-down design'' proclaims that modules should be designed in order
starting with the most general, overall module and working down to
the nitty-gritty modules. Proponents of top-down design have witnessed
shameful wastes of time due to lack of planning, They've learned through
painful experience that trying to correct programs after they've been
written---a practice known as {}``patching''---is like locking the
barn door after the horse has bolted. So they offer as a countermeasure
this official rule of top-down programming:

\begin{quote}
Write no code until you have planned every last detail.
\end{quote}
Because programs are so difficult to change once they've been written,
any design oversight at the preliminary planning stage should be revealed
before the actual code-level modules are written, according to the
top-down design, Otherwise, man-years of effort may be wasted writing
code that cannot be used.


\subsection{Subroutines}
We've been discussing {}``modules'' as abstract entities only. But
all high-level programming languages incorporate techniques that allow
modules of design to be coded as modules of code---discrete units
that can be given names and ``invoked'' by other pieces of code. These
units are called subroutines, procedures, or functions, depending
on the particular high-level language and on how they happen to be
implemented. Suppose we write {}``MAKE-CEREAL'' as a subroutine.
It might look something like this:

\begin{Code}[fontfamily=cmss]
procedure make-cereal
   get clean bowl
   open cereal box
   pour cereal
   open milk
   pour milk
   get spoon
end
\end{Code}

\wepsfigp{img1-010}{Software patches are ugly and conceal structural weaknesses.}

We can also write {}``MAKE-EGGS'' and {}``CLEANUP'' as subroutines.
Elsewhere we can define {}``MAKE-BREAKFAST'' as a simple routine
that invokes, or calls, these subroutines:

\begin{Code}[fontfamily=cmss,commandchars=\&\{\}]
procedure make-breakfast
   var h: boolean (indicates hurried)
   &textit{test for hurried}
   if h = true then
      &textbf{call make-cereal}
   else
      &textbf{call make-eggs}
   end
   &textbf{call cleanup}
end
\end{Code}
\label{make-breakfast}
\medskip

The phrase {}``call make-cereal'' causes the subroutine named {}``make-cereal''
to be executed. When the subroutine has finished being executed, control
returns back to the calling program at the point following the call.
Subroutines obey the rules of structured programming.

As you can see, the effect of the subroutine call is as if the subroutine
code were written out in full within the calling module. But unlike
the code produced by the macro assembler, the subroutine can be compiled
elsewhere in memory and merely referenced. It doesn't necessarily
have to be compiled within the object code of the main program (\Fig{fig1-5}).

\wepsfigb{fig1-5}{A main program and a subroutine in memory.}

Over the years computer scientists have become more forceful in favoring
the use of many small subroutines over long-winded, continuous programs.
Subroutines can be written and tested independently. This makes it
easier to reuse parts of previously written programs, and easier to
assign different parts of a program to different programmers. Smaller
pieces of code are easier to think about and easier to verify for
correctness.

When subroutines are compiled in separate parts of memory and referred
to you can invoke the same subroutine many times throughout a program
without wasting space on repeated object code. Thus the judicious
use of subroutines can also decrease program size. 

Unfortunately, there's a penalty in execution speed when you use a
subroutine. One problem is the overhead in saving registers before
jumping to the subroutine and restoring them afterwards. Even more
time-consuming is the invisible but significant code needed to pass
parameters to and from the subroutine.

Subroutines are also fussy about how you invoke them and particularly
how you pass data to and from them. To test them independently you
need to write a special testing program to call them from.

For these reasons computer scientists recommend their use in moderation.
In practice subroutines are usually fairly large between a half page
to a full page of source code in length.


\subsection{Successive Refinement}
An approach that relies heavily on subroutines is called {}``Successive
Refinement'' \cite{wirth71}. The idea is that you begin by writing a skeletal
version of your program using natural names for procedures for data
structures. Then you write versions of each of the named procedures.
You continue this process to greater levels of detail until the procedures
can only be written in the computer language itself.

At each step the programmer must make decisions about the algorithms
being used and about the data structures they're being used on. Decisions
about the algorithms and associated data structures should be made
in parallel.

If an approach doesn't work out the programmer is encouraged to back
track as far as necessary and start again.

Notice this about successive refinement: You can't actually run any
part of the program until its lowest-level components are written.
Typically this means you can't test the program until after you've
completely designed it. 

Also notice: Successive refinement forces you to work out all details
of control structure on each level before proceeding to the next lower
level.


\subsection{Structured Design}

By the middle of late '70s, the computing industry had tried all the
concepts we've described, and it was still unhappy. The cost of maintaining
software---keeping it functional in the face of change---accounted
for more than half of the total cost of software, in some estimates
as much as ninety percent!

\wepsfigp{img1-013}{Tobias, I think you've carried the successive
refinement of that module far enough.}

Everyone agreed that these atrocities could usually be traced back
to incomplete analysis of the program, or poorly thought-out designs.
Not that there was anything wrong with structured programming \emph{per
se}. When projects came in late, incomplete, or incorrect, the designers
took the blame for not anticipating the unforeseen. 

Scholars naturally responded by placing more emphasis on design. {}``Next
time let's think things out better.''

About this time a new philosophy arose, described in an article called
``Structured Design'' \cite{stevens74}. One of its principles is stated in
this paragraph: 

\begin{quote}
Simplicity is the primary measurement recommended for evaluating alternative
designs relative to reduced debugging and modification time. Simplicity
can be enhanced by dividing the system into separate pieces in such
a way that pieces can be considered, implemented, fixed and changed
with minimal consideration or effect on the other pieces of the system.
\end{quote}
By dividing a problem into simple modules, programs were expected
to be easier to write, easier to change, and easier to understand.

But what is a module, and on what basis does one make the divisions?
{}``Structured Design'' outlines three factors for designing modules.


\subsection{Functional Strength}
One factor is something called {}``functional strength,'' which
is a measure of the uniformity of purpose of all the statements within
a module. If all the statements inside the module collectively can
be thought of as performing a single task, they are functionally bound.

You can generally tell whether the statements in a module are functionally
bound by asking the following questions. First, can you describe its
purpose in one sentence? If not, the module is probably not functionally
bound. Next, ask these four questions about the module:

\begin{enumerate}
\item Does the description have to be a compound sentence?
\item Does it use words involving time such as {}``first'', {}``next'',
{}``then,'' etc.?
\item Does it use a general or nonspecific object following the verb?
\item Does it use words like initialize which imply a lot of different functions
being done at the same time?
\end{enumerate}
If the answer to any of these four questions is {}``yes,'' you're
looking at some less cohesive type of binding than functional binding.
Weaker forms of binding include:

\begin{description}
\item [Coincidental~binding]the statements just happen to appear in the
same module
\item [Logical~binding]the module has several related functions and requires
a flag or parameter to decide which particular function to perform
\item [Temporal~binding]the module contains a group of statements that
happen at the same time, such as initialization but have no other
relationship
\item [Communicational~binding]the module contains a group of statements
that all refer to the same set of data
\item [Sequential~binding]where the output of one statement serves as
input for the next statement
\end{description}
Our {}``MAKE-CEREAL'' module exhibits functional binding, because
it can be thought of as doing one thing, even though it consists of
several subordinate tasks.


\subsection{Coupling}
A second tenet of structured design concerns {}``coupling,'' a measure
of how modules influence the behavior of other modules. Strong coupling
is considered bad form. The worst case is when one module actually
modifies code inside another module. Even passing control flags to
other modules with the intent to control their function is dangerous.

An acceptable form of coupling\index{D!Data coupling}
is {}``data coupling,'' which involves
passing data (not control information) from one module to another.
Even then, systems are easiest to build and maintain when the data
interfaces between modules are as simple as possible.

When data can be accessed by many modules (for instance, global variables),
there's stronger coupling between the modules. If a programmer needs
to change one module, there's a greater danger that the other modules
will exhibit {}``side effects.''

The safest kind of data coupling is the passing of local variables
as parameters from one module to another. The calling module says
to the subordinate module, in effect, {}``I want you to use the data
I've put in these variables named X and Y, and when you re done, I
expect you to have put the answer in the variable named Z. No one
else will use these variables.''

As we said, conventional languages that support subroutines include
elaborate methods of passing arguments from one module to another.


\subsection{Hierarchical Input-Process-Output Designing}
A third precept of structured design concerns the design process.
Designers are advised to use a top-down approach, but to pay less
attention initially to control structures. {}``Decision designing''
can wait until the later, detailed design of modules. Instead, the
early design should focus on the program's hierarchy (which modules
call which modules) and to the passing of data from one module to
another.

To help designers think along these new lines, a graphic representation
was invented, called the {}``structure chart.'' (A slightly different
form is called the {}``HIPO chart,'' which stands for {}``hierarchical
input-process-output.'') Structure charts include two parts, a hierarchy
chart and an input-output chart.

\wtexfigt{fig1-6}{The form of a structured chart, from ``Structured Design,''
\emph{IBM Systems Journal.}}

\Fig{fig1-6} shows these two parts. The main program, called DOIT,
consists of three subordinate modules, which in turn invoke the other
modules shown below them. As you can see, the design emphasizes the
transformation of input to output. 

The tiny numbers of the hierarchy chart refer to the lines on the in-out
chart. At point 1 (the module READ), the output is the value A. At
point 2 (the module TRANSFORM-TO-B), the input is A, and the output
is B. 

Perhaps the greatest contribution of this approach is recognizing
that decisions about control flow should not dominate the emerging
design. As we' ll see, control flow is a superficial aspect of the
problem. Minor changes in the requirements can profoundly change the
program's control structures, and {}``deep-six'' years of work.
But if programs are designed around other concerns, such as the flow
of data, t hen a change in plan won' t have so disastrous an effect.


\subsection{Information-Hiding }
In a paper \cite{parnas72} published back in 1972, Dr. David L. Parnas showed
that the criteria for decomposing modules should not be steps in the
process, but rather pieces of information that might possibly change.
Modules should be used to hide such information.

Let's look at this important idea of ``information- hiding'': Suppose
you are writing a Procedures Manual for your company. Here s a portion:

\begin{quote}
Sales Dept. takes order\\
sends blue copy to Bookkeeping\\
orange copy to Shipping\\
Jay logs the orange copy in the red binder on his desk, and completes
packing slip.
\end{quote}
Everyone agrees that this procedure is correct, and your manual gets
distributed to everyone in the company. Then Jay quits, and Marilyn
takes over. The new duplicate forms have green and yellow sheets,
not blue and orange. The red binder fills up and gets replaced with
a black one. Your entire manual is obsolete. You could have avoided
the obsolescence by using the term {}``Shipping Clerk'' instead
of the name Jay, the terms {}``Bookkeeping Dept. copy and'' {}``Shipping Dept.
copy'' instead of {}``blue'' and {}``orange,'' etc.

This example illustrates that in order to maintain correctness in
the face of a changing environment, arbitrary details should be excluded
from procedures. The details can be recorded elsewhere if necessary.
For instance, every week or so the personnel department might issue
a list of employees and their job titles, so anyone who needed to
know who the shipping clerk was could look it up in this single source.
As the personnel changes, this list would change.

This technique is very important in writing software. Why would a
program ever need to change, once it's running? For any of a million
reasons. You might want to run an old program on new equipment; the
program must be changed just enough to accommodate the new hardware.
The program might not be fast enough, or powerful enough, to suit
the people who are using it. Most software groups find themselves
writing {}``families'' of programs; that is, many versions of related
programs in their particular application field, each a variant on
an earlier program. 

To apply the principle of information-hiding to software, certain
details of the program should be confined to a single location, and
any useful piece of information should be expressed only once. Programs
that ignore this maxim are guilty of redundancy. While hardware redundancy
(backup computers, etc. ) can make a system more secure, redundancy
of information is dangerous.

As any knowledgeable programmer will tell you, a number that a program
uses and that might conceivably change should be made into a {}``constant''
and referred to throughout the program by name, not by value. For
instance, the number of columns representing the width of your computer
paper forms should be expressed as a constant. Even assembly languages
provide {}``EQU''s and labels for associating values such as addresses
and bit-patterns with names.

Any good programmer will also apply the concept of information-hiding
to the development of subroutines, ensuring that each module knows
as little as possible about the insides of other modules. Contemporary
programming languages such as C, Modula 2, and Edison apply this concept
to the architecture of their procedures. But Parnas takes the idea
much further. He suggests that the concept should be extended to algorithms
and data structures. In fact, hiding information---not decision-structure
or calling- hierarchy---should be the primary basis for design!


\section{The Superficiality of Structure}
Parnas proposes two criteria for decomposition:

\begin{enumerate}
\item possible (though currently unplanned) reuse, and
\item possible (though unplanned) change.
\end{enumerate}
This new view of a {}``module'' is different than the traditional
view. This {}``module'' is a collection of routines, usually very
small, which together hide information about some aspect of the problem. 

Two other writers describe the same idea in a different way, using
the term {}``data abstraction'' \cite{liskov75}.\index{A!Abstraction}
Their example is a push-down
stack. The stack {}``module'' consists of routines to initialize
the stack, push a value onto the stack, pop a value from the stack,
and determine whether the stack is empty. This {}``multiprocedure
module'' hides the information of how the stack is constructed from
the rest of the application. The procedures are considered to be a
single module because they are interdependent. You can't change the
method for pushing a value without also changing the method for popping
a value.

The word \emph{uses} plays an important role in this concept. Parnas
writes in a later paper \cite{parnas79}:

\begin{quote}
Systems that have achieved a certain {}``elegance''... have done
so by having parts of the system use other parts...

If such a hierarchical ordering exists then each level offers a testable
and usable subset of the system...

The design of the {}``uses'' hierarchy should be one of the major
milestones in a design effort. The division of the system into independently
callable subprograms has to go in parallel with the decisions about
uses, because they influence each other.
\end{quote}
A design in which modules are grouped according to control flow or
sequence will not readily allow design changes. Structure, in the
sense or control-flow hierarchy, is superficial.

A design in which modules are grouped according to things that may
change can readily accommodate change.


\section{Looking Back, and FORTH}
In this section we'll review the fundamental features of FORTH and
relate them to what we've seen about traditional methodologies. Here's
an example of FORTH code;

\begin{Code}
: BREAKFAST
   HURRIED?  IF  CEREAL  ELSE  EGGS  THEN CLEAN ;
\end{Code}
This is structurally identical to the procedure MAKE-BREAKFAST on
page \pageref{make-breakfast}. (If you're new to FORTH, refer to \Chap{A}
for an explanation.)  The words HURRIED?, CEREAL, EGGS, and CLEAN are
(most likely) also defined. as colon definitions. Up to a point, FORTH
exhibits all the traits we've studied: mnemonic value,
abstraction,\index{A!Abstraction} power, structured control operators,
strong functional binding, limited coupling, and modularity. But
regarding modularity, we encounter what may be FORTH's most
significant breakthrough:

\begin{quote}
The smallest atom of a FORTH program is not a module or a subroutine
or a procedure, but a {}``word.''
\end{quote}
Furthermore, there are no subroutines, main programs, utilities, or
executives, each of which must be invoked differently. Everything
in FORTH is a word. Before we explore the significance of a word-based
environment, let's first study two FORTH inventions that make it possible.


\subsection{Implicit Calls}
First, calls are implicit. You don't have to say CALL CEREAL, you
simply say CEREAL. In FORTH, the definition of CEREAL {}``knows''
what kind of word it is and what procedure to use to invoke itself.

Thus variables and constants, system functions, utilities, as well
as any user-defined commands or data structures can all be {}``called''
simply by name.


\subsection{Implicit Data Passing}
Second, data passing is implicit. The mechanism that produces this
effect is FORTH's data stack. FORTH automatically pushes numbers onto
the stack; words that require numbers as input automatically pop them
off the stack; words that produce numbers as output automatically
push them onto the stack. The words PUSH and POP do not exist in high-level
FORTH.

Thus we can write

\begin{Code}
: DOIT
    GETC  TRANSFORM-TO-DO  PUT-D ;
\end{Code}
confident that GETC will get {}``C'', and leave it on the stack.
TRANSFORM-TO-D will pick up {}``C'' from the stack, transform it,
and leave {}``D'' on the stack. Finally, PUT-D will pick up {}``D''
on the stack and write it. FORTH eliminates the act of passing data
from our code, leaving us to concentrate on the functional steps of
the data's transformation.

Because FORTH uses a stack for passing data, words can nest within
words. Any word can put numbers on the stack and take them off without
upsetting the f1ow of data between words at a higher level (provided,
of course, that the word doesn't consume or leave any unexpected values).
Thus the stack supports structured, modular programming while providing
a simple mechanism for passing local arguments.

FORTH eliminates from our programs the details of \emph{how} words
are invoked and \emph{how} data are passed. What's left? Only the
words that describe our problem.

Having words, we can fully exploit the recommendations of Parnas---to
decompose problems according to things that may change, and have each
{}``module'' consist of many small functions, as many as are needed
to hide information about that module. In FORTH we can write as many
words as we need to do that, no matter how simple each of them may
be.

A line from a typical FORTH application might read:

\begin{Code}
20 ROTATE LEFT TURRET
\end{Code}
Few other languages would encourage you to concoct a subroutine called
LEFT, merely as a modifier, or a subroutine called TURRET, merely
to name part of the hardware. Since a FORTH word is easier to invoke
than a subroutine (simply by being named, not by being called), a
FORTH program is likely to be decomposed into more words than a conventional
program would be into subroutines.


\section{Component Programming}
Having a larger set of simpler words makes it easy to use a technique
we'll call ``component programming.'' To explain, let's first reexamine
these collections we have vaguely described as ``things that may change.''
In a typical system, just about everything is subject to change: I/O
devices such as terminals and printers, interfaces such as UART chips,
the operating system, any data structure or data representation, any
algorithm, etc.

The question is: ``How can we minimize the impact of any such change?
What is the smallest set of other things that must change along with
such a change?''

The answer is: ``The smallest set of interacting data structures and
algorithms that share knowledge about how they collectively work.''
We'll call this unit a ``component.''

A component is a resource. It may be a piece of hardware such as a
UART or a hardware stack. Or the component may be a software resource
such as a queue, a dictionary, or a software stack.

All components involve data objects and algorithms. It doesn't matter
whether the data object is physical (such as a hardware register),
or abstract (such as a stack location or a field in a data base).
It doesn't matter whether the algorithm is described in machine code
or in problem- oriented words such as CEREAL and EGGS.

\wepsfigt{fig1-7}{Structured design vs. component design.}

\Fig{fig1-7} contrasts the results of structured design with the results
of designing by components. Instead of \emph{modules} called READ--RECORD,
EDIT--RECORD, and WRITE--RECORD, we're concerned with components
that describe the structure of records, provide a set of editor commands,
and provide read/write routines to storage.

What have we done? We've inserted a new stage in the development process:
We decomposed by components in our \emph{design}, then we described
the sequence, hierarchy, and input-process-output in our \emph{implementation}.
Yes, it's an extra step, but we now have an extra dimension for decomposition---not
just slicing but \emph{dicing}\index{D!Dicing}.

Suppose that, after the program is written, we need to change the
record structure. In the sequential, hierarchical design, this change
would affect all three modules. In the design by components, the change
would be confined to the record-structure component. No code that
uses this component needs to know of the change.

Aside from maintenance, an advantage to this scheme is that programmers
on a team can be assigned components individually, with less interdependence.
The principle of component programming applies to team management
as well as to software design.


We'll call the set of words which describe a component a {}``lexicon.''
(One meaning of lexicon is {}``a set of words pertaining to a particular
field of interest.'') The lexicon is your interface with the component
from the outside (\Fig{fig1-8}).

In this book, the term {}``lexicon'' refers only to those words
of a component that are used by name outside of a component. A component
may also contain definitions written solely to support the externally
visible lexicon. We'll call the supporting definitions {}``internal''
words.

\wepsfigt{fig1-8}{A lexicon describes a component.}

The lexicon provides the logical equivalents to the data objects and
algorithms in the form of names. The lexicon veils the component's
data structures and algorithms---the {}``how it works.'' It presents
to the world only a {}``conceptual model'' of the component described
in simple words---the {}``what it does.''

These words then become the language for describing the data structures
and algorithms of components written at a a higher level. The {}``what''
of one component becomes the {}``how'' of a higher component.

\wepsfigt{fig1-9}{The entire application consists of components.}

Written in FORTH, an entire application consists of nothing but components.
\Fig{fig1-9} show show a robotics application might be decomposed.

You could even say that each lexicon is a special-purpose compiler,
written solely for the purpose of supporting higher-level application
code in the most efficient and reliable way.

By the way, FORTH itself doesn't support components. It doesn't need
to. Components are the product of the program designer's decomposition.
(FORTH does have {}``screens.'' however---small units of mass storage
for saving source code. A component can usually be written in one
or two screens of FORTH.)

It's important to understand that a lexicon can be used by any and
all of the components at higher levels. Each successive component
does \emph{not} bury its supporting components, as is often the case
with layered approaches to design. Instead, each lexicon is free to
use all of the commands beneath it. The robot-movement command relies
on the root language, with its variables, constants, stack operators,
math operators, and so on, as heavily as any other component.

An important result of this approach is that the entire application
employs a single syntax, which makes it easy to learn and maintain.
This is why I use the term {}``lexicon'' and not {}``language.''
Languages have unique syntaxes.

This availability of commands also makes the process of testing and
debugging a whole lot easier. Because FORTH is interactive, the programmer
can type and test the primitive commands, such as

\begin{Code}
RIGHT SHOULDER 20 PIVOT
\end{Code}
from the {}``outside'' as easily as the more powerful ones like

\begin{Code}
LIFT COFFEE-POT
\end{Code}
At the same time, the programmer can (if he or she wants) deliberately
seal any commands, including FORTH itself, from being accessed by
the end user, once the application is complete.

Now FORTH's methodology becomes clear. FORTH programming consists
of extending the root language toward the application, providing new
commands that can be used to describe the problem at hand.

Programming languages designed especially for particular applications
such as robotics, inventory control, statistics, etc., are known as
{}``application-oriented languages.'' FORTH is a programming environment
for \emph{creating} app\-li\-ca\-tion-oriented languages. (That last sentence
may be the most succinct description of FORTH that you'll find.)

In fact, you shouldn't write any serious application in FORTH; as
a language it's simply not powerful enough. What you \emph{should}
do is write your own language in FORTH (lexicons) to model your understanding
of the problem, in which you can elegantly describe its solution.


\section{Hide From Whom?}
Because modern mainstream languages give a slightly different meaning
to the phrase {}``information-hiding,'' we should clarify. From
what, or whom are we hiding information?

The newest traditional languages (such as Modula 2) bend over backwards
to ensure that modules hide internal routines and data structures
from other modules. The goal is to achieve module independence (a
minimum coupling). The fear seems to be that modules strive to attack
each other like alien antibodies. Or else, that evil bands of marauding
modules are out to clobber the precious family data structures.

This is \emph{not} what we're concerned about. The purpose of hiding
information, as we mean it, is simply to minimize the effects of a
possible design-change by localizing things that might change within
each component.

FORTH programmers generally prefer to keep the program under their
own control and not to employ any techniques to physically hide data
structures. (Nevertheless a brilliantly simple technique for adding
Modula-type modules to FORTH has been implemented, in only three lines
of code, by Dewey Val Shorre \cite{shorre71}.)


\section{Hiding the Construction of Data Structures}
We've noted two inventions of FORTH that make possible the methodology
we've described---implicit calls and implicit data passing. A third
feature allows the data structures within a component to be described
in terms of previously-defined components. This feature is direct
access memory.

Suppose we define a variable called APPLES, like this:

\begin{Code}
VARIABLE APPLES
\end{Code}
We can store a number into this variable to indicate how many apples
we currently have:

\begin{Code}
20 APPLES !
\end{Code}
We can display the contents of the variable:

\begin{Code}
APPLES ? 20 ok
\end{Code}
We can up the count by one:

\begin{Code}
1 APPLES +!
\end{Code}
(The newcomer can study the mechanics of these phrases in Appendix
A.)

The word APPLES has but one function: to put on the stack the \emph{address}
of the memory location where the tally of apples is kept. The tally
can be thought of as a {}``thing,'' while the words that set the
tally, read the tally, or increment the tally can be considered as
{}``actions.''

FORTH conveniently separates {}``things'' from {}``actions'' by
allowing addresses of data structures to be passed on the stack and
providing the {}``fetch'' and {}``store'' commands.

We've discussed the importance of designing around things that may
change. Suppose we've written a lot of code using this variable APPLES.
And now, at the eleventh hour, we discover that we must keep track
of two different kinds of apples, red and green!

We needn't wring our hands, but rather remember the function of APPLES:
to provide an address. If we need two separate tallies, APPLES can
supply two different addresses depending on which kind of apple we're
currently talking about. So we define a more complicated version of
APPLES as follows:

\begin{Code}
VARIABLE COLOR  ( pointer to current tally)
VARIABLE REDS  ( tally of red apples)
VARIABLE GREENS  ( tally of green apples)
: RED  ( set apples-type to RED)  REDS COLOR ! ;
: GREEN  ( set apple-type to GREEN)  GREENS COLOR ! ;
: APPLES  ( -- adr of current apple tally)  COLOR @ ;
\end{Code}

\wepsfiga{fig1-10}{Changing the indirect pointer.}

Here we've redefined APPLES. Now it fetches the contents of a variable
called COLOR. COLOR is a pointer, either to the variable REDS or to
the variable GREENS. These two variables are the real tallies.

If we first say RED, then we can use APPLES to refer to red apples.
If we say GREEN, we can use it to refer to green apples (\Fig{fig1-10}).

We didn't need to change the syntax of any existing code that uses
APPLES. We can still say

\begin{Code}
20 APPLES !
\end{Code}
and

\begin{Code}
1 APPLES +!
\end{Code}
Look again at what we did. We changed the definition of APPLES from
that of a variable to a colon definition, without affecting its usage.
FORTH allows us to hide the details of how APPLES is defined from
the code that uses it. What appears to be {}``thing'' (a variable)
to the original code is actually defined as an {}``action'' (a colon
definition) within the component.

FORTH encourages the use of abstract data types by allowing data structures
to be defined in terms of lower level components. Only FORTH, which
eliminates the CALLs from procedures, which allows addresses and data
to be implicitly passed via the stack, and which provides direct access
to memory locations with @ and !, can offer this level of information-hiding.

FORTH pays little attention to whether something is a data structure
or an algorithm. This indifference allows us programmers incredible
freedom in creating the parts of speech we need to describe our applications.

I tend to think of any word which returns an address, such as APPLES,
as a {}``noun,'' regardless of how it's defined. A word that performs
an obvious action is a {}``verb.''

Words such as RED and GREEN in our example can only be called {}``adjectives''
since they modify the function of APPLES. The phrase

\begin{Code}
RED APPLES ?
\end{Code}
is different from

\begin{Code}
GREEN APPLES ?
\end{Code}
FORTH words can also serve as adverbs and prepositions. There's little
value in trying to determine what part of speech a particular word
is, since FORTH doesn't care anyway. We need only enjoy the ease of
describing an application in natural terms.


\section{But Is It a High-Level Language?}
In our brief technical overview, we noted that traditional high-level
languages broke away from assembly-language by eliminating not only
the \emph{one-for-one} correspondence between commands and machine
operations, but also the \emph{linear} correspondence. Clearly FORTH
lays claim to the first difference; but regarding the second, the
order of words that you use in a definition is the order in which
those commands are compiled. Does this disqualify FORTH from the ranks
of high-level languages? Before we answer, let's explore the advantages
of the FORTH approach.

\wepsfigp{img1-028}{Two points of view.}

\bigskip\blackline{2ex}
Here's what Charles Moore, the inventor of FORTH, has to say:

\begin{tfquot}
You define each word so that the computer knows what it means. The
way it knows is that it executes some code as a consequence of being
invoked. The computer takes an action on every word. It doesn't store
the word away and keep it in mind for later.

In a philosophical sense I think this means that the computer {}``understands''
a word. It understands the word \textbf{DUP}, perhaps more profoundly
than you do, because there's never any question in its mind what \textbf{DUP}
means.

The connection between words that have meaning to you and words that
have meaning to the computer is a profound one. The computer becomes
the vehicle for communication between human being and concept.
\end{tfquot}
\blackline{1ex}
One advantage of the correspondence between source code and machine
execution is the tremendous simplification of the compiler and interpreter.
This simplification improves performance in several ways, as we'll
see in a later section.

From the standpoint of programming methodology, the advantage to the
FORTH approach is that \emph{new} words and \emph{new} syntaxes can
easily be added. FORTH cannot be said to be {}``looking'' for words---it
finds words and executes them. If you add new words FORTH will find
and execute them as well. There's no difference between existing words
and words that you add.

What's more, this {}``extensibility'' applies to all types of words,
not just action-type functions. For instance, FORTH allows you to
add new \emph{compiling} words---like \textbf{IF} and \textbf{THEN}
that provide structured control flow. You can easily add a case statement
or a multiple-exit loop if you need them, or, just as importantly,
take them out if you don't need them.

By contrast, any language that depends on word order to understand
a statement must {}``know'' all legal words and all legal combinations.
Its chances of including all the constructs you'd like are slim. The
language exists as determined by its manufacturer; you can't extend
its knowledge.

Laboratory researchers cite flexibility and extensibility as among
FORTH's most important benefits in their environment. Lexicons can
be developed to hide information about the variety of test equipment
attached to the computer. Once this work has been done by a more experienced
programmer, the researchers are free to use their {}``software toolbox''
of small words to write simple programs for experimentation. As new
equipment appears, new lexicons are added.

Mark Bernstein has described the problem of using an off-the-shelf
special-purpose procedure library in the laboratory \cite{bern83}: {}``The
computer, not the user, dominates the experiment.'' But with FORTH,
he writes, {}``the computer actually encourages scientists to modify,
repair, and improve the software, to experiment with and characterize
their equipment. Initiative becomes once more the prerogative of the
researcher.''

\wepsfigp{img1-030}{Two solutions to the problem of security.}

For those purists who believe FORTH isn't fit to be called a high-level
language, FORTH makes matters even worse. While strong syntax checking
and data typing are becoming one of the major thrusts of contemporary
programming languages, FORTH does almost no syntax checking at all.
In order to provide the kind of freedom and flexibility we have described,
it cannot tell you that you meant to type RED APPLES instead of APPLES
RED. You have just invented syntax!

Yet FORTH more than makes up for its omission by letting you compile
each definition, one at a time, with turnaround on the order of seconds.
You discover your mistake soon enough when the definition doesn't
work. In addition, you can add appropriate syntax checking in your
definitions if you want to.

An artist's paintbrush doesn't notify the artist of a mistake; the
painter will be the judge of that. The chef's skillet and the composer's
piano remain simple and yielding. Why let a programming language try
to out think you?

So is FORTH a high-level language? On the question of syntax checking,
it strikes out. On the question of abstraction\index{A!Abstraction}
and power, it seems to be of \emph{infinite} level---supporting
everything from bit manipulation at an output port to business
applications.

You decide. (FORTH doesn't care.)


\section{The Language of Design}
FORTH is a design language. To the student of traditional computer
science, this statement is self-contradictory. {}``One doesn't design
with a language, one implements with a language. Design precedes implementation.''

Experienced FORTH programmers disagree. In FORTH you can write
abstract,\index{A!Abstraction}
design-level code and still be able to test it at any time by taking
advantage of decomposition into lexicons. A component can easily be
rewritten, as development proceeds, underneath any components that
use it. At first the words in a component may print numbers on your
terminal instead of controlling stepper motors. They may print their
own names just to let you know they've executed. They may do nothing
at all.

Using this philosophy you can write a simple by testable version of
your application, then successively change and refine it until you
reach your goal.

Another factor that makes designing in code possible is that FORTH,
like some of the newer languages, eliminates the {}``batch-compile''
development sequence (edit-compile-test-edit-compile-test). Because
the feedback is instantaneous, the medium becomes a partner in the
creative process. The programmer using a batch-compiler language can
seldom achieve the productive state of mind that artists achieve when
the creative current flows unhindered.

For these reasons, FORTH programmers spend less time planning than
their classical counterparts, who feel righteous about planning. To
them, not planning seems reckless and irresponsible. Traditional environments
force programmers to plan because traditional programming languages
do not readily accommodate change.

Unfortunately, human foresight is limited even under the best conditions.
Too much planning becomes counterproductive.

Of course FORTH doesn't eliminate planning. It allows prototyping.
Constructing a prototype is a more refined way to plan, just as breadboarding
is in electronic design.

As we'll see in the next chapter, experimentation proves more reliable
in arriving at the truth than the guesswork of planning.


\section{The Language of Performance}
Although performance is not the main topic of this book, the newcomer
to FORTH should be reassured that its advantages aren't purely philosophical.
Overall, FORTH outdoes all other high-level languages in speed, capability
and compactness.


\subsection{Speed}
Although FORTH is an interpretive language, it executes compiled code.
Therefore it runs about ten times faster than interpretive BASIC.

FORTH is optimized for the execution of words by means of a technique
known as {}``threaded code'' \cite{bell72}, \cite{dewar}, \cite{kogge82}. The penalty
for modularizing into very small pieces of code is relatively slight.

It does not run as fast as assembler code because the inner interpreter
(which interprets the list of addresses that comprise each colon definition)
may consume up to 50\% of the run time of primitive words, depending
on the processor.

But in large applications, FORTH comes very close to the speed of
assembler. Here are three reasons:

First and foremost, FORTH is simple. FORTH's use of a data stack greatly
reduces the performance cost of passing arguments from word to word.
In most languages, passing arguments between modules is one of the
main reasons that the use of subroutines inhibits performance.

Second, FORTH allows you to define words either in high-level or in
machine language. Either way, no special calling sequence is needed.
You can write a new definition in high level and, having verified
that it is correct, rewrite it in assembler without changing any of
the code that uses it. In a typical application, perhaps 20\% of the
code will be running 80\% of the time. Only the most often used, time-critical
routines need to be machine coded. The FORTH system itself is largely
implemented in machine-code definitions, so you'll have few application
words that need to be coded in assembler.

\wepsfigp{img1-033}{The best top-down designs of mice and young men.}

Third, FORTH applications tend to be better designed than those written
entirely in assembler. FORTH programmers take advantage of the language's
prototyping capabilities and try out several algorithms before settling
on the one best suited for their needs. Because FORTH encourages change,
it can also be called the language of optimization.

FORTH doesn't guarantee fast applications. It does give the programmer
a creative environment in which to design fast applications.


\subsection{Capability}
FORTH\index{C!Capability} can do anything any other language can
do---usually easier.

At the low end, nearly all FORTH systems include assemblers. These
support control-structure operators for writing conditionals and loops
using structured programming techniques. They usually allow you to
write interrupts---you can even write interrupt code in high level
if desired.

Some FORTH systems are multitasked, allowing you to add as many foreground
or background tasks as you want.

FORTH can be written to run on top of any operating system such as
RT-11, CP/M, or MS-DOS---or, for those who prefer it, FORTH can be
written as a self-sufficient operating system including its own terminal
drivers and disk drivers.

With a FORTH cross-compiler or target compiler, you can use FORTH
to recreate new FORTH systems, for the same computer or for different
computers. Since FORTH is written in FORTH, you have the otherwise
unthinkable opportunity to rewrite the operating system according
to the needs of your application. Or you can transport streamlined
versions of applications over to embedded systems. 


\subsection{Size}
There are two considerations here: the size of the root FORTH system,
and the size of compiled FORTH applications.

The FORTH nucleus is very flexible. In an embedded application, the
part of FORTH you need to run your application can fit in as little
as 1K. In a full development environment, a multitasked FORTH system
including interpreter, compiler, assembler, editor, operating system,
and all other support utilities averages 16K. This leaves plenty of
room for applications. (And some FORTHs on the newer processors handle
32-bit addressing, allowing unimaginably large programs.)

Similarly, FORTH compiled applications tend to be very small---usually
smaller than equivalent assembly language programs. The reason, again,
is threaded code. Each reference to a previously defined word, no
matter how powerful, uses only two bytes.

One of the most exciting new territories for FORTH is the production
of FORTH chips such as the Rockwell R65F11 FORTH-based microprocessor
\cite{dumse}. The chip includes not only hardware features but also the
run-time portions of the FORTH language and operating system for dedicated
applications. Only FORTH's architecture and compactness make FORTH-based
micros possible.


\section{Summary}
FORTH has often been characterized as offbeat, totally unlike any
other popular language in structure or in philosophy, On the contrary,
FORTH incorporates many principles now boasted by the most contemporary
languages. Structured design, modularity, and information-hiding are
among the buzzwords of the day.

Some newer languages approach even closer to the spirit of FORTH.
The language C, for instance, lets the programmer define new functions
either in C or in assembly language, as does FORTH. And as with FORTH,
most of C is defined in terms of functions.

But FORTH extends the concepts of modularity and information-hiding
further than any other contemporary language. FORTH even hides the
manner in which words are invoked and the way local arguments are
passed.

The resulting code becomes a concentrated interplay of words, the
purest expression of abstract thought. As a result, FORTH programmers
tend to be more productive and to write tighter, more efficient, and
better maintainable code.

FORTH may not be the ultimate language. But I believe the ultimate
language, if such a thing is possible, will more closely resemble
FORTH than any other contemporary language.


\begin{references}{99}
\bibitem{dahl72} O.J. Dahl, E.W. Dijkstra, and C.A.R. Hoare,
Structured Programming, London, Academic Press, 1972.
\bibitem{wirth71} Niklaus Wirth, {}``Program Development by Stepwise
Refinement,''{} Com- munications of ACM, 14, No. 4 (1971), 221-27.
\bibitem{stevens74} W.P. Stevens, G.J. Myers, and L.L. Constantine,
{}``Structured Design,''{} IBM Systems Journal, Vol. 13, No. 2, 1974.
\bibitem{parnas72} David L. Parnas, {}``On the Criteria To Be Used in
Decomposing Systems into Modules,''{} Communications of the ACM,
December 1972.
\bibitem{liskov75} Barbara H. Liskov and Stephen N. Zilles,
{}``Specification Techniques for Data Abstractions,''{} IEEE
Transactions on Software Engineering, March 1975.
\bibitem{parnas79} David L. Parnas, {}``Designing Software for Ease of
Extension and Contraction,''{} IEEE Transactions on Software
Engineering, March 1979.
\bibitem{shorre71} Dewey Val Shorre, {}``Adding Modules to FORTH,''{}
1980 FORML Proceedings, p. 71.
\bibitem{bern83} Mark Bernstein, ``Programming in the Laboratory,''
  unpublished paper, 1983.
\bibitem{bell72} James R. Bell, ``Threaded Code,'' Communications
of ACM, Vol. 16, No. 6, 370-72.
\bibitem{dewar} Robert B.K. DeWar, ``Indirect Threaded Code,''
Communications of ACM, Vol. 18, No. 6, 331.
\bibitem{kogge82} Peter M. Kogge, ``An Architectural Trail to
Threaded-Code Systems,'' Computer, March, 1982.
\bibitem{dumse} Randy Dumse, ``The R65F11 FORTH Chip,'' FORTH
Dimensions, Vol. 5, No. 2, p. 25.
\end{references}

%% Thinking Forth
%% Copyright (C) 2004 Leo Brodie
%% Initial transcription by Andrew Nicholson, based on OCR output
%% provided by John Hogerhuis
%%

%Page 037 in first edition

%\chapternum{TWO}
\chapter{Analysis}\Chapmark{2}

%Page 038 in first edition

\initialb Anyone who tells you there is some definite number of phases to the
software development cycle is a fool.

Nevertheless \dots{}

\section{The Nine Phases of the Programming Cycle}%

As we've seen, \Forth{} integrates aspects of design with aspects of
implementation and maintenance. As a result, the notion of a ``typical
development cycle'' makes as much sense as a ``typical noise.''

But any approach is better than no approach, and indeed, some
approaches have worked out better than others. Here is a development
cycle that represents an ``average'' of the most successful approaches
used in software projects:%
\index{P!Programming cycle:!phases|(}
\begin{description}
\item[Analysis] \hfill\index{A!Analysis}
    \begin{enumerate}
    \item Discover the Requirements and Constraints
    \item Build a Conceptual Model of the Solution
    \item Estimate Cost/Schedule/Performance
    \end{enumerate}
\item[Engineering] \hfill
    \begin{enumerate}
    \setcounter{enumi}{3}
    \item Preliminary Design
    \item Detailed Design
    \item Implementation
    \end{enumerate}
\item[Usage] \hfill
    \begin{enumerate}
    \setcounter{enumi}{6}
    \item Optimization
    \item Validation and Debugging
    \item Maintenance
    \end{enumerate}
\end{description}
In this book we'll treat the first six stages of the cycle,
focusing on analysis, design, and implementation.

In a \Forth{} project the phases occur on several levels. Looking at a
project from the widest perspective, each of these steps could take a
month or more.  One step follows the next, like seasons.

%Page 039 in first edition

But \Forth{} programmers also apply these same phases toward
defining each word. The cycle then repeats on the order of minutes.%
\index{P!Programming cycle:!phases|)}

Developing an application with this rapid repetition of the
programming cycle is known as using the ``Iterative Approach.''

\section{The Iterative Approach}%
\index{I!Iterative approach|(}%
\index{P!Programming cycle:!iterative approach|(}

The iterative approach was explained eloquently by
\person{Kim Harris} \cite{harris81}. \index{H!Harris, Kim|(}
He begins by describing the scientific method:

\begin{tfquot}
\dots{} a never-ending cycle of discovery and refinement. It first
studies a natural system and gathers observations about its behavior.
Then the observations are modeled to produce a theory about the
natural system. Next, analysis tools are applied to the model, which
produces predictions about the real system's behavior. Experiments
are devised to compare actual behavior to the predicted behavior. The
natural system is again studied, and the model is revised.

\wepsfigb{fig2-1}{The iterative approach to the
software development cycle, from ``The \Forth{} Philosophy,''
by \person{Kim Harris}, \emph{Dr.\@ Dobb's Journal.}}%
\index{H!Harris, Kim|)}

The \emph{goal} of the method is to produce a model which accurately
predicts all observable behavior of the natural system.
\end{tfquot}
\person{Harris} then applies the scientific method to the software
development cycle, illustrated in \Fig{fig2-1}:

\begin{enumerate}
\item A problem is analyzed to determine what functions are required
in the solution.
\item Decisions are made about how to achieve those functions with
the available resources.
\item A program is written which attempts to implement the design.
\item The program is tested to determine if the functions were
implemented correctly.
\end{enumerate}
%Page 040 in first edition
Mr. \person{Harris} adds:

\begin{tfquot}
Software development in \Forth{} seeks first to find the simplest
solution to a given problem. This is done by implementing selected
parts of the problem separately and by ignoring as many constraints as
possible. Then one or a few constraints are imposed and the program is
modified.
\end{tfquot}
An excellent testimonial to the development/testing model of design is
evolution. From protozoa to tadpoles to people, each species along the
way has consisted of functional, living beings. The Creator does not
appear to be a top-down designer.

\begin{tip}
Start simple. Get it running. Learn what you're trying to do. Add
complexity gradually, as needed to fit the requirements and
constraints. Don't be afraid to restart from scratch.
\end{tip}%
\index{I!Iterative approach|)}%
\index{P!Programming cycle:!iterative approach|)}

\section{The Value of Planning}%
\index{P!Planning:!value of|(}%
\index{P!Programming cycle:!value of planning|(}

In the nine phases at the start of this chapter we listed five steps
\emph{before} ``implementation.'' Yet in \Chap{1} we saw that an
overindulgence in planning is both difficult and pointless.

Clearly you can't undertake a significant software
project---regardless of the language---without some degree of planning.
Exactly what degree is appropriate?%
\index{J!Johnson, Dave|(}
\begin{interview}
\noindent More than one \Forth{} programmer has expressed high regard for
\person{Dave Johnson}'s meticulous approach to planning. \person{Johnson}
is supervisor at Moore Products Co. in Springhouse, Pennsylvania. The firm
specializes in industrial instrumentation and process control
applications. Dave has been using \Forth{} since 1978.

He describes his approach:
\begin{tfquot}
Compared with many others that use \Forth{}, I suppose we take a more
formal approach. I learned this the hard way, though. My lack of
discipline in the early years has come back to haunt me.

We use two tools to come up with new products: a functional specification
and a design specification. Our department of Sales \& Applications comes
up with the functional specification, through customer contact.

Once we've agreed on what we're going to do, the functional
specification is turned over to our department. At that point we work
through a design, and come up with the design specification.

Up to this point our approach is no different from programming in any
language. But with \Forth{}, we go about designing somewhat
differently.  With \Forth{} you don't have to work 95\% through your
design before you can start coding, but rather 60\% before you can get
into the iterative process.

%Page 041 in first edition

A typical project would be to add a functional enhancement to one of
our products. For example, we have an intelligent terminal with disk
drives, and we need certain protocols for communicating with another
device. The project to design the protocols, come up with displays,
provide the operator interfaces, etc. may take several months. The
functional specification takes a month; the design specification takes
a month; coding takes three months; integration and testing take
another month.

This is the typical cycle. One project took almost two years, but six
or seven months is reasonable.

When we started with \Forth{} five years ago, it wasn't like that. When I
received a functional specification, I just started coding. I used a
cross between top-down and bottom-up, generally defining a structure,
and as I needed it, some of the lower level, and then returning with
more structure.

The reason for that approach was the tremendous pressure to show
something to management. We wound up never writing down what we were
doing. Three years later we would go back and try to modify the code,
without any documentation. \Forth{} became a disadvantage because it
allowed us to go in too early. It was fun to make the lights flash and
disk drives hum. But we didn't go through the nitty-gritty design
work. As I said, our ``free spirits'' have come back to haunt us.

Now for the new programmers, we have an established requirement: a
thorough design spec that defines in detail all the high-level \Forth{}
words---the tasks that your project is going to do. No more reading a
few pages of the functional specification, answering that, reading a
few more, answering that, etc.

No living programmer likes to document. By ensuring the design ahead
of time, we're able to look back several years later and remember what
we did.

I should mention that during the design phase there is some amount of
coding done to test out certain ideas. But this code may not be part
of the finished product. The idea is to map out your design.
\end{tfquot}%
\index{J!Johnson, Dave|)}
\end{interview}
\person{Johnson} advises us to complete the design specification
before starting to code, with the exception of needed preliminary
tests. The next interview backs up this point, and adds some
additional reasons.%
\begin{interview}
\index{T!Teleska, John|(}
\noindent \person{John Teleska} has been an independent software
consultant since 1976, specializing in custom applications for
academic research environments.  He enjoys providing research tools
``right at the edge of what technology is able to do.''
\person{Teleska} works in Rochester, New York:

\begin{tfquot}
I see the software development process as having two phases. The first is
making sure I understand what the problem is. The second is
implementation, including debugging, verification, etc.

My goal in Phase One is an operational specification. I start with a
problem description, and as I proceed it becomes the operational
specification. My understanding of the problem metamorphoses into a
solution. The better the understanding, the more complete the
solution. I look for closure; a sense of having no more questions that
aren't answered in print.

I've found that on each project I've been putting more time into Phase
One, much to the initial dismay of many of my clients. The limiting
factor is how
%Page 042 in first edition
much I can convince the client it's necessary to spend that time up
front.  Customers generally don't know the specifications for the job
they want done. And they don't have the capital---or don't feel they
do---to spend on good specs. Part of my job is to convince them it
will end up costing more time and money not to.

Some of Phase One is spent on feasibility studies. Writing the spec
unearths uncertainties. I try to be as uncertain about uncertainties
as possible. For instance, they may want to collect 200,000 samples a
second to a certain accuracy. I first need to find out if it's even
possible with the equipment they've got. In this case I've got to test
its feasibility by writing a patch of code.

Another reason for the spec is to cover myself. In case the
application performs to the spec but doesn't fully satisfy the
customer, it's the customer's responsibility. If the customer wants
more, we'll have to renegotiate. But I see it as the designer's
responsibility to do whatever is necessary to generate an operational
specification that will do the job to the customer's satisfaction.

I think there are consultants who bow to client pressure and limit the
time they spend on specs, for fear of losing the job. But in these
situations nobody ends up happy.
\end{tfquot}
\index{T!Teleska, John|)}
\end{interview}%
We'll return to the \person{Teleska} interview momentarily.%
\index{P!Planning:!value of|)}%
\index{P!Programming cycle:!value of planning|)}

\section{The Limitations of Planning}%
\index{P!Planning:!limitations of|(}%
\index{P!Programming cycle:!limitations of planning|(}

Experience has taught us to map out where we're going before we begin
coding. But planning has certain limitations. The following interviews
give different perspectives to the value of planning.%
\begin{interview}
\index{T!Teleska, John|(}
\noindent Despite \person{Teleska}'s preference for a well-planned
project, he suggests that the choice between a top-down and bottom-up
approach may depend on the situation:

\begin{tfquot}
On two recent projects involving a lot of technical interface work, I
did the whole thing bottom-up. I milled around in a bunch of
data-sheets and technical descriptions of little crannies of the
operating system I was dealing with. I felt lost most of the time,
wondering why I ever took the job on.  Then finally I reached a
critical mass of some sort and began putting small programs together
that made small things happen. I continued, bottom-up, until I matched
the target application.

My top-down sense was appalled at this procedure. But I've seen me go
through this process successfully too many times to discount it for
any pedagogical reasons. And there is always this difficult phase
which it seems no amount of linear thinking will penetrate.
Programming seems a lot more intuitive than we, in this business, tell
each other it ought to be.

I think if the application elicits this sense of being lost, I proceed
bottom-up. If the application is in familiar territory then I'll
probably use a more traditional by-the-book approach.
\end{tfquot}
\index{T!Teleska, John|)}
\end{interview}

%Page 043 in first edition

\noindent And here's another view:%
\index{S!Starling, Michael|(}
\begin{interview}
\noindent At the time I interviewed him, \person{Michael Starling} of Union
Carbide was putting the final touches on two applications involving
user-configurable laboratory automation and process control automation
systems. For the pilot plant system, \person{Starling} designed both the
hardware and software to known requirements; on the laboratory
automation system he also defined the requirements himself.

His efforts were extremely successful. On one project, the new system
typically costs only 20\% as much as the equivalent system and
requires days, instead of months, to install and configure.

I asked him what techniques of project management he employed.

\begin{tfquot}
On both of these projects much design was needed. I did not follow the
traditional analysis methods, however. I did employ these steps:

First, I clearly defined the boundaries of the problem.

Second, I determined what the smaller functional pieces, the software
subsystems, had to be.

Third, I did each piece, put them together, and the system ran.

Next, I asked the users ``Does this meet your requirements?''
Sometimes it didn't, and in ways that neither the users nor the
specification designers could have anticipated.

For instance, the designers didn't realize that the original
specification wouldn't produce pleasing, human-oriented graphics
displays. After working with the interactive graphics on the first
version, users were applying arbitrary scales and coming up with
oddball displays.

So even after the basic plot algorithm was designed, we realized we
needed auto-scaling. We went back in and analyzed how human beings
plot data and wrote a first level plot function that evaluates the x
and y data and how much will fit on the graph.

After that, we realized that not all the data taken will be of
interest to experimenters. So we added a zoom capability.

This iterative approach resulted in cleaner code and better thought
out code. We established a baseline set of goals and built a minimal
system to the users' known requirements. Then we'd crank in the
programmer's experience to improve it and determine what the users
forgot they needed when they generated the specs.

The users did not invent most of the new ideas. The programmers did,
and they would bounce these ideas off the users. The problem
definition was a two-way street. In some cases they got things they
didn't know they could do on such a small computer, such as applying
digital filters and signal processing to the data.

One of the things about \Forth{} that makes this approach possible is
that primitives are easily testable. It takes some experience with
\Forth{} to learn how to take advantage of this. Guys from traditional
environments want to write ten pages of code at their desk, then sit
down to type it in and expect it to work.

To summarize my approach: I try to find out from the users what they
need, but at the same time recognizing its incompleteness. Then I keep
%Page 044 in first edition
them involved in the design during the implementation, since they have
the expertise in the application.  When they see the result, they feel
good because they know their ideas were involved.

The iterative approach places highest value on producing a good
solution to the real problem. It may not always give you the most
predictable software costs. The route to a solution may depend upon
your priorities. Remember:
\begin{list}{}{}
\item Good
\item Fast
\item Cheap
\end{list}
Pick any two!
\end{tfquot}
\end{interview}

\noindent As \person{Starling} observes, you don't know completely what
you're doing till you've done it once. In my own experience, the best way
to write an application is to write it twice. Throw away the first
version and chalk it up to experience.%
\index{S!Starling, Michael|)}%
\index{K!Kogge, Peter|(}
\begin{interview}
\noindent \person{Peter Kogge} is Senior Technical Staff in the IBM
Federal Systems Division, Oswego, New York:

\begin{tfquot}
One of the key advantages I find in \Forth{} is that it allows me to
very quickly prototype an application without all the bells and
whistles, and often with significant limitations, but enough to wring
out the ``human interface'' by hands-on trial runs.

When I build such a prototype, I do so with the firm constraint that I
will use not a single line of code from the prototype in the final
program. This enforced ``do-over'' almost always results in far
simpler and more elegant final programs, even when those programs are
written in something other than \Forth{}.
\end{tfquot}
\end{interview}%
\index{K!Kogge, Peter|)}

\noindent Our conclusions? In the \Forth{} environment planning is
necessary. But it should be kept short. Testing and prototyping are
the best ways to discover what is really needed.

A word of caution to project managers: If you're supervising any
experienced \Forth{} programmers, you won't have to worry about them
spending too much time on planning. Thus the following tip has two
versions:

\begin{tip}
For newcomers to \Forth{} (with ``traditional'' backgrounds):
Keep the analysis phase to a minimum.

\medskip
For \Forth{} addicts (without a ``traditional'' background):
Hold off on coding as long as you can possibly stand it.
\end{tip}
%Page 045 in first edition
Or, as we observed in \Chap{1}:

\begin{tip}
Plan for change (by designing components that can be changed).
\end{tip}
Or, simply:

\begin{tip}
Prototype.
\end{tip}%
\index{P!Planning:!limitations of|)}%
\index{P!Programming cycle:!limitations of planning|)}

\section{The Analysis Phase}%
\index{A!Analysis|(}%

In the remainder of this chapter we'll discuss the analysis phase.
Analysis is an organized way of understanding and documenting what
the program should do.%
\index{A!Analysis!defined}

With a simple program that you write for yourself in less than an
hour, the analysis phase may take about 250 microseconds. At the other
extreme, some projects will take many man-years to build. On such a
project, the analysis phase is critical to the success of the entire
project.

We've indicated three parts to the analysis phase:

\begin{enumerate}\parsep=0pt\itemsep=0pt
\item Discovering the requirements and constraints
\item Building a conceptual model of the solution
\item Estimating cost, scheduling, and performance
\end{enumerate}

\noindent Let's briefly describe each part:

\subsection{Discovering the Requirements}%
\index{A!Analysis!requirements|(}%
\index{R!Requirements|(}

The first step is to determine what the application should do. The
customer, or whoever wants the system, should supply a ``requirements
specification.'' This is a modest document that lists the minimum
capabilities for the finished product.

The analyst may also probe further by conducting interviews and
sending out questionnaires to the users.%
\index{A!Analysis!requirements|)}%
\index{R!Requirements|)}

\subsection{Discovering the Constraints}%
\index{A!Analysis!constraints|(}%
\index{C!Constraints|(}

The next step is to discover any limiting factors. How important is
speed? How much memory is available? How soon do you need it?

No matter how sophisticated our technology becomes, programmers will
always be bucking limitations. System capacities inexplicably
%Page 046 in first edition
diminish over time. The double-density disk drives that once were the
answer to my storage prayers no longer fill the bill. The
double-sided, double-density drives I'll get next will seem like a
vast frontier---for a while. I've heard guys with 10-megabyte hard
disks complain of feeling cramped.

Whenever there's a shortage of something---and there always will
be---tradeoffs have to be made. It's best to use the analysis phase to
anticipate most limitations and decide which tradeoffs to make.

On the other hand, you should \emph{not} consider other types of
constraints during analysis, but should instead impose them gradually
during implementation, the way one stirs flour into gravy.

The type of constraint to consider during analysis includes those that
might affect the overall approach. The type to defer includes those
that can be handled by making iterative refinements to the planned
software design.

As we heard in our earlier interviews, finding out about \emph{hardware}
constraints\index{H!Hardware constraints} often requires writing
some test code and trying things out.

Finding out about the \emph{customer's} constraints%
\index{C!Customer constraints}
is usually a matter of asking the customer, or of taking written
surveys. ``How fast do you need such-and-such, on a scale of one to
ten?'', etc.%
\index{A!Analysis!constraints|)}%
\index{C!Constraints|)}

\subsection{Building a Conceptual Model of the Solution}%
\index{A!Analysis!conceptual model|(}%
\index{C!Conceptual model|(}

A conceptual model is an imaginary solution to the problem. It is a
view of how the system \emph{appears} to work. It is an answer to all
the requirements and constraints.%
\index{C!Conceptual model!defined}

\wepsfigp{img2-047}{Refining the conceptual model to meet
requirements and constraints.}

If the requirements definition is for ``something to stand on to paint
the ceiling,'' then a description of the conceptual model is ``a
device that is free-standing (so you can paint the center of the
room), with several steps spaced at convenient intervals (so you can
climb up and down), and having a small shelf near the top (to hold
your paint can).''

A conceptual model is not quite a design, however. A design begins to
describe how the system \emph{really} works. In design, the image of a
step ladder would begin to emerge.

\Forth{} blurs the distinction a little, because all definitions are
written in conceptual terms, using the lexicons of lower level
components. In fact, later in this chapter we'll use \Forth{}
``pseudocode'' to describe conceptual model solutions.

Nevertheless, it's useful to make the distinction. A conceptual model
is more flexible than a design. It's easier to fit the requirements
and constraints into the model than into a design.

\begin{tip}
Strive to build a solid conceptual model before beginning the design.
\end{tip}

%Page 047 in first edition
%img2-47 moved forward
%Page 048 in first edition

\noindent Analysis consists of expanding the requirements definition
into a conceptual model. The technique involves two-way communication
with the customer in successive attempts to describe the model.

Like the entire development cycle, the analysis phase is best approached
iteratively. Each new requirement will tend to suggest something in
your mental model. Your job is to juggle all the requirements and
constraints until you can weave a pattern that fits the bill.

\wepsfigb{fig2-2}{An iterative approach to analysis.}

\Fig{fig2-2} illustrates the iterative approach to the analysis phase.
The final step is one of the most important: show the documented model
to the customer. Use whatever means of communication are
necessary---diagrams, tables, or cartoons---to convey your
understanding to the customer and get the needed feedback. Even if you
cycle through this loop a hundred times, it's worth the effort.

In the next three sections we'll explore three techniques for defining
and documenting the conceptual model:

\begin{enumerate}
\item defining the interfaces
\item defining the rules
\item defining the data structures.
\end{enumerate}%
\index{A!Analysis|)}

\section{Defining the Interfaces}%
\index{A!Analysis!interface definition|(}%
\index{C!Conceptual model!interface definition|(}%
\index{I!Interface definition|(}

\begin{tip}
First, and most importantly, the conceptual model should describe the
system's interfaces.\index{I!Interface definition}
\end{tip}%
\begin{interview}
\index{T!Teleska, John|(}
\person{Teleska}:

\begin{tfquot}
The ``spec'' basically deals with WHAT. In its most glorious form, it
describes what the system would look like to the user---you might call it
the user's manual. I find I write more notes on the human
interaction---what it
%Page 049 in first edition
will look like on the outside---than on the part that gets the job
done. For instance, I'll include a whole error-action listing to show
what happens when a particular error occurs. Oddly, this is the part
that takes the most time to implement anyway.

I'm currently working on a solid-state industrial washing-machine
timer. In this case, the user interface is not that complex. What is
complex is the interface to the washing machine, for which I must
depend on the customer and the documentation they can provide.

The significant interface is whatever is the arms and legs of the
product. I don't make the distinction between hardware and software at
this early stage. They can be interchanged in the implementation.

The process of designing hardware and the process of designing
software are analogous. The way I design hardware is to treat it as a
black box. The front panel is input and output. You can do the same
with software.

I use any techniques, diagrams, etc., to show the customer what the
inputs and outputs look like, using his description of what the
product has to do.  But in parallel, in my own mind, I'm imagining how
it will be implemented.  I'm evaluating whether I can do this
efficiently. So to me it's not a black box, it's a gray box. The
designer must be able to see inside the black boxes.

When I design a system that's got different modules, I try to make the
coupling as rational and as little as possible. But there's always
give and take, since you're compromising the ideal.

\index{D!Data-flow diagrams|(}
For the document itself, I use DFDs {[}data-flow diagrams, which we'll
discuss later{]}, and any other kind of representation that I can show
to my client. I show them as many diagrams as I can to clarify my
understanding.  I don't generally use these once it comes to
implementation. The prose must be complete, even without reference to
the diagrams.
\end{tfquot}%
\index{T!Teleska, John|)}
\end{interview}%
\index{E!Error handling|(}
\begin{tip}
Decide on error- and exception-handling early as part of defining the
interface.
\end{tip}

\noindent It's true that when coding for oneself, a programmer can
often concentrate first on making the code run correctly under
\emph{normal} conditions, then worry about error-handling later. When
working for someone else, however, error-handling should be worked out
ahead of time. This is an area often overlooked by the beginning
programmer.

The reason it's so important to decide on error-handling at this stage
is the wide divergence in how errors can be treated. An error might be:

\begin{itemize}
\item ignored
\item made to set a flag indicating that an error occurred, while
processing continues
\item made to halt the application immediately
\item designed to initiate procedures to correct the problem and keep
the program running.
\end{itemize}
%
%Page 050 in first edition
%
There's room for a serious communications gap if the degree of
complexity required in the error-handling is not nailed down early.
Obviously, the choice bears tremendous impact on the design and
implementation of the application.%
\index{E!Error handling|)}

\begin{tip}
Develop the conceptual model by imagining the data traveling through and
being acted upon by the parts of the model.
\end{tip}%
\index{S!Structured analysis|(}
A discipline called \emph{structured analysis} \cite{weinberg80}
offers some techniques for describing interfaces in ways that your
clients will easily understand.  One of these techniques is called the
``data-flow diagram'' (DFD), which \person{Teleska} mentioned.

\wepsfiga{fig2-3}{A data-flow diagram.}

A data-flow diagram, such as the one depicted in \Fig{fig2-3},
emphasizes what happens to items of data as they travel through the
system.  The circles represent ``transforms,'' functions that act upon
information.  The arrows represent the inputs and outputs of the
transforms.

The diagram depicts a frozen moment of the system in action. It
ignores initialization, looping structures, and other details of
programming that relate to time.

Three benefits are claimed for using DFDs:

First, they speak in simple, direct terms to the customer. If your
%Page 051 in first edition
customer agrees with the contents of your data-flow diagram, you know
you understand the problem.

Second, they let you think in terms of the logical ``whats,'' without
getting caught up in the procedural ``hows,'' which is consistent with
the philosophy of hiding information as we discussed in the last chapter.

Third, they focus your attention on the interfaces to the system and
between modules.

\Forth{} programmers, however, rarely use DFDs except for the customer's
benefit. \Forth{} encourages you to think in terms of the conceptual
model, and \Forth{}'s implicit use of a data stack makes the passing of
data among modules so simple it can usually be taken for granted.
This is because \Forth{}, used properly, approaches a functional language.%
\index{D!Data-flow diagrams|)}%
\index{S!Structured analysis|)}%
\index{P!Pseudocode|(}

For anyone with a few days' familiarity with \Forth{}, simple definitions
convey at least as much meaning as the diagrams:
\begin{Code}
: REQUEST  ( quantity part# -- )
   ON-HAND?  IF  TRANSFER  ELSE  REORDER  THEN ;
: REORDER   AUTHORIZATION?  IF  P.O.  THEN ;
: P.O.   BOOKKEEPING COPY   RECEIVING COPY
   VENDOR MAIL-COPY ;
\end{Code}
This is \Forth{} pseudocode. No effort has been made to determine what
values are actually passed on the stack, because that is an
implementation detail. The stack comment for \forth{REQUEST} is used
only to indicate the two items of data needed to initiate the process.

(If I were designing this application, I'd suggest that the user
interface be a word called \forth{NEED}, which has this syntax:

\begin{Code}
NEED 50 AXLES
\end{Code}

\noindent \forth{NEED} converts the quantity into a numeric value on
the stack, translates the string \forth{AXLES} into a part number,
also on the stack, then calls \forth{REQUEST}. Such a command should
be defined only at the outer-most level.)%
\begin{interview}
\index{J!Johnson, Dave|(}
\person{Johnson} of Moore Products Co. has a few words on \Forth{}
pseudocode:
\begin{tfquot}
IBM uses a rigorously documented PDL (program design language). We use
a PDL here as well, although we call it FDL, for \Forth{} design
language. It's probably worthwhile having all those standards, but
once you're familiar with \Forth{}, \Forth{} itself can be a design
language. You just have to leave out the so-called ``noise'' words:
\forth{C@}, \forth{DUP}, \forth{OVER}, etc., and show only the basic
flow. Most \Forth{} people probably do that informally. We do it
purposefully.
\end{tfquot}%
\index{J!Johnson, Dave|)}%
\end{interview}%
\index{P!Pseudocode|)}
%Page 052 in first edition
\begin{interview*}
During one of our interviews I asked \person{Moore}\index{M!Moore,
Charles|(} if he used diagrams of any sort to plan out the conceptual
model, or did he code straight into \Forth{}? His reply:

\begin{tfquot}
The conceptual model \emph{is} \Forth{}. Over the years I've learned
to think that way.
\end{tfquot}
Can everyone learn to think that way?

\begin{tfquot}
I've got an unfair advantage. I codified my programming style and other
people have adopted it. I was surprised that this happened. And I feel at a
lovely advantage because it is my style that others are learning to emulate.
Can they learn to think like I think? I imagine so. It's just a matter of
practice, and I've had more practice.
\end{tfquot}\index{M!Moore, Charles|)}
\end{interview*}%
\index{A!Analysis!interface definition|)}%
\index{C!Conceptual model!interface definition|)}%
\index{I!Interface definition|)}

\section{Defining the Rules}%
\index{A!Analysis!rule definition|(}%
\index{C!Conceptual model!rule definition|(}%
\index{R!Rule definition|(}%
\program{phone1}

Most of your efforts at defining a problem will center on describing
the interface.%
\index{I!Interface definition}
Some applications will also require that you define the set of
application rules.

All programming involves rules. Usually these rules are so simple it
hardly matters how you express them: ``If someone pushes the button,
ring the bell.''

Some applications, however, involve rules so complicated that they
can't be expressed in a few sentences of English. A few formal techniques
can come in handy to help you understand and document these more
complicated rules.

Here's an example. Our requirements call for a system to compute the
charges on long-distance phone calls. Here's the customer's
explanation of its rate structure. (I made this up; I have no idea how
the phone company actually computes their rates except that they
overcharge.)

\begin{tfquot}
All charges are computed by the minute, according to distance in
hundreds of miles, plus a flat charge. The flat charge for direct dial
calls during weekdays between 8 A.M. and 5 P.M. is .30 for the first
minute, and .20 for each additional minute; in addition, each minute
is charged .12 per 100 miles. The flat charge for direct calls during
weekdays between 5 P.M. and 11 P.M. is .22 for the first minute, and
.15 for each additional minute; the distance rate per minute is .10
per 100 miles. The flat charge for direct calls late during weekdays
between 11 P.M. or anytime on Saturday, Sundays, or holidays is .12
for the first minute, and .09 for each additional minute; the distance
rate per minute is .06 per 100 miles. If the call requires assistance
from the operator, the flat charge increases by .90, regardless of the hour.
\end{tfquot}
This description is written in plain old English, and it's quite a
mouthful.  It's hard to follow and, like an attic cluttered with
accumulated belongings, it may even hide a few bugs.

%Page 053 in first edition
In building a conceptual model for this system, we must describe the
rate structure in an unambiguous, useful way. The first step towards
cleaning up the clutter involves factoring out irrelevant pieces of
information---that is, applying the rules of limited redundancy. We
can improve this statement a lot by splitting it into two statements.
First there's the time-of-day rule:%
\index{A!Analysis!conceptual model|)}

\begin{tfquot}
Calls during weekdays between 8 A.M. and 5 P.M. are charged at ``full'' rate.
Calls during weekdays between 5 P.M. and 11 P.M. are charged at ``lower''
rate. Calls placed during weekdays between 11 P.M. or anytime on Saturday,
Sundays, or holidays are charged at the ``lowest'' rate.
\end{tfquot}
Then there's the rate structure itself, which should be described in
terms of ``first-minute rate,'' ``additional minute rate,'' ``distance
rate,'' and ``operator-assistance rate.''

\begin{tip}
Factor the fruit. (Don't confuse apples with oranges.)
\end{tip}
These prose statements are still difficult to read, however. System
analysts use several techniques to simplify these statements:
structured English, decision trees, and decision tables. Let's study
each of these techniques and evaluate their usefulness in the \Forth{}
environment.

\subsection{Structured English}%
\index{A!Analysis!Structured English|(}%
\index{S!Structured English|(}

Structured English is a sort of structured pseudocode in which our rate
statement would read something like this:

%Page 053/054 in first edition
\begin{Code}[baselinestretch=0.95]
IF full rate
   IF direct-dial
      IF first-minute
	 .30 + .12/100miles
      ELSE ( add'l- minute)
	 .20 + .12/100miles
      ENDIF
   ELSE ( operator )
      IF first-minute
	 1.20 + .12/100miles
      ELSE ( add'l- minute)
	 .20 + .12/100miles
      ENDIF
   ENDIF
ELSE  ( not-full-rate)
   IF lower-rate
      IF direct-dial
	 IF first-minute
	    .22 + .10/100miles
	 ELSE ( add'l- minute)
	    .15 + .10/100miles
	 ENDIF
      ELSE ( operator)
	 IF first-minute
	    1.12 + .10/100miles
	 ELSE ( add'l- minute)
	    .15 + .10/100miles
	 ENDIF
      ENDIF
   ELSE ( lowest-rate)
      IF direct-dial
	 IF first-minute
	    .12 + .06/100miles
	 ELSE ( add'l- minute)
	    .09 + .O6/100miles
	 ENDIF
      ELSE ( operator)
	 IF first-minute
	    1.02 + .O6/100miles
	 ELSE ( add'l- minute)
	    .09 + .06/100miles
	 ENDIF
      ENDIF
   ENDIF
ENDIF
\end{Code}
This is just plain awkward. It's hard to read, harder to maintain, and
hardest to write. And for all that, it's worthless at implementation
time. I don't even want to talk about it anymore.%
\index{A!Analysis!Structured English|)}%
\index{S!Structured English|)}

\subsection{The Decision Tree}%
\index{A!Analysis!decision tree|(}%
\index{D!Decision tree|(}

\wepsfigt{fig2-4}{Example of a decision tree.}

\Fig{fig2-4} illustrates the telephone rate rules by means of a
decision tree.  The decision tree is the easiest method of any to
``follow down'' to determine the result of certain conditions. For
this reason, it may be the best representation to show the customer.

Unfortunately, the decision tree is difficult to ``follow up,'' to
determine which conditions produce certain results. This difficulty
inhibits seeing ways to simplify the problem. The tree obscures the
fact that additional minutes cost the same, whether the operator
assists or not. You can't see the facts for the tree.%
\index{A!Analysis!decision tree|)}%
\index{D!Decision tree|)}

%Page 055 in first edition

% Figure 2-4 was here

\subsection{The Decision Table}%
\index{A!Analysis!decision table|(}%
\index{D!Decision table|(}

The decision table, described next, provides the most usable graphic
representation of compound rules for the programmer, and possibly for
the customer as well. \Fig{fig2-5} shows our rate structure rules in
decision-table form.

\wepsfigt{fig2-5}{The decision table.}

In \Fig{fig2-5} there are three dimensions: the rate discount, whether
an operator intervenes, and initial minute vs. additional minute.

Drawing problems with more than two dimensions gets a little tricky.
As you can see, these additional dimensions can be depicted on
%Page 056 in first edition
paper as subdimensions within an outer dimension. All of the
subdimension's conditions appear within every condition of the outer
dimension.  In software, any number of dimensions can be easily
handled, as we'll see.

All the techniques we've described force you to analyze which
conditions apply to which dimensions. In factoring these dimensions,
two rules apply:

First, all the elements of each dimension must be mutually exclusive.
You don't put ``first minute'' in the same dimension as ``direct
dial,'' because they are not mutually exclusive.

Second, all possibilities must be accounted for within each dimension.
If there were another rate for calls made between 2 A.M. to 2:05 A.M.,
the table would have to be enlarged.

But our decision tables have other advantages all to themselves.  The
decision table not only reads well to the client but actually benefits
the implementor in several ways:

\begin{description}
\item[Transferability to actual code.] This is particularly true in
\Forth{}, where decision tables are easy to implement in a form very
similar to the drawing.

\item[Ability to trace the logic upwards.] Find a condition and see what
factors produced it.

\item[Clearer graphic representation.] Decision tables serve as a better
tool for understanding, both for the implementor and the analyst.
\end{description}

Unlike decision trees, these decision tables group the \emph{results}
together in a graphically meaningful way. Visualization of ideas helps in
understanding problems, particularly those problems that are too
complex to perceive in a linear way.

For instance, \Fig{fig2-5} clearly shows that the charge for
additional minutes does not depend on whether an operator assisted or not.
With this new understanding we can draw a simplified table, as shown
in \Fig{fig2-6}.

\wepsfigt{fig2-6}{A simplified decision table.}

%Page 057 in first edition

It's easy to get so enamored of one's analytic tools that one forgets
about the problem. The analyst must do more than carry out all
possibilities of a problem to the nth degree, as I have seen authors
of books on structured analysis recommend. That approach only
increases the amount of available detail. The problem solver must also
try to simplify the problem.

\begin{tip}
You don't understand a problem until you can simplify it.
\end{tip}

\noindent If the goal of analysis is not only understanding, but
simplification, then perhaps we've got more work to do.

Our revised decision table (\Fig{fig2-6}) shows that the per-mile
charge depends only on whether the rate is full, lower, or lowest. In
other words, it's subject to only one of the three dimensions shown in
the table.  What happens if we split this table into two tables, as in
\Fig{fig2-7}?

\wepsfiga{fig2-7}{The sectional decision table.}

Now we're getting the answer through a combination of table look-up
and calculation. The formula for the per-minute charge can be
expressed as a pseudo\Forth{} definition:

\begin{Code}
: PER-MINUTE-CHARGE ( -- per-minute-charge)
        CONNECT-CHARGE  MILEAGE-CHARGE  + ;
\end{Code}
The ``\forth{+}'' now appears once in the definition,
not nine times in the table.

Taking the principle of calculation one step further, we note (or
remember from the original problem statement) that operator assistance
merely adds a one-time charge of .90 to the total charge. In this
sense, the operator charge is not a function of any of the three
dimensions. It's more
%Page 058 in first edition
appropriately expressed as a ``logical calculation''; that is, a
function that combines logic with arithmetic:

\begin{Code}
: ?ASSISTANCE
   ( direct-dial-charge -- total-charge)
   OPERATOR? IF .90 + THEN ;
\end{Code}
(But remember, this charge applies only to the first minute.)

\wepsfigb{fig2-8}{The decision table without operator involvement depicted.}

This leaves us with the simplified table shown in \Fig{fig2-8}, and an
increased reliance on expressing calculations. Now we're getting
somewhere.

Let's go back to our definition of \forth{PER-MINUTE-CHARGE}:
\begin{Code}
: PER-MINUTE-CHARGE ( -- per-minute-charge)
   CONNECT-CHARGE  MILEAGE-CHARGE  + ;
\end{Code}
Let's get more specific about the rules for the connection charge and for
the mileage charge.

The connection charge depends on whether the minute is the first or
an additional minute. Since there are two kinds of per-minute charges,
perhaps it will be easiest to rewrite \forth{PER-MINUTE-CHARGE} as two
different words.

Let's assume we will build a component that will fetch the appropriate
rates from the table. The word \forth{1MINUTE} will get the rate for
the first minute; \forth{+MINUTES} will get the rate for each
additional minute.  Both of these words will depend on the time of day
to determine whether to use the full, lower, or lowest rates.

Now we can define the pair of words to replace \forth{PER-MINUTE-CHARGE}:

%Page 059 in first edition

\begin{Code}
: FIRST  ( -- charge)
  1MINUTE  ?ASSISTANCE   MILEAGE-CHARGE + ;
: PER-ADDITIONAL  ( -- charge)
   +MINUTES  MILEAGE-CHARGE + ;
\end{Code}
What is the rule for the mileage charge? Very simple. It is the rate
(per hundred miles) times the number of miles (in hundreds). Let's
assume we can define the word \forth{MILEAGE-RATE}, which will fetch
the mileage rate from the table:

\begin{Code}
: MILEAGE-CHARGE  ( -- charge)
   #MILES @  MILEAGE-RATE * ;
\end{Code}
Finally, if we know the total number of minutes for a call, we can now
calculate the total direct-dial charge:

\begin{Code}
: TOTAL   ( -- total-charge)
   FIRST                        ( first minute rate)
   ( #minutes) 1-               ( additional minutes)
      PER-ADDITIONAL *          ( times the rate)
   +  ;                         ( added together)
\end{Code}
We've expressed the rules to this particular problem through a
combination of simple tables and logical calculations.

(Some final notes on this example: We've written something very close
to a running \Forth{} application. But it is only pseudocode. We've
avoided stack manipulations by assuming that values will somehow be on
the stack where the comments indicate. Also, we've used hyphenated
names because they might be more readable for the customer. Short
names are preferred in real code---see \Chap{5}.)\program{phone2}

We'll unveil the finished code for this example in \Chap{8}.%
\index{A!Analysis!decision table|)}%
\index{A!Analysis!rule definition|)}%
\index{C!Conceptual model!rule definition|)}%
\index{D!Decision table|)}%
\index{R!Rule definition|)}

\section{Defining the Data Structures}%
\index{A!Analysis!conceptual model|(}%
\index{A!Analysis!data structure definition|(}%
\index{C!Conceptual model|)}%
\index{C!Conceptual model!data structure definition|(}%
\index{D!Data structures:!defining|(}

After defining the interfaces, and sometimes defining the rules,
occasionally you'll need to define certain data structures as well.
We're not referring here to the implementation of the data structures,
but rather to a description of their conceptual model.

If you're automating a library index, for instance, a crucial portion
of your analysis will concern developing the logical data structure.
You'll have to decide what information will be kept for each book:
title, author, subject, etc. These ``attributes'' will comprise an
``entity'' (set of related records) called \forth{BOOKS}. Then you'll
have to determine what other data structures will be required to let
the users search the \forth{BOOKS} efficiently.

%Page 060 in first edition

\wepsfigp{img2-060}{Given two adequate solutions,
the correct one is the simpler.}

%Page 061 in first edition

%%! cut suggested by Leo
%You may need another entity consisting of authors' names in
%alphabetical order, along with ``attribute pointers'' to the books
%each author has written.

Certain constraints will also affect the conceptual model of the data
structure. In the library index example, you need to know not only
\emph{what} information the users need, but also how long they're
willing to \emph{wait} to get it.

For instance, users can request listings of topics by year of
publication---say everything on ladies' lingerie between 1900 and
1910.  If they expect to get this information in the snap of a girdle,
you'll have to index on years and on topics. If they can wait a day,
you might just let the computer search through all the books in the
library.%
\index{A!Analysis!conceptual model|)}%
\index{A!Analysis!data structure definition|)}%
\index{C!Conceptual model!data structure definition|)}%
\index{D!Data structures:!defining|)}

\section{Achieving Simplicity}%
\index{A!Analysis!simplicity|(}%
\index{S!Simplicity|(}

\begin{tip}
Keep it simple.
\end{tip}

\noindent While you are taking these crucial first steps toward
understanding the problem, keep in mind the old saying:

\begin{tfquot}
Given two solutions to a problem, the correct one is the simpler.
\end{tfquot}
This is especially true in software design. The simpler solution is often
more difficult to discover, but once found, it is:

\begin{itemize}
\item easier to understand
\item easier to implement
\item easier to verify and debug
\item easier to maintain
\item more compact
\item more efficient
\item more fun
\end{itemize}

\begin{interview}
One of the most compelling advocates of simplicity is
\person{Moore}\index{M!Moore, Charles|(}:

\begin{tfquot}
You need a feeling for the size of the problem. How much code should
it take to implement the thing? One block? Three? I think this is a
very useful design tool. You want to gut-feel whether it's a trivial
problem or a major problem, how much time and effort you should spend
on it.

When you're done, look back and say, ``Did I come up with a solution
that is reasonable?'' If your solution fills six screens, it may seem
you've used a sledgehammer to kill a mosquito. Your mental image is
out of proportion to the significance of the problem.

I've seen nuclear physics programs with hundreds of thousands of lines
of FORTRAN. Whatever that code does, it doesn't warrant hundreds of
%Page 062 in first edition
thousands of lines of code. Probably its writers have overgeneralized
the problem. They've solved a large problem of which their real needs
are a subset. They have violated the principle that the solution
should match the problem.
\end{tfquot}\index{M!Moore, Charles|)}
\end{interview}
\begin{tip}
Generality usually involves complexity. Don't generalize your solution any
more than will be required; instead, keep it changeable.
\end{tip}
\begin{interview}
\person{Moore} continues:\index{M!Moore, Charles|(}
\begin{tfquot}

Given a problem, you can code a solution to it. Having done that, and found
certain unpleasantnesses to it, you can go back and change the problem,
and end up with a simpler solution.

There's a class of device optimization---minimizing the number of gates in a
circuit-where you take advantage of the ``don't care'' situation. These
occur either because a case won't arise in practice or because you really
don't care. But the spec is often written by people who have no appreciation
for programming. The designer may have carefully specified all the cases,
but hasn't told you, the programmer, which cases are really important.

If you are free to go back and argue with him and take advantage of the
``don't cares,'' you can come up with a simpler solution.

Take an engineering application, such as a 75-ton metal powder press,
stamping out things. They want to install a computer to control the
valves in place of the hydraulic control previously used. What kind of
spec will you get from the engineer? Most likely the sensors were
placed for convenience from an electromechanical standpoint. Now they
could be put somewhere else, but the engineer has forgotten. If you
demand explanations, you can come closer to the real world and further
from their model of the world.

Another example is the PID (proportional integration and
differentiation) algorithm for servos. You have one term that
integrates, another term that differentiates, and a third term that
smooths. You combine those with 30\% integration, 10\%
differentiation, or whatever. But it's only a digital filter. It used
to be convenient in analog days to break out certain terms of the
digital filter and say, ``This is the integrator and this is the
differentiator. I'll make this with a capacitor and I'll make that
with an inductor.''

Again the spec writers will model the analog solution which was
modeling the electromechanical solution, and they're several models
away from reality. In fact, you can replace it all with two or three
coefficients in a digital filter for a much cleaner, simpler and more
efficient solution.
\end{tfquot}\index{M!Moore, Charles|)}
\end{interview}

\begin{tip}
Go back to what the problem was before the customer tried to solve it.
Exploit the ``don't cares.''
\end{tip}

%Page 063 in first edition

\wepsfigp{img2-063}{An overgeneralized solution.}

%Page 064 in first edition

\begin{interview}
\person{Moore} continues:\index{M!Moore, Charles|(}

\begin{tfquot}
Sometimes the possibilities for simplification aren't immediately
obvious.

There's this problem of zooming in a digitized graphics display, such
as CAD systems. You have a picture on the screen and you want to zoom
in on a portion to see the details.

I used to implement it so that you move the cursor to the position of
interest, then press a button, and it zooms until you have a window of
the desired size. That was the way I've always done it. Until I
realized that that was stupid. I never needed to zoom with such fine
resolution.

So instead of moving the cursor a pixel at a time, I jump the cursor
by units of, say, ten. And instead of increasing the size of box, I
jump the size of the box. You don't have a choice of sizes. You zoom
by a factor of four. The in-between sizes are not interesting. You can
do it as many times as you like.

By quantizing things fairly brutally, you make it easier to work with,
more responsive, and simpler.
\end{tfquot}\index{M!Moore, Charles|)}
\end{interview}

\begin{tip}
To simplify, quantize.\index{Q!Quantization}
\end{tip}

\begin{interview}
\person{Moore} concludes:\index{M!Moore, Charles|(}%
\begin{tfquot}
It takes arrogance to go back and say ``You didn't really mean this,''
or ``Would you mind if I took off this page and replaced it with this
expression?'' They get annoyed. They want you to do what they told you
to do.

\index{S!Stuart, LaFarr|(}
\person{LaFarr Stuart} took this attitude when he redesigned \Forth{}
\cite{stuart80}. He didn't like the input buffer, so he implemented
\Forth{} without it, and discovered he didn't really need an input buffer.%
\index{S!Stuart, LaFarr|)}

If you can improve the problem, it's a great situation to get into.
It's much more fun redesigning the world than implementing it.
\end{tfquot}\index{M!Moore, Charles|)}
\end{interview}

\noindent Effective programmers learn to be tactful and to couch their
approaches in non-threatening ways: ``What would be the consequences
of replacing that with this?'' etc.

Yet another way to simplify a problem is this:

\begin{tip}
To simplify, keep the user out of trouble.
\end{tip}
\noindent Suppose you're designing part of a word processor that
displays a directory of stored documents on the screen, one per line.
You plan that the
%Page 065 in first edition
user can move the cursor next to the name of any document, then type a
one-letter command indicating the chosen action: ``p'' for print,
``e'' for edit, etc.

Initially it seems all right to let the user move the cursor anywhere
on the screen. This means that those places where text already appears
must be protected from being overwritten. This implies a concept of
``protected fields'' and special handling. A simpler approach confines
the cursor to certain fields, possibly using reverse video to let the
user see the size of the allowable field.

Another example occurs when an application prompts the user for a
numeric value. You often see such applications that don't check input
until you press ``return,'' at which time the system responds with an
error message such as ``invalid number.'' It's just as easy---probably
easier---to check each key as it's typed and simply not allow
non-numeric characters to appear.

\begin{tip}
To simplify, take advantage of what's available.
\end{tip}%
\index{L!LaManna, Michael|(}%
\begin{interview}
\person{Michael LaManna}, a \Forth{} programmer in Long Island, New York,
comments:

\begin{tfquot}
I always try to design the application on the most powerful processor
I can get my hands on. If I have a choice between doing development on
a 68000-based system and a 6809-based system, I'd go for the
68000-based system. The processor itself is so powerful it takes care
of a lot of details I might otherwise have to solve myself.

If I have to go back later and rewrite parts of the application for a
simpler processor, that's okay. At least I won't have wasted my time.
\end{tfquot}
\end{interview}%
\index{L!LaManna, Michael|)}

\noindent A word of caution: If you're using an existing component to
simplify your prototype, don't let the component affect your design.
You don't want the design to depend on the internals of the component.%
\index{A!Analysis!simplicity|)}%
\index{S!Simplicity|)}

\section{Budgeting and Scheduling}%
\index{B!Budgeting|(}%
\index{A!Analysis!budgeting|(}%
\index{A!Analysis!scheduling|(}%
\index{S!Scheduling|(}

Another important aspect of the analysis phase is figuring the price
tag.  Again, this process is much more difficult than it would seem.
If you don't know the problem till you solve it, how can you possibly
know how long it will take to solve it?

Careful planning is essential, because things always take longer than
you expect. I have a theory about this, based on the laws of
probability:

%Page 066 in first edition

\wepsfigp{img2-066}{Conventional wisdom reveres complexity.}

%Page 067 in first edition
\begin{tip}
The mean time for making a ``two-hour'' addition to an application is
approximately 12 hours.
\end{tip}

\noindent Imagine the following scenario: You're in the middle of
writing a large application when suddenly it strikes you to add some
relatively simple feature. You think it'll take about two hours, so
without further planning, you just do it. Consider: That's two hours
coding time. The design time you don't count because you perceived the
need---and the design---in a flash of brilliance while working on the
application. So you estimate two hours.

But consider the following possibilities:

\begin{enumerate}
\item Your implementation has a bug. After two hours it doesn't work.
So you spend another two hours recoding. (Total 4.)
\item OR, before you implemented it, you realized your initial design
wouldn't work. You spend two hours redesigning. \emph{These} two hours
count. Plus another two hours coding it. (Total 4.)
\item OR, you implement the first design before you realize the design
wouldn't work. So you redesign (two more hours) and reimplement (two
more). (Total 6.)
\item OR, you implement the first design, code it, find a bug, rewrite
the code, find a design flaw, redesign, recode, find a bug in the new
code, recode again. (Total 10.)
\suspend{enumerate}
You see how the thing snowballs?
\resume{enumerate}
\item Now you have to document your new feature. Add two hours to the
above. (Total 12.)
\item After you've spent anywhere from 2 to 12 hours installing and
debugging your new feature, you suddenly find that element Y of your
application bombs out. Worst yet, you have no idea why. You spend two
hours reading memory dumps trying to divine the reason. Once you do,
you spend as many as 12 additional hours redesigning element Y. (Total
26.) Then you have to document the syntax change you made to element
Y. (Total 27.)
\end{enumerate}

\noindent That's a total of over three man-days. If all these mishaps
befell you at once, you'd call for the men with the little white
coats. It rarely gets that bad, of course, but the odds are decidedly
\emph{against} any project being as easy as you think it will be.

How can you improve your chances of judging time requirements
correctly? Many fine books have been written on this topic, notably
\emph{The Mythical Man-Month} by \person{Frederick P. Brooks}, Jr.
\cite{brooks75}.%
\index{B!Brooks,Fredrick P., Jr.}%
\index{M!Mythical@\emph{Mythical Man-Month, The} Brooks}
I have little to add to this body of knowledge except for some
personal observations.

\begin{enumerate}

\item Don't guess on a total. Break the problem up into the smallest
possible pieces, then estimate the time for each piece. The sum of the
pieces is always greater than what you'd have guessed the total would
be. (The whole appears to be less than the sum of the parts.)

%Page 068 in first edition
\item In itemizing the pieces, separate those you understand well
enough to hazard a guess from those you don't. For the second
category, give the customer a range.

\item A bit of psychology: always give your client some options.
Clients \emph{like} options. If you say, ``This will cost you \$6,000,''
the client will probably respond ``I'd really like to spend \$4,000.''
This puts you in the position of either accepting or going without a job.

But if you say, ``You have a choice: for \$4,000 I'll make it
\emph{walk} through the hoop; for \$6,000 I'll make it \emph{jump}
through the hoop. For \$8,000 I'll make it \emph{dance} through the
hoop waving flags, tossing confetti and singing `Roll Out the Barrel.'\,''

Most customers opt for jumping through the hoop.
\end{enumerate}

\begin{tip}
Everything takes longer than you think, including thinking.
\end{tip}%
\index{A!Analysis!budgeting|)}%
\index{A!Analysis!scheduling|)}%
\index{B!Budgeting|)}%
\index{S!Scheduling|)}

\section{Reviewing the Conceptual Model}%
\index{A!Analysis!conceptual model|(}%
\index{C!Conceptual model|(}

The final box on our iterative analytic wheel is labeled ``Show Model
to Customer.'' With the tools we've outlined in this chapter, this job
should be easy to do.

In documenting the requirements specification, remember that specs are
like snowmen. They may be frozen now, but they shift, slip, and melt
away when the heat is on. Whether you choose data-flow diagrams or
straight \Forth{} pseudocode, prepare yourself for the great thaw by
remembering to apply the concepts of limited redundancy.

Show the documented conceptual model to the customer. When the
customer is finally satisfied, you're ready for the next big step: the
design!%
\index{A!Analysis!conceptual model|)}%
\index{C!Conceptual model|)}

\begin{references}{9}
\bibitem{harris81} \person{Kim Harris}, ``The \Forth{} Philosophy,''
  \emph{Dr.\@ Dobb's Journal,} Vol. 6, Iss. 9, No. 59 (Sept. 81),
  pp. 6-11.
\bibitem{weinberg80} \person{Victor Weinberg}, \emph{Structured Analysis,}
  Englewood Cliffs, N.J.: Prentice-Hall, Inc., 1980.
\bibitem{stuart80} \person{LaFarr Stuart}, ``LaFORTH,''
  \emph{1980 FORML Proceedings,} p. 78.
\bibitem{brooks75} \person{Frederick P. Brooks}, Jr., \emph{The Mythical
  Man-Month,} Reading, Massachusetts, Addison-Wesley, 1975.
\end{references}

% end of chapter two

%% Thinking Forth
%% Copyright (C) 2004 Leo Brodie
%% Initial transcription by Joseph Knapka, based on OCR output
%% provided by John Hogerhuis
%%
%\chapternum{THREE}
\chapter{Preliminary Design/\hy Decomposition}\Chapmark{3}%
\index{P!Preliminary Design|(}%

\initial Assuming you have some idea of what your program should
accomplish, it's time to begin the design. The first stage,
preliminary design, focuses on shrinking your mountainous problem into
manageable molehills.

In this chapter we'll discuss two ways to decompose your \Forth{}
application.

\section{Decomposition by Component}%
\index{C!Components:!decomposition by|(}%
\index{D!Decomposition!by component|(}%
\index{P!Preliminary Design!decomposition by component|(}%

Has this sort of thing ever happened to you? You've been planning for
three months to take a weekend vacation to the mountains. You've been
making lists of what to bring, and daydreaming about the slopes.

Meanwhile you're deciding what to wear to your cousin's wedding
next Saturday. They're informal types, and you don't want to overdress.
Still, a wedding's a wedding. Maybe you should rent a tuxedo anyway.

For all this planning, it's not until Thursday that you realize the
two events coincide. You have expletives for such moments.

How is such a mental lapse possible, in one so intelligent as
yourself? Apparently the human mind actually makes links between
memories. New ideas are somehow added onto existing paths of related
thoughts.

\wepsfiga{fig3-1}{Pools of thought not yet linked}

In the mishap just described, no connection was ever made between the
two separately-linked pools of thought until Thursday.  The conflict
probably occurred when some new input (something as trivial as hearing
Saturday's weather report) got linked into both pools of thought. A
lightning flash of realization arced between the pools, followed
inexorably by thunderous panic.

A simple tool has been invented to avoid such disasters. It's called a
calendar. If you had recorded both plans in the same calendar, you would
have seen the other event scheduled, something your brain failed to do
for all its intricate magnificence.

\begin{tip}
To see the relationship between two things, put them close
together. To remind yourself of the relationship, \emph{keep} them
together.
\end{tip}
These truisms apply to software design, particularly to the preliminary
design phase. This phase is traditionally the one in which the designer
dissects a large application into smaller, programmer-sized modules.

In \Chap{1} we discovered that applications can be
conveniently decomposed into components.

\begin{tip}
The goal of preliminary design is to determine what components are
necessary to accomplish the requirements.
\end{tip}
For instance, you might have an application in which events must occur
according to some predetermined schedule. To manage the scheduling,
you might first design a few words to constitute a ``schedule-building
lexicon.'' With these words you'll be able to describe the order of events
that must occur within your application.

Thus within a single component, you'll not only share information,
but also work out potential conflicts. The wrong approach would be to let
each functional module ``know'' things about its schedule that could
potentially conflict with another module's schedule.

How can you know, in designing a component, what commands the
using components will need? Admittedly, this is something of a ``chicken
vs. egg'' problem. But \Forth{} programmers handle it the same way
chickens and eggs do: iteratively.

If the component is well-designed, completeness doesn't matter. In
fact, a component need only suffice for the current iteration's design. No
component should be considered a ``closed book'' until the application
has been completed---which, in the case of maintained applications, is
never.

As an example, imagine that your product needs to ``talk'' to other
machines in the outside world via a universal I/O chip that is part of your
system. This particular chip has a ``control register'' and a ``data
register.'' In a badly designed application, pieces of code throughout the
program would access the communication chip by simply invoking the
\code{OUT} instruction to put an appropriate command byte into the command
register. This makes the entire application needlessly dependent on that
particular chip---very risky.

Instead, \Forth{} programmers would write a component to control the I/O
chip. These commands would have logical names and a convenient
interface (usually \Forth{}'s stack) to allow usage by the rest of the
application.

For any iteration of your product's design, you would implement only
the commands needed so far---not all the valid codes that could be
sent to the ``control register.'' If later in the project cycle you
realize that you need an additional command, say one to change the
baud rate, the new command would be added to the I/O chip lexicon, not
to the code needed to set the baud rate. There's no penalty for
making this change except the few minutes (at most) it takes to edit
and recompile.

\begin{tip}
Within each component, implement only the commands needed for the
current iteration. (But don't preclude future additions.)
\end{tip}
What goes on inside a component is pretty much its own business. It's
not necessarily bad style for definitions within the component to share
redundant information.

For instance, a record in a certain data structure is fourteen bytes
long. One definition in the component advances a pointer 14 bytes to
point to the next record; another definition decrements the pointer 14
bytes.

As long as that number 14 remains a ``secret'' to the component and
won't be used elsewhere, you don't need to define it as constant. Just use
the number 14 in both definitions:

\begin{Code}
: +RECORD 14 RECORDS +! ;
: -RECORD -14 RECORD# +! ;
\end{Code}
On the other hand, if the value will be needed outside of the component,
or if it's used several times within the component and there's a good
chance that it will change, you're better off hiding it behind a name:

\begin{Code}
14 CONSTANT /RECORD
: +RECORD /RECORD RECORD# +! ;
: -RECORD /RECORD NEGATE RECORD# +! ;
\end{Code}
(The name \forth{/RECORD}, by convention, means ``bytes per record.'')

\section{Example: A Tiny Editor}\program{editor1}

Let's apply decomposition by component to a real problem. It would be
nice to design a large application right here in \Chap{3}, but alas, we
don't have the room and besides, we'd get sidetracked in trying to
understand the application.

Instead, we'll take a component from a large application that has
already been decomposed. We'll design this component by decomposing it
further, into subcomponents.

Imagine that we must create a tiny editor that will allow users to
change the contents of input fields on their terminal screen. For instance,
the screen might look like this:

%!! no bf cmtt font (use Courier?)
\begin{quote}
\textsf{\textbf{Name of Member}\quad \fbox{Justine Time\underline{~}\quad}}
\end{quote}
The editor will provide three modes for users to change the contents of
the input field:

%!! original style: label in italics, no indentation
\begin{description}

\item[Overwrite.] Typing ordinary characters overwrites any characters
that were there before.

\item[Delete.] Pressing the combination of keys ``Ctrl D'' deletes the
character under the cursor and slides the remaining characters
leftwards.

\item[Insert.] Pressing the combination of keys ``Ctrl I'' switches the
editor into ``Insert Mode,'' where subsequently typing ordinary
characters inserts them at the cursor position, sliding the remaining
characters rightwards.

\end{description}
As part of the conceptual model we should also consider the error or
exception-handling; for instance, what is the limit of the field? what
happens in insert mode when characters spill off the right? etc.

That's all the specification we have right now. The rest is up to us.

Let's try to determine what components we'll need. First, the editor
will react to keys that are typed at the keyboard. Therefore we'll
need a keystroke interpreter---some kind of routine that awaits
keystrokes and matches them up with a list of possible operations. The
keystroke interpreter is one component, and its lexicon will consist
of a single word.  Since that word will allow the editing of a field,
let's call the word \forth{EDIT}.

The operations invoked by the keystroke interpreter will comprise a
second lexicon. The definitions in this lexicon will perform the
various functions required. One word might be called \forth{DELETE}, another
\forth{INSERT}, etc. Since each of these commands will be invoked by the
interpreter, each of them will process a single keystroke.

Below these commands should lie a third component, the set of
words that implement the data structure to be edited.

\wepsfiga{fig3-2}{Generalized decomposition of the Tiny Editor problem.}

Finally, we'll need a component to display the field on the video
screen. For the sake of simplicity, let's plan on creating one word only,
\forth{REDISPLAY}, to redisplay the entire field after each key is pressed.

\begin{Code}
: EDITOR  BEGIN  KEY  REVISE  REDISPLAY  ... UNTIL ;
\end{Code}
This approach separates revising the buffer from updating the display.
For now, we'll only concentrate on revising the buffer.

Let's look at each component separately and try to determine the words
each will need. We can begin by considering the events that must occur
within the three most important editing functions: overwriting,
deleting, and inserting. We might draw something like the following on
the back of an old pizza menu (we won't pay much attention to
exception-handling in the present discussion):

\begin{description}
%!!small layout differences between the '84 and '94 edition; using '94 layout
\item[To Overwrite:]~\\[\smallskipamount]
\begin{tabular}{@{}p{2.1in}>{\ttfamily}p{2.1in}}
Store new character into byte pointer to by pointer.

Advance pointer (unless at end of field).
& \parbox[t]{2.1in}{
  F U N \fwbox{\pointto{K}} T I O N A L I T Y\\
  F U N \fwbox{\pointto{C}} T I O N A L I T Y\\
  F U N C \pointto{T} I O N A L I T Y}\\
\end{tabular}

\item[To Delete:]~\\[\smallskipamount]
\begin{tabular}{@{}p{2.1in}>{\ttfamily}p{2.1in}}
Copy leftwards, by one place, the string
beginning one place to the right of the pointer.

Store a ``blank'' into the last
position on the line.
& \parbox[t]{2.1in}{
  F U N C T I O N \pointto{S} \fwbox{A L I T Y}\\
  F U N C T I O N \fwbox{\pointto{A} L I T Y} Y\\
  F U N C T I O N \pointto{A} L I T Y \fwbox{\newlength{\charheight}\settoheight{\charheight}{K}\rule{0cm}{\charheight} }}\\
\end{tabular}

% Note, pointers-and-boxes omitted in this diagram for some reason.
\item[To Insert:]~\\[\smallskipamount]
\begin{tabular}{@{}p{2.1in}>{\ttfamily}p{2.1in}}
Copy rightwards, by one place,
the string beginning at the pointer.
%!!'84 ed: the string beginning one place\\ to the right of the pointer.

Store new character into
byte pointed to by pointer.

Advance pointer (unless at end of field).
& \parbox[t]{2.1in}{
  F U N \fwbox{\pointto{T} I O N A L I T Y}\\
  F U N \pointto{T} \fwbox{T I O N A L I T Y}\\
  F U N \fwbox{\pointto{C}} T I O N A L I T Y\\
  F U N C \pointto{T} I O N A L I T Y} \\

\end{tabular}
\end{description}
We've just developed the algorithms for the problem at hand.

Our next step is to examine these three essential procedures, looking
for useful ``names''---that is procedures or elements which can
either:

\begin{enumerate}
\item possibly be reused, or

\item possibly change
\end{enumerate}
We discover that all three procedures use something called a ``pointer.''
We need two procedures:

\begin{enumerate}
\item to get the pointer (if the pointer itself is relative, this
function will perform some computation).

\item to advance the pointer

%!!next line not part of enumeration (not indented in '94 edition).
\suspend{enumerate}
Wait, three procedures:
\resume{enumerate}
\item to move the pointer backwards
\end{enumerate}
because we will want ``cursor keys'' to move the cursor forward and back
without editing changes.

These three operators will all refer to a physical pointer somewhere
in memory. Where it is kept and how (relative or absolute) should be
hidden within this component.

Let's attempt to rewrite these algorithms in code:

\begin{Code}
: KEY#  ( returns value of key last pressed )  ... ;
: POSITION  ( returns address of character pointed-to)  ;
: FORWARD  ( advance pointer, stopping at last position)  ;
: BACKWARD  ( decrement pointer, stopping at first position)  ;
: OVERWRITE   KEY# POSITION C!  FORWARD ;
: INSERT   SLIDE>  OVERWRITE ;
: DELETE   SLIDE<  BLANK-END ;
\end{Code}
To copy the text leftwards and rightwards, we had to invent two new
names as we went along, \forth{SLIDE<} and \forth{SLIDE>} (pronounced
``slide-backwards'' and ``slide-forwards'' respectively). Both of them
will certainly use \forth{POSITION}, but they also must rely on an element
we've deferred considering: a way to ``know'' the length of the
field. We can tackle that aspect when we get to writing the third
component.  But look at what we found out already: we can describe
``Insert'' as simply ``\forth{SLIDE> OVERWRITE}''.

In other words, ``Insert'' actually \emph{uses} ``Overwrite'' even though
they appear to exist on the same level (at least to a Structured
Programmer).

Instead of probing deeper into the third component, let's lay out
what we know about the first component, the key interpreter. First we
must solve the problem of ``insert mode.'' It turns out that ``insert'' is
not just something that happens when you press a certain key, as
delete is. Instead it is a \emph{different way of interpreting}
some of the possible keystrokes.

For instance in ``overwrite'' mode, an ordinary character gets stored
into the current cursor position; but in ``insert mode'' the remainder
of the line must first be shifted right. And the backspace key works
differently when the editor is in Insert Mode as well.

Since there are two modes, ``inserting'' and ``not-inserting,'' the
keystroke interpreter must associate the keys with two possible sets of
named procedures.

We can write our keystroke interpreter as a decision table (worrying
about the implementation later):

\bigskip
\begin{tabular}{>{\ttfamily}l>{\ttfamily}l>{\ttfamily}l}
\emph{\textrm{Key}} &\emph{\textrm{Not-inserting}}& \emph{\textrm{Inserting}}\\
Ctrl-D        & DELETE              & INSERT-OFF\\
Ctrl-I        & INSERT-ON           & INSERT-OFF\\
backspace      & BACKWARD            & INSERT<	 \\
left-arrow     & BACKWARD            & INSERT-OFF\\
right-arrow    & FORWARD             & INSERT-OFF\\
return         & ESCAPE              & INSERT-OFF\\
any printable  & OVERWRITE           & INSERT    \\
\end{tabular}
\bigskip

\noindent We've placed the possible types of keys in the left column,
what they do normally in the middle column, and what they do in
``insert mode'' in the right column.

To implement what happens when ``backspace'' is pressed while in
Insert Mode, we add a new procedure:

\begin{Code}
: INSERT<   BACKWARD  SLIDE< ;
\end{Code}
(move the cursor backwards on top of the last character typed, then slide
everything to the right leftward, covering the mistake).

This table seems to be the most logical expression of the problem at
the current level. We'll save the implementation for later (\Chap{8}).

Now we'll demonstrate the tremendous value of this approach in
terms of maintainability. We'll throw ourselves a curve---a major change
of plans!

\section{Maintaining a Component-based Application}

How well will our design fare in the face of change? Envision the following
scenario:

We originally assumed that we could refresh the video display simply
by retyping the field every time a key is pressed. We even implemented
the code on our personal computer, with its memory-mapped video that
refreshes an entire line in the blink of a scan cycle. But now our
customer wants the application to run on a telephone-based network,
with all I/O being done at a not-so-fast baud rate. Since some of our
input fields are almost as wide as the video screen, maybe 65
characters, it just takes too long to refresh the entire line on every
key stroke.

We've got to change the application so that we only refresh that
part of the field that actually changes. In ``insert'' and ``delete,'' this
would mean the text to the right of the cursor. In ``overwrite'' it would
mean changing just the single character being overwritten.

This change is significant. The video refresh function, which we
cavalierly relegated to the key interpreter, now must depend on which
editing functions occur. As we've discovered, the most important names
needed to implement the key interpreter are:

\begin{Code}[fontfamily=cmss]
FORWARD
BACKWARD
OVERWRITE
INSERT
DELETE
INSERT<
\end{Code}
None of their descriptions make any reference to the video refresh
process, because that was originally assumed to happen later.

But things aren't as bad as they seem. Looking at it now, the process
\forth{OVERWRITE} could easily include a command to type the new
character where the terminal's cursor is. And \forth{SLIDE<} and
\forth{SLIDE>} could include commands to type everything to the right
of, and including, \forth{POSITION}, then reset the terminal's cursor
to its current position.

Here are our revised procedure names. The commands just added
are in boldface:

%!! "KEY# EMIT" and "RETYPE" in boldface
\begin{Code}
: OVERWRITE KEY# POSITION C! KEY# EMIT FORWARD ;
: RETYPE ( type from current position to
   end of field and reset cursor) ;
: INSERT SLIDE> RETYPE OVERWRITE ;
: DELETE SLIDE< BLANK-END RETYPE ;
\end{Code}
Since these are the only three functions that change memory, they are
the only three functions that need to refresh the screen. This idea is
critical. We must be able to make such assertions to assure program
correctness. The assertion is intrinsic to the nature of the problem.

Note that the additional problem of video refresh adds an additional
``pointer'': the current cursor position on the screen. But decomposition
by component has encouraged us to view the \forth{OVERWRITE} process as
changing both the data field and the video vision of it; similarly with
\forth{SLIDE<} and \forth{SLIDE>}. For this reason it seems natural
now to maintain only one real pointer---a relative one---from which we
can compute either the data address in memory, or the column number on
the screen.

Since the nature of the pointer is wholly hidden within the three
processes \forth{POSITION}, \forth{FORWARD}, and \forth{BACKWARD}, we
can readily accommodate this approach, even if it wasn't our first
approach.

This change may have seemed simple enough here---even obvious. If so,
it's because the technique ensures flexible design. If we had used a
traditional approach---if we had designed according to structure, or
according to data transformation through sequential processes---our
brittle design would have been shattered by the change.

To prove this assertion, we'll have to start all over again from
scratch.

\section{Designing and Maintaining a Traditional Application}

Let's pretend we haven't studied the Tiny Editor problem yet, and
we're back with a minimal set of specs. We'll also start with our
initial assumption, that we can refresh the display by retyping the
entire field after each keystroke.

According to the dictum of top-down design, let's take the
widest-angle view possible and examine the problem.  \Fig{fig3-3}
depicts the program in its simplest terms. Here we've realized that
the editor is actually a loop which keeps getting keystrokes and
performing some editing function, until the user presses the return
key.

\wepsfigb{fig3-3}{The traditional approach: view from the top.}
Inside the loop we have three modules: getting a character from the
keyboard, editing the data, and finally refreshing the display to match
the data.

Clearly most of the work will go on inside ``Process a Keystroke.''

Applying the notion of successive refinement, \fig{fig3-4} shows the
editor problem redrawn with ``Process a Keystroke'' expanded. We find it
takes several attempts before we arrive at this configuration. Designing
this level forces us to consider many things at once that we had deferred
till later in the previous try.

\wepsfigt{fig3-4}{A structure for ``Process a Keystroke.''}

For instance, we must determine all the keys that might be pressed.
More significantly, we must consider the problem of ``insert mode.'' This
realization forces us to invent a flag called \forth{INSERT-MODE} which gets
toggled by the ``Ctrl I'' key. It's used within several of the structural
lines to determine how to process a type of key.

A second flag, called \forth{ESCAPE}, seems to provide a nice structured
way of escaping the editor loop if the user presses the return key while
not in insert mode.

Having finished the diagram, we're bothered by the multiple tests
for Insert Mode. Could we test for Insert Mode once, at the beginning?
Following this notion, we draw yet another chart (\fig{fig3-5}).

As you can see, this turns out even more awkward than the first
figure. Now we're testing for each key twice. It's interesting though,
how the two structures are totally different, yet functionally
equivalent. It's enough to make one wonder whether the control
structure is terribly relevant to the problem.

\wepsfigt{fig3-5}{Another structure for ``Process a Keystroke.''}

Having decided on the first structure, we've finally arrived at the
most important modules---the ones that do the work of overwriting,
inserting, and deleting. Take another look at our expansion of
``Process a Character'' in \fig{fig3-4}. Let's consider just one of the
seven possible execution paths, the one that happens if a printable
character is pressed.

In \fig{fig3-6}(a) we see the original structural path for a printable
character.

Once we figure out the algorithms for overwriting and inserting
characters, we might refine it as shown in \fig{fig3-6}(b). But look at
that embarrassing redundancy of code (circled portions). Most
competent structured programmers would recognize that this redundancy
is unnecessary, and change the structure as shown in \fig{fig3-6}(c).
Not too bad so far, right?

\subsection{Change in Plan}

Okay, everyone, now act surprised. We've just been told that this
application won't run on a memory-mapped display. What does this
change do to our design structure?

\wepsfigt{fig3-6}{The same section, ``refined'' and ``optimized.''}

Well, for one thing it destroys ``Refresh Display'' as a separate
module. The function of ``Refresh Display'' is now scattered among the
various structural lines inside ``Process a Keystroke.'' The structure of
our entire application has changed. It's easy to see how we might have
spent weeks doing top-down design only to find we'd been barking down
the wrong tree.

What happens when we try to change the program? Let's look again
at the path for any printable character.

\fig{fig3-7} (a) shows what happens to our first-pass design when we
add refresh. Part (b) shows our ``optimized'' design with the refresh
modules expanded. Notice that we're now testing the Insert flag twice
within this single leg of the outer loop.

But worse, there's a bug in this design. Can you find it?

In both cases, overwriting and inserting, the pointer is incremented
\emph{before} the refresh. In the case of overwrite, we're displaying
the new character in the wrong position. In the case of insert, we're
typing the remainder of the line but not the new character.

Granted, this is an easy problem to fix. We need only move the refresh
modules up before ``Increment Pointer.'' The point here is: How did we
miss it? By getting preoccupied with control flow structure, a
superficial element of program design.

\wepsfigt{fig3-7}{Adding refresh.}

In contrast, in our design by components the correct solution fell
out naturally because we ``used'' the refresh component inside the editing
component. Also we used \forth{OVERWRITE} inside \forth{INSERT}.

By decomposing our application into components which use one another,
we achieved not only \emph{elegance} but a more direct path to
\emph{correctness}.\program{editor2}

\section{The Interface Component}%
\index{I!Interface component|(}

In computer science terminology, interfacing between modules has two
aspects. First, there's the way other modules invoke the module; this is
the control interface. Second, there's the way other modules pass and
receive data to and from the module; this is the data interface.

Because of \Forth{}'s dictionary structure, control is not an issue.
Definitions are invoked by being named. In this section, when we use the
term ``interface'' we're referring to data.

When it comes to data interfaces between modules, traditional wisdom
says only that ``interfaces should be carefully designed, with a
minimum of complexity.'' The reason for the care, of course, is that
each module must implement its own end of the interface (\fig{fig3-8}).

This means the presence of redundant code. As we've seen, redundant
code brings at least two problems: bulky code and poor maintainability.
A change to the interface of one module will affect the interface
of the opposite module.

\wepsfiga{fig3-8}{Traditional view of the interface as a junction.}

There's more to good interface design than that. Allow me to introduce
a design element which I call the ``interface component.'' The purpose
an interface component is to implement, and
\emph{hide information about},%
\index{I!Information-hiding}
the data interface between two or more other components
(\fig{fig3-9}).

\wepsfiga{fig3-9}{Use of the interface component.}

\begin{tip}
Both data structures and the commands involved in the communication of
data between modules should be localized in an interface component.
\end{tip}
Let me give an example from my own recent experience. One of my hobbies
is writing text formatter/editors. (I've written two of them, including
the one on which I am writing this book.)

In my latest design the formatter portion contains two components.
The first component reads the source document and decides where to
make line and page breaks, etc. But instead of sending the text directly to
the terminal or printer, it saves up a line's worth at a time in a ``line
buffer.''

Similarly, instead of sending printer-control commands---for
bold-facing, underlining, etc.---as the text is being formatted, it
defers these commands until the text is actually sent. To defer the
control commands, I have a second buffer called the ``attribute
buffer.'' It corresponds, byte-for-byte, with the line buffer, except
that each byte contains a set of flags that indicate whether the
corresponding character should be underlined, boldfaced, or whatever.

The second component displays or prints the contents of the line
buffer. The component knows whether it is transmitting to the terminal
or to the printer, and outputs the text according to the attributes
indicated by the attribute buffer.

Here we have two well-defined components---the line-formatter and
the output component---each one shouldering part of the function of the
formatter as a whole.

The data interface between these two components is fairly complex.
The interface consists of two buffers, a variable that indicates the
current number of valid characters, and finally a ``knowledge'' of what
all those attribute patterns mean.

In \Forth{} I've defined these elements together in a single screen.  The
buffers are defined with \forthb{CREATE}, the count is an ordinary
\forthb{VARIABLE}, and the attribute patterns are defined as
\forthb{CONSTANT}s, such as:
\begin{Code}
1 CONSTANT UNDERNESS ( bit mask for underlining)
2 CONSTANT BOLDNESS ( bit mask for boldface)
\end{Code}
The formatting component uses phrases like \forth{UNDERNESS SET-FLAG}
to set bits in the attribute buffer. The output component uses phrases
like \forth{UNDERNESS AND} to read the attribute buffer.

\subsection{A Design Mistake}

In designing an interface component, you should ask yourself ``What is
the set of structures and commands that must be shared by the
communicating components?'' It's important to determine what elements
belong to the interface and what elements should remain within a
single component.

In writing my text formatter, I failed to answer this question fully
and found myself with a bug. The problem was this:

I allow different type widths to be used: condensed, double width,
etc. This means not only sending different signals to the printer, but
changing the number of characters allowed per line.

I keep a variable, called \forth{WALL}, for the formatter.
\forth{WALL} indicates the right margin: the point beyond which no
more text can be set. Changing to a different type width means
changing the value of \forth{WALL} proportionately.  (Actually, this
turns out to be a mistake in itself. I should be using a finer unit of
measurement, the number of which remains constant for the line.
Changing type widths would mean changing the number of units per
character. But getting back to the mistake at hand\dots)

Alas, I was also using \forth{WALL} inside the output component to
determine how many characters to display. My reasoning was that this value
would change depending on what type-width I was using.

I was right---99\% of the time. But one day I discovered that, under
a certain condition, a line of condensed text was being somehow cut
short. The final couple of words were just missing. The reason turned out
to be that \forth{WALL} was getting changed before the output component had
a chance to use it.

Originally I had seen nothing wrong with letting the output component
blithely use the formatter's \forth{WALL} as well. Now I realized that the
formatter had to leave a separate variable for the output component,
to indicate how many valid characters were in the buffers. This would
leave any subsequent font commands free to change \forth{WALL}.

It was important that the two buffers, the attribute commands, and the
new variable were the \emph{only} elements that could be shared
between the two modules. Reaching into either module from the other
one spells trouble.

The moral of this story is that we must distinguish between data
structures that are validly used only within a single component and those
that may be shared by more than one component.

A related point:

\begin{tip}
Express in objective units any data to be shared by components.
\end{tip}
For example:

\begin{itemize}
\item	Module A measures the temperature of the oven.

\item	Module B controls the burner.

\item	Module C makes sure the door is locked if the oven is too hot.
\end{itemize}
The information of global interest is the temperature of the oven,
expressed objectively in degrees. While Module A might receive a value
representing the voltage from a heat sensor, it should convert this
value to degrees before presenting it to the rest of the application.%
\index{C!Components:!decomposition by|)}%
\index{D!Decomposition!by component|)}%
\index{I!Interface component|)}%
\index{P!Preliminary Design!decomposition by component|)}

\section{Decomposition by Sequential Complexity}%
\index{D!Decomposition!by sequential complexity|(}%
\index{P!Preliminary Design!decomposition by sequential complexity|(}%
\index{S!Sequential complexity, decomposition by|(}

We've been discussing one way to do decomposition: according to
components. The second way is according to sequential complexity.

One of \Forth{}'s rules is that a word must already have been defined to
be invoked or referred to. Usually the sequence in which words are
defined parallels the order of increasing capabilities which the words
must possess. This sequence leads to a natural organization of the
source listing. The powerful commands are simply added on top of the
elementary application (\fig{fig3-10}a).

Like a textbook, the elementary stuff comes first. A newcomer to
the project would be able to read the elementary parts of the code before
moving on the advanced stuff.

\wepsfigt{fig3-10}{Two ways to add advanced capabilities.}

But in many large applications, the extra capabilities are best
implemented as an enhancement to some private, root function in the
elementary part of the application (\fig{fig3-10}b). By being able to
change the root's capability, the user can change the capability of all the
commands that use the root.

Returning to the word processor for an example, a fairly primitive
routine is the one that starts a new page. It's used by the word that
starts a new line; when we run out of lines we must start a new
page. The word that starts a new line, in turn, is used by the routine
that formats words on the line; when the next word won't fit on the
current line, we invoke \forth{NEWLINE}. This ``uses'' hierarchy demands that
we define \forth{NEWPAGE} early in the application.

The problem? One of the advanced components includes a routine that
must be invoked by \forth{NEWPAGE}. Specifically, if a figure or table
appears in the middle of text, but at format time won't fit on what's
left of the page, the formatter defers the figure to the next page
while continuing with the text. This feature requires somehow
``getting inside of'' \forth{NEWPAGE}, so that when \forth{NEWPAGE} is
next executed, it will format the deferred figure at the top of the
new page:

\begin{Code}
: NEWPAGE  ... ( terminate page with footer)
   ( start new page with header)  ...  ?HOLDOVER ... ;
\end{Code}
How can \forth{NEWPAGE} invoke \forth{?HOLDOVER}, if \forth{?HOLDOVER} is
not defined until much later?

While it's theoretically possible to organize the listing so that the
advanced capability is defined before the root function, that approach is
bad news for two reasons.

First, the natural organization (by degree of capability) is
destroyed. Second, the advanced routines often use code that is defined
amid the elementary capabilities. If you move the advanced routines to
the front of the application, you'll also have to move any routines they
use, or duplicate the code. Very messy.

\index{V!Vectored execution|(}
You can organize the listing by degree of complexity using a technique
called ``vectoring.'' You can allow the root function to invoke (point
to) any of various routines that have been defined after the root
function itself. In our example, only the \emph{name} of the routine
\forth{?HOLDOVER} need be created early; its definition can be given later.%
\index{V!Vectored execution|)}

\Chap{7} treats the subject of vectoring in \Forth{}.%
\index{D!Decomposition!by sequential complexity|)}%
\index{P!Preliminary Design!decomposition by sequential complexity|)}%
\index{S!Sequential complexity, decomposition by|)}

\section{The Limits of Level Thinking}%
\index{L!Level@``Level'' thinking, limits of|(}

Most of us are guilty of over-emphasizing the difference between
``high-level'' and ``low-level.'' This notion is an arbitrary one. It
limits our ability to think clearly about software problems.

``Level'' thinking, in the traditional sense, distorts our efforts in
three ways:

\begin{enumerate}
\item It implies that the order of development should follow a
hierarchical structure

\item It implies that levels should be segregated from each
other, prohibiting the benefits of reusability

\item It fosters syntactical differences between levels (e.g.,
assembler vs. ``high-level'' languages) and a belief that the
nature of programming somehow changes as we move further from
machine code.
\end{enumerate}
Let's examine each of these misconceptions one by one.

\subsection{Where to Begin?}

\begin{interview}
I asked \person{Moore}\index{M!Moore, Charles|(} how he would go about
developing a particular application, a game for children. As the child
presses the digits on the numeric keypad, from zero to nine, that same
number of large boxes would appear on the screen.

	\person{Moore}:

\begin{tfquot}
	I don't start at the top and work down. Given that exact
	problem, I would write a word that draws a box. I'd start at
	the bottom, and I'd end up with a word called \forth{GO},
	which monitored the keyboard.
\end{tfquot}


\noindent How much of that is intuitive?

\begin{tfquot}
	Perhaps some degree of it. I know where I'm going so I don't
	have to start there. But also it's more fun to draw boxes than
	to program a keyboard. I'll do the thing that's most fun in
	order to get into the problem. If I have to clean up all those
	details later, that's the price I pay.
\end{tfquot}

\noindent Are you advocating a ``fun-down'' approach?%
\index{F!Fundown@``Fun-down'' approach|(}

\begin{tfquot}
	Given that you're doing it in a free-spirit fashion, yes. If
	we were giving a demonstration to a customer in two days, I'd
	do it differently. I would start with the most visible thing,
	not the most fun thing. But still not in that hierarchical
	sequence, top down. I base my approach on more immediate
	considerations such as impressing the customer, getting
	something to work, or showing other people how it's going to
	work to get them interested.

	If you define a level as ``nesting,'' then yes, it's a good
	way to decompose a problem. But I've never found the notion of
	``level'' useful. Another aspect of levels is languages,
	metalanguages, meta-metalanguages. To try and split hairs as
	to which level you are on---assembler level, first integration
	level, last integration level---it's just tedious and not
	helpful. My levels get all mixed up hopelessly.
\end{tfquot}\index{M!Moore, Charles|)}%
\index{F!Fundown@``Fun-down'' approach|)}
\end{interview}
Designing by components makes where you start less important. You
could start with the key interpreter, for instance. Its goal is to
receive keystrokes and convert them to numbers, passing these numbers
to an internally invoked word. If you substitute the \Forth{} word
\forth{.} (``dot,'' which prints a number from the stack), then we can
implement the key interpreter, test it, and debug it without using
routines that have anything to do with drawing squares.
%!! original had a typo: s/it without using/it. without using/
%   that was fixed in the '94 edition and here

On the other hand, if the application required hardware support (such as a
graphics package) that we didn't have on hand, we might want to
substitute something available, such as displaying an asterisk, just to
get into the problem. Thinking in terms of lexicons is like painting a
huge mural that spans several canvases. You work on all the canvases at
once, first sketching in the key design elements, then adding splashes of
color here and there\dots{} until the entire wall is complete.

\begin{tip}
In deciding where to start designing, look for:\medskip
\begin{itemize}
\item areas where the most creativity is required (the areas where change
is most likely)
\item areas that give the most satisfying feedback (get the juices
flowing)
\item areas in which the approach decided upon will greatly affect other
areas, or which will determine whether the stated problem can be solved at
all
\item things you should show the customer, for mutual understanding
\item things you can show the investors, if necessary for the rent.
\end{itemize}
\end{tip}

\subsection{No Segregation Without Representation}%
\index{O!Object|(}%
The second way in which levels can interfere with optimal solutions is
by encouraging segregation of the levels. A popular design construct
called the ``object'' typifies this dangerous
philosophy.\ifeightyfour\else\footnote{Editor's note: But see the
recant in the 1994 Preface on page \pageref{preface94}\ifofour, and
the clairification in the 2004 Preface on page
\pageref{preface2004}\fi. Think of something like Windows COM
``objects'' or CORBA.}\fi

An object is a portion of code that can be invoked by a single name,
but that can perform more than one function. To select a particular
function you have to invoke the object and pass it a parameter or a
group of parameters. You can visualize the parameters as representing
a row of buttons you can push to make the object do what you want.

The benefit of designing an application in terms of objects is that,
like a component, the object hides information from the rest of the
application, making revision easier.

There are several problems, though. First, the object must contain a
complicated decision structure to determine which function it must
perform. This increases object size and decreases performance. A
lexicon, on the other hand, provides all usable functions by name for
you to invoke directly.

Second, the object is usually designed to stand alone. It can't take
advantage of tools provided by supporting components. As a result, it
tends to duplicate code inside itself that will appear elsewhere in
the application. Some objects are even required to parse text in order
to interpret their parameters. Each may even use its own syntax. A
shameless waste of time and energy!

%!! should be captionless but framed
\wepsfigp{no-scrambled}{``No scrambled?''}

Finally, because the object is constructed to recognize a finite set
of possibilities, it's difficult to make additions to the row of
buttons when a new function is needed. The tools inside the object
have not been designed for reuse.%
\index{O!Object|)}

The idea of levels pervades the design of the IBM Personal
Computer. Besides the processor itself (with its own machine
instruction set, of course), there are these software levels:

\begin{itemize}
\item the set of utilities written in assembler and burned into the
system's ROM
\item the disk operating system, which invokes the utilities
\item the high-level language of choice, which invokes the operating
system and the utilities
\item and finally, any application using the language.
\end{itemize}

\noindent The ROM utilities provide the hardware-dependent routines:
those that handle the video screen, disk drives, and keyboard. You
invoke them by placing a control code in a certain register and
generating the appropriate software interrupt.

For instance, software interrupt \code{10H} causes entry to the video
routines. There are 16 of these routines. You load register \code{AH} with
the number of the video routine you want.

Unfortunately, in all 16 routines there is not one that displays a text
string. To do that, you must repeat the process of loading registers and
generating a software interrupt, which in turn must make a decision
about which routine you want, and do a few other things you don't
need---for \emph{every single character}.

Try writing a text editor in which the entire screen may need to be
refreshed with each keystroke. Slow as mail! You can't improve the speed
because you can't reuse any of the information within the video routines
except for what's provided on the outside. The stated reason for this is to
``insulate'' the programmer from device addresses and other details of
the hardware. After all, these could change with future upgrades.

The only way to efficiently implement video I/O on this machine is
to move strings directly into video memory. You can do this easily,
because the reference manual tells you the address at which video
memory starts. But this defeats the intent of the system's designers.
Your code may no longer survive a hardware revision.

By supposedly ``protecting'' the programmer from details, segregation
has defeated the purpose of information hiding. Components, in
contrast, are not segregated modules but rather cumulative additions
to the dictionary. A video lexicon would, at the very least, give a
name for the address of video memory.

It's not that anything's wrong with the concept of a bit-switch
function interface between components, when it's necessary. The
problem here is that this video component was incompletely
designed. On the other hand, if the system had been fully
integrated---operating system and drivers written in \Forth{}---the video
component would not \emph{have} to be designed to suit all needs. An
application programmer could either rewrite the driver or write an
extension to the driver using available tools from the video lexicon.

\begin{tip}
Don't bury your tools.
\end{tip}

\subsection{The Tower of Babble}

The final deception perpetrated by level thinking is that programming
languages should become qualitatively different the ``higher'' you go. We
tend to speak of high-level code as something rarefied, and low-level code
as something grubby and profane.

To some degree these distinctions have validity, but this is only the
result of certain arbitrary architectural constraints that we all accept as
the norm. We've grown accustomed to assemblers with terse mnemonics
and unnatural syntactical rules, because they're ``low-level.''

The component concept rebels against the polarity of high-level vs.
low-level. All code should look and feel the same. A component is simply
a set of commands that together transform data structures and
algorithms into useful functions. These functions can be used without
knowledge of the structures and/or algorithms within.

The distance of these structures from actual machine code is
irrelevant. The code written to toggle bits in an output port should,
in theory, look no more intimidating than the code to format a report.

Even machine code should be readable. A true \Forth{}-based
engine would enjoy a syntax and dictionary identical and continuous
with the ``high-level'' dictionary we know today.%
\index{L!Level@``Level'' thinking, limits of|)}

\section{Summary}

In this chapter we've seen two ways that applications can be
decomposed: into components, and according to sequential complexity.

Special attention should be paid to those components that serve as
interfaces between other components.

Now, if you've done preliminary design correctly, your problem is
lying at your feet in a heap of manageable pieces. Each piece represents a
problem to solve. Grab your favorite piece and turn to the next chapter.%
\index{P!Preliminary Design|)}

\section{For Further Thinking}

\emph{(Answers appear in \App{D}.)}

%!! use enumeration?
\begin{enumerate}
\item Below are two approaches to defining an editor's keyboard
interpreter.  Which would you prefer? Why?

\begin{enumerate}
\item \begin{Code}
( Define editor keys )
HEX
72 CONSTANT UPCURSOR
80 CONSTANT DOWNCURSOR
77 CONSTANT RIGHTCURSOR
75 CONSTANT LEFTCURSOR
82 CONSTANT INSERTKEY
83 CONSTANT DELETEKEY
DECIMAL
( Keystroke interpreter)
: EDITOR
   BEGIN  MORE WHILE  KEY   CASE
      UPCURSOR     OF  CURSOR-UP     ENDOF
      DOWNCURSOR   OF  CURSOR-DOWN   ENDOF
      RIGHTCURSOR  OF  CURSOR>       ENDOF
      LEFTCURSOR   OF  CURSOR<       ENDOF
      INSERTKEY    OF  INSERTING     ENDOF
      DELETEKEY    OF  DELETE        ENDOF
   ENDCASE  REPEAT ;
\end{Code}
\medbreak\item \begin{Code}
( Keystroke interpreter)
: EDITOR
   BEGIN  MORE WHILE  KEY   CASE
      72 OF  CURSOR-UP     ENDOF
      80 OF  CURSOR-DOWN   ENDOF
      77 OF  CURSOR>       ENDOF
      75 OF  CURSOR<       ENDOF
      82 OF  INSERTING     ENDOF
      83 OF  DELETE        ENDOF
   ENDCASE  REPEAT ;
\end{Code}
\end{enumerate}

\item This problem is an exercise in information hiding.
%%!!! The original stops the enumerate here. The question however goes
%%!!! to the end of the chapter.

Let's suppose we have a region of memory outside of the \Forth{}
dictionary which we want to allocate for data structures (for whatever
reason). The region of memory begins at \forth{HEX} address \code{C000}.
We want to define a series of arrays which will reside in that memory.

 We might do something like this:

\begin{Code}
HEX
C000 CONSTANT FIRST-ARRAY  ( 8 bytes)
C008 CONSTANT SECOND-ARRAY  ( 6 bytes)
C00C CONSTANT THIRD ARRAY  ( 100 bytes)
\end{Code}
Each array-name defined above will return the starting address of the
appropriate array. But notice we had to compute the correct starting
address for each array, based on how many bytes we had already
allocated. Let's try to automate this, by keeping an ``allocation
pointer,'' called \forth{>RAM}, showing where the next free byte is. We first
set the pointer to the beginning of the RAM space:

\begin{Code}
VARIABLE >RAM
C000 >RAM !
\end{Code}
Now we can define each array like this:
x\begin{Code}
>RAM @ CONSTANT FIRST-ARRAY    8 >RAM +!
>RAM @ CONSTANT SECOND-ARRAY   6 >RAM +!
>RAM @ CONSTANT THIRD-ARRAY  100 >RAM +!
\end{Code}
Notice that after defining each array, we increment the pointer by the
size of the new array to show that we've allocated that much additional
RAM.

\goodbreak
To make the above more readable, we might add these two
definitions:

\begin{Code}
: THERE ( -- address of next free byte in RAM)
     >RAM @ ;
: RAM-ALLOT ( #bytes to allocate -- )  >RAM +! ;
\end{Code}
We can now rewrite the above equivalently as:

\begin{Code}
THERE CONSTANT FIRST-ARRAY    8 RAM-ALLOT
THERE CONSTANT SECOND-ARRAY   6 RAM-ALLOT
THERE CONSTANT THIRD-ARRAY  100 RAM-ALLOT
\end{Code}
(An advanced \Forth{} programmer would probably combine these operations
into a single defining word, but that whole topic is not germane to
what I'm leading up to.)

Finally, suppose we have 20 such array definitions scattered
throughout our application.

Now, the problem: Somehow the architecture of our system changes and
we decide that we must allocate this memory such that it \emph{ends}
at \forth{HEX} address \forth{EFFF}. In other words, we must start at the
end, allocating arrays backwards. We still want each array name to return
its \emph{starting} address, however.

To do this, we must now write:

\begin{Code}
F000 >RAM ! ( EFFF, last byte, plus one)
: THERE ( -- address of next free byte in RAM)
     >RAM @ ;
: RAM-ALLOT  ( #bytes to allocate)  NEGATE >RAM +! ;
  8 RAM-ALLOT  THERE CONSTANT FIRST-ARRAY
  6 RAM-ALLOT  THERE CONSTANT SECOND-ARRAY
100 RAM-ALLOT  THERE CONSTANT THIRD-ARRAY
\end{Code}
This time \forth{RAM-ALLOT} \emph{decrements} the pointer. That's okay,
it's easy to add \forth{NEGATE} to the definition of \forth{RAM-ALLOT}.
Our present concern is that each time we define an array we must
\forth{RAM-ALLOT} \emph{before} defining it, not after. Twenty places in
our code need finding and correcting.

The words \forth{THERE} and \forth{RAM-ALLOT} are nice and friendly, but
they didn't succeed at hiding \emph{how} the region is allocated. If they
had, it wouldn't matter which order we invoked them in.

At long last, our question: What could we have done to \forth{THERE} and
\forth{RAM-ALLOT} to minimize the impact of this design change? (Again, the
answer I'm looking for has nothing to do with defining words.)
\end{enumerate}

%% Thinking Forth
%% Copyright (C) 2004 Leo Brodie
%% Initial transcription by Nils M Holm
%% Based on OCR scans by John Hogerhuis
%% 
%% Chapter: Four - Detailed Design/Problem Solving

\chapter{Detailed~Design/\penalty 0\allowhyphens Problem~Solving}\Chapmark{4}
\index{D!Detailed design|(}%
\index{D!Detailed design!problem solving techniques|(}%
\begin{tfquot}
\emph{Trivial:} I can see how to do this. I just don't know how long it
will take.\\
\emph{Non-trivial:} I haven't a \emph{clue} how to do this!

%!! right-justify paragraph
\begin{flushright}
---\emph{Operating philosophy developed at the Laboratory\\
Automation and Instrumentation Design Group,\\
Chemistry Dept., Virginia Polytechnic Institute and State University}
\end{flushright}
\end{tfquot}
\initial Once you've decided upon the components in your application, your next
step is to design those components. In this chapter we'll apply
problem-solving techniques to the detailed design of a \Forth{}
application.  This is the time for pure invention, the part that many of
us find the most fun. There's a special satisfaction in going to the mat
with a non-trivial problem and coming out the victor.

In English it's difficult to separate an idea from the words used to
express the idea. In writing a \Forth{} application it's difficult to
separate the detailed design phase from implementation, because we tend to
design in \Forth{}. For this reason, we'll get a bit ahead of ourselves in
this chapter by not only presenting a problem but also designing a
solution to it, right on through to the coded implementation.

\section{Problem-Solving Techniques}%
\index{P!Problem-solving techniques|(}%

Even neophytes can solve programming problems without devoting any
conscious thought to problem solving techniques. So what's the point in
studying techniques of problem solving? To quicken the process. By
thinking about the \emph{ways} in which we solve problems, apart from
the problems themselves, we enrich our subconscious storehouse of
techniques.

\person{G. Polya} has written several books on problem solving, especially of
the mathematical problem. The most accessible of these is
\emph{How to Solve It} \cite{polya}.%
\index{P!Polya, G.}%
\index{H!How to@\emph{How to Solve It} (Polya)}
Although solving a mathematical problem isn't quite the same as
solving a software problem, you'll find some valuable suggestions there.

The following series of tips summarize several techniques recommended by
the science of problem solving:

\begin{tip}
Determine your goal.
\end{tip}
Know what you're trying to accomplish. As we saw in \Chap{2}, this
step can be detailed further:

\index{I!Interface definition}
Determine the data interfaces: Know what data will be required to
accomplish the goal, and make sure those data are available (input).
Know what data the function is expected to produce (output). For a single
definition, this means writing the stack-effect comment.

Determine the rules;\index{R!Rule definition}
review all the facts that you know. In Chapter Two we described the
rates for computing the cost of a phone call along with the rules for
applying the rates.

\begin{tip}
Picture the problem as a whole.
\end{tip}
In the \emph{analysis} phase we separated the problem into its parts, to
clarify our understanding of each piece. We are now entering the
\emph{synthesis} phase. We must visualize the problem as a whole.

Try to retain as much information about the problem in your mind
as possible. Use words, phrases, figures and tables, or any kind of graphic
representation of the data and/or rules to help you see the maximum
information at a glance. Fill your mind to bursting with the requirements
of the problem you need to solve, the way you might fill your lungs with
air.

Now hold that mental image, the way you might hold your breath.

One of two things will happen:

You may see the solution in a flash of insight. Great! Exhale a sigh
of relief and proceed directly to implementation. Or\dots{}, the problem is
too complex or too unfamiliar to be solved so easily. In this case, you'll
have to turn your attention to analogies and partial solutions. As you do
so, it's important that you have already concentrated on the problem's
requirements all at once, engraving these requirements on your mental
retina.

\begin{tip}
Develop a plan.
\end{tip}
If the solution didn't come at a glance, the next step is to determine the
approach that you will take to solve it. Set a course for action and avoid
the trap of fumbling about aimlessly.

The following tips suggest several approaches you might consider.

\begin{tip}
Think of an analogous problem.
\index{A!Analogous problems}
\end{tip}
Does this problem sound familiar? Have you written a definition like it
before? Figure out what parts of the problem are familiar, and in what
ways this problem might differ. Try to remember how you solved it
before, or how you solved something like it.

\begin{tip}
Work forward.
\end{tip}
The normal, obvious way to attack a problem is by beginning with the
known, and proceeding to the unknown. In deciding which horse to bet
on, you'd begin with their recent histories, their current health, and so on,
apply weights to these various factors and arrive at a favorite.

\begin{tip}
Work backward.
\end{tip}
More complicated problems present many possible ways to go with the
incoming data. How do you know which route will take you closer to the
solution? You don't. This class of problem is best solved by working
backward (\Fig{fig4-1}).

\wepsfiga{fig4-1}{A problem that is easier to solve backward than
forward.}

%!! include Figure 4-1 here

\begin{tip}
Believe.
\end{tip}
Belief is a necessary ingredient for successfully working backward. We'll
illustrate with a famous mathematical problem. Suppose we have two
containers. The containers have no graduation marks, but one holds nine
gallons and the other holds four gallons. Our task is to measure out exactly
six gallons of water from the nearby stream in one of the containers
(\Fig{fig4-2}).

\wepsfigt{fig4-2}{Two containers.}

%!! include Figure 4-2 here


Try to solve this on your own before reading further.

How can we get a ``six'' out of a ``nine'' and a ``four''? We can
start out working forward\index{W!Working forwards}, by mentally
transferring water from one container to the other. For example, if we
fill the large container twice from the small container, we'll get
eight gallons. If we fill the nine-gallon container to the brim, then
empty enough water to fill the four-gallon container, we'll have
exactly five gallons in the large container.

These ideas are interesting, but they haven't gotten us six gallons.
And it's not clear how they will get us six gallons.

Let's try working backward.\index{W!Working backwards|(}
We assume we've measured six gallons
of water, and it's sitting in the large container (it won't fit in the
small one!). Now, how did we get it there? What was the state of our
containers one step previously?

There are only two possibilities (\Fig{fig4-3}):
%!! there's certainly some better way to do ordered lists in TeX
\begin{enumerate}
\item The four-gallon container was full, and we just added it to the large
   container. This implies that we already had two gallons in the large
   container. Or\dots{}
%!! line break? % enumerate suspend/resume?
\item The nine-gallon container was full, and we just poured off three gallons
   into the small container.
\end{enumerate}
Which choice? Let's make a guess. The first choice requires a two-gallon
measurement, the second requires a three-gallon measurement. In our initial
playing around, we never saw a unit like two. But we did see a difference
of one, and one from four is three. Let's go with version b.

Now comes the real trick. We must make ourselves \emph{believe} without
doubt that we have arrived at the situation described. We have just
poured off three gallons into the small container. Suspending all disbelief,
we concentrate on how we did it.

How can we pour off three gallons into the small container? If there
had already been one gallon in the small container! Suddenly we're over
the hump. The simple question now is, how do we get one gallon in the
small container? We must have started with a full nine-gallon container,
poured off four gallons twice, leaving one gallon. Then we transferred the
one gallon to the small container.

\wepsfigt{fig4-3}{Achieving the end result.}

%!! include Figure 4-3 here

%!! include cartoon of page 103 here

\wepsfigp{img4-103}{Intent on a complicated problem.}


Our final step should be to check our logic by running the problem
forwards again.

Here's another benefit of working backward:\index{W!Working backwards|)}
If the problem is unsolvable,
working backward helps you quickly prove that it has no solution.

\begin{tip}
Recognize the auxiliary problem.\index{A!Auxiliary problems}
\end{tip}
Before we've solved a problem, we have only a hazy notion of what
steps---or even how many steps---may be required. As we become more
familiar with the problem, we begin to recognize that our problem
includes one or more subproblems that somehow seem different from the
main outline of the proposed procedure.

In the problem we just solved, we recognized two subproblems: filling
the small container with one gallon and then filling the large container
with six gallons.

Recognizing these smaller problems, sometimes called
``auxiliary problems,''\index{A!Auxiliary problems}
is an important problem-solving technique. By identifying
the subproblem, we can assume it has a straightforward solution.
Without stopping to determine what that solution might be, we forge
ahead with our main problem.

(\Forth{} is ideally suited to this technique, as we'll see.)

\begin{tip}
Step back from the problem.
\end{tip}
It's easy to get so emotionally attached to one particular solution that we
forget to keep an open mind.

The literature of problem solving often employs the example of the
nine dots. It stumped me, so I'll pass it along. We have nine dots arranged
as shown in \Fig{fig4-4}. The object is to draw straight lines that
touch or pass through all nine dots, without lifting the pen off the paper.
The constraint is that you must touch all nine dots with only four lines.

\wepsfiga{fig4-4}{The nine dots problem.}

%!! include Figure 4-4 here

You can sit a good while and do no better than the almost-right
\Fig{fig4-5}. If you concentrate really hard, you may eventually conclude
that the problem is a trick---there's no solution.

\wepsfiga{fig4-5}{Not quite right.}

%!! include Figure 4-5 here

But if you sit back and ask yourself,
%!! indent paragraph
\begin{tfquot}
``Am I cheating myself out a useful tack by being narrow-minded? Am I
assuming any constraints not specified in the problem? What constraints
might they be?''
\end{tfquot}
then you might think of extending some of the lines beyond the perimeter
of the nine dots.

\begin{tip}
Use whole-brain thinking.\index{W!Whole-brain thinking}
\end{tip}
When a problem has you stumped and you seem to be getting nowhere,
relax, stop worrying about it, perhaps even forget about it for a while.

Creative people have always noted that their best ideas seem to
come out of the blue, in bed or in the shower. Many books on problem
solving suggest relying on the subconscious for the really difficult
problems.

Contemporary theories on brain functions explore the differences
between rational, conscious thought (which relies on the manipulation of
symbols) and subconscious thought (which correlates perceptions to
previously stored information, recombining and relinking knowledge in
new and useful ways).

\person{Leslie Hart}\index{H!Hart, Leslie} \cite{hart75} explains the
difficulty of solving a large problem by means of logic:

%!! begin indented paragraph
\begin{tfquot}
A huge load is placed on that one small function of the brain that can be
brought into the attention zone for a period. The feat is possible, like
the circus act, but it seems more sensible to\dots{} use the full
resources of our glorious neocortex\dots{} the multibillion-neuron
capacity of the brain.

\dots{} The work aspect lies in providing the brain with raw input, as in
observing, reading, collecting data, and reviewing what others have
achieved.  Once in, [subconscious] procedures take over, simultaneously,
automatically, outside of the attention zone.

\dots{} It seems apparent\dots{} that a search is going on during the
interval, though not necessarily continuously, much as in a large
computer. I would hazard the guess that the search ramifies, starts and
stops, reaches dead ends and begins afresh, and eventually assembles an
answer that is evaluated and then popped into conscious attention---often
in astonishingly full-blown detail.
\end{tfquot}
%!! end indented paragraph

\begin{tip}
Evaluate your solution. Look for other solutions.
\end{tip}
You may have found one way of skinning the cat. There may be other
ways, and some of them may be better.

Don't invest too much effort in your first solution without asking
yourself for a second opinion.
\index{P!Problem-solving techniques|)}%
\index{D!Detailed design!problem solving techniques|)}

%!! include cartoon of page 106 here
\wepsfigp{img4-106}{``I'm not just sleeping. I'm using my neocortex.''}

\section{Interview with a Software Inventor}

\index{B!Burgess, Donald A.|(}
\begin{interview}
\person{Donald A. Burgess}, owner and president of Scientek Instrumentation,
Inc.:

%!! begin indented paragraph
\begin{tfquot}
I have a few techniques I've found useful over the years in designing
anything, to keep myself flexible.
My first rule is, ``Nothing is impossible.''
My second rule is, ``Don't forget, the object is to make a buck.''

First examine the problem, laying out two or three approaches on paper.
Then try the most appealing one, to see if it works. Carry it through. Then
deliberately go all the way back to the beginning, and start over.

Starting over has two values. First, it gives you a fresh approach. You
either gravitate back to the way you started, or the way you started
gravitates toward the new way.

Second, the new approach may show all kinds of powerful possibilities. Now
you have a benchmark. You can look at both approaches and compare the
advantages of both. You're in a better position to judge.

Getting stuck comes from trying too hard to follow a single approach.
Remember to say, ``I want this kumquat crusher to be different. Let's
reject the traditional design as not interesting. Let's try some crazy
ideas.''

The best thing is to start drawing pictures. I draw little men. That keeps
it from looking like ``data'' and interfering with my thinking process. The
human mind works exceptionally well with analogies. Putting things in
context keeps you from getting stuck within the confines of any language,
even \Forth{}.

When I want to focus my concentration, I draw on little pieces of paper.
When I want to think in broad strokes, to capture the overall flow, I draw
on great big pieces of paper. These are some of the crazy tricks I use to keep
from getting stagnant.

When I program in \Forth{}, I spend a day just dreaming, kicking around
ideas. Usually before I start typing, I sketch it out in general terms. No
code, just talk. Notes to myself.

Then I start with the last line of code first. I describe what I would like
to do, as close to English as I can. Then I use the editor to slide this
definition towards the bottom of the screen, and begin coding the internal
words. Then I realize that's a lousy way to do it. Maybe I split my top word
into two and transfer one of them to an earlier block so I can use it earlier.
I run the hardware if I have it; otherwise I simulate it.

\Forth{} requires self-discipline. You have to stop diddling with the
keyboard. \Forth{} is so willing to do what I tell it to, I'll tell it to
do all kinds of ridiculous things that have nothing to do with where I'm
trying to go. At those times I have to get away from the keyboard.

\Forth{} lets you play. That's fine, chances are you'll get some ideas. As
long as you keep yourself from playing as a habit. Your head is a whole lot
better than the computer for inventing things.
\end{tfquot}
\end{interview}%
\index{B!Burgess, Donald A.|)}

\section{Detailed Design}%
\index{L!Lexicons|(}

We're now at the point in the development cycle at which we've decided
we need a component (or a particular word). The component will consist
of a number of words, some of which (those that comprise the lexicon) will
be used by other components and some of which (the internal words) will
be only used within this component.

Create as many words as necessary to obey the following tip:

\begin{tip}
Each definition should perform a simple, well-defined task.
\end{tip}%
\index{D!Detailed design!steps in|(}
Here are the steps generally involved in designing a component:
%!! there's certainly some better way to do ordered lists in TeX
\begin{enumerate}
\item Based on the required functions, decide on the names and syntax for the
   external definitions (define the interfaces).
\item Refine the conceptual model by describing the algorithm(s) and data
   structure(s).
\item Recognize auxiliary definitions.
\item Determine what auxiliary definitions and techniques are already available.
\item Describe the algorithm with pseudocode.
\item Implement it by working backwards from existing definitions to the inputs.
\item Implement any missing auxiliary definitions.
\item If the lexicon contains many names with strong elements in common,
   design and code the commonalities as internal definitions, then implement
   the external definitions.
\end{enumerate}%
\index{D!Detailed design!steps in|)}
We'll discuss the first two steps in depth. Then we'll engage in an
extended example of designing a lexicon.

\section{\Forth{} Syntax}%
\index{D!Detailed design!forth@\Forth{} syntax|(}%
\index{F!forth@\Forth{}!syntax|(}%
\index{S!Syntax \Forth{}|(}

At this point in the development cycle you must decide how the words in
your new lexicon will be used in context. In doing so, keep in mind how
the lexicon will be used by subsequent components.

\begin{tip}
In designing a component, the goal is to create a lexicon that will make your
later code readable and easy to maintain.
\end{tip}
Each component should be designed with components that use it in mind.
You must design the syntax of the lexicon so that the words make sense
when they appear in context. Hiding interrelated information within the
component will ensure maintainability, as we've seen.

At the same time, observe \Forth{}'s own syntax. Rather than insisting
on a certain syntax because it seems familiar, you may save
yourself from writing a lot of unnecessary code by choosing a syntax that
\Forth{} can support without any special effort on your part.

Here are some elementary rules of \Forth{}'s natural syntax:

\begin{tip}
Let numbers precede names.
\index{N!Numbers-precede-names syntax rule}
\end{tip}
%!! include cartoon of page 110 here
Words that require a numeric argument will naturally expect to find that
number on the stack. Syntactically speaking, then, the number should
precede the name. For instance, the syntax of the word \forth{SPACES}, which
emits ``$n$'' number of spaces, is

\begin{Code}
20 SPACES
\end{Code}
Sometimes this rule violates the order that our ear is accustomed to
hearing. For instance, the \Forth{} word \forth{+} expects to be preceded
by both arguments, as in

\begin{Code}
3 4 +
\end{Code}
\index{P!Postfix notation|(}
This ordering, in which values precede operators, is called ``postfix.''

\Forth{}, in its magnanimity, won't \emph{insist} upon postfix notation.
You could redefine \forth{+} to expect one number in the input stream,
like this:

\begin{Code}
3 + 4
\end{Code}
by defining it so:
\begin{Code}
: +   BL WORD  NUMBER DROP  + ;
\end{Code}
(where \forthb{WORD} is 79/83 Standard, returning an address, and
\forthb{NUMBER} returns a double-length value as in the 83 Standard
Uncontrolled Reference Words).

Fine. But you wouldn't be able to use this definition inside other colon
definitions or pass it arguments, thereby defeating one of \Forth{}'s
major advantages.

Frequently, ``noun'' type words pass their addresses (or any type of
pointer) as a stack argument to ``verb'' type words. The \Forth{}-like
syntax of
\begin{quote}
%!! sans-serif
{\sf ``noun'' ``verb''}
\end{quote}
\wepsfigxx{img4-110}
will generally prove easiest to implement because of the stack.\medbreak

In some cases this word order\index{W!Words:!ordering} sounds
unnatural.  For instance, suppose we have a file named
\forth{INVENTORY}. One thing we can do with that file is \forth{SHOW}
it; that is, format the information in pretty columns. If
\forth{INVENTORY} passes a pointer to \forth{SHOW}, which acts upon
it, the syntax becomes
\begin{Code}
INVENTORY SHOW
\end{Code}
If your spec demands the English word-order,\index{W!Words:!ordering}
\Forth{} offers ways to
achieve it. But most involve new levels of complexity. Sometimes the
best thing to do is to choose a better name. How about
\begin{Code}
INVENTORY REPORT
\end{Code}
(We've made the ``pointer'' an adjective, and the ``actor'' a noun.)

If the requirements insist on the syntax
\begin{Code}
SHOW INVENTORY
\end{Code}
we have several options. \forth{SHOW} might set a flag and
\forth{INVENTORY} would act according to the flag. Such an approach has
certain disadvantages, especially that \forth{INVENTORY} must be ``smart''
enough to know all the possible actions that might be taken on it. (We'll
treat these problems in Chapters \ref{chapter-7} and \ref{chapter-8}.)%
\index{P!Postfix notation|)}

Or, \forth{SHOW} might look ahead at the next word in the input stream.
We'll discuss this approach in a tip, ``Avoid expectations,'' later in this
chapter.

Or, the recommended approach, \forth{SHOW} might set an ``execution
variable'' that \forth{INVENTORY} will then execute. (We'll discuss vectored
execution in \Chap{7}.)

\begin{tip}
Let text follow names.
\index{T!Text-follows-names syntax rule}
\end{tip}
If the \Forth{} interpreter finds a string of text that is neither a number
nor a predefined word, it will abort with an error message. For this
reason, an undefined string must be preceded by a defined word.

An example is \forth{."} (dot-quote), which precedes the text it will
later print. Another example is \forthb{CREATE} (as well as all defining
words), which precedes the name that is, at the moment, still undefined.

The rule also applies to defined words that you want to refer to, but
not execute in the usual way. An example is \forthb{FORGET}, as in

\begin{Code}
FORGET TASK
\end{Code}
Syntactically, \forthb{FORGET} must precede \forth{TASK} so that
\forth{TASK} doesn't execute.

\begin{tip}
Let definitions consume their arguments.
\end{tip}
This syntax rule is more a convention of good \Forth{} programming
than a preference of \Forth{}.
\index{D!Definitions-consume-arguments syntax rule}

Suppose you're writing the word \forth{LAUNCH}, which requires the
number of a launch pad and fires the appropriate rocket. You want the
definition to look roughly like this:
\begin{Code}
: LAUNCH  ( pad#)  LOAD  AIM  FIRE ;
\end{Code}
Each of the three internal definitions will require the same argument, the
launch pad number. You'll need two \forthb{DUP}s somewhere. The question
is where? If you put them inside \forth{LOAD} and \forth{AIM}, then you
can keep them out of \forth{LAUNCH}, as in the definition above. If you
leave them out of \forth{LOAD} and \forth{AIM}, you'll have to define:
\begin{Code}
: LAUNCH  ( pad#)  DUP LOAD  DUP AIM  FIRE ;
\end{Code}
By convention, the latter version is preferable, because \forth{LOAD} and
\forth{AIM} are cleaner. They do what you expect them to do. Should you
have to define \forth{READY}, you can do it so:
\begin{Code}
: READY  ( pad#)  DUP LOAD  AIM ;
\end{Code}
and not
\begin{Code}
: READY  ( pad#)  LOAD  AIM  DROP ;
\end{Code}

\begin{tip}
Use zero-relative numbering.
\end{tip}
By habit we humans number things starting with one: ``first, second,
third,'' etc. Mathematical models, on the other hand, work more naturally
\index{Z!Zero-relative numbering|(}
when starting with zero. Since computers are numeric processors, software
becomes easier to write when we use zero-relative numbering.

To illustrate, suppose we have a table of eight-byte records. The first
record occupies the first eight bytes of the table. To compute its
starting address, we add ``0'' to \forth{TABLE}. To compute the starting
address of the ``second'' record, we add ``8'' to \forth{TABLE}.

\wepsfiga{fig4-6}{A table of 8-byte records.}

%!! include Figure 4-6 here
\medbreak
It's easy to derive a formula to achieve these results:
%!! should be two columns, right-hand side sans-serif

\bigskip
\begin{tabular}{@{}l@{ }l@{}r}
\sf first record starts at:  &  $\mathsf{0 \times 8} = {}$ & \sf 0  \\
\sf second record starts at: &  $\mathsf{1 \times 8} = {}$ & \sf 8  \\
\sf third record starts at:  &  $\mathsf{2 \times 8} = {}$ & \sf 16 \\
\end{tabular}
\bigskip

\noindent
We can easily write a word which converts a record\# into the address
where that record begins:

\begin{Code}
: RECORD  ( record# -- adr )
   8 *  TABLE + ;
\end{Code}
Thus in computer terms it makes sense to call the ``first record'' the 0th
record.

If your requirements demand that numbering start at one, that's
fine. Use zero-relative numbering throughout your design and then, only
in the ``user lexicons'' (the set of words that the end-user will use)
include the conversion from zero-to one-relative numbering:
\begin{Code}
: ITEM  ( n -- adr)  1- RECORD ;
\end{Code}
\index{Z!Zero-relative numbering|)}

\begin{tip}
Let addresses precede counts.\index{A!Address-precedes-counts syntax rule}
\end{tip}
Again, this is a convention, not a requirement of \Forth{}, but such
conventions are essential for readable code. You'll find examples of this
rule in the words \forthb{TYPE}, \forthb{ERASE}, and \forthb{BLANK}.

\begin{tip}
Let sources precede destinations.
\index{S!Source-precede-destination syntax rule}
\end{tip}
Another convention for readability. For instance, in some systems, the
phrase
\begin{Code}
22 37 COPY
\end{Code}
copies Screen 22 to Screen 37. The syntax of \forth{CMOVE} incorporates both
this convention and the previous convention:
\begin{Code}[commandchars=\&\{\}]
source destination count &poorbf{CMOVE}
\end{Code}

\index{E!Expectations-in-input-stream syntax rule|(}
\begin{tip}
Avoid expectations (in the input stream).
\end{tip}
Generally try to avoid creating words that presume there will be other
words in the input stream.

Suppose your color computer represents blue with the value 1, and
light-blue with 9. You want to define two words: \forth{BLUE} will return
1; \forth{LIGHT} may precede \forth{BLUE} to produce 9.

In \Forth{}, it would be possible to define \forth{BLUE} as a constant, so
that when executed it always returns 1.\program{color1}

\begin{Code}
1 CONSTANT BLUE
\end{Code}
And then define \forth{LIGHT} such that it looks for the next word in the
input stream, executes it, and ``ors'' it with 8 (the logic of this will
become apparent when we visit this example again, later in the book):
\begin{Code}
: LIGHT  ( precedes a color)  ( -- color value)
     ' EXECUTE  8 OR ;
\end{Code}
(in fig-\Forth{}:
\begin{Code}[commandchars=\&\{\}]
: LIGHT [COMPILE] '  CFA EXECUTE  8 OR ;&textrm{)}
\end{Code}
%!! The closing paren at the end of this code is a bit misleading.
%!! At least a leading space and a different font should be used.
\noindent (For novices: The apostrophe in the definition of \forth{LIGHT}
is a \Forth{} word called ``tick.'' Tick is a dictionary-search word; it
takes a name and looks it up in the dictionary, returning the address
where the definition resides. Used in this definition, it will find the
address of the word following \forth{LIGHT}---for instance,
\forth{BLUE}---and pass this address to the word \forthb{EXECUTE}, which
will execute \forth{BLUE}, pushing a one onto the stack.  Having ``sucked
up'' the operation of \forth{BLUE}, \forth{LIGHT} now ``or''s an 8 into
the 1, producing a 9.)

This definition will work when invoked in the input stream, but special
handling is required if we want to let \forth{LIGHT} be invoked within a
colon definition, as in:
\begin{Code}
: EDITING   LIGHT BLUE BORDER ;
\end{Code}
Even in the input stream, the use of \forth{EXECUTE} here will cause a
crash if \forth{LIGHT} is accidentally followed by something other than a
defined word.

The preferred technique, if you're forced to use this particular
syntax, is to have \forth{LIGHT} set a flag, and have \forth{BLUE}
determine whether that flag was set, as we'll see later on.

There will be times when looking ahead in the input stream is desirable,
even necessary. (The proposed \forth{TO} solution is often implemented
this way \cite{rosen82}.)

But generally, avoid expectations. You're setting yourself up for
disappointment.
\index{E!Expectations-in-input-stream syntax rule|)}

\begin{tip}
Let commands perform themselves.
\end{tip}%
\index{C!Compilers|(}%
This rule is a corollary to ``Avoid expectations.'' It's one of
\Forth{}'s philosophical quirks to let words do their own work. Witness
the \Forth{} compiler (the function that compiles colon definitions),
caricatured in \Fig{fig4-7}. It has very few rules:

\wepsfigt{fig4-7}{The traditional compiler vs. the \Forth{} compiler.}

%!! include Figure 4-7 here

%!! this paragraph should be an indented, unordered list
\begin{itemize}
\item Scan for the next word in the input stream and look it up in the
dictionary.

\item If it's an ordinary word, \emph{compile} its address.

\item If it's an ``immediate'' word, \emph{execute} it.

\item If it's not a defined word, try to convert it to a number and
compile it as a literal.

\item If it's not a number, abort with an error message.
\end{itemize}
Nothing is mentioned about compiling-words such as \forthb{IF},
\forthb{ELSE}, \forthb{THEN}, etc. The colon compiler doesn't know about
these words. It merely recognizes certain words as ``immediate'' and
executes them, letting them do their own work. (See \emph{Starting
\Forth{}}, Chapter Eleven, ``How to Control the Colon Compiler.'')

The compiler doesn't even ``look for'' semicolon to know when to
stop compiling. Instead it executes semicolon, allowing it to do the work
of ending the definition and shutting off the compiler.

There are two tremendous advantages to this approach. First, the
compiler is so simple it can be written in a few lines of code. Second,
there's no limit on the number of compiling words you can add at any
time, simply by making them immediate. Thus, even \Forth{}'s colon
compiler is extensible!

\index{I!Interpreters|(}
\Forth{}'s text interpreter and \Forth{}'s address interpreter also
adhere to this same rule.

The following tip is perhaps the most important in this chapter:


\begin{tip}
Don't write your own interpreter/compiler when you can use \Forth{}'s.
\end{tip}
One class of applications answers a need for a special purpose
language---a self-contained set of commands for doing one particular
thing. An example is a machine-code assembler. Here you have a large
group of commands, the mnemonics, with which you can describe the
instructions you want assembled. Here again, \Forth{} takes a radical
departure from mainstream philosophy.

Traditional assemblers\index{A!Assemblers} are special-purpose
interpreters---that is, they are complicated programs that scan the
assembly-language listing looking for recognized mnemonics such as
\code{ADD}, \code{SUB}, \code{JMP}, etc., and
assemble machine instructions correspondingly. The \Forth{} assembler,
however, is merely a lexicon of \Forth{} words that themselves assemble
machine instructions.

There are many more examples of the special purpose language,
each specific to individual applications. For instance:

%!! begin ordered list
\begin{enumerate}
\item If you're building an Adventure-type game, you'd want to write a
language that lets you create and describe monsters and rooms, etc. You
might create a defining word called \forth{ROOM} to be used like this:

\begin{Code}
ROOM DUNGEON
\end{Code}
Then create a set of words to describe the room's attributes by building
unseen data structures associated with the room:

\begin{Code}
EAST-OF DRAGON-LAIR
WEST-OF BRIDGE
CONTAINING POT-O-GOLD
etc.
\end{Code}
%!! 'etc.' should not be part of the code, should it?
The commands of this game-building language can simply be \Forth{}
words, with \Forth{} as the interpreter.
% typo: originally, ``WORDS'' was upper case. Doesn't make sense here.

\item If you're working with Programmable Array Logic (PAL) devices,
\index{P!Programmable Array Logic (PAL)}
you'd like a form of notation that lets you describe the behavior of
the output pins in logical terms, based on the states of the input
pins. A PAL programmer was written with wonderful simplicity in
\Forth{} by \person{Michael Stolowitz} \cite{stolowitz82}.
\index{S!Stolowitz, Michael}

\item If you must create a series of user menus to drive your application,
you might want to first develop a menu-compiling language. The words of
this new language allow an application programmer to quickly program the
needed menus---while hiding information about how to draw borders, move
the cursor, etc.
\end{enumerate}
All of these examples can be coded in \Forth{} as lexicons, using the
normal \Forth{} interpreter, without having to write a special-purpose
interpreter or compiler.
\index{C!Compilers|)}

\index{M!Moore, Charles|(}
\begin{interview}
\person{Moore}:
\begin{tfquot}
A simple solution is one that does not obscure the problem with
irrelevancies.  It's conceivable that something about the problem requires
a unique interpreter. But every time you see a unique interpreter, it
implies that there is something particularly awkward about the problem.
And that is almost never the case.

If you write your own interpreter, the interpreter is almost certainly the
most complex, elaborate part of your entire application. You have switched
from solving a problem to writing an interpreter.

I think that programmers like to write interpreters. They like to do these
elaborate difficult things. But there comes a time when the world is going
to have to quit programming keypads and converting numbers to binary,
and start solving problems.
\end{tfquot}
\end{interview}%
\index{M!Moore, Charles|)}%
\index{D!Detailed design!forth@\Forth{} syntax|)}%
\index{F!forth@\Forth{}!syntax|)}%
\index{S!Syntax \Forth{}|)}%
\index{I!Interpreters|)}%
\index{L!Lexicons|)}

\section{Algorithms and Data Structures}%
\index{A!Algorithms|(}%
\index{D!Detailed design!algorithms|(}

In \Chap{2} we learned how to describe a problem's requirements in terms
of interfaces and rules. In this section we'll refine the conceptual model
for each component into clearly defined algorithms and data structures.

An algorithm is a procedure, described as a finite number of rules, for
accomplishing a certain task. The rules must be unambiguous and guaranteed
to terminate after a finite number of applications. (The word is named for
the ninth century Persian mathematician \person{al-Khowarizimi}.)

An algorithm lies halfway between the imprecise directives of human
speech, such as ``Please sort these letters chronologically,'' and the
precise directives of computer language, such as ``\forth{BEGIN 2DUP < IF}
\dots{}'' etc. The algorithm for sorting letters chronologically might be
this:
\begin{enumerate}
\item Take an unsorted letter and note its date.
\item Find the correspondence folder for that month and year.
\item Flip through the letters in the folder, starting from the front, until
   you find the first letter dated later than your current letter.
\item Insert your current letter just in front of the letter dated later.
   (If the folder is empty, just insert the letter.)
\end{enumerate}
There may be several possible algorithms for the same job. The algorithm
given above would work fine for folders containing ten or fewer letters,
but for folders with a hundred letters, you'd probably resort to a more
efficient algorithm, such as this:
\begin{enumerate}
\item (same)
\item (same)
\item If the date falls within the first half of the month, open the
folder a third of the way in. If the letter you find there is dated later
than your current letter, search forward until you find a letter dated the
same or before your current letter. Insert your letter at that point. If
the letter you find is dated earlier than your current letter, search
backward\dots{}
\end{enumerate}
\dots{} You get the point. This second algorithm is more complicated than
the first. But in execution it will require fewer steps on the average
(because you don't have to search clear from the beginning of the folder
every time) and therefore can be performed faster.

\index{D!Data structures:!defined|(}
A data structure is an arrangement of data or locations for data,
organized especially to match the problem. In the last example, the file
cabinet containing folders and the folders containing individual letters
can be thought of as data structures.
\index{D!Data structures:!defined|)}%

The new conceptual model includes the filing cabinets and folders
(data structures) plus the steps for doing the filing (algorithms).
\index{A!Algorithms|)}%
\index{D!Detailed design!algorithms|)}

\section{Calculations vs. Data Structures vs. Logic}%
\index{C!Calculations|(}%
\index{D!Data structures:!calculation@vs. calculation vs. logic|(}%
\index{D!Detailed design!calculations|(}%
\index{D!Detailed design!data structures|(}%
\index{D!Detailed design!logic|(}%
\index{L!Logic|(}

We've stated before that the best solution to a problem is the simplest
adequate one; for any problem we should strive for the simplest
approach.

Suppose we must write code to fulfill this specification:

%!! sans-serif
\begin{Code}[fontfamily=cmss]
if the input argument is 1, the output is 10
if the input argument is 2, the output is 12
if the input argument is 3, the output is 14
\end{Code}
There are three approaches we could take:

\begin{description}
\item[Calculation]~
\begin{Code}
( n)  1-  2*  10 +
\end{Code}
\medbreak
\item[Data Structure]~
\begin{Code}
CREATE TABLE  10 C,  12 C,  14 C,
( n)  1- TABLE + C@
\end{Code}
\medbreak
\item[Logic]~
\begin{Code}
( n)  CASE
         1 OF 10 ENDOF
         2 OF 12 ENDOF
         3 OF 14 ENDOF  ENDCASE
\end{Code}
\end{description}
In this problem, calculation is simplest. Assuming it is also adequate
(speed is not critical), calculation is best.

The problem of converting angles to sines and cosines can be implemented
more simply (at least in terms of lines of code and object size)
by calculating the answers than by using a data structure. But for many
applications requiring trig, it's faster to look up the answer in a table
stored in memory. In this case, the simplest \emph{adequate} solution is
using the data structure.

In \Chap{2} we introduced the telephone rate problem. In that
problem the rates appeared to be arbitrary, so we designed a data
structure:

\bigskip
%!! this should be a table; see pg 119
\begin{tabular}{lccc}
          & \emph{Full Rate} & \emph{Lower Rate} & \emph{Lowest Rate} \\ \hline
First Min.  &       .30        &        .22        &          .12 \\ \hline
Add'1 Mins. &       .12        &        .10        &          .06 \\ \hline
\end{tabular}\program{phone3}
\bigskip

\noindent Using a data structure was simpler than trying to invent a
formula by which these values could be calculated. And the formula might
prove wrong later. In this case, table-driven code is easier to maintain.

In \Chap{3} we designed a keystroke interpreter for our Tiny
Editor using a decision table:\program{editor3}

%!! this should be a table; see pg 119
%%typo: The original has here ``Cntrl''. Usually, it's Ctrl, and other chapters
%%use that name, too.
\medskip
\begin{tabular}{lll}
\emph{Key}         & \emph{Not-Inserting} & \emph{Inserting} \\
\texttt{Ctrl-D}    & \texttt{DELETE}      & \texttt{INSERT-OFF} \\
\texttt{Ctrl-I}    & \texttt{INSERT-ON}   & \texttt{INSERT-OFF} \\
\texttt{backspace} & \texttt{BACKWARD}    & \texttt{INSERT<} \\
etc. & &
\end{tabular}
\medskip

\medbreak\noindent
We could have achieved this same result with logic:
\begin{Code}
CASE
   CTRL-D     OF  'INSERTING @  IF
      INSERT-OFF   ELSE DELETE     THEN   ENDOF
   CTRL-I     OF  'INSERTING @  IF
      INSERT-OFF   ELSE INSERT-ON  THEN   ENDOF
   BACKSPACE  OF  'INSERTING @  IF
      INSERT<      ELSE BACKWARD   THEN   ENDOF
ENDCASE
\end{Code}
but the logic is more confusing. And the use of logic to express such a
multi-condition algorithm gets even more convoluted when a table was
not used in the original design.

The use of logic becomes advisable when the result is not calculable,
or when the decision is not complicated enough to warrant a decision
table. \Chap{8} is devoted to the issue of minimizing the use of logic
in your programs.

\begin{tip}
In choosing which approach to apply towards solving a problem, give
preference in the following order:
\medskip
%!! there's certainly some better way to do ordered lists in TeX
\begin{enumerate}
\item calculation (except when speed counts)
\item data structures
\item logic
\end{enumerate}
\end{tip}
Of course, one nice feature of modular languages such as \Forth{} is that
the actual implementation of a component---whether it uses calculation,
data structures, or logic---doesn't have to be visible to the rest of the
application.
\index{C!Calculations|)}%
\index{D!Data structures:!calculation@vs. calculation vs. logic|)}%
\index{D!Detailed design!calculations|)}%
\index{D!Detailed design!data structures|)}%
\index{D!Detailed design!logic|)}%
\index{L!Logic|)}

\section{Solving a Problem: Computing Roman Numerals}%
\index{D!Detailed design!demonstration of|(}%
\index{R!Roman numeral computation problem|(}%
\program{roman1}

In this section we'll attempt to demonstrate the process of designing a
lexicon. Rather than merely present the problem and its solution, I'm
hoping we can crack this problem together. (I kept a record of my
thought processes as I solved this problem originally.) You'll see
elements of the problem-solving guidelines previously given, but you'll
also see them being applied in a seemingly haphazard order---just as they
would be in reality.

Here goes: The problem is to write a definition that consumes a
number on the stack and displays it as a Roman numeral.

This problem most likely represents a component of a larger system. We'll
probably end up defining several words in the course of solving this
problem, including data structures. But this particular lexicon will
include only one name, \forth{ROMAN}, and it will take its argument from
the stack. (Other words will be internal to the component.)

Having thus decided on the external syntax, we can now proceed to
devise the algorithms\index{A!Algorithms} and data structures.

We'll follow the scientific method---we'll observe reality, model a
solution, test it against reality, modify the solution, and so on. We'll
begin by recalling what we know about Roman numerals.

Actually, we don't remember any formal rules about Roman
numerals. But if you give us a number, we can make a Roman numeral
out of it. We know how to do it---but we can't yet state the procedure as
an algorithm.

So, let's look at the first ten Roman numerals:

%!! paragraph should be sans-serif + indented
\medskip
{\sf\begin{tabular}{r}
I \\
II \\
III \\
IV \\
V \\
VI \\
VII \\
VIII \\
IX \\
X \\
\end{tabular}}
\medskip

\noindent We make a few observations. First, there's the idea of a
tally, where we represent a number by making that many marks (3 =
III). On the other hand, special symbols are used to represent groups
(5 = V). In fact, it seems we can't have more than three I's in a row
before we use a larger symbol.

Second, there's a symmetry around five. There's a symbol for five (V), and
a symbol for ten (X). The pattern I, II, III repeats in the second half,
but with a preceding V.

One-less-than-five is written IV, and one-less-than-ten is written IX.
It seems that putting an ``I'' in front of a larger-value symbol is like
saying ``one-less-than\dots{}''

These are vague, hazy observations. But that's alright. We don't
have the whole picture yet.

Let's study what happens above ten:

%!! paragraph should be sans-serif + indented
\medskip
{\sf\begin{tabular}{r}
XI \\
XII \\
XIII \\
XIV \\
XV \\
XVI \\
XVII \\
XVIII \\
XIX \\
XX \\
\end{tabular}}

\medskip
\noindent This is exactly the pattern as before, with an extra ``X''
in front. So there's a repeating cycle of ten, as well.

If we look at the twenties, they're the same, with two ``X''s; the
thirties with three ``X''s. In fact, the number of ``X'' is the same as the
number in the tens column of the original decimal number.

This seems like an important observation: we can decompose our decimal
number into decimal digits, and treat each digit separately. For instance,
37 can be written as

%!! sans-serif + indented
\begin{quote}
{\sf XXX (thirty)}
\end{quote}

\noindent followed by

%!! sans-serif + indented
\begin{quote}
{\sf VII (seven)}
\end{quote}

\noindent It may be premature, but we can already see a method by
which \Forth{} will let us decompose a number into decimal digits---with
modulo division by ten. For instance, if we say

\begin{Code}
37 10 /MOD
\end{Code}
we'll get a 7 and a 3 on the stack (the three---being the quotient---is on
top.)

But these observations raise a question: What about below ten, where there
is no decimal digit? Is this a special case? Well, if we consider that
each ``X'' represents ten, then the absence of ``X'' represents zero.  So
it's \emph{not} a special case. Our algorithm works, even for numbers less
than ten.

Let's continue our observations, paying special attention to the
cycles of ten. We notice that forty is ``XL.'' This is analogous to 4 being
``IV,'' only shifted by the value of ten. The ``X'' before the ``L'' says
``ten-less-than-fifty.'' Similarly,

%!! paragraph should be sans-serif + indented
\bigskip
{\sf\begin{tabular}{rrcrr}
   L &  (50) &   is analogous to &     V & (5) \\
  LX &  (60) &          "        &    VI & (6) \\
 LXX &  (70) &          "        &   VII & (7) \\
LXXX &  (80) &          "        &  VIII & (8) \\
  XC &  (90) &          "        &    IX & (9) \\
   C & (100) &          "        &     X & (10) \\
\end{tabular}}
\bigskip

\noindent Apparently the same patterns apply for any decimal digit---only
the symbols themselves change. Anyway, it's clear now that we're dealing
with an essentially decimal system.

If pressed to do so, we could even build a model for a system to display
Roman numerals from 1 to 99, using a combination of algorithm and
data structure.

\subsection{Data Structure}

%!! see table on pg 123
\bigskip
\begin{tabular}{lcrclcr}
\multicolumn{3}{c}{\underline{\emph{Ones' Table}}} &
\hbox to 3em{\hss} & \multicolumn{3}{c}{\underline{\emph{Tens' Table}}} \\
& 0 &          & &    & 0 & \\
& 1 &       I  & &    & 1 &      X \\
& 2 &      II  & &    & 2 &     XX \\
& 3 &     III  & &    & 3 &    XXX \\
& 4 &      IV  & &    & 4 &     XL \\
& 5 &       V  & &    & 5 &      L \\
& 6 &      VI  & &    & 6 &     LX \\
& 7 &     VII  & &    & 7 &    LXX \\
& 8 &    VIII  & &    & 8 &   LXXX \\
& 9 &      IX  & &    & 9 &     XC \\
\end{tabular}
\bigskip

\subsection{Algorithm}

Divide $n$ by 10. The quotient is the tens' column digit; the remainder is
the ones' column digit. Look up the ten's digit in the tens' table and
print the corresponding symbol pattern. Look up the ones' digit in the
one's table and print that corresponding symbol pattern.

For example, if the number is 72, the quotient is 7, the remainder is 2. 7
in the tens' table corresponds to ``LXX,'' so print that. 2 in the ones'
column corresponds to ``II,'' so print that. The result:

%!! sans-serif + indented
\begin{quote}
{\sf LXXII}
\end{quote}
We've just constructed a model that works for numbers from one to 99.
Any higher number would require a hundreds' table as well, along with
an initial division by 100.

The logical model just described might be satisfactory, as long as it
does the job. But somehow it doesn't seem we've fully solved the problem.
We avoided figuring out how to produce the basic pattern by storing
all possible combinations in a series of tables. Earlier in this chapter we
observed that calculating an answer, if it's possible, can be easier than
using a data structure.

Since this section deals with devising algorithms\index{A!Algorithms},
let's go all the way.
Let's look for a general algorithm for producing any digit, using
only the elementary set of symbols. Our data structure should contain
only this much information:

%!! sans-serif + indented
\bigskip
{\sf\begin{tabular}{c@{\hspace{1em}}c}
I & V \\
X & L \\
C & D \\
M & \\
\end{tabular}}
\bigskip

\noindent In listing the symbols, we've also \emph{organized} them in a
way that seems right. The symbols in the left column are all multiples of
ten; the symbols in the right column are multiples of five. Furthermore,
the symbols in each row have ten times the value of the symbols directly
above them.

Another difference, the symbols in the first column can all be combined in
multiples, as ``XXXIII.'' But you can't have multiples of any of the
right-column symbols, such as VVV. Is this observation useful? Who knows?

Let's call the symbols in the left column \forth{ONERS} and in the right
column \forth{FIVERS}. The \forth{ONERS} represent the values 1, 10, 100,
and 1,000; that is, the value of one in every possible decimal place. The
\forth{FIVERS} represent 5, 50, and 500; that is, the value of five in
every possible decimal place.

Using these terms, instead of the symbols themselves, we should be able to
express the algorithm\index{A!Algorithms} for producing any digit. (We've
factored out the actual symbols from the \emph{kind} of symbols.) For
instance, we can state the following preliminary algorithm:

%!! indent paragraph
\begin{tfquot}
For any digit, print as many \forth{ONERS} as necessary to add up to the
value.
\end{tfquot}
Thus, for 300 we get ``CCC,'' for 20 we get ``XX'' for one we get ``I.''
And for 321 we get ``CCCXXI.''

This algorithm works until the digit is 4. Now we'll have to expand
our algorithm to cover this exception:

\begin{tfquot}
Print as many \forth{ONERS} as necessary to add up to the value, but if
the digit is 4, print a \forth{ONER} then a \forth{FIVER}.
%!! vertical half space
Hence, 40 is ``XL''; 4 is ``IV.''
\end{tfquot}
This new rule works until the digit is 5. As we noticed before, digits of
five and above begin with a \forth{FIVER} symbol. So we expand our rule
again:

\begin{tfquot}
If the digit is 5 or more, begin with a \forth{FIVER} and subtract five
from the value; otherwise do nothing. Then print as many \forth{ONERS} as
necessary to add up to the value. But if the digit is 4, print only a
\forth{ONER} and a \forth{FIVER}.
\end{tfquot}
This rule works until the digit is 9. In this case, we must print a
\forth{ONER} preceding a---what? A \forth{ONER} from the next higher
decimal place (the next row below). Let's call this a \forth{TENER}. Our
complete model, then is:

\begin{tfquot}
If the digit is 5 or more, begin with a \forth{FIVER} and subtract five
from the value; otherwise do nothing. Then, print as many \forth{ONERS} as
necessary to add up to the value. But if the digit is 4, print only a
\forth{ONER} and a \forth{FIVER}, or if it's 9, print only a \forth{ONER}
and a \forth{TENER}.
\end{tfquot}
\index{A!Algorithms|(}%
We now have an English-language version of our algorithm.  But we still
have some steps to go before we can run it on our computer.

In particular, we have to be more specific about the exceptions. We
can't just say,

\begin{tfquot}
Do a, b, and c. \emph{But} in such and such a case, do something different.
\end{tfquot}
because the computer will do a, b, and c before it knows any better.
\medbreak

Instead, we have to check whether the exceptions apply \emph{before} we
do anything else.

\begin{tip}
In devising an algorithm, consider exceptions last. In writing code, handle
exceptions first.
\end{tip}
This tells us something about the general structure of our digit-producing
word. It will have to begin with a test for the 4/9 exceptions. In
either of those cases, it will respond accordingly. If neither exception
applies, it will follow the ``normal'' algorithm. Using pseudocode, then:

\begin{Code}
: DIGIT  ( n )  4-OR-9? IF  special cases
   ELSE  normal case  THEN ;
\end{Code}
An experienced \Forth{} programmer would not actually write out this
pseudocode, but would more likely form a mental image of the structure
for eliminating the special cases. A less experienced programmer might
find it helpful to capture the structure in a diagram, or in code as we've
done here.
\index{A!Algorithms|)}

In \Forth{} we try to minimize our dependence on logic. But in this case
we need the conditional \forthb{IF} because we have an exception we need
to eliminate. Still, we've minimized the complexity of the control
structure by limiting the number of \forthb{IF }\forthb{THEN}s in this
definition to one.

Yes, we still have to distinguish between the 4-case and the 9-case,
but we've deferred that structural dimension to lower-level
definitions---the test for 4-or-9 and the ``special case'' code.

What our structure really says is that either the 4-exception or the
9-exception must prohibit execution of the normal case. It's not enough
merely to test for each exception, as in this version:

\begin{Code}
: DIGIT  ( n )  4-CASE? IF  ONER FIVER  THEN
                9-CASE? IF  ONER TENER  THEN
                normal case... ;
\end{Code}
because the normal case is never excluded. (There's no way to put an
\forthb{ELSE} just before the normal case, because \forthb{ELSE} must
appear between \forthb{IF} and \forthb{THEN}.)

If we insist on handling the 4-exception and the 9-exception separately,
we could arrange for each exception to pass an additional flag, indicating
that the exception occurred. If either of these flags is true, then we can
exclude the normal case:

%!! warning: \textbf{} in verbatim block
\begin{Code}[commandchars=\&\{\}]
: DIGIT  ( n )  4-CASE? &poorbf{DUP} IF  ONER FIVER  THEN
                9-CASE? &poorbf{DUP} IF  ONER TENER  THEN
                OR  NOT IF normal case THEN ;
\end{Code}
But this approach needlessly complicates the definition by adding new
control structures. We'll leave it like it was.

Now we have a general idea of the structure of our main definition.

We stated, ``If the digit is 5 or more, begin with a \forth{FIVER} and
subtract five from the value; otherwise do nothing. Then, print as many
\forth{ONERS} as necessary to add up to the value.''

A direct translation of these rules into \Forth{} would look like this:

\begin{Code}
( n)  DUP  4 > IF  FIVER 5 -  THEN  ONERS
\end{Code}
This is technically correct, but if we're familiar with the technique of
modulo division, we'll see this as a natural situation for modulo division
by 5. If we divide the number by five, the quotient will be zero (false)
when the number is less than five, and one (true) when it's between 5 and
9. We can use it as the boolean flag to tell whether we want the leading
\forth{FIVER}:

\begin{Code}
( n )  5 / IF FIVER THEN ...
\end{Code}
The quotient/flag becomes the argument to \forth{IF}.

Furthermore, the remainder of modulo 5 division is always a number between
0 and 4, which means that (except for our exception) we can use the
remainder directly as the argument to \forth{ONERS}. We revise our phrase
to

\begin{Code}[commandchars=\&\{\}]
( n )  5 &poorbf{/MOD} IF FIVER THEN  &poorbf{ONERS}
\end{Code}
Getting back to that exception, we now see that we can test for both 4 and
9 with a single test---namely, if the remainder is 4. This suggests that
we can do our 5 \forthb{/MOD} first, then test for the exception.
Something like this:

%!! warning: \textbf{} in verbatim block
\begin{Code}[commandchars=\&\{\}]
: DIGIT  ( n )
     5 /MOD  &poorbf{OVER 4 =  IF  special case  ELSE}
     IF FIVER THEN  ONERS  THEN ;
\end{Code}
(Notice that we \forthb{OVER}ed the remainder so that we could compare it with
4 without consuming it.)

So it turns out we \emph{do} have a doubly-nested \forthb{IF THEN}
construct after all. But it seems justified because the \forthb{IF THEN}
is handling the special case. The other is such a short phrase,
``\forth{IF FIVER THEN},'' it's hardly worth making into a separate
definition. You could though. (But we won't.)

Let's focus on the code for the special case. To state its algorithm: ``If
the digit is four, print a \forth{ONER} and a \forth{FIVER}. If the digit
is nine, print a \forth{ONER} and a \forth{TENER}.''

We can assume that the digit will be one or the other, or else we'd never
be executing this definition. The question is, how do we tell which one?

Again, we can use the quotient of division by five. If the quotient is
zero, the digit must have been four; otherwise it was nine. So we'll play
the same trick and use the quotient as a boolean flag. We'll write:

\begin{Code}
: ALMOST  ( quotient )
     IF  ONER TENER  ELSE  ONER FIVER  THEN ;
\end{Code}
In retrospect, we notice that we're printing a \forth{ONER} either way. We
can simplify the definition to:

\begin{Code}
: ALMOST  ( quotient )
     ONER  IF TENER ELSE FIVER THEN ;
\end{Code}
We've assumed that we have a quotient on the stack to use. Let's go back
to our definition of \forth{DIGIT} and make sure that we do, in fact:

%!! warning: \textbf{} in verbatim block
\begin{Code}[commandchars=\&\{\}]
: DIGIT  ( n )
     5 /MOD  OVER 4 =  IF  &poorbf{ALMOST}  ELSE
     IF FIVER THEN  ONERS  THEN ;
\end{Code}
It turns out that we have not only a quotient, but a remainder underneath
as well. We're keeping both on the stack in the event we execute the
\forthb{ELSE} clause. The word \forthb{ALMOST}, however, only needs the
quotient. So, for symmetry, we must \forthb{DROP} the remainder like this:

%!! warning: \textbf{} in verbatim block
\begin{Code}[commandchars=\&\{\}]
: DIGIT  ( n )
     5 /MOD  OVER 4 =  IF  ALMOST  &poorbf{DROP}  ELSE
     IF FIVER THEN  ONERS  THEN ;
\end{Code}
There we have the complete, coded definition for producing a single digit
of a Roman numeral. If we were desperate to try it out before writing the
needed auxiliary definitions, we could very quickly define a lexicon of
words to print one group of symbols, say the \forth{ONES} row:

%!! there is a small mistake in the original book:
%!! : TENER   ." X"; <-- missing space
%!! Fixed that.
\begin{Code}
: ONER    ." I" ;
: FIVER   ." V" ;
: TENER   ." X" ;
: ONERS  ( # of oners -- )
     ?DUP IF 0 DO  ONER  LOOP  THEN ;
\end{Code}
before loading our definitions of \forth{ALMOST} and \forth{DIGIT}.

But we're not that desperate. No, we're anxious to move on to the
problem of defining the words \forth{ONER}, \forth{FIVER}, and \forth{TENER} so that their
symbols depend on which decimal digit we're formatting.

Let's go back to the symbol table we drew earlier:

%!! sans-serif + indented, heading in italics
\bigskip{\sf
\begin{tabular}{rcc}
          & \it \forth{ONER}s & \it \forth{FIVER}s \\
     ones & I  &  V \\
     tens & X  &  L \\
 hundreds & C  &  D \\
thousands & M  & \\
\end{tabular}}
\bigskip

\noindent
We've observed that we also need a ``\forth{TENER}''---which is the
\forth{ONER} in the next row below. It's as if the table should really
be written:

%!! Here comes above table again, but this time with
%!! some circles and lines in it. I have no clue how
%!! to render this. See pg 128.
\bigskip
{\sf\begin{tabular}{rccc}
          & \it \forth{ONER}s & \it \forth{FIVER}s & \it \forth{TENER}s \\
     ones & \circlenode[linecolor=white]{I2}{\boxto{}I}  &  V & \circlenode{X1}{\boxto{}X} \\
     tens & \circlenode{X2}{\boxto{}X}  &  L & \circlenode{C1}{\boxto{}C} \\
 hundreds & \circlenode{C2}{\boxto{}C}  &  D & \circlenode{M1}{\boxto{}M} \\
thousands & \circlenode{M2}{\boxto{}M}  & \\
\end{tabular}
\ncline{X1}{X2}
\ncline{C1}{C2}
\ncline{M1}{M2}}
\bigskip

\noindent But that seems redundant. Can we avoid it? Perhaps if we try
a different model, perhaps a linear table, like this:

\bigskip
{\sf\begin{tabular}{rc}
ones      & I \\
          & V \\
tens      & X \\
          & L \\
hundreds  & C \\
          & D \\
thousands & M \\
\end{tabular}}
\bigskip

\noindent Now we can imagine that each column name (``ones,'' ``tens,''
etc.) points to the \forth{ONER} of that column. From there we can also
get each column's \forth{FIVER} by reaching down one slot below the
current \forth{ONER}, and the \forth{TENER} by reaching down two slots.

It's like building an arm with three hands. We can attach it to the
\forth{ONES} column, as in \Fig{fig4-8}a, or we can attach it to the tens'
column, as in \Fig{fig4-8}b, or to any power of ten.

\wepsfigt{fig4-8}{A mechanical representation: accessing the data
structure.}

%!! include Figure 4-8 here

An experienced \Forth{} programmer is not likely to imagine arms,
hands, or things like that. But there must be a strong mental image---the
stuff of right-brain thinking---before there's any attempt to construct the
model with code.

Beginners who are learning to think in this right-brain way might
find the following tip helpful:

\begin{tip}
If you have trouble thinking about a conceptual model,
\index{C!Conceptual model}
visualize it---or draw it---as a mechanical device.
\end{tip}
Our table is simply an array of characters. Since a character requires
only a byte, let's make each ``slot'' one byte. We'll call the table
\forth{ROMANS}:

\begin{Code}
CREATE ROMANS    ( ones)  ASCII I  C,   ASCII V  C,
                 ( tens)  ASCII X  C,   ASCII L  C,
             ( hundreds)  ASCII C  C,   ASCII D  C,
            ( thousands)  ASCII M  C,
\end{Code}
Note: This use of \forthb{ASCII} requires that \forthb{ASCII} be
``\forthb{STATE}-dependent'' (see \App{C}). If you don't have \forthb{ASCII},
or if it is not state-dependent, use:

\begin{Code}
CREATE ROMANS  73 C,  86 C,  88 C,  76 C,
   67 C,  68 C,  77 C,
\end{Code}
We can select a particular symbol from the table by applying two
different offsets at the same time. One dimension represents the decimal
place: ones, tens, hundreds, etc. This dimension is made ``current,'' that
is, its state stays the same until we change it.

The other dimension represents the kind of symbol we want---\forth{ONER},
\forth{FIVER}, \forth{TENER}---within the current decimal column. This
dimension is incidental, that is, we'll specify which symbol we want each
time.

Let's start by implementing the ``current'' dimension. We need
some way to point to the current decimal column. Let's create a variable
called \forth{COLUMN\#} (pronounced ``column-number'') and have it contain an
offset into the table:

\begin{Code}
VARIABLE COLUMN#  ( current offset)
: ONES        O COLUMN# ! ;
: TENS        2 COLUMN# ! ;
: HUNDREDS    4 COLUMN# ! ;
: THOUSANDS   6 COLUMN# ! ;
\end{Code}
Now we can find our way to any ``arm position'' by adding the contents of
\forth{COLUMN\#} to the beginning address of the table, given by
\forth{ROMANS}:

\begin{Code}
: COLUMN  ( -- adr of column)  ROMANS  COLUMN# @  + ;
\end{Code}
Let's see if we can implement one of the words to display a symbol. We'll
start with \forth{ONER}.

The thing we want to do in \forth{ONER} is \forthb{EMIT} a character.

\begin{Code}
: ONER                   EMIT ;
\end{Code}
Working backward, \forthb{EMIT} requires the ASCII character on the stack.
How do we get it there? With \forthb{C@}.

\begin{Code}
: ONER                C@ EMIT ;
\end{Code}
\forthb{C@} requires the \emph{address} of the slot that contains the
symbol we want. How do we get that address?

The \forth{ONER} is the first ``hand'' on the movable arm---the position
that \forth{COLUMN} is already pointing to. So, the address we want is
simply the address returned by \forth{COLUMN}:

\begin{Code}
: ONER   COLUMN       C@ EMIT ;
\end{Code}
Now let's write \forth{FIVER}. It computes the same slot address, then
adds one to get the next slot, before fetching the symbol and emitting it:

\begin{Code}
: FIVER  COLUMN 1+    C@ EMIT ;
\end{Code}
And \forth{TENER} is:

\begin{Code}
: TENER  COLUMN 2+    C@ EMIT ;
\end{Code}
These three definitions are redundant. Since the only difference between
them is the incidental offset, we can factor the incidental offset out
from the rest of the definitions:

\begin{Code}
: .SYMBOL  ( offset)  COLUMN +  C@ EMIT ;
\end{Code}
Now we can define:

\begin{Code}
: ONER    O .SYMBOL ;
: FIVER   1 .SYMBOL ;
: TENER   2 .SYMBOL ;
\end{Code}
All that remains for us to do now is to decompose our complete decimal
number into a series of decimal digits. Based on the observations we've
already made, this should be easy. \Fig{fig4-9} shows our completed
listing.

Voila! From problem, to conceptual model, to code.

Note: this solution is not optimal. The present volume does not address
the optimization phase.

One more thought: Depending on who uses this application, we may
want to add error-checking. Fact is, the highest symbol we know is M; the
highest value we can represent is 3,999, or MMMCMXCIX.

\goodbreak
We might redefine ROMAN as follows:

%!! I have added an IF...THEN around ABORT" Too large".
%!! Without it, it does not seem to make sense, does it?
%!! (Pg 131)
\begin{Code}
: ROMAN  ( n)
   DUP  3999 >  ABORT" Too large"  ROMAN ;
\end{Code}

\begin{interview}
\person{Moore}:\index{M!Moore, Charles|(}

\begin{tfquot}
There's a definite sense of rightness when you've done it right. It may be
that feeling that distinguishes \Forth{} from other languages, where you
never feel you've really done well. In \Forth{}, it's the ``Aha!''
reaction. You want to run off and tell somebody.

Of course, nobody will appreciate it like you do.
\end{tfquot}\index{M!Moore, Charles|)}
\end{interview}

\begin{figure*}[tttt]
\caption{Roman numerals, solved.}
\labelfig{fig4-9}
\vspace{1ex}
\setcounter{screen}{20}
\begin{Screen}
\ Roman numerals                                         8/18/83
CREATE ROMANS    ( ones)  ASCII I  C,   ASCII V  C,
                 ( tens)  ASCII X  C,   ASCII L  C,
             ( hundreds)  ASCII C  C,   ASCII D  C,
            ( thousands)  ASCII M  C,
VARIABLE COLUMN#  ( current offset)
: ONES       O COLUMN# ! ;
: TENS       2 COLUMN# ! ;
: HUNDREDS   4 COLUMN# ! ;
: THOUSANDS  6 COLUMN# ! ;

: COLUMN ( -- address of column)  ROMANS  COLUMN# @  + ;

\end{Screen}

\begin{Screen}
\ Roman numerals cont'd                                  8/18/83
: .SYMBOL  ( offset -- )  COLUMN +  C@ EMIT ;
: ONER    O .SYMBOL ;
: FIVER   1 .SYMBOL ;
: TENER   2 .SYMBOL ;

: ONERS  ( # of oners -- )
   ?DUP  IF  O DO  ONER  LOOP  THEN ;
: ALMOST  ( quotient of 5 / -- )
   ONER  IF  TENER  ELSE  FIVER  THEN ;
: DIGIT  ( digit -- )
   5 /MOD  OVER  4 = IF  ALMOST  DROP  ELSE  IF FIVER THEN
   ONERS THEN ;

\end{Screen}

\begin{Screen}
\ Roman numerals cont'd                                  8/18/83
: ROMAN  ( number --)  1000 /MOD  THOUSANDS DIGIT
                        100 /MOD   HUNDREDS DIGIT
                         10 /MOD       TENS DIGIT
                                       ONES DIGIT  ;

\end{Screen}
\vspace{1ex}
\label{fig-fig4-9}
\end{figure*}%
\index{R!Roman numeral computation problem|)}%
\program{roman2}

\section{Summary}

In this chapter we've learned to develop a single component, starting
first with deciding on its syntax, then proceeding with determining its
algorithm(s)\index{A!Algorithms}
and data structure(s), and concluding with an implementation in \Forth{}.

With this chapter we complete our discussion of design. The remainder
of the book will discuss style and programming techniques.
\index{D!Detailed design|)}%
\index{D!Detailed design!demonstration of|)}

%!! the ordered list thing again
\begin{references}{9}
\bibitem{polya} \person{G. Polya}, \emph{How To Solve It: A New Aspect of
   Mathematical Method}, (Princeton, New Jersey, Princeton University
   Press).
%!! (C) should be a nicer copyright sign
\bibitem{hart75} \person{Leslie A. Hart}, \emph{How the Brain Works},
   \copyright{} 1975 by \person{Leslie A. Hart}, (New York, Basic Books, Inc.,
   1975).
\bibitem{rosen82} \person{Evan Rosen}, ``High Speed, Low Memory Consumption
   Structures,'' 1982 \emph{FORML Conference Proceedings}, p. 191.
\bibitem{stolowitz82} \person{Michael Stolowitz}, ``A Compiler for Programmable
   Logic in FORTH,'' 1982 \emph{FORML Conference Proceedings}, p. 257.
\end{references}

\section{For Further Thinking}

Design the components and describe the algorithm(s) necessary to simulate
shuffling a deck of cards. Your algorithm will produce an array of
numbers, 0--51, arranged in random order.

The special constraint of this problem, of course, is that no one card
may appear twice in the array.

You may assume you have a random-number generator called
CHOOSE. It's stack argument is ``$n$''; it produces a random number
between zero and $n-1$ inclusive. (See the Handy Hint, Chapter Ten,
\emph{Starting \Forth{}}.)

Can you design the card-shuffling algorithm so that it avoids the
time-consuming burden of checking some undetermined number of slots
on each pass of the loop? Can you do so using only the one array?

%% Thinking Forth
%% Copyright (C) 2004 Leo Brodie
%% Initial transcription by Albert van der Horst
%%
%% Chapter: Implementation: Elements of FORTH Style
%%
%% ``%%!!!te'' marks mark typesetting errors in the original
%% We don't want to retain typesetting errors.

\chapter{Implementation: Elements~of \Forth{}~Style}\Chapmark{5}%
\index{I!Implementation|(}%

\initial Badly written \Forth{} has been accused of looking like ``code that
went through a trash compactor.'' It's true, \Forth{} affords more
freedom in the way we write applications.  But that freedom also gives
us a chance to write exquisitely readable and easily maintainable
code, provided we consciously employ the elements of good \Forth{}
style.

In this chapter we'll delve into \Forth{} coding convention
including:

\begin{itemize}%%originally: no bullet points
\item listing organization
\item screen layout, spacing and indentation
\item commenting
\item choosing names
\end{itemize}
I wish I could recommend a list of hard-and-fast conventions for
everyone.  Unfortunately, such a list may be inappropriate in many
situations.  This chapter merges many widely-adopted conventions with
personal preferences, commented with alternate ideas and the reasons
for the preferences.  In other words:
\begin{Code}
: TIP  VALUE JUDGEMENT ;
\end{Code}
I'd especially like to thank \person{Kim Harris},\index{H!Harris, Kim} who
proposed many of the conventions described in this chapter, for his
continuing efforts at unifying divergent views on good \Forth{} style.

\section{Listing Organization}%
\index{L!Listing organization|(}%
\index{I!Implementation!listing organization|(}

A well-organized book has clearly defined chapters, with clearly defined
sections, and a table of contents to help you see the organization at a
glance.  A well-organized book is easy to read.  A badly organized book
makes comprehension more difficult, and makes finding information
later on nearly impossible.

\wepsfigp{fig5-1}{I still don't see how these programming conventions
enhance readability.}

The necessity for good organization applies to an application listing
as well.  Good organization has three aspects:

\begin{enumerate}\parsep=0pt\itemsep=0pt
\item Decomposition
\item Composition
\item Disk partitioning
\end{enumerate}

\subsection{Decomposition}%
\index{D!Decomposition|(}%
\index{L!Lexicons|(}
As we've already seen, the organization of a listing should follow the
decomposition of the application into lexicons.  Generally these
lexicons should be sequenced in ``uses'' order.  Lexicons being
\emph{used} should precede the lexicons which \emph{use} them.%
\index{L!Lexicons|)}

On a larger scale, elements in a listing should be organized by degree
of complexity, with the most complex variations appearing towards the
end.  It's best to arrange things so that you can leave off the
lattermost screens (i.e., not load them) and still have a
self-sufficient, running application, working properly except for the
lack of the more advanced features.

We discussed the art of decomposition extensively in \Chap{3}.%
\index{D!Decomposition|)}

\subsection{Composition}%
\index{C!Composition|(}%
\index{L!Load screens|(}%

Composition is the putting together of pieces to create a whole.  Good
composition requires as much artistry as good decomposition.

One of \Forth{}'s present conventions is that source code resides in
``screens,'' which are 1K units of mass storage.  (The term ``screen''
refers to a block used specifically for source code.) It's possible in
\Forth{} to chain every screen of code to the next, linking the entire
listing together linearly like a lengthy parchment scroll.  This is
not a useful approach.  Instead:

\begin{tip}
Structure your application listing like a book: hierarchically.
\end{tip}
An application may consist of:

\begin{description}\parsep=0pt\itemsep=0pt
\item[Screens:] the smallest unit of \Forth{} source
\item[Lexicons:]\index{L!Lexicons} one to three screens, enough to
implement a component
\item[Chapters:]\index{C!Chapters} a series of related lexicons, and
\item[Load screens:] analogous to a table of contents, a screen that
loads the chapters in the proper sequence.
\end{description}

%Page 138 in first edition.

\begin{figure*}[tttt]
\caption{Example of an application-load screen.}
\labelfig{fig5-1}
\setcounter{screen}{1}
\begin{Screen}
\ QTF+ Load Screen                                      07/09/83
: RELEASE#   ." 2.01" ;
  9 LOAD  \ compiler tools, language primitives
 12 LOAD  \ video primitives
 21 LOAD  \ editor
 39 LOAD  \ line display
 48 LOAD  \ formatter
 69 LOAD  \ boxes
 81 LOAD  \ deferring
 90 LOAD  \ framing
 96 LOAD  \ labels, figures, tables
102 LOAD  \ table of contents generator




\end{Screen}
\end{figure*}%
\index{C!Composition|)}

\subsection{Application-load Screen}%
\index{L!Load screens!application|(}

\Fig{fig5-1} is an example of an application-load screen.  Since it
resides in Screen 1, you can load this entire application by entering
\begin{Code}
1 LOAD
\end{Code}
The individual load commands within this screen load the chapters of
the application.  For instance, Screen 12 is the load screen for the
video primitives chapter.

As a reference tool, the application-load screen tells you where to
find all of the chapters.  For instance, if you want to look at the
routines that do framing, you can see that the section starts at
Screen 90.

Each chapter-load screen in turn, loads all of the screens comprising
the chapter.  We'll study some formats for chapter-load screens shortly.

The primary benefit of this hierarchical scheme is that you can load
any section, or any screen by itself, without having to load the
entire application.  Modularity of the source code is one of the
reasons for \Forth{}'s quick turnaround time for editing, loading, and
testing (necessary for the iterative approach).  Like pages of a book,
each screen can be accessed individually and quickly.  It's a ``random
access'' approach to source-code maintenance.

You can also replace any passage of code with a new, trial version by
simply changing the screen numbers in the load screen.  You don't have to
move large passages of source code around within a file.

In small applications, there may not be such things as chapters.  The
application-load screen will directly load all the lexicons.  In
larger applications, however, the extra level of hierarchy can improve
maintainability.
%Page 139 in first edition.

A screen should either be a load-screen or a code-screen, not a mixture.
Avoid embedding a \forthb{LOAD} or \forthb{THRU} command in the middle of a
screen containing definitions just because you ``need something'' or
because you ``ran out of room.''

\subsection{Skip Commands}%
\index{S!Skip commands}
Two commands make it easy to control what gets loaded in each screen
and what gets ignored.  They are:

\medbreak
\begin{description}
\item[\emph{\forthb{\bs}}]~
\item[\emph{\forthb{\bs S}}] also called \forthb{EXIT}\index{E!EXIT}
\end{description}
\forth{\bs} is pronounced ``skip-line.'' It causes the
\Forth{} interpreter to ignore everything to the right of it on the
same line. (Since \forth{\bs} is a \Forth{} word, it must be followed
by a space.) It does not require a delimiter.

In \Fig{fig5-1}, you see \forth{\bs} used in two ways: to begin the
screen-comment line (Line 0), and to begin comments on individual
lines which have no more code to the right of the comment.

During testing, \forth{\bs} also serves to temporarily ``paren out''
lines that already contain a right parenthesis in a name or comment.
For instance, these two ``skip-line''s keep the definition of
\forth{NUTATE} from being compiled without causing problems in
encountering either right parenthesis:
\begin{Code}
\ : NUTATE  ( x y z )
\   SWAP ROT  (NUTATE) ;
\end{Code}
\forthb{\bs S} is pronounced ``skip-screen.'' It causes the \Forth{}
interpreter to stop interpreting the screen entirely, as though there
were nothing else in the screen beyond \forthb{\bs S}.

In many \Forth{} systems, this function is the same as
\forthb{EXIT},\index{E!EXIT} which is the run-time routine for
semicolon.  In these systems the use of \forthb{EXIT} is acceptable.
Some \Forth{} systems, however, require for internal reasons a
different routine for the ``skip-screen'' function.

Definitions for \forthb{\bs} and \forthb{\bs S} can be found in \App{C}.%
\index{L!Load screens!application|)}

\subsection{Chapter-load Screens}%
\index{L!Load screens!chapter|(}%
\index{C!Chapter-load screens|(}

\Fig{fig5-2} illustrates a typical chapter-load screen.  The screens
loaded by this screen are referred to relatively, not absolutely as
they were in the application-load screen.

This is because the chapter-load screen is the first screen of the
contiguous range of screens in the chapter.  You can move an entire
chapter forward or backward within the listing; the relative pointers
in the chapter-load screen are position-independent.  All you have to
change is the single number in the application-load screen that points
to the beginning of the chapter.
%Page 140 in first edition.

\begin{figure*}
\caption{Example of a chapter-load screen.}
\labelfig{fig5-2}

\setcounter{screen}{100}
\begin{Screen}
\ GRAPHICS                 Chapter load                 07/11/83

 1 FH LOAD            \ dot-drawing primitive
 2 FH 3 FH THRU       \ line-drawing primitives
 4 FH 7 FH THRU       \ scaling, rotation
 8 FH LOAD            \ box
 9 FH 11 FH THRU      \ circle



CORNER  \ initialize relative position to low-left corner





\end{Screen}
\end{figure*}

\index{R!Relative loading|(}%
\begin{tip}
Use absolute screen numbers in the application-load screen.  Use
relative screen numbers in the chapter- or section-load screens.
\end{tip}
There are two ways to implement relative loading.
The most common is to define:
\begin{Code}
: +LOAD  ( offset -- )  BLK @ +  LOAD ;
\end{Code}
and
\begin{Code}
: +THRU  ( lo-offset hi-offset -- )
     1+ SWAP DO  I +LOAD  LOOP ;
\end{Code}
My own way, which I submit as a more useful factoring, requires a
single word, \forthb{FH} (see \App{C} for its definition).

The phrase
\begin{Code}
1 FH LOAD
\end{Code}
is read ``1 from here \forth{LOAD},'' and is equivalent
to \forth{1 +LOAD}.

Similarly,
\begin{Code}
2 FH   5 FH THRU
\end{Code}
is read ``2 from here, 5 from here \forth{THRU}.''

Some programmers begin each chapter with a dummy word; e.g.,
\begin{Code}
: VIDEO-IO ;
\end{Code}
%Page 141 in first edition.
and list its name in the comment on the line where the chapter is
loaded in the application-load screen.  This permits selectively
\forth{FORGET}ting any chapter and reloading from that point on
without having to look at the chapter itself.

Within a chapter the first group of screens will usually define those
variables, constants, and other data structures needed globally within
the chapter.  Following that will come the lexicons, loaded in
``uses'' order.  The final lines of the chapter-load screen normally
invoke any needed initialization commands.

\index{M!Moore Products Company|(}%
Some of the more style-conscious \Forth{}wrights begin each chapter
with a ``preamble'' that discusses in general terms the theory of
operation for the components described in the chapter.  \Fig{fig5-3}
is a sample preamble screen which demonstrates the format required at
Moore Products Co.

\begin{figure*}[tttt]
\caption{Moore Products Co.'s format for chapter preambles.}
\labelfig{fig5-3}
\begin{Screen}
CHAPTER 5  -  ORIGIN/DESTINATION - MULTILOOP BIT ROUTINES

DOCUMENTS - CONSOLE STRUCTURE CONFIGURATION
        DESIGN SPECIFICATION
        SECTIONS - 3.2.7.5.4.1.2.8
                   3.2.7.5.4.1.2.10

ABSTRACT  -  File control types E M T Q and R can all
             originate from a Regional Satellite or a
             Data Survey Satellite.  These routines allow
             the operator to determine whether the control
             originated from a Regional Satellite or not.




\end{Screen}

\begin{Screen}
CHAPTER NOTES - Whether or not a point originates from
                a Regional Satellite is determined by
                the Regional bit in BITS, as follows:

                  1 = Regional Satellite
                  2 = Data Survey Satellite

                 For the location of the Regional bit
                 in BITS, see the Design Specification
                 Section - 3.2.7.5.4.1.2.10

HISTORY  -




\end{Screen}
\end{figure*}%
\index{M!Moore Products Company|)}

%Page 142 in first edition.
\begin{interview}
\person{Charles Moore}\index{M!Moore, Charles|(} (no relation to Moore
Products Co.) places less importance on the well-organized
hierarchical listing than I do. \person{Moore}:

\begin{tfquot}
I structure \emph{applications} hierarchically, but not necessarily
\emph{listings.}  My listings are organized in a fairly sloppy way,
not at all hierarchically in the sense of primitives first.

I use \forthb{LOCATE} {[}also known as \forthb{VIEW}; see the Handy Hint
in \emph{Starting \Forth{},} Chapter Nine{]}.  As a result, the
listing is much less carefully organized because I have \forthb{LOCATE}
to find things for me.  I never look at listings.
\end{tfquot}\index{M!Moore, Charles|)}
\end{interview}


\subsection{--\,--> vs.\ THRU}%
\index{T!THRU|(}

\index{N!Next block|(}%
On the subject of relative loading, one popular way to load a series
of adjacent screens is with the word \forth{-{}->} (pronounced ``next
block'').  This word causes the interpreter to immediately cease
interpreting the current screen and begin interpreting the next
(higher-numbered) screen.

If your system provides \forth{-{}->}, you must choose between using
the \forthb{THRU} command in your chapter-load screen to load each
series of screens, or linking each series together with the arrows and
\forth{LOAD}ing only the first in the series.  (You can't do both;
you'd end up loading most of the screens more than once.)%
\index{T!THRU|)}

The nice thing about the arrows is this: suppose you change a screen
in the middle of a series, then reload the screen.  The rest of the
series will automatically get loaded.  You don't have to know what the
last screen is.

That's also the nasty thing about the arrows: There's no way to stop
the loading process once it starts.  You may compile a lot more
screens than you need to test this one screen.%
\index{N!Next block|)}

To get analytical about it, there are three things you might want to
do after making the change just described:

\begin{enumerate}
\item load the one screen only, to test the change,
\item load the entire section in which the screen appears,
\suspend{enumerate}
or
\resume{enumerate}
\item load the entire remainder of the application.
\end{enumerate}
The use of \forthb{THRU} seems to give you the greatest control.

Some people consider the arrow to be useful for letting definitions
cross screen boundaries.  In fact \forth{-{}->} is the only way to
compile a high-level (colon) definition that occupies more than one
screen, because \forth{-{}->} is ``immediate.'' But it's \emph{never}
good style to let a colon definition cross screen boundaries.  (They
should never be that long!)

On the other hand, an extremely complicated and time-critical piece
of assembler coding might occupy several sequential screens.  In this
case, though, normal \forthb{LOAD}ing will do just as well, since the
assembler does not use compilation mode, and therefore does not
require immediacy.

Finally, the arrow wastes an extra line of each source screen.  We
don't recommend it.%
\index{R!Relative loading|)}%
\index{L!Load screens!chapter|)}%
\index{C!Chapter-load screens|)}%
\index{L!Load screens|)}%
\index{L!Listing organization|)}

%Page 143 in first edition.
\section{An Alternative to Screens: Source in Named Files}%
\index{N!Named files, storing source code in|(}

\index{F!File-based system|(}%
Some \Forth{} practitioners advocate storing source code in
variable-length, named text files, deliberately emulating the approach
used by traditional compilers and editors.  This approach may become
more and more common, but its usefulness is still controversial.

Sure, it's nice not to have to worry about running out of room in a
screen, but the hassle of writing in a restricted area is compensated
for by retaining control of discrete chunks of code.  In developing an
application, you spend a lot more time loading and reloading screens
than you do rearranging their contents.

``Infinite-length'' files allow sloppy, disorganized thinking and bad
factoring.  Definitions become longer without the discipline imposed
by the 1K block boundaries.  The tendency becomes to write a 20K file,
or worse: a 20K definition.

Perhaps a nice compromise would be a file-based system that allows
nested loading, and encourages the use of very small named files.
Most likely, though, the more experienced \Forth{} programmers would
not use named files longer than 5K to 10K.  So what's the benefit?%
\index{F!File-based system|)}

Some might answer that rhetorical question: ``It's easier to remember
names than numbers.'' If that's so, then predefine those block numbers
as constants, e.g.:
\begin{Code}
90 CONSTANT FRAMING
\end{Code}
Then to load the ``framing'' section, enter
\begin{Code}
FRAMING LOAD
\end{Code}
Or, to list the section's load block, enter
\begin{Code}
FRAMING LIST
\end{Code}
(It's a convention that names of sections end in ``ING.'')

Of course, to minimize the hassle of the screen-based approach you
need good tools, including editor commands that move lines of source
from one screen to another, and words that slide a series of screens
forward or back within the listing.%
\index{N!Named files, storing source code in|)}

\subsection{Disk Partitioning}%
\index{D!Disk partitioning|(}
The final aspect of the well-organized listing involves standardizing
an arrangement for what goes where on the disk.  These standards must
be set by each shop, or department, or individual programmer,
depending on the nature of the work.
%Page 144 in first edition.

\begin{figure*}
\caption{Example of a disk-partitioning scheme within one department.}
\labelfig{fig5-4}

\begin{description}
\item[Screen 0] is the title screen, showing the name of the
    application, the current release number, and primary author.
\item[Screen 1] is the application-load block.
\item[Screen 2] is reserved for possible continuation from Screen 1
\item[Screen 4 and 5] contain system messages.
\item[Screens 9 thru 29] incorporate general utilities needed
    in, but not restricted to, this application.
\item[Screen 30] begins the application screens.
\end{description}
\end{figure*}

\Fig{fig5-4} shows a typical department's partitioning scheme.

In many \Forth{} shops it's considered desirable to begin sections of
code on screen numbers that are evenly divisible by three.  Major
divisions on a disk should be made on boundaries evenly divisible by
thirty.%
\index{D!Disk partitioning|)}

The reason? By convention, \Forth{} screens are printed three to a
page, with the top screen always evenly divisible by three.  Such a
page is called a ``triad;'' most \Forth{} systems include the word
\forth{TRIAD} to produce it, given as an argument the number of any of
the three screens in the triad.  For instance, if you type
\begin{Code}
77 TRIAD
\end{Code}
you'll get a page that includes 75, 76, and 77.

The main benefit of this convention is that if you change a single
screen, you can slip the new triad right into your binder containing
the current listing, replacing exactly one page with no overlapping
screens.

Similarly, the word \forth{INDEX} lists the first line of each screen,
60 per page, on boundaries evenly divisible by 60.\index{I!INDEX}

\index{S!Screens!numbering|(}%
\begin{tip}
Begin sections or lexicons on screen numbers evenly divisible by three.
Begin applications or chapters on screen numbers evenly divisible by
thirty.
\end{tip}%
\index{S!Screens!numbering|)}

\subsection{Electives}%
\index{E!Electives|(}
Vendors of \Forth{} systems have a problem.  If they want to include
every command that the customer might expect---words to control
graphics, printers, and other niceties---they often find that the system
has swollen to more than half the memory capacity of the computer,
leaving less room for serious programmers to compile their
applications.
%Page 145 in first edition.

The solution is for the vendor to provide the bare bones as a
precompiled nucleus, with the extra goodies provided in \emph{source} form.
This approach allows the programmer to pick and choose the special
routines actually needed.

These user-loadable routines are called ``electives.'' Double-length
arithmetic, date and time support, \forth{CASE} statements and the
\forth{DOER}/\hy \forth{MAKE} construct (described later) are some of
the features that \Forth{} systems should offer as electives.%
\index{E!Electives|)}%
\index{I!Implementation!listing organization|)}

\section{Screen Layout}%
\index{S!Screens!layout|(}%
\index{I!Implementation!screen layout|(}

In this section we'll discuss the layout of each source screen.

\index{C!Comment line|(}%
\begin{tip}
Reserve Line 0 as a ``comment line.''
\end{tip}
The comment line serves both as a heading for the screen, and also as
a line in the disk \forth{INDEX}.  It should describe the purpose of
the screen (not list the words defined therein).

The comment line minimally contains the name of the screen.  In larger
applications, you may also include both the chapter name and screen
name.  If the screen is one of a series of screens implementing a
lexicon, you should include a ``page number'' as well.

\index{S!Stamp|(}%
\index{D!Dates, representation of|(}
The upper right hand corner is reserved for the ``stamp.'' The stamp
includes the date of latest revision and, when authorship is
important, the programmer's initials (three characters to the left of
the date); e.g.:
\begin{Code}
( Chapter name        Screen Name -- pg #      JPJ 06/10/83)
\end{Code}
Some \Forth{} editors will enter the stamp for you at the press of a key.

A common form for representing dates is
\begin{Code}
mm-dd-yy
\end{Code}
that is, February 6, 1984 would be expressed
\begin{Code}
02-06-84
\end{Code}
An increasingly popular alternative uses
\begin{Code}
ddMmmyy
\end{Code}
where ``Mmm'' is a three-letter abbreviation of the month.  For instance:
\begin{Code}
22Oct84
\end{Code}
%Page 146 in first edition.
This form requires fewer characters than
\begin{Code}
10-22-84
\end{Code}
and eliminates possible confusion between dates and months.%
\index{S!Stamp|)}%
\index{D!Dates, representation of|)}

If your system has \forth{\bs} (``skip-line''---see \App{C}), you can
write the comment line like this:
\begin{Code}
\ Chapter name        Screen Name -- pg.#       JPJ 06/10/83
\end{Code}
As with all comments, use lower-case or a mixture of lower- and
upper-case text in the comment line.

One way to make the index of an application reveal more about the
organization of the screens is to indent the comment line by three
spaces in screens that continue a lexicon.  \Fig{fig5-5} shows a
portion of a list produced by \forthb{INDEX} in which the comment lines
for the continuing screens are indented.%
\index{C!Comment line|)}

\begin{figure*}[bbbb]
\caption{The output of \forth{INDEX} showing indented comment lines.}
\labelfig{fig5-5}
\begin{Code}
 90 \ Graphics           Chapter load               JPJ 06/10/83
 91    \ Dot-drawing primitives                     JPJ 06/10/83
 92 \ Line-drawing primitives                       JPJ 06/11/83
 93    \ Line-drawing primitives                    JPJ 06/10/83
 94    \ Line-drawing primitives                    JPJ 09/02/83
 95 \ Scaling, rotation                             JPJ 06/10/83
 96    \ Scaling, rotation                          JPJ 02/19/84
 97    \ Scaling, rotation                          JPJ 02/19/84
 98    \ Scaling, rotation                          JPJ 02/19/84
 99 \ Boxes                                         JPJ 06/10/83
100 \ Circles                                       JPJ 06/10/83
101    \ Circles                                    JPJ 06/10/83
102    \ Circles                                    JPJ 06/10/83
\end{Code}
\end{figure*}

\begin{tip}
Begin all definitions at the left edge of the screen, and define only
one word per lines.
\end{tip}
%! There should be no page breaks between the text and code
\noindent \emph{Bad:}
\begin{Code}
: ARRIVING   ." HELLO" ;   : DEPARTING   ." GOODBYE" ;
\end{Code}
\noindent \emph{Good:}
\begin{Code}
: ARRIVING   ." HELLO" ;
: DEPARTING   ." GOODBYE" ;
\end{Code}
This rule makes it easier to find a definition in the listing.  (When
definitions continue for more than one line, the subsequent lines
should always be indented.)
%Page 147 in first edition.

\forthb{VARIABLE}s and \forthb{CONSTANT}s should also be defined one per
line.  (See ``Samples of Good Commenting Style'' in \App{E}.) This
leaves room for an explanatory comment on the same line.  The
exception is a large ``family'' of words (defined by a common
defining-word) which do not need unique comments:
\begin{Code}
0 HUE BLACK     1 HUE BLUE      2 HUE GREEN
3 HUE CYAN      4 HUE RED       5 HUE MAGENTA
\end{Code}
\begin{tip}
Leave lots of room at the bottom of the screen for later additions.
\end{tip}
On your first pass, fill each screen no more than half with code.  The
iterative approach demands that you sketch out the components of your
application first, then iteratively flesh them out until all the
requirements are satisfied.  Usually this means adding new commands,
or adding special-case handling, to existing screens.  (Not
\emph{always,} though.  A new iteration may see a simplification of
the code. Or a new complexity may really belong in another component
and should be factored out, into another screen.)

Leaving plenty of room at the outset makes later additions more
pleasant.  One writer recommends that on the initial pass, the screen
should contain about 20--40 percent code and 80--60 percent
whitespace \cite{stevenson81}.

Don't skip a line between each definition.  You may, however, skip a
line between \emph{groups} of definitions.%
\index{D!DECIMAL|(}\index{B!BASE|(}
\begin{tip}
All screens must leave \forthb{BASE} set to \forthb{DECIMAL}.
\end{tip}
Even if you have three screens in a row in which the code is written
in \forthb{HEX} (three screens of assembler code, for instance), each
screen must set \forth{BASE} to \forthb{HEX} at the top, and restore
base to \forthb{DECIMAL} at the bottom.  This rule ensures that each
screen could be loaded separately, for purposes of testing, without
mucking up the state of affairs.  Also, in reading the listing you
know that values are in decimal unless the screen explicitly
says \forthb{HEX}.

Some shops take this rule even further.  Rather than brashly resetting
base to \forthb{DECIMAL} at the end, they reset base to \emph{whatever
it was at the beginning.} This extra bit of insurance can be
accomplished in this fashion:
\begin{Code}
BASE @       HEX    \ save original BASE on stack
0A2 CONSTANT BELLS
0A4 CONSTANT WHISTLES
... etc. ...
BASE !              \ restore it
\end{Code}
%Page 148 in first edition.
\noindent Sometimes an argument is passed on the stack from screen to
screen, such as the value returned by \forthb{BEGIN} or \forthb{IF} in a
multiscreen assembler definition, or the base address passed from one
defining word to another---see ``Compile-Time Factoring'' in \Chap{6}.
In these cases, it's best to save the value of \forth{BASE} on the
return stack like this:
\begin{Code}
BASE @ >R     HEX
... etc. ...
R> BASE !
\end{Code}
Some folks make it a policy to use this approach on any screen that
changes \forthb{BASE}, so they don't have to worry about it.%
\index{B!BASE|)}

\person{Moore}\index{M!Moore, Charles} prefers to define \forthb{LOAD}
to invoke \forthb{DECIMAL} after loading.  This approach simplifies the
screen's contents because you don't have to worry about resetting.%
\index{D!DECIMAL|)}

\subsection{Spacing and Indentation}%
\index{S!Spacing|(}

\begin{tip}
Spacing and indentation are essential for readability.
\end{tip}
The examples in this book use widely accepted conventions of spacing
and indenting style.  Whitespace, appropriately used, lends
readability.  There's no penalty for leaving space in source screens
except disk memory, which is cheap.

For those who like their conventions in black and white, Table
\ref{tab-5-1} is a list of guidelines.  (But remember, \Forth{}'s
interpreter couldn't care less about spacing or indentation.)

\begin{table}[bbbb]
\caption{Indentation and spacing guidelines}
\label{tab-5-1}
\medskip\blackline{0pt}\medskip
\begin{minipage}{\textwidth}
\begin{quote}
1 space between the colon and the name\\
2 spaces between the name and the comment\footnote{An often-seen
alternative calls for 1 space between the name and comment and 3
between the comment and the definition.  A more liberal technique uses
3 spaces before and after the comment.  Whatever you choose, be
consistent.}\\
2 spaces, or a carriage return, after the comment and
before the definition\footnotemark[1]\\
3 spaces between the name and definition if no comment is used\\
3 spaces indentation on each subsequent line (or multiples
of 3 for nested indentation)\\
1 space between words/numbers within a phrase\\
2 or 3 spaces between phrases\\
1 space between the last word and the semicolon\\
1 space between semicolon and \forthb{IMMEDIATE} (if invoked)
\end{quote}
No blank lines between definitions, except to separate distinct groups
of definitions

%% Note: AH: Again type writer style footnotes.
\end{minipage}
\medskip\blackline{0pt}
\end{table}

%Page 149 in first edition.

The last position of each line should be blank except for:
%%!!!te original used a) b) enumeration here.
\ifeightyfour
\begin{list}{\alph{enumi})}{\usecounter{enumi}}
\else
\begin{enumerate}
\fi
\item quoted strings that continue onto the next line, or
\item the end of a comment.
\ifeightyfour
\end{list}
\else
\end{enumerate}
\fi
A comment that begins with \forth{\bs} may continue right to the end
of the line.  Also, a comment that begins with \forth{(} may have its
delimiting right parenthesis in the last column.

\index{I!Indentation|(}%
Here are some common errors of spacing and indentation:

\goodbreak\bigskip\noindent
\emph{Bad} (name not separated from the body of the definition):
\begin{Code}
: PUSH HEAVE HO ;
\end{Code}
\emph{Good:}
\begin{Code}
: PUSH   HEAVE HO ;
\end{Code}
\emph{Bad} (subsequent lines not indented three spaces):
\begin{Code}
: RIDDANCE  ( thing-never-to-darken-again -- )
DARKEN  NEVER AGAIN ;
\end{Code}
\emph{Good:}
\begin{Code}
: RIDDANCE  ( thing-never-to-darken-again -- )
   DARKEN  NEVER AGAIN ;
\end{Code}
\emph{Bad} (lack of phrasing):
\begin{Code}
: GETTYSBURG   4 SCORE 7 YEARS + AGO ;
\end{Code}
\goodbreak\noindent
\emph{Good:}
\begin{Code}
: GETTYSBURG   4 SCORE   7 YEARS +   AGO ;
\end{Code}
Phrasing is a subjective art;
I've yet to see a useful set of formal rules.\\
Simply strive for readability.%
\index{S!Spacing|)}%
\index{S!Screens!layout|)}%
\index{I!Implementation!screen layout|)}%
\index{I!Indentation|)}

\section{Comment Conventions}%
\index{C!Comment conventions|(}%
\index{I!Implementation!comment conventions|(}

Appropriate commenting is essential.  There are five types of comments:
stack-effect comments, data-structure comments, input-stream comments,
purpose comments and narrative comments.
%%!!!te Original didn't put this into a description
\begin{description}
\item[\emph{A} stack-effect comment]\index{S!Stack-effect comment}%
shows the arguments that the definition consumes from the stack, and
the arguments it returns to the stack, if any.
%Page 150 in first edition.

\item[\emph{A} data-structure comment]\index{D!Data-structure comment}%
indicates the position and meaning of elements in a data structure.
For instance, a text buffer might contain a count in the first byte,
and 63 free bytes for text.

\item[\emph{An} input-stream comment]\index{I!Input-stream comment}%
indicates what strings the word expects to see in the input stream.
For example, the \Forth{} word \forth{FORGET} scans for the name of a
dictionary entry in the input stream.

\item[\emph{A} purpose comment]\index{P!Purpose comment}%
describes, in as few words possible, what the definition does.  How
the definition works is not the concern of the purpose comment.

\item[\emph{A} narrative comment]\index{N!Narrative comments}%
appears amidst a definition to explain what is going on, usually
line-by-line.  Narrative comments are used only in the ``vertical
format,'' which we'll describe in a later section.
\end{description}

Comments are usually typed in lower-case letters to distinguish them
from source code.  (Most \Forth{} words are spelled with upper-case
letters, but lower-case spellings are sometimes used in special cases.)

In the following sections we'll summarize the standardized formats
for these types of comments and give examples for each type.

\subsection{Stack Notation}%
\index{S!Stack notation|(}
\begin{tip}
Every colon or code definition that consumes and/or returns any arguments
on the stack must include a stack-effect comment.
\end{tip}
``Stack notation'' refers to conventions for representing what's on
the stack.  Forms of stack notation include ``stack pictures,''
``stack effects,'' and ``stack-effect comments.''%
\index{S!Stack notation|)}%
\index{P!Purpose comment}

\subsection{Stack Picture}%
\index{S!Stack picture|(}

A stack picture depicts items understood to be on the stack at a given
time.  Items are listed from left to right, with the leftmost item
representing the bottom of the stack and the rightmost item
representing the top.

For instance, the stack picture
\begin{Code}
nl n2
\end{Code}
indicates two numbers on the stack, with n2 on the top (the most
accessible position).

This is the same order that you would use to type these values in;
i.e., if n1 is 100 and n2 is 5000, then you would type
\begin{Code}
100 5000
\end{Code}
to place these values correctly on the stack.

%Page 151 in first edition.
A stack picture can include either abbreviations, such as ``n1,'' or
fully spelled-out words.  Usually abbreviations are used.  Some
standard abbreviations appear in Table \ref{tab-5-2}.  Whether
abbreviations or fully spelled-out words are used, each stack item
should be separated by a space.

If a stack item is described with a phrase (such as
``address-of-latest-link''), the words in the phrase should be joined
by hyphens.  For example, the stack picture:
\begin{Code}
address current-count max-count
\end{Code}
shows three elements on the stack.%
\index{S!Stack picture|)}

\subsection{Stack Effect}%
\index{S!Stack effect|(}

A ``stack effect'' shows two stack pictures: one picture of any items
that may be \emph{consumed} by a definition, and another picture of
any items \emph{returned} by the definition.  The ``before'' picture
comes first, followed by two hyphens, then the ``after'' picture.

For instance, the stack effect for \Forth{}'s addition operator,
\forth{+} is
\begin{Code}
n n -- sum
\end{Code}
where \forth{+} consumes two numbers and returns their sum.

Remember that the stack effect describes only the \emph{net result} of
the operation on the stack.  Other values that happen to reside on the
stack beneath the arguments of interest don't need to be shown.  Nor
do values that may appear or disappear while the operation is
executing.

If the word returns any input arguments unchanged, they
should be repeated in the output picture; e.g.,
\begin{Code}
3rd 2nd top-input -- 3rd 2nd top-output
\end{Code}
Conversely, if the word changes any arguments, the stack comment must
use a different descriptor:
\begin{Code}
nl -- n2
n -- n'
\end{Code}
A stack effect might appear in a formatted glossary.%
\index{S!Stack effect|)}

\subsection{Stack Effect Comment}%
\index{S!Stack-effect comment|(}

A ``stack-effect comment'' is a stack effect that appears in source
code surrounded by parentheses.  Here's the stack-effect comment for
the word \forth{COUNT}:
\begin{Code}
( address-of-counted-string -- address-of-text count)
\end{Code}
%Page 152 in first edition.
or:
\begin{Code}
( 'counted-string -- 'text count)
\end{Code}
(The ``count'' is on top of the stack after the word has executed.)

If a definition has no effect on the stack (that is, no effect the
user is aware of, despite what gyrations occur within the definition),
it needs no stack-effect comment:
\begin{Code}
: BAKE   COOKIES OVEN ! ;
\end{Code}
On the other hand, you may want to use an empty stack comment---i.e.,
\begin{Code}
: BAKE   ( -- )  COOKIES OVEN ! ;
\end{Code}
to emphasize that the word has no effect on the stack.

If a definition consumes arguments but returns none, the double-hyphen
is optional.  For instance,
\begin{Code}
( address count -- )
\end{Code}
can be shortened to
\begin{Code}
( address count)
\end{Code}
The assumption behind this convention is this: There are many more
colon definitions that consume arguments and return nothing than
definitions that consume nothing and return arguments.%
\index{S!Stack-effect comment|)}

\subsection{Stack Abbreviation Standards}%
\index{S!Stack abbreviation standards|(}

Abbreviations used in stack notation should be consistent. Table
\ref{tab-5-2} lists most of the commonly used abbreviations.  (This
table reappears in \App{E}.) The terms ``single-length,''
``double-length,'' etc. refer to the size of a ``cell'' in the
particular \Forth{} system.  (If the system uses a 16-bit cell, ``n''
represents a 16-bit number; if the system uses a 32-bit cell, ``n''
represents a 32-bit number.)%
\index{S!Stack abbreviation standards|)}

\subsection{Notation of Flags}

Table \ref{tab-5-2} shows three ways to represent a boolean flag.  To
illustrate, here are three versions of the same stack comment for the
word \forth{-TEXT}:
\begin{Code}
( at u a2 -- ?)
( at u a2 -- t=no-match)
( at u a2 -- f=match)
\end{Code}
%Page 153 in first edition.
\begin{table}[hhhh]
\caption{Stack-comment abbreviations.}
\label{tab-5-2}
\vspace{1ex}
\blackline{1ex}
\begin{tabular}{@{\hspace{2.5em}}ll}
n             &  single-length signed number \\
d             &  double-length signed number \\
u             &  single-length unsigned number \\
ud            &  double-length unsigned number \\
t             &  triple-length \\
q             &  quadruple-length \\
c             &  7-bit character value \\
b             &  8-bit byte \\
?             &  boolean flag; or; \\
\quad t=         &  true \\
\quad f=         &  false \\
a or adr      &  address \\
acf           &  address of code field \\
apf           &  address of parameter field \\
'             &  (as prefix) address of \\
s d           &  (as a pair) source destination \\
lo hi         &  lower-limit upper-limit (inclusive) \\
\#            &   count \\
0             &  offset \\
i             &  index \\
m             &  mask \\
x             &  don't care (data structure notation) \\
\end{tabular}
\smallskip

An ``offset'' is a difference expressed in absolute units, such as bytes.\\
An ``index'' is a difference expressed in logical units, such as
elements or records.
\vspace{0ex}
\blackline{0ex}
\end{table}
The equal sign after the symbols ``t'' and ``f'' equates the flag
outcome with its meaning.  The result-side of the second version would
be read ``true means no match.''

\subsection{Notation of Variable Possibilities}%
\index{V!Variable possibilities, notation of|(}
Some definitions yield a different stack effect under different
circumstances.

If the number of items on the stack remains the same under all
conditions, but the items themselves change, you can use the vertical
bar ( \forth{|} ) to mean ``or.'' The following stack-effect comment
describes a word that returns either the address of a file or, if the
requested file is not found, zero:
\begin{Code}
( -- address|O=undefined-file)
\end{Code}
If the number of items in a stack picture can vary---in either the
``before'' or ``after'' picture---you must write out both versions of
the entire stack
%Page 154 in first edition.
picture, along with the double-hyphen, separated by the ``or'' symbol.
For instance:
\begin{Code}
-FIND   ( -- apf len t=found | -- f=not-found )
\end{Code}
This comment indicates that if the word is found, three arguments are
returned (with the flag on top); otherwise only a false flag is
returned.

Note the importance of the second ``-\,-''.  Its omission would indicate
that the definition always returned three arguments, the top one being
a flag.

\index{S!Stack-effect comment|(}%
If you prefer, you can write the entire stack effect twice, either on
the same line, separated by three spaces:
\begin{Code}
?DUP   \ if zero: ( n -- n)    if non-zero:( n -- n n)
\end{Code}
or listed vertically:
\begin{Code}
-FIND  \     found:( -- apf len t )
       \ not-found:( -- f )
\end{Code}
\index{S!Stack-effect comment|)}%
\index{V!Variable possibilities, notation of|)}

\subsection{Data-Structure Comments}%
\index{D!Data-structure comment|(}

A ``data-structure comment'' depicts the elements in a data structure.
For example, here's the definition of an insert buffer called
\forth{|INSERT} :

\begin{Code}
CREATE |INSERT  64 ALLOT  \  { 1# | 63text }
\end{Code}
The ``faces'' (curly-brackets) begin and end the structure comment;
the bars separate the various elements in the structure; the numbers
represent bytes per element.  In the comment above, the first byte
contains the count, and the remaining 63 bytes contain the text.

A ``bit comment'' uses the same format as a data-structure comment to
depict the meaning of bits in a byte or cell.  For instance, the bit
comment
\begin{Code}
{ 1busy? | 1acknowledge? | 2x | 6input-device |
   6output-device }
\end{Code}
describes the format of a 16-bit status register of a communications
channel.  The first two bits are flags,\index{F!Flags}
the second two bits are unused, and the final pair of six-bit fields
indicate the input and output devices which this channel is connected
to.

If more than one data structure employs the same pattern of elements,
write out the comment only once (possibly in the preamble), and
%Page 155 in first edition.
give a name to the pattern for reference in subsequent screens.  For
instance, if the preamble gives the above bit-pattern the name
``status,'' then ``status'' can be used in stack comments to indicate
values with that pattern:
\begin{Code}
: STATUS?  ( -- status) ... ;
\end{Code}
If a \forthb{2VARIABLE} contains one double-length value, the comment
should be a stack picture that indicates the contents:
\begin{Code}
2VARIABLE PRICE  \ price in cents
\end{Code}
If a \forthb{2VARIABLE} contains two single-length data elements, it's
given a stack picture showing what would be on the stack after a
\forth{2@}. Thus:
\begin{Code}
2VARIABLE MEASUREMENTS  ( height weight )
\end{Code}
This is different from the comment that would be used if
\forth{MEASUREMENTS} were defined by \forthb{CREATE}.
\begin{Code}
CREATE MEASUREMENTS  4 ALLOT    \ { 2weight | 2height }
\end{Code}
(While both statements produce the same result in the dictionary, the
use of \forthb{2VARIABLE} implies that the values will normally be
``2-fetched'' and ``2-stored'' together-thus we use a \emph{stack}
comment.  The high-order part, appearing on top of the stack, is
listed to the right.  The use of \forthb{CREATE} implies that the
values will normally be fetched and stored separately--thus we use a
data structure comment.  The item in the 0th position is listed to the
left.)%
\index{D!Data-structure comment|)}

\subsection{Input-stream Comments}%
\index{I!Input-stream comment|(}

The input-stream comment indicates what words and/or strings are
presumed to be in the input stream. Table \ref{tab-5-3} lists the
designations used for input stream arguments.

\begin{table*}[hhhh]
\caption{Input-stream comment designations.}
\label{tab-5-3}
\vspace{1ex}\blackline{1ex}
\centerline{\begin{tabular}{ll}
c     &  single character, blank-delimited \\
name  &  sequence of characters, blank delimited \\
text  &  sequence of characters, delimited by non-blank \\
\end{tabular}}

\medskip
Follow ``text'' with the actual delimiter required; e.g.:
\forth{text"} or \forth{text)}
\vspace{1ex}\blackline{0ex}
\end{table*}
%% Note: AH : typewriter style footnote with *
%% Note: AH : text" and text) are literal examples,
%%    and should be slanted or such.
The input-stream comment appears \emph{before} the stack comment, and
is \emph{not} encapsulated between its own pair of parentheses, but
simply surrounded
%Page 156 in first edition.
by three spaces on each side.  For instance, here's one way to comment
the definition of \forth{'} (tick) showing first the input-stream comment,
then the stack comment:
\begin{Code}
: '   \ name   ( -- a)
\end{Code}
If you prefer to use \forth{(} , the comment would look like this:
\begin{Code}
: '   ( name   ( -- a)
\end{Code}
\index{S!String input, receiving|(}%
Incidentally, there are three distinct ways to receive string input.
To avoid confusion, here are the terms:

%%!!!te Original had bullet points here + description
\begin{description}
\item[Scanning-for]\index{S!Scanning-for}%
means looking ahead in the input stream, either for a word or number
as in the case of tick, or for a delimiter as in the case of
\forth{."} and \forth{(} .
\item[Expecting]\index{E!Expecting}%
means waiting for. \forth{EXPECT} and \forth{KEY}, and definitions
that invoke them, are ones that ``expect'' input.
\item[Presuming]\index{P!Presuming}%
indicates that in normal usage something will follow.  The word:
``scans-for'' the name to be defined, and ``presumes'' that a
definition will follow.
\end{description}
The input-stream comment is only appropriate for input being scanned-for.%
\index{S!String input, receiving|)}%
\index{I!Input-stream comment|)}

\subsection{Purpose Comments}%
\index{P!Purpose comment|(}

\begin{tip}
Every definition should bear a purpose comment unless:
\medskip
%%!!!te Original had a. b. enumeration
\begin{enumerate}
\item its purpose is clear from its name or its stack-effect comment, or
\item if it consists of three or fewer words.
\end{enumerate}
\end{tip}
The purpose comment should be kept to a minimum-never more than a full
line.  For example:
\begin{Code}
: COLD   \ restore system to start condition
    ... ;
\end{Code}
Use the imperative mood: ``set Foreground color,'' not ``sets
Foreground color.''

On the other hand, a word's purpose can often be described in terms of
its stack-effect comment.  You rarely need both a stack comment and a
purpose comment.  For instance:
\begin{Code}
: SPACES  ( #)
\end{Code}
%Page 157 in first edition.
or
\begin{Code}
: SPACES  ( #spaces-to-type -- )   ... ;
\end{Code}
This definition takes as its incoming argument a number that
represents the number of spaces to type.
\begin{Code}
: ELEMENT  ( element# -- 'element)  2*  TABLE + ;
\end{Code}
This definition converts an index, which it consumes, into an address
within a table of 2-byte elements corresponding to the indexed element.
\begin{Code}
: PAD  ( -- 'scratch-pad)  HERE  80 + ;
\end{Code}
This definition returns an address of a scratch region of memory.

Occasionally, readability is best served by including both types of
comment. In this case, the purpose comment should appear last.  For
instance:
\begin{Code}
: BLOCK  ( n -- a)  \   ensure block n in buffer at a
\end{Code}

\begin{tip}
Indicate the type of comment by ordering: input-stream comments first,
stack-effect comments second, purpose comments last.
\end{tip}
For example:
\begin{Code}
: GET   \   name   ( -- a)   get first match
\end{Code}
If you prefer to use \forth{(}, then write:
\begin{Code}
: GET   (   name  ( -- a)    ( get first match)
\end{Code}
If necessary, you can put the purpose comment on a second line:
\begin{Code}
: WORD   \   name   ( c -- a)
   \ scan for string delimt'd by "c"; leave at a
   ...  ;
\end{Code}
\index{P!Purpose comment|)}%

\subsection{Comments for Defining Words}%
\index{D!Defining words:!comments for|(}
The definition of a defining word involves two behaviors:

\begin{itemize}
\item that of the defining word as it defines its ``child''
(compile-time behavior), and
\item that of the child itself (run-time behavior).
\end{itemize}
%Page 158 in first edition.
These two behaviors must be commented separately.

\index{C!Compiling words, comments for|(}%
\begin{tip}
Comment a defining word's compile-time behavior in the usual way;
comment its run-time behavior separately, following the word
\forthb{DOES>} (or \forthb{;CODE}).
\end{tip}
For instance,
\begin{Code}
: CONSTANT  ( n )
   DOES>  ( -- n)  @ ;
\end{Code}
The stack-effect comment for the run-time (child's) behavior
represents the net stack effect for the child word.  Therefore it does
not include the address returned by \forthb{DOES>,} even though this
address is on the stack when the run-time code begins.

\bigskip\noindent
\emph{Bad} (run-time comment includes apf):
\begin{Code}
: ARRAY   \  name  ( #cells)
   CREATE 2* ALLOT
   DOES>   ( i apf -- 'cell)  SWAP  2* + ;
\end{Code}
\goodbreak\noindent
\emph{Good:}
\begin{Code}
: ARRAY   \  name  ( #cells)
   CREATE 2* ALLOT
    DOES>  ( i -- 'cell)  SWAP  2* + ;
\end{Code}
%% Note: AH: CELLS instead of 2* is indicated
Words defined by this word \forth{ARRAY} will exhibit the stack effect:
\begin{Code}
( i -- 'cell)
\end{Code}
If the defining word does not specify the run-time behavior, there still
exists a run-time behavior, and it may be commented:
\begin{Code}
: VARIABLE   (  name  ( -- )  CREATE  2 ALLOT ;
   \ does>   ( -- adr )
\end{Code}
\index{D!Defining words:!comments for|)}%

\subsection{Comments for Compiling Words}
As with defining words, most compiling words involve two behaviors:
\begin{enumerate}
\item That of the compiling word as the definition in which it appears
is compiled
\item That of the run-time routine which will execute when we invoke
the word being defined.  Again we must comment each behavior
separately.
\end{enumerate}
%Page 159 in first edition.

\begin{tip}
Comment a compiling word's run-time behavior in the usual way; comment
its compile-time behavior separately, beginning with the label
``Compile:''.
\end{tip}
For instance:
\begin{Code}
: IF   ( ? -- ) ...
\ Compile:   ( -- address-of-unresolved-branch)
   ... ; IMMEDIATE
\end{Code}
In the case of compiling words, the first comment describes the
run-time behavior, which is usually the \emph{syntax for using} the
word.  The second comment describes what the word \emph{actually does}
in compiling (which is of less importance to the user).

Other examples:
\begin{Code}
: ABORT"  ( ? -- )
\ Compile:   text"   ( -- )
\end{Code}
Occasionally a compiling word may exhibit a different behavior when it
is invoked \emph{outside} a colon definition.  Such words (to be
fastidious about it) require three comments.  For instance:
\begin{Code}
: ASCII  ( -- c)
\ Compile:   c   ( -- )
\ Interpret:   c   ( -- c )
     ... ; IMMEDIATE
\end{Code}
\App{E} includes two screens showing good commenting style.%
\index{C!Comment conventions|)}%
\index{C!Compiling words, comments for|)}%
\index{I!Implementation!comment conventions|)}

\section{Vertical Format vs.\ Horizontal Format}%
\index{V!Vertical format vs. horizontal format|(}%
\index{I!Implementation!vertical format vs. horizontal format|(}
The purpose of commenting is to allow a reader of your code to easily
determine what's going on.  But how much commenting is necessary? To
determine the level of commenting appropriate for your circumstances,
you must ask yourself two questions:

%%!!! this didn't have any bullet points or whatever in the first edition
\begin{itemize}
\expandafter\item\ifeightyfour[]\fi Who will be reading my code?
\expandafter\item\ifeightyfour[]\fi How readable are my definitions?
\end{itemize}%
\index{N!Narrative comments|(}
There are two basic styles of commenting to choose from.  The first
style, often called the ``vertical format,'' includes a step-by-step
description of the process, in the manner of a well-commented assembly
language listing.  These line-by-line comments are called ``narrative
comments.''

%Page 160 in first edition.
\begin{Code}
\ CRC Checksum                                      07/15/83
: ACCUMULATE   ( oldcrc char -- newcrc)
   256 *               \ shift char to hi-order byte
   XOR                 \ & xor into previous crc
   8 0 DO              \ Then for eight repetitions,
       DUP 0< IF       \ if hi-order bit is "1"
          16386 XOR    \ xor it with mask and
          DUP +        \ shift it left one place
          1+           \ set lo-order bit to "1"
              ELSE     \ otherwise, i.e. hi-order bit is "0"
          DUP +        \ shift it left one place
              THEN
       LOOP ;          \ complete the loop
\end{Code}
The other approach does not intersperse narrative comments between
code phrases.  This is called the ``horizontal format.''
\begin{Code}
: ACCUMULATE  ( oldcrc char -- newcrc)
   256 *  XOR  8 0 DO  DUP 0< IF
      16386 XOR  DUP +  1+  ELSE  DUP +  THEN  LOOP ;
\end{Code}
The vertical format is preferred when a large team of programmers are
coding and maintaining the application.  Typically, such a team will
include several junior-level programmers responsible for minor
corrections.  In such an environment, diligent commenting can save a
lot of time and upset.  As \person{Johnson}\index{J!Johnson, Dave} of Moore
Products Co.  says: ``When maintaining code you are usually interested
in just one small section, and the more information written there the
better your chances for a speedy fix.''

Here are several pertinent rules required of the \Forth{} programmers
at Moore Products Co. (I'm paraphrasing):

\begin{enumerate}
\item A vertical format will be used.  Comments will appear to the
right of the source code, but may continue to engulf the next line
totally if needed.
\item There should be more comment characters than source characters.
(The company encourages long descriptive names, greater than ten
characters, and allows the names to be counted as comment characters.)
\item Any conditional structure or application word should appear on a
separate line.  ``Noise words'' can be grouped together.  Indentation
is used to show nested conditionals.
\end{enumerate}
There are some difficulties with this format, however.%
\index{L!Line-by-line commenting|(}
For one thing, line-by-line commenting is time-consuming, even with a
good screen editor.  Productivity can be stifled, especially when
stopping to write the comments breaks your chain of thought.

Also, you must also carefully ensure that the comments are up-to-date.
Very often code is corrected, the revision is tested, the change
%Page 161 in first edition.
works---and the programmer forgets to change the comments.  The more
comments there are, the more likely they are to be wrong.  If they're
wrong, they're worse than useless.

This problem can be alleviated if the project supervisor carefully
reviews code and ensures the accuracy of comments.

Finally, line-by-line commenting can allow a false sense of security.
Don't assume that because each \emph{line} has a comment, the
\emph{application} is well-com\-men\-ted.  Line-by-line commenting doesn't
address the significant aspects of a definition's operation.  What,
for instance, is the thinking behind the checksum algorithm used? Who
knows, from the narrative comments?%
\index{L!Line-by-line commenting|)}%
\index{N!Narrative comments|)}

To properly describe, in prose, the implications of a given procedure
usually requires many paragraphs, not a single phrase.  Such
descriptions properly belong in auxiliary documentation or in the
chapter preamble.

Despite these cautions, many companies find the vertical format
necessary.  Certainly a team that is newly exposed to \Forth{} should
adopt it, as should any very large team.

What about the horizontal format? Perhaps it's an issue of art vs.\@
practicality, but I feel compelled to defend the horizontal format as
equally valid and in some ways superior.

If \Forth{} code is really well-written, there should be nothing
ambiguous about it.  This means that:

\begin{itemize}
\item supporting lexicons have a well-designed syntax
\item stack inputs and outputs are commented
\item the purpose is commented (if it's not clear from the name or
stack comment)
\item definitions are not too long
\item not too many arguments are passed to a single definition via the
stack (see ``The Stylish Stack'' in \Chap{7}).
\end{itemize}
\Forth{} is simply not like other languages, in which line-by-line
commenting is one of the few things you can do to make programs more
readable.

Skillfully written \Forth{} code is like poetry, containing precise
meaning that both programmer and machine can easily read.  Your
\emph{goal} should be to write code that does not need commenting,
even if you choose to comment it.  Design your application so that the
code, not the comments, conveys the meaning.

If you succeed, then you can eliminate the clutter of excessive
commenting, achieving a purity of expression without redundant
explanations.

\wepsfigp{fig5-2}{Wiggins, proud of his commenting technique.}

\begin{tip}
The most-accurate, least-expensive documentation
is self-documenting code.
\end{tip}
%Page 162 in first edition.
Unfortunately, even the best programmers, given the pressure of a
deadline, may write working code that is not easily readable without
comments.  If you are writing for yourself, or for a small group with
whom you can verbally communicate, the horizontal format is ideal.
Otherwise, consider the vertical format.%
\index{V!Vertical format vs. horizontal format|)}%
\index{I!Implementation!vertical format vs. horizontal format|)}

\section{Choosing Names: The Art}%
\index{I!Implementation!choosing names|(}%
\index{N!Names:!choosing|(}

\begin{tfquot}
Besides a mathematical inclination, an exceptionally good mastery of
one's native tongue is the most vital asset of a competent programmer
\emph{(Prof.  \person{Edsger W. Dijkstra} \cite{dijkstra82}).}
\end{tfquot}%
\index{W!Words:!choosing names|(}
We've talked about the significance of using names to symbolize ideas
and objects in the application.  The choosing of names turns out to be
an important part of the design process.

Newcomers tend to overlook the important of names.  ``After all,''
they think, ``the computer doesn't care what names I choose.''

But good names are essential for readability.  Moreover, the mental
exercise of summoning a one-word description bears a synergistic effect
on your perceptions of what the entity should or should not do.

\smallbreak
Here are some rules for choosing good names:
\begin{tip}
Choose names according to ``what,'' not ``how.''
\end{tip}
A definition should hide the complexities of implementation from other
definitions which invoke it.  The name, too, should hide the details
of the procedure, and instead should describe the outward appearance
or net effect.

For instance, the \Forth{} word \forthb{ALLOT} simply increments the
dictionary pointer (called \forthb{DP} or \forthb{H} in most systems).
But the name \forthb{ALLOT} is better than \forth{DP+!} because the
user is thinking of reserving space, not incrementing a pointer.

The '83 Standard adopted the name \forthb{CMOVE>} instead of the
previous name for the same function, \forthb{<CMOVE}.  The operation
makes it possible to copy a region of memory \emph{forward} into
overlapping memory.  It accomplishes this by starting with the last
byte and working \emph{backward}.  In the new name, the forwardness of
the ``what'' supersedes the backwardness of the ``how.''

\begin{tip}
Find the most expressive word.\index{E!Expressive words}
\end{tip}
%Page 164 in first edition.

\begin{tfquot}
A powerful agent is the right word.  Whenever we come upon one of
those intensely right words in a book or a newspaper the resulting
effect is physical as well as spiritual, and electrically prompt
\emph{(\person{Mark Twain}).}

The difference between the right word and the almost-right word is
like the difference between lightning and the lightning bug
\emph{(\person{Mark Twain}).}

Suit the action to the word, the word to the action
\emph{(\person{Shakespeare}, Hamlet, Act~III).}
\end{tfquot}
\person{Henry Laxen},\index{L!Laxen, Henry}
a \Forth{} consultant and author, suggests that the most important
\Forth{} development tool is a good thesaurus \cite{laxen}.%
\index{T!Thesaurus}

Sometimes you'll think of an adequate word for a definition, but it
doesn't feel quite right.  It may be months later before you realize
that you fell short of the mark.  In the Roman numeral example in
\Chap{4}, there's a word that handles the exception case: numbers that
are one-less-than the next symbol's value.  My first choice was
\forth{4-0R-9}.  That's awkward, but it was much later that I thought
of \forth{ALMOST}.

Most fig-\Forth{} systems include the word \forth{VLIST}, which lists
the names of all the words in the current vocabulary.  After many
years someone realized that a nicer name is \forth{WORDS}.  Not only
does \forth{WORDS} sound more pleasant by itself, it also works nicely
with vocabulary names.  For instance:
\begin{Code}
EDITOR WORDS
\end{Code}
or
\begin{Code}
ASSEMBLER WORDS
\end{Code}
On the other hand, \person{Moore}\index{M!Moore, Charles} points out
that inappropriate names can become a simple technique for encryption.%
\index{E!Encryption}
If you need to provide security when you're forced to distribute
source, you can make your code very unreadable by deliberately
choosing misleading names.  Of course, maintenance becomes impossible.

\index{P!Phrases|(}%
\begin{tip}
Choose names that work in phrases.
\end{tip}
Faced with a definition you don't know what to call, think about how the
word will be used in context.  For instance:
%%!!!te this is how it looks like in the original
%%!!!te This really seems to be a description environment
\ifeightyfour
\begin{Code}[fontfamily=cmss,commandchars=\&\{\}]
SHUTTER OPEN
  OPEN is the appropriate name for a word that sets a
  bit in an I/O address identified with the name
  SHUTTER.&medskip
3 BUTTON DOES IGNITION
  DOES is a good choice for a word that vectors the
  address of the function IGNITION into a table of
  functions, so that IGNITION will be executed when
  Button 3 is pushed.&medskip
SAY HELLO
  SAY is the perfect choice for vectoring HELLO into an
  execution variable.  (When I first wrote this example
  for Starting &Forth{}, I called it VERSION. &person{Moore}
  reviewed the manuscript and suggested SAY, which is
  clearly much better.)&medskip
I'M HARRY
  The word I'M seems more natural than LOGON HARRY,
  LOGIN HARRY or SESSION HARRY, as often seen.
\end{Code}
\else
\begin{description}
\item[\forth{SHUTTER OPEN}]~

  \forth{OPEN} is the appropriate name for a word that sets a
  bit in an I/O address identified with the name
  \forth{SHUTTER}.
\item[\forth{3 BUTTON DOES IGNITION}]~

  \forth{DOES} is a good choice for a word that vectors the
  address of the function \forth{IGNITION} into a table of
  functions, so that \forth{IGNITION} will be executed when
  Button 3 is pushed.
\item[\forth{SAY HELLO}]~

  \forth{SAY} is the perfect choice for vectoring \forth{HELLO} into an
  execution variable.  (When I first wrote this example
  for Starting \Forth{}, I called it \forth{VERSION}. \person{Moore}
  reviewed the manuscript and suggested \forth{SAY}, which is
  clearly much better.)
\item[\forth{I'M HARRY}]~

  The word \forth{I'M} seems more natural than \forth{LOGON HARRY},
  \forth{LOGIN HARRY} or \forth{SESSION HARRY}, as often seen.
\end{description}
\fi

\index{M!Moore, Charles|(}%
\begin{interview}
The choice of \forth{I'M} is another invention of \person{Moore}, who
says:

\begin{tfquot}
I detest the word \forth{LOGON}. There is no such word in English.  I
was looking for a word that said, ``I'm \dots{}'' It was a natural.
I just stumbled across it.  Even though it's clumsy with that
apostrophe, it has that sense of rightness.

All these little words are the nicest way of getting the ``Aha!''
reaction.  If you think of the right word, it is \emph{obviously} the
right word.

If you have a wide recall vocabulary, you're in a better position to
come up with the right word.
\end{tfquot}
\end{interview}\index{M!Moore, Charles|)}%
Another of \person{Moore}'s favorite words is \forth{TH}, which he
uses as an array indexing word.  For instance, the phrase
\begin{Code}
5 TH
\end{Code}
returns the address of the ``fifth'' element of the array.%
\index{P!Phrases|)}

\index{S!Spelled-out names|(}%
\begin{tip}
Spell names in full.
\end{tip}
I once saw some \Forth{} code published in a magazine in which the
author seemed hell-bent on purging all vowels from his names,
inventing such eyesores as \forth{DSPL-BFR} for ``display buffer.''
Other writers seem to think that three characters magically says it
all, coining \forth{LEN} for ``length.'' Such practices reflect
thinking from a bygone age.

\index{A!Abbreviations|(}%
\Forth{} words should be fully spelled out.  Feel proud to type every
letter of \forth{INITIALIZE} or \forth{TERMINAL} or \forth{BUFFER}.
These are the words you mean.
%Page 166 in first edition.

The worst problem with abbreviating a word is that you forget just
how you abbreviated it.  Was that \forth{DSPL} or \forth{DSPLY?}

Another problem is that abbreviations hinder readability.
Any programming language is hard enough to read without compounding
the difficulty.

Still, there are exceptions.  Here are a few:

\begin{enumerate}
\item Words that you use extremely frequently in code. \Forth{}
employs a handful of commands that get used over and over, but have
little or no intrinsic meaning:
\begin{Code}
:   ;   @   !   .   ,
\end{Code}
But there are so few of them, and they're used so often, they become
old friends.  I would never want to type, on a regular basis,
\begin{Code}
DEFINE   END-DEFINITION   FETCH   STORE
     PRINT   COMPILE#
\end{Code}
(Interestingly, most of these symbols don't have English counterparts.
We use the phrase ``\emph{colon} definition'' because there's no other
term; we say ``\emph{comma} a number into the dictionary'' because
it's not exactly compiling, and there's no other term.)
\item Words that a terminal operator might use frequently to control
an operation.  These words should be spelled as single letters, as are
line editor commands.
\item Words in which familiar usage implies that they be abbreviated.
\Forth{} assembler mnemonics are typically patterned after the
manufacturer's suggested mnemonics, which are abbreviations (such as
\forth{JMP} and \forth{MOV}).
\end{enumerate}%
\index{S!Spelled-out names|)}
Your names should be pronounceable; otherwise you may regret it when
you try to discuss the program with other people.  If the name is
symbolic, invent a pronunciation (e.g., \forth{>R} is called ``to-r'';
\forth{R>} is called ``r-from'').

\begin{tip}
Favor short words.
\end{tip}
Given the choice between a three-syllable word and a one-syllable word
that means the same thing, choose the shorter.  \forth{BRIGHT} is a
better name than \forth{INTENSE}.  \forth{ENABLE} is a better name
than \forth{ACTIVATE}; \forth{GO}, \forth{RUN}, or \forth{ON} may be
better still.

Shorter names are easier to type.  They save space in the source
screen.  Most important, they make your code crisp and clean.

\index{H!Hyphenated names|(}%
\begin{tip}
Hyphenated names may be a sign of bad factoring.
\end{tip}%
\index{A!Abbreviations|)}
%Page 168 in first edition.

\begin{interview}
\person{Moore}:\index{M!Moore, Charles|(}

\begin{tfquot}
There are diverging programming styles in the \Forth{} community.  One
uses hyphenated words that express in English what the word is doing.
You string these big long words together and you get something that is
quite readable.

But I immediately suspect that the programmer didn't think out the words
carefully enough, that the hyphen should be broken and the words defined
separately.  That isn't always possible, and it isn't always advantageous.
But I suspect a hyphenated word of mixing two concepts.%
\end{tfquot}\index{M!Moore, Charles|)}
\end{interview}
Compare the following two strategies for saying the same thing:
\begin{Code}
ENABLE-LEFT-MOTOR        LEFT MOTOR ON
ENABLE-RIGHT-MOTOR       RIGHT MOTOR ON
DISABLE-LEFT-MOTOR       LEFT MOTOR OFF
DISABLE-RIGHT-MOTOR      RIGHT MOTOR OFF
ENABLE-LEFT-SOLENOID     LEFT SOLENOID ON
ENABLE-RIGHT-SOLENOID    RIGHT SOLENOID ON
DISABLE-LEFT-SOLENOID    LEFT SOLENOID OFF
DISABLE-RIGHT-SOLENOID   RIGHT SOLENOID OFF
\end{Code}
The syntax on the left requires eight dictionary entries; the syntax
on the right requires only six-and some of the words are likely to be
reused in other parts of the application.  If you had a \forth{MIDDLE}
motor and solenoid as well, you'd need only seven words to describe
sixteen combinations.%
\index{H!Hyphenated names|)}

\begin{tip}
Don't bundle numbers into names.
\end{tip}
Watch out for a series of names beginning or ending with numbers, such
as \forth{1CHANNEL}, \forth{2CHANNEL}, \forth{3CHANNEL}, etc.

\index{B!Bundling of names and numbers|(}%
This bundling of names and numbers may be an indication of bad
factoring.  The crime is similar to hyphenation, except that what
should be factored out is a number, not a word.  A better factoring of
the above would be
\begin{Code}
1 CHANNEL
2 CHANNEL
3 CHANNEL
\end{Code}
In this case, the three words were reduced to one.

Often the bundling of names and numbers indicates fuzzy naming.
%Page 168 in first edition.
In the above case, more descriptive names might indicate the purpose
of the channels, as in
\begin{Code}
VOICE , TELEMETRY , GUITAR
\end{Code}
%% Note: AH: the commas be better left out here.
\index{B!Bundling of names and numbers|)}%
We'll amplify on these ideas in the next chapter on ``Factoring.''%
\index{N!Names:!choosing|)}

\section{Naming Standards: The Science}%
\index{N!Naming conventions|(}

\begin{tip}
Learn and adopt \Forth{}'s naming conventions.
\end{tip}
In the quest for short, yet meaningful names, \Forth{} programmers
have adopted certain naming conventions.  \App{E} includes a list of
the most useful conventions developed over the years.

An example of the power of naming conventions is the use of ``dot''
to mean ``print'' or ``display.'' \Forth{} itself uses
\begin{Code}
.   D.   U.R
\end{Code}
for displaying various types of numbers in various formats.  The
convention extends to application words as well.  If you have a
variable called \forth{DATE,} and you want a word that displays the
date, use the name
\begin{Code}
.DATE
\end{Code}
\index{S!Suffixes|(}%
\index{P!Prefixes|(}%
A caution: The overuse of prefixes and suffixes makes words uglier and
ultimately less readable.  Don't try to describe everything a word
does by its name alone.  After all, a name is a symbol, not a
shorthand for code.  Which is more readable and natural sounding?:

\begin{tfquot}
Oedipus complex
\end{tfquot}
(which bears no intrinsic meaning), or
\begin{tfquot}
subconscious-attachment-to-parent-of-opposite-sex complex
\end{tfquot}
Probably the former, even though it assumes you know the play.

\begin{tip}
Use prefixes and suffices to differentiate between like words rather
than to cram details of meaning into the name itself.
\end{tip}
%Page 169 in first edition.
For instance, the phrase
\begin{Code}
... DONE IF CLOSE THEN ...
\end{Code}
is just as readable as
\begin{Code}
... DONE? IF CLOSE THEN ...
\end{Code}
and cleaner as well.  It is therefore preferable, unless we need an
additional word called \forth{DONE} (as a flag, for instance).

\medbreak
A final tip on naming:\index{W!Words:!choosing names|)}
\begin{tip}
Begin all hex numbers with ``0'' (zero) to avoid potential collisions
with names.
\end{tip}
For example, write \forth{0ADD}, not \forth{ADD}.%
\index{S!Suffixes|)}%
\index{P!Prefixes|)}

By the way, don't expect your \Forth{} system to necessarily conform
to the above conventions.  The conventions are meant to be used in
new applications.

\Forth{} was created and refined over many years by people who used it
as a means to an end.  At that time, it was neither reasonable nor
possible to impose naming standards on a tool that was still growing
and evolving.

Had \Forth{} been designed by committee, we would not love it so.%
\index{N!Naming conventions|)}%
\index{I!Implementation!choosing names|)}

\section{More Tips for Readability}

Here are some final suggestions to make your code more readable.
(Definitions appear in \App{C}.)

One constant that pays for itself in most applications is
\forth{BL}\index{B!Blank space (BL)}
(the ASCII value for ``blank-space'').

The word \forthb{ASCII} is used primarily within colon definitions to
free you from having to know the literal value of an ASCII character.
For instance, instead of writing:
\begin{Code}
: (    41 WORD  DROP ;  IMMEDIATE
\end{Code}
where 41 is the ASCII representation for right-parenthesis, you can
write
\begin{Code}
: (    ASCII ) WORD  DROP ;  IMMEDIATE
\end{Code}
\index{T!TRUE|(}%
\index{F!FALSE|(}%
A pair of words that can make dealing with booleans more readable are
\forthb{TRUE} and \forthb{FALSE}.  With these additions you can write
phrases such as

%Page 170 in first edition.
\begin{Code}
TRUE 'STAMP? !
\end{Code}
to set a flag or
\begin{Code}
FALSE 'STAMP? !
\end{Code}
to clear it.

(I once used \forthb{T} and \forthb{F}, but the words are
needed so rarely I now heed the injunction against abbreviations.)

As part of your application (not necessarily part of your \Forth{}
system), you can take this idea a step further and define:
\begin{Code}
: ON   ( a)  TRUE SWAP ! ;
: OFF   ( a)  FALSE SWAP ! ;
\end{Code}
\index{T!TRUE|)}\index{F!FALSE|)}\index{O!ON}\index{O!OFF}%
These words allow you to write:
\begin{Code}
'STAMP? ON
\end{Code}
or
\begin{Code}
'STAMP? OFF
\end{Code}
Other names for these definitions include \forth{SET} and
\forth{RESET}, although \forth{SET} and \forth{RESET} most commonly
use bit masks to manipulate individual bits.

An often-used word is \forthb{WITHIN},\index{W!WITHIN} which determines
whether a given value lies within two other values.  The syntax is:
\begin{Code}
n  lo hi WITHIN
\end{Code}
where ``n'' is the value to be tested and ``lo'' and ``hi'' represent
the range. \forthb{WITHIN} returns true if ``n'' is
\emph{greater-than or equal-to} ``lo'' and \emph{less-than}
``hi.'' This use of the non-inclusive upper limit parallels the syntax
of \forthb{DO }\forthb{LOOP}s.

\person{Moore}\index{M!Moore, Charles} recommends the word \forthb{UNDER+}.
It's useful for adding a value to the number just under the top stack
item, instead of to the top stack item.  It could be implemented in
high level as:
\begin{Code}
: UNDER+  ( a b c -- a+c b )  ROT +  SWAP ;
\end{Code}

\section{Summary}
Maintainability requires readability.  In this chapter we've enumerated
various ways to make a source listing more readable.  We've assumed a
%Page 171 in first edition.
policy of making our code as self-documenting as possible.  Techniques
include listing organization, spacing and indenting, commenting, name
choices, and special words that enhance clarity.

We've mentioned only briefly auxiliary documentation, which includes
all documentation apart from the listing itself.  We won't discuss
auxiliary documentation further in this volume, but it remains an
integral part of the software development process.%
\index{I!Implementation|)}

\begin{references}{9}
\bibitem{stevenson81} \person{Gregory Stevenson}, ``Documentation Priorities,''
\emph{1981 FORML Conference Proceedings,} p. 401.
\bibitem{lee81} \person{Joanne Lee}, ``Quality Assurance in a \Forth{}
Environment,'' (Appendix A), \emph{1981 FORML Proceedings,} p. 363.
\bibitem{dijkstra82} \person{Edsger W. Dijkstra}, \emph{Selected Writings on
Computing: A Personal Perspective,} New York, Springer Verlag, Inc.,
1982.
\bibitem{laxen} \person{Henry Laxen}, ``Choosing Names,'' \emph{\Forth{} Dimensions,}
vol. 4, no.\ 4, \Forth{} Interest Group.
\end{references}

%Page 172 in first edition.

%% Thinking Forth
%% Copyright (C) 2004 Leo Brodie
%% Initial transcription by <Josef Gabriel>
%% 
%% Chapter: <6>

\chapter{Factoring}\Chapmark{6}

In this chapter we'll continue our study of the implementations phase,
this time focusing on factoring.

Decomposition and factoring are chips off the same block. Both involve
dividing and organizing. Decomposition occurs during preliminary
design; factoring occurs during detailed design and implementation.

Since every colon definition reflects decisions of factoring, an
understanding of good factoring technique is perhaps the most
important skill for a \Forth{} programmer.

What is factoring? Factoring means organizing code into useful
fragments. To make a fragment useful, you often must separate reusable
parts from non-reusable parts. The reusable parts become new
definitions. The non-reusable parts become arguments or parameters to the
definitions.

Making this separation is usually referred to as ``factoring out.''
The first part of this chapter will discuss various ``factoring-out''
techniques.

Deciding how much should go into, or stay out of, a definition is
another aspect of factoring. The second section will outline the
criteria for useful factoring.

\section{Factoring Techniques}

\begin{tfquot}
If a module seems almost, but not quite, useful form a second place in
the system, try to identify and isolate the useful subfunction. The
remainder of the module might be incorporated in its original caller
(from ``\emph{Structured Design}'' \cite{stevens74-6}).
\end{tfquot}
The ``useful subfunction'' of course becomes the newly factored
definition.What about the part that ``isn't quite useful''? That
depends on what it is.

\subsection{Factoring Out Data}
\index{D!Data, factoring out|(}
The simplest thing to factor out is data thanks to \Forth{}'s data
stack. For instance, to compute two-thirds of 1,000, we write

\begin{Code}
1000 2 3 */
\end{Code}
To define a word that computes two-thirds of \emph{any} number,
we factor out the argument from the definition:

\begin{Code}
: TWO-THIRDS  ( n1 -- n2)  2 3 */ ;
\end{Code}
When the datum comes in the \emph{middle} of the useful phrase, we
have to use stack manipulation. For instance, to center a piece of
text ten characters long on an 80-column screen, we would write:

\begin{Code}
80  10 -   2/ SPACES
\end{Code}
But text isn't always 10 characters long. To make the phrase useful
for any string, you'd factor out the length by writing:

\begin{Code}
: CENTER (length -- ) 80  SWAP -  2/ SPACES ;
\end{Code}
The data stack can also be used to pass addresses. Therefore what's
factored out may be a \emph{pointer} to data rather than the data
themselves. The data can be numbers or even strings, and still be
factored out through use of the stack.

Sometimes the difference appears to be a function, but you can factor
it out simply as a number on the stack. For instance:

\begin{description}
\item[Segment 1:] \forth{WILLY NILLY\ \ PUDDIN' PIE AND}
\item[Segment 2:] \forth{WILLY NILLY\ \ 8 *\ \ PUDDIN' PIE AND}
\end{description}

\noindent How can you factor out the ``\forth{8 *}'' operation?
By including ``\forth{*}'' in the factoring and passing it a one or eight:

\begin{Code}
: NEW  ( n )  WILLY NILLY  *  PUDDIN' PIE AND ;
\end{Code}

\begin{description}
\item[Segment 1:] \forth{1 NEW}
\item[Segment 2:] \forth{8 NEW}
\end{description}

\noindent 
(Of course if \forth{WILLY NILLY} changes the stack, you'll need to
add appropriate stack-manipulation operators.)

If the operation involves addition, you can nullify it by passing a
zero.

\begin{tip}
For simplicity, try to express the difference between similar
fragments as a numeric difference (values or addresses), rather than
as a procedural difference.
\end{tip}
\index{D!Data, factoring out|)}

\subsection{Factoring Out Functions}
On the other hand, the difference sometimes \emph{is} a function. Witness:

%ARN - this gets messed up Code doesn't indent properly
\begin{description}
\item[Segment 1:]~~
\begin{minipage}[t]{.6\hsize}
\begin{Code}[commandchars=\&\{\}]
BLETCH-A  BLETCH-B   &poorbf{BLETCH-C}
         BLETCH-D  BLETCH-E  BLETCH-F
\end{Code}
\end{minipage}
\item[Segment 2:]~~
\begin{minipage}[t]{.6\hsize}
\begin{Code}[commandchars=\&\{\}]
BLETCH-A  BLETCH-B  &poorbf{PERVERSITY}
         BLETCH-D  BLETCH-E  BLETCH-F
\end{Code}
\end{minipage}
\end{description}

\noindent Wrong approach:

\begin{Code}[commandchars=\&\{\}]
: BLETCHES  ( t=do-BLETCH-C | f=do-PERVERSITY -- ) 
   BLETCH-A  BLETCH-B  IF  &poorbf{BLETCH-C}  ELSE  &poorbf{PERVERSITY}
      THEN  BLETCH-D BLETCH-E BLETCH-F ;
\end{Code}

\begin{description}
\item[Segment 1:]~~ \forth{TRUE BLETCHES}
\item[Segment 2:]~~ \forth{FALSE BLETCHES}
\end{description}

\noindent A better approach:

\begin{Code}
: BLETCH-AB   BLETCH-A BLETCH-B ;
: BLETCH-DEF   BLETCH-D BLETCH-E BLETCH-F ;
\end{Code}

\begin{description}
\item[Segment 1:]~~ \forth{BLETCH-AB ~}\forthb{BLETCH-C}\forth{~ BLETCH-DEF}
\item[Segment 2:]~~ \forth{BLETCH-AB ~}\forthb{PERVERSITY}\forth{~ BLETCH-DEF}
\end{description}

\index{C!Control flags|(}
\index{F!Flags|(}
\begin{tip}
Don't pass control flags downward.
\end{tip}
Why not? First, you are asking your running application to make a
pointless decision-one you knew the answer to while
programming-thereby reducing efficiency. Second, the terminology
doesn't match the conceptual model. What are \forth{TRUE BLETCHES}
as opposed to \forth{FALSE BLETCHES}?
\index{C!Control flags|)}
\index{F!Flags|)}

\subsection{Factoring Out Code from Within Control Structures}
\index{C!Code:!factoring out from within control structures|(}
\index{C!Control structures:!factoring out|(}
\index{C!Control structures:!factoring out code from within|(}

Be alert to repetitions on either side of an \forthb{IF} \forthb{ELSE} \forthb{THEN}
statement. For instance:

\index{E!EMIT|(}
\begin{Code}
... ( c)  DUP  BL 127 WITHIN
       IF  EMIT  ELSE
       DROP  ASCII . EMIT   THEN ...
\end{Code}
This fragment normally emits an ASCII character, but if the character
is a control code, it emits a dot. Either way, an \forthb{EMIT} is
performed. Factor \forthb{EMIT} out of the conditional structure, like
this:

\begin{Code}
... ( c)  DUP  BL 127 WITHIN NOT
       IF  DROP  ASCII .  THEN  EMIT  ...
\end{Code}
\index{E!EMIT|)}
The messiest situation occurs when the difference between two
definitions is a function within a structure that makes it impossible
to factor out the half-fragments. In this case, use stack arguments,
variables, or even vectoring. We'll see how vectoring can be used in a
section of \Chap{7} called ``Using DOER/MAKE.''

Here's a reminder about factoring code from out of a \forthb{DO} \forthb{LOOP}:

\begin{tip}
In factoring out the contents of a \forthb{DO} \forthb{LOOP} into a new
definition, rework the code so that \forth{I} (the index) is not
referenced within the new definition, but rather passed as a stack
argument to it.
\end{tip}
\index{C!Code:!factoring out from within control structures|)}
\index{C!Control structures:!factoring out|)}

\subsection{Factoring Out Control Structures Themselves}

Here are two definitions whose differences lies within a \forthb{IF}
\forthb{THEN} construct:

\begin{Code}
: ACTIVE    A B OR  C AND  IF  TUMBLE JUGGLE JUMP THEN ;
: LAZY      A B OR  C AND  IF   SIT  EAT  SLEEP   THEN ;
\end{Code}
The condition and control structure remain the same; only the event
changes. Since you can't factor the \forthb{IF} into one word and the
\forthb{THEN} into another, the simplest thing is to factor the
condition:

\begin{Code}
: CONDITIONS? ( -- ?) A B OR C AND ;
: ACTIVE    CONDITIONS? IF TUMBLE JUGGLE JUMP THEN ;
: LAZY      CONDITIONS? IF    SIT  EAT  SLEEP THEN ;
\end{Code}

\noindent 
Depending on the number of repetitions of the same condition and
control structure, you may even want to factor out both. Watch this:

\begin{Code}
: CONDITIONALLY   A B OR  C AND NOT IF  R> DROP   THEN ;
: ACTIVE   CONDITIONALLY   TUMBLE JUGGLE JUMP ;
: LAZY   CONDITIONALLY  SIT  EAT  SLEEP ;
\end{Code}
The word \forthb{CONDITIONALLY} may---depending on the condition---alter
the control flow so that the remaining words in each definition will be
skipped. This approach has certain disadvantages as well. We'll
discuss this technique---pros and cons---in \Chap{8}.

More benign examples of factoring-out control structures include case
statements,\index{C!Case statements} which eliminate nested
\forthb{IF} \forthb{ELSE} \forthb{THEN}s, and multiple exit loops (the
\forthb{BEGIN} \forthb{WHILE} \forthb{WHILE} \forthb{WHILE}
\forthb{\dots{}} \forthb{REPEAT} construct). We'll also discuss these
topics in \Chap{8}.  \index{C!Control structures:!factoring out code
from within|)}

\subsection{Factoring Out Names}
It's even good to factor out names, when the names seem almost, but
not quite, the same. Examine the following terrible example of code,
which is meant to initialize three variables associated with each of
eight channels:

\begin{Code}
VARIABLE OSTS       VARIABLE 1STS       VARIABLE 2STS 
VARIABLE 3STS       VARIABLE 4STS       VARIABLE 5STS
VARIABLE 6STS       VARIABLE 7STS       VARIABLE 0TNR
VARIABLE 1TNR       VARIABLE 2TNR       VARIABLE 3TNR
VARIABLE 4TNR       VARIABLE 5TNR       VARIABLE 6TNR
VARIABLE 7TNR       VARIABLE OUPS       VARIABLE 1UPS
VARIABLE 2UPS       VARIABLE 3UPS       VARIABLE 4UPS
VARIABLE 5UPS       VARIABLE 6UPS       VARIABLE 7UPS
\end{Code}

\begin{Code} 
: INIT-CHO   0 0STS !  1000 0TNR !  -1 OUPS ! ; 
: INIT-CHO   0 1STS !  1000 1TNR !  -1 1UPS ! ; 
: INIT-CH2   0 2STS !  1000 2TNR !  -1 2UPS ! ; 
: INIT-CH3   0 3STS !  1000 3TNR !  -1 3UPS ! ; 
: INIT-CH4   0 4STS !  1000 4TNR !  -1 4UPS ! ; 
: INIT-CH5   0 5STS !  1000 5TNR !  -1 5UPS ! ; 
: INIT-CH6   0 6STS !  1000 6TNR !  -1 6UPS ! ; 
: INIT-CH7   0 7STS !  1000 7TNR !  -1 7UPS ! ; 
\end{Code}

\begin{Code} 
: INIT-ALL-CHS    INIT-CHO  INIT-CH1  INIT-CH2  INIT-CH3
   INIT-CH4  INIT-CH5  INIT-CH6  INIT-CH7 ;
\end{Code}
First there's a similarity among the names of the variables; then
there's a similarity in the code used in all the \forth{INIT-CH}
words.

Here's an improved rendition. The similar variable names have been
factored into three data structures, and the lengthy recital of
\forth{INIT-CH} words has been factored into a \forthb{DO} \forthb{LOOP}:

\begin{Code}
: ARRAY  ( #cells -- )  CREATE  2* ALLOT
   DOES> ( i -- 'cell)  SWAP  2* + ; 
8 ARRAY STATUS  ( channel# -- adr)
8 ARRAY TENOR   (        "       )
8 ARRAY UPSHOT  (        "       )
: STABLE   8 0 DO  0 I STATUS !  1000 I TENOR ! 
   -1 I UPSHOT !  LOOP ;
\end{Code}
That's all the code we need.

\index{C!Counted strings|(}
Even in the most innocent cases, a little data structure can eliminate
extra names. By convention \Forth{} handles text in ``counted
strings'' (i.e., with the count in the first byte). Any word that
returns the ``address of a string'' actually returns this beginning
address, where the count is. not only does use of this two-element
data structure eliminate the need for separate names for string and
count, it also makes it easier to move a string in memory, because you
can copy the string \emph{and} the count with a single \forthb{CMOVE}.
\index{C!Counted strings|)}

When you start finding the same awkwardness here and there, you can
combine things and make the awkwardness go away.

\subsection{Factoring Out Functions into Defining Words}

\begin{tip}
If a series of definitions contains identical functions, with
variation only in data, use a defining word.
\end{tip}
Examine the structure of this code (without worrying about its
purpose---you'll see the same example later on):

\begin{Code}
: HUE  ( color -- color') 
   'LIGHT? @  OR  0 'LIGHT? ! ;
: BLACK   0 HUE ;
: BLUE   1 HUE ;
: GREEN   2 HUE ;
: CYAN   3 HUE ;
: RED   4 HUE ;
: MAGENTA   5 HUE ;
: BROWN   6 HUE ;
: GRAY   7 HUE ;
\end{Code}

\noindent The above approach is technically correct, but less
memory-efficient than the following approach using defining words:

\begin{Code}
: HUE   ( color -- )  CREATE ,
   DOES>  ( -- color )  @ 'LIGHT? @  OR  0 'LIGHT? ! ;
 0 HUE BLACK         1 HUE BLUE          2 HUE GREEN
 3 HUE CYAN          4 HUE RED           5 HUE MAGENTA
 6 HUE BROWN         7 HUE GRAY
\end{Code}
(Defining words are explained in \emph{Starting \Forth{}}, Chapter Eleven).

By using a defining word, we save memory because each compiled colon
definition needs the address of \forthb{EXIT} to conclude the
definition.(In defining eight words, the use of a defining word saves
14 bytes on a 16-bit \Forth{}.) Also, in a colon definition each
reference to a numeric literal requires the compilation of
\forthb{LIT} (or \forthb{literal}), another 2 bytes per
definition. (If 1 and 2 are predefined constants, this costs another
10 bytes---24 total.)

In terms of readability, the defining word makes it absolutely clear
that all the colors it defines belong to the same family of words.

The greatest strength of defining words, however, arises when a series
of definitions share the same \emph{compile-time} behavior. This topic
is the subject of a later section, ``Compile-Time Factoring.''

\section{Factoring Criteria}
Armed with an understanding of factoring techniques, let's now discuss
several of the criteria for factoring \Forth{} definitions. They
include:

\begin{enumerate}
\item Limiting the size of definitions
\item Linting repetition of code
\item Nameability
\item Information hiding
\item Simplifying the command interface
\end{enumerate}

\index{D!Defining words:!length|(}
\begin{tip}
Keep definitions short.
\end{tip}

\begin{interview}
We asked \person{Moore},\index{M!Moore, Charles|(}
``How long should a \Forth{} definition be?''

\begin{tfquot}
A word should be a line long. That's the target.

When you have a whole lot of words that are all useful in their own
right---perhaps in debugging or exploring, but inevitably there's a
reason for their existence---you feel you've extracted the essence of
the problem and that those words have expressed it.

Short words give you a good feeling.
\end{tfquot}\index{M!Moore, Charles|)}
\end{interview}
An informal examination of one of \person{Moore}'s applications shows
that he averages seven references, including both words and numbers,
per definition. These are remarkably short definitions. (Actually his
code was divided about 50--50 between one-line and two-line
definitions.)

Psychological tests have shown that the human mind can only focus its
conscious attention on seven things, plus or minus two, at a time
\cite{miller56}. Yet all the while, day and night, the vast amounts of
data, making connections and associations and solving problems.

Even if out subconscious mind knows each part of an application inside
out, our narrow-viewed conscious mind can only correlate seven
elements of it at once. Beyond that, our grasp wavers. Short
definitions match our mental capabilities.

Something that many \Forth{} programmers to write overly long
definitions is the knowledge that headers take space in the
dictionary. The coarser the factoring, the fewer the names, and the
less memory that will be wasted.

It's true that more memory will be used, but it's hard to say that
anything that helps you test, debug and interact with your code is a
``waste.'' If your application is large, try using a default width of
three, with the ability to switch to a full-length name to avoid a
specific collision. (``Width'' refers to a limit on the number of
characters stored in the name field of each dictionary header.)

If the application is still too big, switch to a \Forth{} with
multiple dictionaries on a machine with extended memory, or better
yet, a 32-bit \Forth{} on a machine with 32-bit addressing.

A related fear is that over-factoring will decrease performance due to
the overhead of \Forth{}'s inner interpreter. Again, it's true that
there is some penalty for each level of nesting. But ordinarily the
penalty for extra nesting due to proper factoring will not be
noticeable. If you timings are that tight, the real solution is to
translate something into assembler.

\begin{tip}
Factor at the point where you feel unsure about your code (where
complexity approaches the conscious limit).
\end{tip}
Don't let your ego take over with an ``I can lick this!'' attitude.
\Forth{} code should never feel uncomfortably complex. Factor!

\begin{interview}
\person{Moore}:\index{M!Moore, Charles|(}

\begin{tfquot}
Feeling like you might have introduced a bug is one reason for
factoring. Any time you see a doubly-nested \forthb{DO} \forthb{LOOP}, that's a
sign that something's wrong because it will be hard to debug. Almost
always take the inner \forthb{DO} \forthb{LOOP} and make a word.

And having factored out a word for testing, there's no reason for
putting it back. You found it useful in the first place. There's no
guarantee you won't need it again.
\end{tfquot}\index{M!Moore, Charles|)}
\end{interview}
Here's another facet of the same principle:

\begin{tip}
Factor at the point where a comment seems necessary
\end{tip}
Particularly if you feel a need to remind yourself what's on the
stack, this may be a good time to ``make a break.''

\goodbreak
Suppose you have
\begin{Code}
... BALANCE  DUP xxx xxx xxx xxx xxx xxx xxx xxx xxx
     xxx xxx xxx xxx xxx xxx   ( balance) SHOW  ...
\end{Code}
which begins by computing the balance and ends by displaying it. In
the meantime, several lines of code use the balance for purposes of
their own. Since it's difficult to see that the balance is still on
the stack when \forth{SHOW} executes, the programmer has interjected a
stack picture.

This solution is generally a sign of bad factoring. Better to write:
\begin{Code}
: REVISE  ( balance -- )  xxx xxx xxx xxx xxx xxx xxx
     xxx xxx xxx xxx xxx xxx xxx ;
... BALANCE  DUP REVISE  SHOW  ...
\end{Code}
No narrative stack pictures are needed. Furthermore, the programmer
now has a reusable, testable subset of the definition.

\begin{tip}
Limit repetition of code.
\end{tip}
\index{C!Code:!repetition of|(}%
The second reason for factoring, to eliminate repeated fragments of
code, is even more important than reducing the size of definitions.
\index{D!Defining words:!length|)}

\begin{interview}
\person{Moore}:\index{M!Moore, Charles|(}

\begin{tfquot}
When a word is just a piece of something, it's useful for clarity or
debugging, but not nearly as good as a word that is used multiple
times. Any time a word is used only once you want to question its
value.

Many times when a program has gotten too big I will go back through it
looking for phrases that strike my eye as candidates for factoring.
The computer can't do this; there are too many variables.
\end{tfquot}\index{M!Moore, Charles|)}
\end{interview}
In looking over your work, you often find identical phrases or short
passages duplicated several times. In writing an editor I found this
phrase repeated several times:

\begin{Code}
FRAME  CURSOR @ +
\end{Code}
Because it appeared several times I factored it into a new word called
\forth{AT}.

It's up to you to recognize fragments that are coded differently but
functionally equivalent, such as:

\begin{Code}
FRAME  CURSOR @ 1-  +
\end{Code}
The \forth{1-} appears to make this phrase different from the one defined as
\forth{AT}. But in fact, it can be written

\begin{Code}
AT 1-
\end{Code}
On the other hand:

\begin{tip}
When factoring out duplicate code, make sure the factored code serves
a single purpose.
\end{tip}
Don't blindly seize upon duplications that may not be useful. For
instance, in several places in one application I used this phrase:

\begin{Code}
BLK @ BLOCK  >IN @ +  C@
\end{Code}
I turned it into a new word and called it \forth{LETTER}, since it
returned the letter being pointed to by the interpreter.

In a later revision, I unexpectedly had to write:

\begin{Code}
BLK @ BLOCK  >IN @ +  C!
\end{Code}
I could have used the existing \forth{LETTER} were it not for its
\forth{C@} at the end. Rather than duplicated the bulk of the phrase
in the new section, I chose to refactor \forth{LETTER} to a finer
resolution, taking out the \forth{C@}.  The usage was then either
\forth{LETTER C@} or \forth{LETTER C!}. This change required me to
search through the listing changing all instances of \forth{LETTER} to
\forth{LETTER C@}.  But I should have done that in the first place,
separating the computation of the letter's address from the operation
to be performed on the address.
\index{C!Code:!repetition of|)}

Similar to our injunction against repetition of code:

\begin{tip}
Look for repetition of patterns.
\end{tip}
If you find yourself referring back in the program to copy the pattern
of previously-used words, then you may have mixed in a general idea
with a specific application. The part of the pattern you are copying
perhaps can be factored out as an independent definition that can be
used in all the similar cases.

\begin{tip}
Be sure you can name what you factor.
\end{tip}

\medbreak
\begin{interview}
\person{Moore}:\index{M!Moore, Charles|(}
\begin{tfquot}
If you have a concept that you can't assign a single name to , not a
hyphenated name, but a name, it's not a well-formed concept. The
ability to assign a name is a necessary part of decomposition.
Certainly you get more confidence in the idea.
\end{tfquot}\index{M!Moore, Charles|)}
\end{interview}
Compare this view with the criteria for decomposing a module espoused
by structured design in \Chap{1}. According to that method, a module
should exhibit ``functional binding,'' Which can be verified by
describing its function in a single, non-compound, \emph{sentence}.
\Forth{}'s ``atom,'' a \emph{name}, is an order of magnitude more
refined.

\begin{tip}
Factor definitions to hide details that may change.
\end{tip}
We've seen the value of information hiding in earlier chapters,
especially with regard to preliminary design. It's useful to remember
this criterion during the implementation stage as well.

Here's a very short definition that does little except hide information:

\begin{Code}
: >BODY  ( acf -- apf )  2+ ;
\end{Code}
This definition allows you to convert an acf (address of code field) to
an apf (address of parameter field) without depending on the actual
structure of a dictionary definition. If you were to use \forthb{2+} instead of
the word \forthb{>BODY}, you would lose transportability if you ever
converted to a \Forth{} system in which the heads were separated from
the bodies. (This is one of a set of words suggested by Kim Harris,
and included as an Experimental Proposal in the \Forth{}-83 Standard
\cite{harris83}.)

Here's a group of definitions that might be used in writing an editor:

\begin{Code}
: FRAME  ( -- a)  SCR @ BLOCK ;
: CURSOR  ( -- a)  R# ;
: AT  ( -- a)  FRAME  CURSOR @ + ;
\end{Code}
These three definitions can form the basis for all calculations of
addresses necessary for moving text around. Use of these three
definitions completely separates your editing algorithms from a
reliance on \Forth{} blocks.

What good is that? If you should decide, during development, to create
an editing buffer to protect the user from making errors that destroy
a block, you merely have to redefine two of these words, perhaps like
this:

\begin{Code}
CREATE FRAME  1024 ALLOT
VARIABLE CURSOR
\end{Code}
The rest of your code can remain intact. 

\begin{tip}
Factor calculations algorithms out of definitions that display results.
\end{tip}
This is really a question of decomposition.

Here's an example. The word defined below, pronounced
``people-to-paths,'' computes how many paths of communication there are
between a given number of people in a group. (This is a good thing for
managers of programmer teams to know---the number of communication
paths increases drastically with each new addition to the team.)
\begin{Code}
: PEOPLE>PATHS  ( #people -- #paths )  DUP 1-  *  2/ ;
\end{Code}
This definition does the calculation only. Here's the ``user definition''
that invokes \forth{PEOPLE>PATHS} to perform the calculation, and then
displays the result:
\begin{Code}
: PEOPLE  ( #people)
    ." = "  PEOPLE>PATHS  .  ." PATHS " ;
\end{Code}
This produces:
%% I don't know how to underline these.
\begin{Code}[commandchars=\&\{\}]
2 PEOPLE&underline{ = 1 PATHS}
3 PEOPLE&underline{ = 3 PATHS}
5 PEOPLE&underline{ = 10 PATHS}
10 PEOPLE&underline{ = 45 PATHS}
\end{Code}
Even if you think you're going to perform a particular calculation
only once, to display it in a certain way, believe me, you're wrong.
You will have to come back later an factor out the calculation part.
Perhaps you'll need to display the information in a right-justified
column, or perhaps you'll want to record the results in a data
base---you never know. But you'll always have to factor it, so you
might as well do it right the first time. (The few times you might get
away with it aren't worth the trouble.)

The word \forth{.} (dot) is a prime example. Dot is great 99\% of the
time, but occasionally it does too much. Here's what it does, in fact
(in \Forth{}--83):

\begin{Code}
: .   ( n )  DUP ABS 0 <# #S  ROT SIGN  #> TYPE SPACE ;
\end{Code}
But suppose you want to convert a number on the stack into an ASCII
string and store it in a buffer for typing later. Dot converts it, but
also types it. Or suppose you want to format playing cards in the form
\forth{10C} (for ``ten of clubs''). You can't use dot to display the 10
because it prints a final space.

Here's a better factoring found in some \Forth{} systems:

\begin{Code}
: (.)  (n -- a #)  DUP ABS 0  <# #S  ROT SIGN  #> ;
: .  ( n)  (.) TYPE SPACE ;
\end{Code}
We find another example of failing to factor the output function from
the calculation function in our own Roman numeral example in \Chap{4}.
Given our solution, we can't store a Roman numeral i8n a buffer or
even center it in a field. (A better approach would have been to use
\forthb{HOLD} instead of \forthb{EMIT}.)

Information hiding can also be a reason \emph{not} to factor. For
instance, if you factor the phrase

\begin{Code}
SCR @ BLOCK
\end{Code}
into the definition

\begin{Code}
: FRAME   SCR @ BLOCK ;
\end{Code}
remember you are doing so only because you may want to change the
location of the editing frame. Don't blindly replace all occurrences
of the phrase with the new word \forth{FRAME,} because you may change the
definition of \forth{FRAME} and there will certainly be times when you really
want \forthb{SCR} \forthb{@} \forthb{BLOCK}.

\begin{tip}
If a repeated code fragment is likely to change in some cases but not
others, factor out only those instances that might change. If the
fragment is likely to change in more than one definition.
\end{tip}
Knowing when to hide information requires intuition and experience.
Having made many design changes in your career, you'll learn the hard
way which things will be most likely to change in the future.

You can never predict everything, though. It would be useless to try,
as we'll see in the upcoming section called ``The Iterative Approach
in Implementation.''

\index{C!Commands, reducing number of|(}
\begin{tip}
Simplify the command interface by reducing the number of commands.
\end{tip}
It may seem paradoxical, but good factoring can often yield
\emph{fewer} names. In \Chap{5} we saw how six simple names
(\forth{LEFT}, \forth{RIGHT}, \forth{MOTOR}, \forth{SOLENOID},
\forth{ON}, and \forth{OFF}) could do the work of eight
badly-factored, hyphenated names.

As another example, I found two definitions circulating in one
department in which \Forth{} had recently introduced. Their purpose
was purely instructional, to remind the programmer which vocabulary
was \forth{CURRENT}, and which was \forth{CONTEXT}:

\begin{Code}
: .CONTEXT   CONTEXT @  8 -  NFA  ID.   ;
: .CURRENT   CURRENT @  8 -  NFA  ID.  ;
\end{Code}
\goodbreak
\noindent If you typed

\begin{Code}
.CONTEXT
\end{Code}
the system would respond

\begin{Code}[commandchars=\&\{\}]
.CONTEXT&underline{ FORTH}
\end{Code}
(They worked---at least on the system used there---by backing up to the
name field of the vocabulary definition, and displaying it.)

The obvious repetition of code struck my eye as a sign of bad
factoring. It would have been possible to consolidate the repeated
passage into a third definition:

\begin{Code}
: .VOCABULARY   ( pointer )  @  8 -  NFA  ID. ;
\end{Code}
shortening the original definitions to:

\begin{Code}
: .CONTEXT   CONTEXT .VOCABULARY ;
: .CURRENT   CURRENT .VOCABULARY ;
\end{Code}
But in this approach, the only difference between the two definitions
was the pointer to be displayed. Since part of good factoring is to
make fewer, not more definitions, it seemed logical to have only one
definition, and let it take as an argument either the word
\forth{CONTEXT} or the word \forth{CURRENT}.

\medbreak
Applying the principles of good naming, I suggested:

\begin{Code}
: IS  ( adr)   @  8 -  NFA  ID. ;
\end{Code}
allowing the syntax

\begin{Code}[commandchars=\&\{\}]
CONTEXT IS&underline{ ASSEMBLER ok}
\end{Code}
or

\begin{Code}[commandchars=\&\{\}]
CURRENT IS&underline{ FORTH ok}
\end{Code}
The initial clue was repetition of code, but the final result came
from attempting to simplify the command interface.

Here's another example. The IBM PC has four modes four displaying text
only:

\begin{quote}\sf
40 column monochrome

40 column color

80 column monochrome

80 column color
\end{quote}
The word \forth{MODE} is available in the \Forth{} system I use.
\forth{MODE} takes an argument between 0 and 3 and changes the mode
accordingly. Of course, the phrase \forth{0 MODE} or \forth{1 MODE}
doesn't help me remember which mode is which.

Since I need to switch between these modes in doing presentations, I
need to have a convenient set of words to effect the change. These
words must also set a variable that contains the current number of
columns---40 or 80.

Here's the most straightforward way to fulfill the requirements:

\begin{Code}
: 40-B&W       40 #COLUMNS !  0 MODE ;
: 40-COLOR     40 #COLUMNS !  1 MODE ;
: 80-B&W       80 #COLUMNS !  2 MODE ;
: 80-COLOR     80 #COLUMNS !  3 MODE ;
\end{Code}
By factoring to eliminate the repetition, we come up with this version:

\begin{Code}
: COL-MODE!     ( #columns mode )  MODE  #COLUMNS ! ;
: 40-B&W       40 0 COL-MODE! ;
: 40-COLOR     40 1 COL-MODE! ;
: 80-B&W       80 2 COL-MODE! ;
: 80-COLOR     80 3 COL-MODE! ;
\end{Code}
But by attempting to reduce the number of commands, and also by
following the injunctions against numerically-prefixed and hyphenated
names, we realize that we can use the number of columns as a stack
argument, and \emph{calculate} the mode:

\begin{Code}
: B&W    ( #cols -- )  DUP #COLUMNS !  20 /  2-     MODE ;
: COLOR  ( #cols -- )  DUP #COLUMNS !  20 /  2-  1+ MODE ;
\end{Code}
This gives us this syntax:

\begin{Code}
40 B&W
80 B&W
40 COLOR
80 COLOR
\end{Code}
We've reduced the number of commands from four to two.

Once again, though, we have some duplicate code. If we factor out this
code we get:

\begin{Code}
: COL-MODE!  ( #columns chroma?)
   SWAP DUP #COLUMNS !  20 / 2-  +  MODE ;
: B&W    ( #columns -- )  0 COL-MODE! ;
: COLOR  ( #columns -- )  1 COL-MODE! ;
\end{Code}
Now we've achieved a nicer syntax, and at the same time greatly
reduced the size of the object code. With only two commands, as in
this example, the benefits may be marginal. But with larger sets of
commands the benefits increase geometrical.

Our final example is a set of words to represent colors on a
particular system. Names like \forth{BLUE} and \forth{RED} are nicer
than numbers. One solution might be to define:

\begin{Code}
 0 CONSTANT BLACK                 1 CONSTANT BLUE
 2 CONSTANT GREEN                 3 CONSTANT CYAN
 4 CONSTANT RED                   5 CONSTANT MAGENTA
 6 CONSTANT BROWN                 7 CONSTANT GRAY
 8 CONSTANT DARK-GRAY             9 CONSTANT LIGHT-BLUE
10 CONSTANT LIGHT-GREEN          11 CONSTANT LIGHT-CYAN
12 CONSTANT LIGHT-RED            13 CONSTANT LIGHT-MAGENTA
14 CONSTANT YELLOW               15 CONSTANT WHITE
\end{Code}
These colors can be used with words such as \forth{BACKGROUND},
\forth{FOREGROUND}, and \forth{BORDER}:

\begin{Code}
WHITE BACKGROUND  RED FOREGROUND  BLUE BORDER
\end{Code}
But this solution requires 16 names, and many of them are hyphenated.
Is there a way to simplify this?

We notice that the colors between 8 and 15 are all ``lighter''
versions of the colors between 0 and 7. (In the hardware, the only
difference between these two sets is the setting of the ``intensity
bit.'') If we factor out the ``lightness,'' we might come up with this
solution:

\begin{Code}
VARIABLE 'LIGHT?  ( intensity bit?)
: HUE  ( color)  CREATE ,
   DOES>  ( -- color )  @  'LIGHT? @  OR  0 'LIGHT? ! ;
 0 HUE BLACK         1 HUE BLUE           2 HUE GREEN
 3 HUE CYAN          4 HUE RED            5 HUE MAGENTA
 6 HUE BROWN         7 HUE GRAY
: LIGHT   8 'LIGHT? ! ;
\end{Code}
With this syntax, the word

\begin{Code}
BLUE
\end{Code}
by itself will return a ``1'' on the stack, but the phrase

\begin{Code}
LIGHT BLUE
\end{Code}
will return a ``9.'' (The adjective \forth{LIGHT} sets flag which is
used by the hues, then cleared.)

\goodbreak
If necessary for readability, we still might want to define:

\begin{Code}
8 HUE DARK-GRAY
14 HUE YELLOW
\end{Code}
Again, through this approach we've achieved a more pleasant syntax and
shorter object code.
\index{C!Commands, reducing number of|)}

\begin{tip}
Don't factor for the sake of factoring. Use clich�s.
\index{C!Cliches|(}
\end{tip}
The phrase

\begin{Code}
OVER + SWAP
\end{Code}
may seem commonly in certain applications. (I converts an address and
count into an ending address and starting address appropriate for a
\forthb{DO LOOP}.)\index{D!DO LOOP}

Another commonly seen phrase is

\begin{Code}
1+ SWAP
\end{Code}
(It rearranges a first-number and last-number into the
last-number-plus-one and first-number order required by \forthb{DO}.)

It's a little tempting to seize upon these phrases and turn them into
words, such as (for the first phrase) \forth{RANGE}.

\begin{interview}
\person{Moore}:\index{M!Moore, Charles|(}

\begin{tfquot}
That particular phrase [\forthb{OVER + SWAP}] is one that's right on
the margin of being a useful word. Often, though if you define
something as a word, it turns out you use it only once. If you name
such a phrase, you have trouble knowing exactly what \forth{RANGE}
does. You can't see the manipulation in your mind. \forthb{OVER +
SWAP} has greater mnemonic value than \forth{RANGE}.
\end{tfquot}
\index{M!Moore, Charles|)}
\end{interview}
%% Could not get that special e in cliches.
I call these phrases ``clich�s.'' They stick together as meaningful
functions. You don't have to remember how the phrase works, just what
it does. And you don't have to remember an extra name.
\index{C!Cliches|)}

\section{Compile-Time Factoring}
\index{C!Compile-time factoring|(}

\noindent
In the last section we looked at many techniques for organizing code
and data to reduce redundancy.

We can also apply limited redundancy during compilation, by letting
\Forth{} do some of out dirty work.

\begin{tip}
For maximum maintainability, limit redundancy even at compile time.
\end{tip}
Suppose in our application we must draw nine boxes as shown in
\Fig{fig6-1}.

\begin{figure*}[hhhh]
\caption{What we're supposed to display}
\labelfig{fig6-1}
\begin{center}
\begin{BVerbatim}[baselinestretch=0.7]
********     ********     ********
********     ********     ********
********     ********     ********
********     ********     ********
********     ********     ********


********     ********     ********
********     ********     ********
********     ********     ********
********     ********     ********
********     ********     ********


********     ********     ********
********     ********     ********
********     ********     ********
********     ********     ********
********     ********     ********
\end{BVerbatim}
\end{center}
\end{figure*}

In our design we need to have constants that represent values such as
the size of each box, the size of the gap between boxes, and the
left-most and top-most coordinates of the first box.

Naturally we can define:

\begin{Code}
8 CONSTANT WIDE
5 CONSTANT HIGH
4 CONSTANT AVE
2 CONSTANT STREET
\end{Code}
(Streets run east and west; avenues run north and south.)

%% I changed "We want" to "we want"
Now, to define the left margin, we might compute it mentally, We want
to center all these boxes on a screen 80 columns wide. To center
something, we subtract its width from 80 and divide by two to
determine the left margin. To figure the total width of all the boxes,
we add
\begin{eqnarray}
8 + 4 + 8 + 4 + 8 = 32\nonumber
\end{eqnarray}
(three widths and two avenues). $(80-31) / 2 = 24$.

\goodbreak
So we could, crudely, define:

\begin{Code}
24 CONSTANT LEFTMARGIN
\end{Code}
and use the same approach for \forth{TOPMARGIN}.

But what if we should later redesign the pattern, so that the width
changed, or perhaps the gap between the boxes? We'd have to recompute
the left margin ourselves.

In the \Forth{} environment, we can use the full power of \Forth{}
even when we're compiling. Why not let \Forth{} do the figuring?

\begin{Code}
WIDE 3 *  AVE 2 *  +  80 SWAP -  2/ CONSTANT LEFTMARGIN
HIGH 3 *  STREET 2 * +  24 SWAP -  2/ CONSTANT TOPMARGIN
\end{Code}

\begin{tip}
If a constant's value depends on the value of an earlier constant, use
\Forth{} to calculate the value of the second.
\end{tip}
None of these computations are performed when the application is
running, so run-time speed is not affected.

Here's another example. \Fig{fig6-2} shows the code for a word that
draws shapes. The word \forth{DRAW} emits a star at every x--y
coordinate listed in the table called \forth{POINTS}. (Note: the word
\forth{XY} positions the cursor to the ( x y ) coordinate on the
stack.)

Notice the line immediately following the list of points:

\begin{Code}
HERE POINTS -  ( /table)  2/  CONSTANT #POINTS
\end{Code}
%% I don't know how to offset the Figure 6-2.
\begin{figure*}[bbbb]
\caption{Another example of limiting compile-time redundancy.}
\labelfig{fig6-2}
\begin{center}
\begin{BVerbatim}
: P  ( x y -- )  C, C, ;
CREATE POINTS
   10 10 P     10 11 P     10 12 P     10 13 P     10 14 P
   11 10 P     12 10 P     13 10 P     14 10 P
   11 12 P     12 12 P     13 12 P     14 12 P
HERE POINTS -  ( /table)  2/  CONSTANT #POINTS
: @POINTS  ( i -- x y)  2* POINTS + DUP 1+ C@  SWAP C@ ;
: DRAW  #POINTS 0 DO  I @POINTS  XY  ASCII * EMIT  LOOP ;
\end{BVerbatim}
\end{center}
\end{figure*}
The phrase ``\forth{HERE POINTS -}'' computes the number of x--y
coordinates in the table: this value becomes the constant
\forth{\#POINTS}, used as the limit in \forth{DRAW}'s \forthb{DO} \forthb{LOOP}.
\index{D!DO LOOP}

This construct lets you add or subtract points from the table without
worrying about the number of points there are. \Forth{} computes this
for you.

\subsection{Compile-Time Factoring through Defining Words}%
\index{W!Words:!defining|see{Defining words}}%
\index{D!Defining words:!compile-time factoring through}
Let's examine a series of approaches to the same problem---defining a
group of related addresses. Here's the first try:

\begin{Code}
HEX 01A0 CONSTANT BASE.PORT.ADDRESS
BASE.PORT.ADDRESS CONSTANT SPEAKER
BASE.PORT.ADDRESS 2+ CONSTANT FLIPPER-A
BASE.PORT.ADDRESS 4 + CONSTANT FLIPPER-B
BASE.PORT.ADDRESS 6 + CONSTANT WIN-LIGHT
DECIMAL
\end{Code}
The idea is right, but the implementation is ugly. The only elements
that change from port to port are the numeric offset and the name of
the port being defined; everything else repeats. This repetition
suggests the use of a defining word.

The following approach, which is more readable, combines all the
repeated code into the ``does'' part  of a defining word:

\begin{Code}
: PORT  ( offset -- )  CREATE ,
   \ does>  ( -- 'port) @ BASE.PORT.ADDRESS + ;
0 PORT SPEAKER
2 PORT FLIPPER-A
4 PORT FLIPPER-B
6 PORT WIN-LIGHT
\end{Code}
In this solution we're performing the offset calculation at
\emph{run}-time, every time we invoke one of these names. It would be
more efficient to perform the calculation at compile time, like this:

\begin{Code}
: PORT  ( offset -- )  BASE.PORT.ADDRESS + CONSTANT ;
   \ does>  ( -- 'port)
0 PORT SPEAKER
2 PORT FLIPPER-A
4 PORT FLIPPER-B
6 PORT WIN-LIGHT
\end{Code}
Here we've created a defining word, \forth{PORT}, that has a unique
\emph{compile}-time behavior, namely adding the offset to
\forth{BASE.PORT.ADDRESS} and defining a \forth{CONSTANT}.

We might even go one step further. Suppose that all port addresses are
two bytes apart. In this case there's no reason we should have to
specify these offsets. The numeric sequence
\begin{quote}
0 2 4 6
\end{quote}
is itself redundant.

In the following version, we begin with the \forth{BASE.PORT.ADDRESS}
on the stack. The defining word \forth{PORT} duplicates this address,
makes a constant out of it, then adds 2 to the address still on the
stack, for the next invocation of \forth{PORT}.

\begin{Code}
: PORT   ( 'port -- 'next-port)  DUP CREATE ,  2+ ;
   \ does>  ( -- 'port)
BASE.PORT.ADDRESS
  PORT SPEAKER
  PORT FLIPPER-A
  PORT FLIPPER-B
  PORT WIN-LIGHT
DROP ( port.address)
\end{Code}
Notice we must supply the initial port address on the stack before
defining the first port, then invoke \forthb{DROP} when we've finished
defining all the ports to get rid of the port address that's still on
the stack.

One final comment. The base-port address is very likely to change, and
therefor should be defined in only one place. This does \emph{not}
mean it has to be defined as a constant. Provided that the base-port
address won't be used outside of this lexicon of port names, it's just
as well to refer to it by number here.

\begin{Code}
HEX 01A0  ( base port adr)
  PORT SPEAKER
  PORT FLIPPER-A
  PORT FLIPPER-B
  PORT WIN-LIGHT
DROP
\end{Code}
\index{C!Compile-time factoring|)}

\section{The Iterative Approach in Implementation}
\index{I!Iterative approach|(}
Earlier in the book we discussed the iterative approach, paying
particular attention to its impact on the design phase. Now that we're
talking about implementation, let's see how the approach is actually
used in writing code.

\begin{tip}
Work on only one aspect of a problem at a time.
\end{tip}
Suppose we're entrusted with the job of coding a word to draw or erase
a box at a given x--y coordinate. (This is the same problem we
introduced in the section called ``Compile-Time Factoring.'')

At first we focus our attention on the problem of drawing a
box---never mind erasing it. We might come up with this:

\begin{Code}
: LAYER   WIDE  0 DO  ASCII * EMIT  LOOP ;
: BOX   ( upper-left-x  upper-left-y -- )
   HIGH  0 DO  2DUP  I +  XY LAYER  LOOP  2DROP ;
\end{Code}
Having tested this to make sure it works correctly, we turn now to the
problem of using the same code to \emph{un}draw a box. The solution is
simple: instead of hard-coding the \forthb{ASCII *} we'd like to change
the emitted character from an asterisk to a blank. This requires the
addition of a variable, and some readable words for setting the
contents of the variable. So:

\begin{Code}
VARIABLE INK
: DRAW   ASCII *  INK ! ;
: UNDRAW   BL  INK ! ;
: LAYER   WIDTH  0 DO  INK @  EMIT  LOOP ;
\end{Code}
The definition of \forth{BOX}, along with the remainder of the application,
remains the same.

This approach allows the syntax

\begin{Code}
( x y ) DRAW BOX
\end{Code}
or

\begin{Code}
( x y ) UNDRAW BOX
\end{Code}
By switching from an explicit value to a variable that contains a
value, we've added a level of indirection. In this case, we've added
indirection ``backwards,'' adding a new level of complexity to the
definition of \forth{LAYER} without substantially lengthening the definition.

By concentrating on one dimension of the problem at a time, you can
solve each dimension more efficiently. If there's an error in your
thinking, the problem will be easier to see if it's not obscured by yet
another untried, untested aspect of your code.

\begin{tip}
Don't change too much at once.
\end{tip}
While you're editing your application---adding a new feature or fixing
something---it's often tempting to go and fix several other things at
the same time. Our advice: Don't.

Make as few changes as you can each time you edit-compile. Be sure to
test the results of each revision before going on. You'd be amazed how
often you can make three innocent modifications, only to recompile and
have nothing work!

Making changes one at a time ensures that when it stops working, you
know why.

\begin{tip}
Don't try to anticipate ways to factor too early.
\end{tip}%
\index{A!Arrays|(}
Some people wonder why most \Forth{} systems don't include the
definition word \forth{ARRAY}. This rule is the reason.
\begin{interview}
\person{Moore}:\index{M!Moore, Charles|(}
\begin{tfquot}
I often have a class of things called arrays. The simplest array
merely adds a subscript to an address and gives you back an
address. You can define an array by saying
\begin{Code}
CREATE X   100 ALLOT
\end{Code}
then saying
\begin{Code}
X +
\end{Code}
Or you can say
\begin{Code}
: X   X + ;
\end{Code}
One of the problems that's most frustrating for me is knowing whether
it's worth creating a defining word for a particular data structure.
Will I have enough instances to justify it?

I rarely know in advance if I'm going to have more than one array. So
I don't define the word \forth{ARRAY}.

After I discover I need two arrays, the question is marginal.

If I need three then it's clear. Unless they're different. And odds
are they will be different. You may want it to fetch it for you. You
may want a byte array, or a bit array. You may want to do bounds
checking, or store its current length so you can add things to the
end.

I grit my teeth and say, ``Should I make the byte array into a cell
array, just to fit the data structure into the word I already have
available?''

The more complex the problem, the less likely it will be that you'll
find a universally applicable data structure. The number of instances
in which a truly complex data structure has found universal use is
very small. One example of a successful complex data structure is the
\Forth{} dictionary. Very firm structure, great versatility. It's used
everywhere in \Forth{}. But that's rare.

If you choose to define the word \forth{ARRAY}, you've done a
decomposition step. You've factored out the concept of an array from
all the words you'll later back in. And you've gone to another level
of abstraction.

Building levels of abstraction is a dynamic process, not one you can
predict.
\end{tfquot}\index{M!Moore, Charles|)}
\end{interview}%
\index{A!Arrays|)}
\begin{tip}
Today, make it work. Tomorrow, optimize it.
\end{tip}
\begin{interview}
Again \person{Moore}.\index{M!Moore, Charles|(} On the day of this
interview, \person{Moore} had been completing work on the design of a
board-level \Forth{} computer, using commercially available ICs. As
part of his toolkit for designing the board, he created a simulator in
\Forth{}, to test the board's logic:

\begin{tfquot}
This morning I realized I've been mixing the descriptions of the chips
with the placement of the chips on the board. This perfectly
convenient for my purposes at the moment, but when I come up with
another board that I want to use the same chips for, I have arranged
things very badly.

I should have factored it with the descriptions here and the uses
there. I would then have had a chip description language. Okay. At the
time I was doing this I was not interested in that level of
optimization.

Even if the thought had occurred to me then, I probably would have
said, ``All right, I'll do that later,'' then gone right ahead with
what I was doing. Optimization wasn't the most important thing to me
at the time.

Of course I try to factor things well. But if there doesn't seem to be
a good way to do something, I say, ``Let's just make it work.''

My motivation isn't laziness, it's knowing that there are other things
coming down the pike that are going to affect this decision in ways I
can't predict. Trying to optimize this now is foolish. Until I get the
whole picture in front of me, I can't know what the optimum is.
\end{tfquot}\index{M!Moore, Charles|)}
\end{interview}
The observations in this section shouldn't contradict what's been said
before about information hiding and about anticipating elements that
may change. A good programmer continually tries to balance the expense
of building-in changeability against the expense of changing things
later if necessary.

These decisions take experience. But as a general rule:

\begin{tip}
Anticipate things-that-may-change by organizing information, not by
adding complexity. Add complexity only as necessary to make the
current iteration work.
\end{tip}
\index{I!Iterative approach|)}

\subsection{Summary}
In this chapter we've discussed various techniques and criteria for
factoring. We also examined how the iterative approach applies to the
implementation phase.

\begin{references}{9}
\bibitem{stevens74-6} \person{W.P. Stevens}, \person{G.J. Myers}, and \person{L.L. Constantine},
\emph{ IBM Systems Journal}, vol. 13, no. 2, 1974, Copyright 1974 by
International Business Machines Corporation.
\bibitem{miller56} \person{G.A. Miller}, ``The Magical Number Seven, Plus or
Minus Two: Some Limits on our Capacity for Processing Information,''
\emph{Psychol. Rev}., vol. 63, pp. 81-97, Mar. 1956.
\bibitem{harris83} \person{Kim R. Harris}, ``Definition Field Address
Conversion Operators,'' \emph{\Forth{}--83 Standard}, \Forth{} Standards
Team.
\end{references}


%% Thinking Forth
%% Copyright (C) 2004 Leo Brodie
%% Initial transcription by Nils M Holm
%% Based on OCR scans by Steve Fisher
%% 
%% Chapter: Seven - Handling Data: Stacks and States

%!! horizontal line
%SEVEN
%!! horizontal line

\chapter{Handling~Data: Stacks~and~States}

%!! horizontal line
FORTH handles data in one of two ways: either on the stack or in data
structures. When to use which approach and how to manage both the
stack and data structures are the topics of this chapter.

\section{The Stylish Stack}

The simplest way for FORTH words to pass arguments to each other is
via the stack. The process is ``simple'' because all the work of pushing
and popping values to and from the stack is implicit.

%!! horizontal line
\bigskip\blackline{2ex}
\noindent Moore:

\begin{tfquot}
The data stack uses this idea of ``hidden information.'' The arguments
being passed between subroutines are not explicit in the calling sequence.
The same argument might ripple through a whole lot of words quite invisibly,
even below the level of awareness of the programmer, simply because it
doesn't have to be referred to explicitly.
\end{tfquot}
%!! horizontal line
\blackline{1ex}

One important result of this approach: Arguments are unnamed. They
reside on the stack, not in named variables. This effect is one of the
reasons for FORTH's elegance. At the same time it's one of the reasons
badly written FORTH code can be unreadable. Let's explore this
paradox.

The invention of the stack is analogous to that of pronouns in
English. Consider the passage:

%!! indent paragraph
\begin{tfquot}
Take this gift, wrap it in tissue paper and put it in a box.
\end{tfquot}
Notice the word ``gift'' is mentioned only once. The gift is referred to
henceforth as ``it.''

The informality of the ``it'' construct makes English more readable
(provided the reference is unambiguous). So with the stack, the implicit
passing of arguments makes code more readable. We emphasize the
\emph{processes}, not the \emph{passing of arguments} to the processes.

Our analogy to pronouns suggests why bad FORTH can be so unreadable.
The spoken language gets confusing when too many things are
referred to with pronouns.

%!! indent paragraph
\begin{tfquot}
Take off the wrapping and open the box. Remove the gift and throw it
away.
\end{tfquot}
The problem with this passage is that we're using ``it'' to refer to too
many things at once. There are two solutions to this error. The easiest
solution is to supply a real name instead of ``it'':

%!! indent paragraph
\begin{tfquot}
Remove the wrapping and open the box. Take out the gift and throw
\emph{the box} away.
\end{tfquot}
Or we can introduce the words ``former'' and ``latter.'' But the best
solution is to redesign the passage:

%!! indent paragraph
\begin{tfquot}
Remove the wrapping and open the present. Throw away the box.
\end{tfquot}
So in FORTH we have analogous observations:

\begin{tip}
Simplify code by using the stack. But don't stack too deeply within any
single definition. Redesign, or, as a last resort, use a named variable.
\end{tip}
Some newcomers to FORTH view the stack the way a gymnast views a
trampoline: as a fun place to bounce around on. But the stack is meant for
data-passing, not acrobatics.

So how deep is ``too deep?'' Generally, three elements on the stack is
the most you can manage within a single definition. (In double-length
arithmetic, each ``element'' occupies two stack positions but is logically
treated as a single element by operators such as \textbf{2DUP}, \textbf{20VER},
etc.)

In your ordinary lexicon of stack operators, \textbf{ROT} is the only one
that gives you access to the third stack item. Aside from \textbf{PICK} and
\textbf{ROLL} (which we'll comment on soon), there's no easy way to get at
anything below that.

To stretch our analogy to the limit, perhaps three elements on the
stack corresponds to the three English pronouns ``this,'' ``that,'' and
``th'other.''

\subsection{Redesign}

Let's witness a case where a wrong-headed approach leads to a messy
stack problem. Suppose we're trying to write the definition of +THRU
(see Chapter Five, ``Listing Organization'' section, ``Relative Loading''
subsection). We've decided that our loop body will be

\begin{Code}
...  DO  I LOAD  LOOP ;
\end{Code}
that is, we'll put \textbf{LOAD} in a loop, then arrange for the index and limit
to correspond to the absolute screens being loaded.

On the stack initially we have:

\begin{Code}
lo hi
\end{Code}
where ``lo'' and ``hi'' are the \emph{offsets} from \textbf{BLK}.

We need to permute them for \textbf{DO}, like this:

\begin{Code}
hi+1+blk lo+blk
\end{Code}
Our biggest problem is adding the value of \textbf{BLK} to both offsets.

We've already taken a wrong turn but we don't know it yet. So let's
proceed. We try:

\begin{Code}
lo hi
               BLK @
lo hi blk
               SWAP
lo blk hi
               OVER
lo blk hi blk
               +
lo blk hi+blk
               1+
lo blk hi+blk+1
               ROT ROT
hi+blk+1 lo blk
               +
hi+blk+1 lo+blk
\end{Code}
We made it, but what a mess!

If we're gluttons for punishment, we might make two more stabs at
it arriving at:

\begin{Code}
BLK @  DUP ROT + 1+  ROT ROT +
\end{Code}
and

\begin{Code}
BLK @  ROT OVER +  ROT ROT + 1+  SWAP
\end{Code}
All three sequences do the same thing, but the code seems to be getting
blurrier, not better.

With experience we learn to recognize the combination ROT ROT as
a danger sign: the stack is too crowded. Without having to work out the
alternates, we recognize the problem: once we make two copies of ``blk,''
we have four elements on the stack.

At this point, the first resort is usually the return stack:

\begin{Code}
BLK @  DUP >R  + 1+  SWAP R> + +
\end{Code}
(See ``The Stylish Return Stack,'' coming up next.) Here we've \textbf{DUP}ed
``blk,'' saving one copy on the return stack and adding the other copy to
``hi.''

Admittedly an improvement. But readable?

Next we think, ``Maybe we need a named variable.'' Of course, we
have one already: \textbf{BLK}. So we try:

\begin{Code}
BLK @  + 1+  SWAP BLK @ +
\end{Code}
Now it's more readable, but it's still rather long, and redundant too.
\textbf{BLK @ +} appears twice.

``\textbf{BLK @ +}''? That sounds famililar. Finally our neurons connect.

We look back at the source for +LOAD just defined:

\begin{Code}
: +LOAD  ( offset -- )  BLK @ +  LOAD ;
\end{Code}
This word, +LOAD, should be doing the work. All we have to write is:

\begin{Code}
: +THRU  ( lo hi )  1+ SWAP  DO  I +LOAD  LOOP ;
\end{Code}
We haven't created a more efficient version here, because the work of
\textbf{BLK @ +} will be done on every pass of the loop. But we have created a
cleaner, conceptually simpler, and more readable piece of code. In this
case, the inefficiency is unnoticeable because it only occurs as each block
is loaded.

Redesigning, or rethinking the problem, was the path we should
have taken as soon as things got ugly.

\subsection{Local Variables}

Most of the time problems can be arranged so that only a few arguments
are needed on the stack at any one time. Occasionally, however, there's
nothing you can do.

Here's an example of a worst case. Assume you have a word called
LINE which draws a line between any two points, specified as coordinates
in this order:

\begin{Code}
( x1 y1 x2 y2)
\end{Code}
where x1,y1 represent the x,y coordinates for the one end-point, and
x2,y2 represent the opposite end-point.

Now you have to write a box-drawing word called [BOX] which
takes four arguments in this order:

\begin{Code}
( x1 y1 x2 y2)
\end{Code}
where x1 y1 represent the x,y coordinates for the upper left-hand corner
of the box, and x2 y2 represent the lower right-hand corner coordinates.
Not only do you have four elements on the stack, they each have to
be referred to more than once as you draw lines from point to point.

Although we're using the stack to get the four arguments, the algorithm
for drawing a box doesn't lend itself to the nature of the stack. If you're
in a hurry, it would probably be best to take the easy way out:

\begin{Code}
VARIABLE TOP         ( y coordinates top of box)
VARIABLE LEFT        ( x     "       left side)
VARIABLE BOTTOM      ( y     "       bottom)
VARIABLE RIGHT       ( x     "       right side)
: [BOX]   ( x1 y1 x2 y2)   BOTTOM !  RIGHT !  TOP !  LEFT !
   LEFT @ TOP @  RIGHT @ TOP @  LINE
   RIGHT @ TOP @  RIGHT @ BOTTOM @  LINE
   RIGHT @ BOTTOM @  LEFT @ BOTTOM @  LINE
   LEFT @ BOTTOM @  LEFT @ TOP @  LINE ;
\end{Code}
What we've done is create four named variables, one for each coordinate.
The first thing [BOX] does is fill these variables with the arguments from
the stack. Then the four lines are drawn, referencing the variables.
Variables such as these that are used only within a definition (or in some
cases, within a lexicon) are called ``local variables.''

I've been guilty many times of playing hotshot, trying to do as
much as possible on the stack rather than define a local variable. There
are three reasons to avoid this cockiness.

First, it's a pain to code that way. Second, the result is unreadable.
Third, all your work becomes useless when a design change becomes
necessary, and the order of two arguments changes on the stack. The
\textbf{DUP}s, \textbf{OVER}s and \textbf{ROT}s weren't really solving the problem,
just jockeying things into position.

With this third reason in mind, I recommend the following:

\begin{tip}
Especially in the design phase, keep on the stack only the arguments you're
using immediately. Create local variables for any others. (If necessary,
eliminate the variables during the optimization phase.)
\end{tip}
Fourth, if the definition is extremely time-critical, those tricky stack
manipulators, (e.g., \textbf{ROT ROT}) can really eat up clock cycles. Direct
access to variables is faster.

If it's \emph{really} time-critical, you may need to convert to assembler
anyway. In this case, all your stack problems fly out the door, because all
your data will be referenced either in registers or indirectly through
registers. Luckily, the definitions with the messiest stack arguments are
often the ones written in code. Our [BOX] primitive is a case in point.
\textbf{CMOVE>} is another.

The approach we took with [BOX] certainly beats spending half an
hour juggling items on the stack, but it is by no means the best solution.
What's nasty about it is the expense of creating four named variables,
headers and all, solely for use within this one routine.

(If you're target compiling an application that will not require
headers in the dictionary, the only loss will be the 8 bytes in RAM for the
variables. In FORTH systems of the future, headers may be separated
into other pages of memory anyway; again the loss will be only 8 bytes.)
Let me repeat: This example represents a worst-case situation, and
occurs rarely in most FORTH applications. If words are well-factored,
then each word is designed to do very little. Words that do little generally
require few arguments.

In this case, though, we are dealing with two points each represented
by two coordinates.

Can we change the design? First, LINE may be \emph{too} primitive a
primitive. It requires four arguments because it can draw lines between
any two points, diagonally, if necessary.

In drawing our box, we may only need perfectly vertical and horizontal
lines. In this case we can write the more powerful, but less specific,
words VERTICAL and HORIZONTAL to draw these lines. Each requires
only \emph{three} arguments: the starting position's x and y and the
length. This factoring of function simplifies the definition of [BOX].

Or we might discover that this syntax feels more natural to the
user:

\begin{Code}
10 10 ORIGIN! 30 30 BOX
\end{Code}
where ORIGIN! sets a two-element pointer to the ``origin,'' the place
where the box will start (the upper left-hand corner). Then ``30 30 BOX''
draws a box 30 units high and 30 units wide, relative to the origin.

This approach reduces the number of stack arguments to BOX as
part of the design.

\begin{tip}
When determining which arguments to handle via data structures rather
than via the stack, choose the arguments that are the more permanent or
that represent a current state.
\end{tip}
\subsection{On PICK and ROLL}

Some folks like the words \textbf{PICK} and \textbf{ROLL}. They use these words to
access elements from any level on the stack. We don't recommend them.
For one thing, \textbf{PICK} and \textbf{ROLL} encourage the programmer to think
of the stack as an array, which it is not. If you have so many elements on
the stack that you need \textbf{PICK} and \textbf{ROLL}, those elements should be in
an array instead.

Second, they encourage the programmer to refer to arguments that
have been left on the stack by higher-level, calling definitions without
being explicitly \emph{passed} as arguments. This makes the definition
dependent on other definitions. That's unstructured---and dangerous.

Finally, the position of an element on the stack depends on what's
above it, and the number of things above it can change constantly. For
instance, if you have an address at the fourth stack position down, you can
write

\begin{Code}
4 PICK @
\end{Code}
to fetch its contents. But you must write

\begin{Code}
( n) 5 PICK !
\end{Code}
because with ``n'' on the stack, the address is now in the fifth position.
Code like this is hard to read and harder to modify.

\subsection{Make Stack Drawings}

When you do have a cumbersome stack situation to solve, it's best to work
it out with paper and pencil. Some people even make up forms, such as the
one in \Fig{fig7-1}. Done formally like this (instead of on the back of your
phone bill) stack commentaries serve as nice auxiliary documentation.

\subsection{Stack Tips}

\begin{tip}
Make sure that stack effects balance out under all possible control flows.
\end{tip}
In the stack commentary for \textbf{CMOVE>} in \Fig{fig7-1}, the inner brace
represents the contents of the \textbf{DO LOOP}. The stack depth upon exiting
the loop is the same as upon entering it: one element. Within the outer
braces, the stack result of the \textbf{IF} clause is the same as that of the
\textbf{ELSE} clause: one element left over. (What that leftover element
represents doesn't matter, as symbolized by the ``x'' next to \textbf{THEN}.)

\wepsfiga{fig7-1}{Example of a stack commentary.}

%!! include \Fig{fig7-1} here

\begin{tip}
When doing two things with the same number, perform the function that
will go underneath first.
\end{tip}
For example:

\begin{Code}
: COUNT  ( a -- a+1 # )  DUP C@  SWAP 1+  SWAP ;
\end{Code}
(where you first get the count) is more efficiently written:

\begin{Code}
: COUNT  ( a -- a+1 # )  DUP 1+  SWAP C@ ;
\end{Code}
(where you first compute the address).

\begin{tip}
Where possible, keep the number of return arguments the same in all
possible cases.
\end{tip}
You'll often find a definition which does some job and, if something goes
wrong, returns an error-code indentifying the problem. Here's one way
the stack interface might be designed:

\begin{Code}
( -- error-code f | -- t)
\end{Code}
If the flag is true, the operation was successful. If the flag is false, it was
unsuccessful and there's another value on the stack to indicate the nature
of the error.

You'll find stack manipulation easier, though, if you redesign the interface
to look like this:

\begin{Code}
( -- error-code | O=no-error)
\end{Code}
One value serves both as a flag and (in case of an error) the error code.
Note that reverse-logic is used; non-zero indicates an error. You can use
any values for the error codes except zero.

\section{The Stylish Return Stack}

What about this use of the return stack to hold temporary arguments? Is
it good style or what?

Some people take great offense to its use. But the return stack
offers the simplest solution to certain gnarly stack jams. Witness the
definition of \textbf{CMOVE>} in the previous section.

If you decide to use the return stack for this purpose, remember
that you are using a component of FORTH for a purpose other than that
intended. (See the section called ``Sharing Components,'' later in this
chapter.)

Here's some suggestions to keep you from shooting yourself in the
foot:

\begin{tip}
%!! there's certainly some better way to do ordered lists in TeX
\begin{enumerate}
\item Keep return stack operators symmetrical.
\item Keep return stack operators symmetrical under all control flow
   conditions.
\item In factoring definitions, watch out that one part doesn't contain one
   return stack operator, and the other its counterpart.
\item If used inside a \textbf{DO LOOP}, return stack operators must by symmetrical
   within the loop, and \textbf{I} is no longer valid in code bounded by
   \textbf{>R} and \textbf{R>}.
\end{enumerate}
\end{tip}
For every \textbf{>R} there must be a \textbf{R>} in the same definition. Sometimes
the operators will appear to be symmetrical, but due to the control structure
they aren't. For instance:

\begin{Code}
... BEGIN ... >R ... WHILE ... R> ... REPEAT
\end{Code}
If this construction is used in the outer loop of your application,
everything will run fine until you exit (perhaps hours later) when you'll
suddenly blow up. The problem? The last time through the loop, the
resolving \textbf{R>} has been skipped.

\section{The Problem With Variables}

Although we handle data of immediate interest on the stack, we depend
on much information tucked away in variables, ready for recurring access.
A piece of code can change the contents of a variable without
necessarily having to know anything about how that data will be used,
who will use it, or when and if it will be used. Another piece of code can
fetch the contents of a variable and use it without knowing where that
value came from.

For every word that pushes a value onto the stack, another word
must consume that value. The stack gives us point-to-point communication,
like the post office.

Variables, on the other hand, can be set by any command and accessed
any number of times---or not at all---by any command. Variables
are available for anyone who cares to look---like graffiti.

Thus variables can be used to reflect the current state of affairs.

Using currentness can simplify problems. In the Roman numeral example
of Chapter Four, we used the variable COLUMN\# to represent the
current decimal-place; the words ONER, FIVER, and TENER depended
on this information to determine which type of symbol to display. We
didn't have to specify both descriptions every time, as in TENS ONER,
TENS FIVER, etc.

On the other hand, currentness adds a new level of complexity. To
make something current we must first define a variable or some type of
data structure. We also must remember to initialize it, if there's any
chance that part of our code will refer to it before another part has had a
chance to set it.

A more serious problem with variables is that they are not ``reentrant.''
On a multi-tasked FORTH system, each task which requires
local variables must have its own copies. FORTH's \textbf{USER} variables
serve this purpose. (See \emph{Starting FORTH}, Chapter Nine, ``FORTH
Geography.'')

Even within a single task, a definition that refers to a variable is
harder to test, verify, and reuse in a different situation than one in which
arguments are passed via the stack.

Suppose we are implementing a word-processor editor. We need a
routine that calculates the number of characters between the current cursor
position and the previous carriage-return/line-feed sequence. So we
write a word that employs a \textbf{DO LOOP} starting at the current position
(CURSOR @) and ending at the zeroth position, searching for the line
feed character.

Once the loop has found the character sequence, we subtract its
relative address from our current cursor position

\begin{Code}
its-position CURSOR @  SWAP -
\end{Code}
to determine the distance between them.

Our word's stack effect is:

\begin{Code}
( -- distance-to-previous-cr/lf)
\end{Code}
But in later coding we find we need a similar word to compute the
distance from an arbitrary character---\emph{not} the current cursor
position---to the first previous line-feed character. We end up factoring
out the ``CURSOR @'' and allowing the starting address to be passed as an
argument on the stack, resulting in:

\begin{Code}
( starting-position -- distance-to-previous-cr/lf)
\end{Code}
By factoring-out the reference to the variable, we made the definition
more useful.

\begin{tip}
Unless it involves cluttering up the stack to the point of unreadability, try
to pass arguments via the stack rather than pulling them out of variables.
\end{tip}
%!! horizontal line
\bigskip
\blackline{2ex}
\noindent Kogge:
\begin{tfquot}
Most of the modularity of FORTH comes from designing and treating
FORTH words as ``functions'' in the mathematical sense. In my experience
a FORTH programmer usually tries quite hard to avoid defining any but
the most essential global variables (I have a friend who has the sign ``Help
stamp out variables'' above his desk), and tries to write words with what is
called ``referential transparency,'' i.e., given the same stack inputs a word
will always give the same stack ouputs regardless of the more global
context in which it is executed.

In fact this property is exactly what we use when we test words in isolation.
Words that do not have this property are significantly harder to test. In a
sense a ``named variable'' whose value changes frequently is the next worst
thing to the now ``forbidden'' GOTO.
\end{tfquot}
%!! horizontal line
\blackline{1ex}

%!! include cartoon of page 211 here
\wepsfigp{img7-211}{``Shot from a cannon on a fast-moving train,
hurtling between the blades of a windmill, and expecting to grab a
trapeze dangling from a hot-air balloon\dots{} I told you Ace, there were
too many variables!''}

Earlier we suggested the use of local variables especially during the
design phase, to eliminate stack traffic. It's important to note that in
doing so, the variables were referred to only within the one definition. In
our example, [BOX] receives four arguments from the stack and immediately
loads them into local variables for its own use. The four variables are not
referred to outside of this definition, and the word behaves safely as a
function.

Programmers unaccustomed to a language in which data can be
passed implicitly don't always utilize the stack as fully as they should.
Michael Ham suggests the reason may be that beginning FORTH users
don't trust the stack [1]. He admits to initially feeling safer about storing
values into variables than leaving them on the stack. ``No telling \emph{what}
might happen with all that thrashing about on the stack,'' he felt.

It took some time for him to appreciate that ``if words keep properly
to themselves, using the stack only for their expected input and output
and cleaning up after themselves, they can be looked upon as sealed
systems \dots{} I could put the count on the stack at the beginning of the
loop, go through the complete routine for each group, and at the end the
count would emerge, back on top of the stack, not a hair out of place.''

\section{Local and Global Variables/Initialization}

As we saw earlier, a variable that is used exclusively within a single
definition (or single lexicon), hidden from other code, is called a local
variable. A variable used by more than one lexicon is called a global
variable. As we've seen in an earlier chapter, a set of global variables that
collectively describe a common interface between several lexicons is
called an ``interface lexicon.''

FORTH makes no distinction between local and global variables.
But FORTH programmers do.

%!! horizontal line
\bigskip\blackline{2ex}
\noindent Moore:

\begin{tfquot}
We should be writing for the reader. If something is referred to only locally,
a temporary variable just for accumulating a sum in, we should define it
locally. It's handier to define it in the block where it's used, where you can
see its comment.

If it's used globally, we should collect things according to their logical
function, and define them together on a separate screen. One per line with a
comment.

The question is, where do you initialize them? Some say on the same line,
immediately following its definition. But that messes up the comments,
and there isn't room for any decent comment. And it scatters the
initialization all over the application.

I tend to do all my initialization in the load screen. After I've loaded all my
blocks, I initialize the things that have to be initialized. It might also set
up color lookup tables or execute some initialization code.

If your program is destined to be target compiled, then it's easy to write a
word at the point that encompasses all the initialization.

It can get much more elaborate. I've defined variables in ROM where the
variables were all off in an array in high memory, and the initial values are
in ROM, and I copy up the initial values at initialization time. But usually
you're only initializing a few variables to anything other than zero.
\end{tfquot}
%!! horizontal line
\blackline{1ex}

\section{Saving and Restoring a State}

Variables have the characteristic that when you change their contents,
you clobber the value that was there before. Let's look at some of the
problems this can create, and some of the things we can do about them.

\textbf{BASE} is a variable that indicates the current number radix for all
numeric input and output. The following words are commonly found in
FORTH systems:

\begin{Code}
: DECIMAL   10 BASE ! ;
: HEX   16 BASE ! ;
\end{Code}
Suppose we've written a word that displays a ``dump'' of memory. Ordinarily,
we work in decimal mode, but we want the dump in hexadecimal. So we write:

\begin{Code}
: DUMP  ( a # )
   HEX   ...   ( code for the dump) ... DECIMAL ;
\end{Code}
This works---most of the time. But there's a presumption that we want to
come back to decimal mode. What if it had been working in hexadecimal,
and wants to come back to hexadecimal? Before we change the base to
\textbf{HEX}, we have to save its current value. When we're done dumping, we
restore it.

This means we have to tuck away the saved value temporarily,
while we format the dump. The return stack is one place to do this:

\begin{Code}
: DUMP  ( a # )
   BASE @ >R  HEX   ( code for dump)  R> BASE ! ;
\end{Code}
If things get too messy, we may have to define a temporary variable:

\begin{Code}
VARIABLE OLD-BASE
: DUMP  ( a # )
   BASE @  OLD-BASE !  HEX ( code for dump )
   OLD-BASE @  BASE ! ;
\end{Code}
How quickly things get complicated.

In this situation, if both the current and the old version of a variable
belong only to your application (and not part of your system), and if this
same situation comes up more than once, apply a technique of factoring:

\begin{Code}
: BURY  ( a)  DUP 2+  2 CMOVE ;
: EXHUME  ( a)  DUP 2+  SWAP 2 CMOVE ;
\end{Code}
Then instead of defining two variables, such as CONDITION and
OLD-CONDITION, define one double-length variable:

\begin{Code}
2VARIABLE CONDITION
\end{Code}
Use BURY and EXHUME to save and restore the original value:

\begin{Code}
: DIDDLE    CONDITION BURY  17 CONDITION !  ( diddle )
   CONDITION EXHUME ;
\end{Code}
BURY saves the ``old'' version of condition at CONDITION 2+.

You still have to be careful. Going back to our \textbf{DUMP} example, suppose
you decided to add the friendly feature of letting the user exit the
dump at any time by pressing the ``escape'' key. So inside the loop you
build the test for a key being pressed, and if so execute \textbf{QUIT}. But what
happens?

The user starts in decimal, then types \textbf{DUMP}. He exits \textbf{DUMP} midway
through and finds himself, strangely, in hexadecimal.

In the simple case at hand, the best solution is to not use \textbf{QUIT}, but
rather a controlled exit from the loop (via \textbf{LEAVE}, etc.) to the end of the
definition where \textbf{BASE} is reset.

In very complex applications a controlled exit is often impractical,
yet many variables must somehow be restored to a natural condition.

%!! horizontal line
\bigskip
\blackline{2ex}
\noindent Moore responds to this example:

\begin{tfquot}
You really get tied up in a knot. You're creating problems for yourself. If I
want a hex dump I say \textbf{HEX DUMP}. If I want a decimal dump I say
\textbf{DECIMAL DUMP}. I don't give \textbf{DUMP} the privilege of messing around
with my environment.

There's a philosophical choice between restoring a situation when you
finish and establishing the situation when you start. For a long time I felt
you should restore the situation when you're finished. And I would try to
do that consistently everywhere. But it's hard to define ``everywhere.'' So
now I tend to establish the state before I start.

If I have a word which cares where things are, it had better set them. If
somebody else changes them, they don't have to worry about resetting
them.

There are more exits than there are entrances.
\end{tfquot}
%!! horizontal line
\blackline{1ex}

In cases in which I need to do the resetting before I'm done, I've found it
useful to have a single word (which I call PRISTINE) to perform this
resetting. I invoke PRISTINE:

%!! '*'s should be bullets
\begin{itemize}
\item at the normal exit point of the application
\item at the point where the user may deliberately exit (just before QUIT)
\item at any point where a fatal error may occur, causing an abort.
\end{itemize}

Finally, when you encounter this situation of having to save/restore a
value, make sure it's not just a case of bad factoring. For example,
suppose we have written:

\begin{Code}
: LONG   18 #HOLES ! ;
: SHORT   9 #HOLES ! ;
: GAME   #HOLES @  O DO  I HOLE PLAY  LOOP ;
\end{Code}
The current GAME is either LONG or SHORT.

Later we decide we need a word to play \emph{any} number of holes. So we
invoke GAME making sure not to clobber the current value of \#HOLES:

\begin{Code}
: HOLES  ( n)  #HOLES @  SWAP #HOLES !  GAME  #HOLES ! ;
\end{Code}
Because we needed HOLES after we'd defined GAME, it seemed to be of
greater complexity; we built HOLES around GAME. But in fact---perhaps
you see it already---rethinking is in order:

\begin{Code}
: HOLES ( n)  O DO  I HOLE PLAY  LOOP ;
: GAME   #HOLES @ HOLES ;
\end{Code}
We can build GAME around HOLES and avoid all this saving/restoring
nonsense.

\section{Application Stacks}

In the last section we examined some ways to save and restore a single
previous value. Some applications require \emph{several} values to be saved and
restored. You may often find the best solution to this problem in defining
your own stack.

Here is the code for a user stack including very simple error checking
(an error clears the stack):

\begin{Code}
CREATE STACK  12 ALLOT  \  ( 2tos-pointer | 10stack [5 cells] )
HERE CONSTANT STACK>
: INIT-STACK   STACK STACK ! ;   INIT-STACK
: ?BAD  ( ?)   IF ." STACK ERROR "  INIT-STACK  ABORT  THEN ;
: PUSH  ( n)   2 STACK +!  STACK @  DUP  STACK> = ?BAD  ! ;
: POP  ( -- n)  STACK @ @  -2 STACK +!  STACK @ STACK < ?BAD ;
\end{Code}
The word PUSH takes a value from off of your data stack and ``pushes''
it onto this new stack. POP is the opposite, ``popping'' a value from off
the new stack, and onto FORTH's data stack.

In a real application you might want to change the names PUSH
and POP to better match their conceptual purposes.

\othersidetrue\section{Sharing Components}
\othersidefalse
\begin{tip}
It's legal to use a component for an additional purpose besides its intended
one, provided:

%!! there's certainly some better way to do ordered lists in TeX
\begin{enumerate}
\item All uses of the component are mutually exclusive
\item Each interrupting use of the component restores the component to
   its previous state when finished.
\end{enumerate}

Otherwise you need an additional component or level of complexity.
\end{tip}
We've seen a simple example of this principle with the return stack. The
return stack is a component of the FORTH system designed to hold
return addresses, and thereby serve as an indication of where you've been
and where you're going. To use the return stack as a holder for temporary
values is possible, and in many cases desirable. Problems occur when one
of the above restrictions is ignored.

In my text formatter the output can go invisible. This feature has
two purposes: (1) for looking ahead to see whether something will fit, and
(2) for formatting the table of contents (the entire document is formatted
and page numbers are calculated without anything actually being
displayed).

It was tempting to think that once having added the ability to make
the output invisible, I could use this feature to serve both purposes.
Unfortunately, the two purposes are not mutually exclusive.

Let's see what would happen if I tried to violate this rule. Imagine
that the word DISPLAY does the output, and it's smart enough to know
whether to be visible or invisible. The words VISIBLE and INVISIBLE
set the state respectively.

My code for looking ahead will first execute INVISIBLE, then
test-format the upcoming text to determine its length, and finally execute
VISIBLE to restore things to the normal state.

This works fine.

Later I add the table-of-contents feature. First the code executes
INVISIBLE, then runs through the document determining page
numbers etc.; then finally executes VISIBLE to restore things to normal.

The catch? Suppose I'm running a table of contents and I hit one of
those places where I look ahead. When I finish looking ahead, I execute
VISIBLE. Suddenly I start printing the document when I was supposed
to be running the table of contents.

The solution? There are several.

One solution views the problem as being that the lookahead code is
clobbering the visible/invisible flag, which may have been preset by
table-of-contents. Therefore, the lookahead code should be responsible for
saving, and later restoring, the flag.

Another solution involves keeping two separate variables---one to
indicate we're looking ahead, the other to indicate we're printing the
table of contents. The word DISPLAY requires that both flags be false in
order to actually display anything.

There are two ways to accomplish the latter approach, depending on
how you want to decompose the problem. First, we could nest one condition
within the other:

\begin{Code}
: [DISPLAY]  ...
( the original definition, always does the output) ... ;
VARIABLE 'LOOKAHEAD?  ( t=looking-ahead)
: <DISPLAY>   'LOOKAHEAD? @ NOT IF  [DISPLAY]  THEN ;
VARIABLE 'TOC?  ( t=setting-table-of-contents)
: DISPLAY   'TOC? @ NOT IF  <DISPLAY>  THEN ;
\end{Code}
DISPLAY checks that we're not setting the table of contents and invokes
<DISPLAY>, which in turn checks that we're not looking ahead
and invokes [DISPLAY].

In the development cycle, the word [DISPLAY] that always does
the output was originally called DISPLAY. Then a new DISPLAY was
defined to include the lookahead check, and the original definition was
renamed [DISPLAY], thus adding a level of complexity backward
without changing any of the code that used DISPLAY.

Finally, when the talbe-of-contents feature was added, a new DISPLAY
was defined to include the table-of-contents check, and the
previous DISPLAY was renamed <DISPLAY>.

That's one approach to the use of two variables. Another is to include
both tests within a single word:

\begin{Code}
: DISPLAY   'LOOKAHEAD? @  'TOC @ OR  NOT IF [DISPLAY] THEN ;
\end{Code}
But in this particular case, yet another approach can simplify the whole
mess. We can use a single variable not as a flag, but as a counter.

We define:

\begin{Code}
VARIABLE 'INVISIBLE?  ( t=invisible)
: DISPLAY   'INVISIBLE? @  O= IF [DISPLAY] THEN ;
: INVISIBLE   1 'INVISIBLE? +! ;
: VISIBLE    -1 'INVISIBLE? +! ;
\end{Code}
The lookahead code begins by invoking INVISIBLE which bumps the
counter up one. Non-zero is ``true,'' so DISPLAY will not do the output.
After the lookahead, the code invokes VISIBLE which decrements the
counter back to zero (``false'').

The table-of-contents code also begins with VISIBLE and ends with
INVISIBLE. If we're running the table of contents while we come upon a
lookahead, the second invocation of VISIBLE raises the counter to two.

The subsequent invocation of INVISIBLE decrements the counter to
one, so we're still invisible, and will remain invisible until the table of
contents has been run.

(Note that we must substitute \textbf{0=} for \textbf{NOT}. The '83 Standard has
changed \textbf{NOT} to mean one's complement, so that \textbf{1 NOT} yields true. By
the way, I think this was a mistake.)

This use of a counter may be dangerous, however. It requires parity
of command usage: two VISIBLEs yields invisible. That is, unless VISIBLE
clips the counter:

\begin{Code}
: VISIBLE   'INVISIBLE? @  1-  O MAX  'INVISIBLE? ! ;
\end{Code}
\section{The State Table}

A single variable can express a single condition, either a flag, a value, or
the address of a function.

A collection of conditions together represent the \emph{state} of the
application or of a particular component [2]. Some applications require the
ability to save a current state, then later restore it, or perhaps to have a
number of alternating states.

\begin{tip}
When the application requires handling a group of conditions
simultaneously, use a state table, not separate variables.
\end{tip}
\begin{figure*}[hhhh]
\caption{A collection of related variables.}
\labelfig{fig7-2}
\begin{center}
\begin{BVerbatim}
VARIABLE TOP
VARIABLE BOTTOM
VARIABLE LEFT
VARIABLE RIGHT
VARIABLE INSIDE
VARIABLE OUT
\end{BVerbatim}
\end{center}
\end{figure*}

The simple case requires saving and restoring a state. Suppose we initially
have six variables representing the state of a particular component, as
shown in \Fig{fig7-2}.

Now suppose that we need to save all of them, so that further processing
can take place, and later restore all of them. We could define:

\begin{Code}
: @STATE ( -- top bottom left right inside out)
   TOP @  BOTTOM @  LEFT @  RIGHT @  INSIDE @  OUT @ ;
: !STATE ( top bottom left right inside out -- )
   OUT !  INSIDE !  RIGHT !  LEFT !  BOTTOM !  TOP ! ;
\end{Code}
thereby saving all the values on the stack until it's time to restore them.
Or, we might define alternate variables for each of the variables above, in
which to save each value separately.

But a preferred technique involves creating a table, with each
element of the table referred to by name. Then creating a second table of
the same length. As you can see in \Fig{fig7-3}, we can save the state by
copying the table, called POINTERS, into the second table, called
SAVED.

\wepsfiga{fig7-3}{Conceptual model for saving a state table.}

%!! include \Fig{fig7-3} here

We've implemented this approach with the code in \Fig{fig7-4}.

\begin{figure*}[tttt]
\caption{Implementation of save/restorable state table.}
\labelfig{fig7-4}
\begin{center}
\begin{BVerbatim}
O CONSTANT POINTERS  \ address of state table PATCHED LATER
: POSITION   ( o -- o+2 ) CREATE DUP ,  2+
   DOES>  ( -- a )  @  POINTERS + ;
O  \ initial offset
POSITION TOP
POSITION BOTTOM
POSITION LEFT
POSITION RIGHT
POSITION INSIDE
POSITION OUT
CONSTANT /POINTERS   \ final computed offset

HERE ' POINTERS >BODY !  /POINTERS ALLOT  \ real table
CREATE SAVED  /POINTERS ALLOT  \ saving place
: SAVE     POINTERS  SAVED  /POINTERS CMOVE ;
: RESTORE  SAVED  POINTERS  /POINTERS CMOVE ;
\end{BVerbatim}
\end{center}
\end{figure*}

Notice in this implementation that the names of the pointers, TOP,
BOTTOM, etc., always return the same address. There is only one location
used to represent the current value of any state at any time.

Also notice that we define POINTERS (the name of the table) with
CONSTANT, not with CREATE, using a dummy value of zero. This is
because we refer to POINTERS in the defining word POSITION, but it's
not until after we've defined all the field names that we know how big the
table must be and can actually ALLOT it.

As soon as we create the field names, we define the size of the table
as a constant /POINTERS. At last we reserve room for the table itself,
patching its beginning address (HERE) into the constant POINTERS.
(The word \textbf{>BODY} converts the address returned by tick into the
address of the constant's value.) Thus POINTERS returns the address of
the table allotted later, just as a name defined by CREATE returns the
address of a table allotted directly below the name's header.

Although it's valid to patch the value of a CONSTANT at compile
time, as we do here, there is a restriction of style:

\begin{tip}
A CONSTANT's value should never be changed once the application is
compiled.
\end{tip}
The case of alternating states is slightly more involved. In this situation
we need to alternate back and forth between two (or more) states, never
clobbering the conditions in each state when we jump to the other state.
\Fig{fig7-5} shows the conceptual model for this kind of state table.

\wepsfiga{fig7-5}{Conceptual model for alternating-states tables.}

%!! include \Fig{fig7-5 here

In this model, the names TOP, BOTTOM, etc., can be made to point
into either of two tables, REAL or PSEUDO. By making the REAL table
the current one, all the pointer names reference addresses in the REAL
table; by making the PSEUDO table current, they address the PSEUDO
table.

The code in \Fig{fig7-6} implements this alternating states
mechanism. The words WORKING and PRETENDING change the
pointer appropriately. For instance:

\begin{figure*}
\caption{Implementation of alternating-states mechanism.}
\labelfig{fig7-6}
\begin{center}
\begin{BVerbatim}
VARIABLE 'POINTERS  \ pointer to state table
: POINTERS ( -- adr of current table)   'POINTERS @ ;
: POSITION   ( o -- o+2 ) CREATE DUP ,  2+
   DOES>  ( -- a )  @ POINTERS + ;
O  \ initial offset
POSITION TOP
POSITION BOTTOM
POSITION LEFT
POSITION RIGHT
POSITION INSIDE
POSITION OUT
CONSTANT /POINTERS  \ final computed offset
CREATE REAL    /POINTERS ALLOT  \ real state table
CREATE PSEUDO  /POINTERS ALLOT  \ temporary state table
: WORKING      REAL 'POINTERS ! ;     WORKING
: PRETENDING   PSEUDO 'POINTERS ! ;
\end{BVerbatim}
\end{center}
\end{figure*}

\begin{Code}[commandchars=&\{\}]
WORKING
10 TOP !
TOP &underline{? 10}
PRETENDING
20 TOP !
TOP &underline{? 20}
WORKING
TOP &underline{? 10}
PRETENDING
TOP &underline{? 20}
\end{Code}
The major difference with this latter approach is that names go through
an extra level of indirection (POINTERS has been changed from a
constant to a colon definition). The field names can be made to point to
either of two state tables. Thus each name has slightly more work to do.
Also, in the former approach the names refer to fixed locations; a
\textbf{CMOVE} is required each time we save or restore the values. In this
approach, we have only to change a single pointer to change the current
table.

\section{Vectored Execution}

Vectored execution extends the ideas of currentness and indirection
beyond data, to functions. Just as we can save values and flags in
variables, we can also save functions, because functions can be referred to
by address.

The traditional techniques for implementing vectored execution are
described in \emph{Starting FORTH}, Chapter Nine. In this section we'll
discuss a new syntax which I invented and which I think can be used in many
circumstances more elegantly than the traditional methods.

The syntax is called DOER/MAKE. (If your system doesn't include
these words, refer to Appendix B for code and implementation details.) It
works like this: You define the word whose behavior will be vectorable
with the defining word \textbf{DOER}, as in

\begin{Code}
DOER PLATFORM
\end{Code}
Initially, the new word PLATFORM does nothing. Then you can write
words that change what PLATFORM does by using the word MAKE:

\begin{Code}
: LEFTWING   MAKE PLATFORM  ." proponent " ;
: RIGHTWING  MAKE PLATFORM  ." opponent " ;
\end{Code}
When you invoke LEFTWING, the phrase MAKE PLATFORM changes
what PLATFORM will do. Now if you type PLATFORM, you'll see:

\begin{Code}[commandchars=\&\{\}]
LEFTWING ok
PLATFORM &underline{proponent ok}
\end{Code}
RIGHTWING will make PLATFORM display ``opponent.'' You can use
PLATFORM within another definition:

\begin{Code}
: SLOGAN   ." Our candidate is a longstanding " PLATFORM
   ." of heavy taxation for business. " ;
\end{Code}
The statement

\begin{Code}
LEFTWING SLOGAN
\end{Code}
will display one campaign statement, while

\begin{Code}
RIGHTWING SLOGAN
\end{Code}
will display another.

The ``MAKE'' code can be any FORTH code, as much or as long as
you want; just remember to conclude it with semicolon. The semicolon at
the end of LEFTWING serves for both LEFTWING and for the bit of
code after MAKE. When MAKE redirects execution of the \textbf{DOER} word,
it also \emph{stops} execution of the word in which it appears.

When you invoke LEFTWING, for example, \textbf{MAKE} redirects
PLATFORM and exits. Invoking LEFTWING does not cause ``proponent''
to be printed. \Fig{fig7-7} demonstrates this point, using a
conceptualized illustration of the dictionary.

If you want to \emph{continue} execution, you can use the word \textbf{;AND} in
place of semicolon. \textbf{;AND} terminates the code that the \textbf{DOER} word

\wepsfigt{fig7-7}{DOER and MAKE.}
\wepsfigt{fig7-8}{Multiple MAKEs in parallel using ;AND.}
\wepsfigt{fig7-9}{Multiple MAKEs in series.}

%!! include \Fig{fig7-7 here

points to, and resumes execution of the definition in which it appears, as
you can see in \Fig{fig7-8}.

Finally, you can chain the ``making'' of \textbf{DOER} words in series by
not using \textbf{;AND}. \Fig{fig7-9} explains this better than I could write about
it.

\section{Using DOER/MAKE}

There are many occasions when the DOER/MAKE construct proves
beneficial. They are:

%!! there's certainly some better way to do ordered lists in TeX
%!! This list has seven items and spans several pages;
%!! I have marked its end using ---End-Of-List---
\begin{enumerate}
\item To change the state of a function (when external testing of the state is not
necessary). The words LEFTWING and RIGHTWING change the state of
the word PLATFORM.

\item To factor out internal phrases from similar definitions, but within control
structures such as loops.

Consider the definition of a word called DUMP, designed to reveal the
contents of a specified region of memory.

%!! include \Fig{fig7-8 here

%!! does \textbf{} work inside verbatim blocks? It better should...
\begin{Code}[commandchars=\&\{\}]
: DUMP  ( a # )
   O DO  I 16 MOD O= IF  CR  DUP I +  5 U.R  2 SPACES  THEN
   DUP I +  &poorbf{@ 6} U.R  &poorbf{2 +LOOP}  DROP ;
\end{Code}
The problem arises when you write a definition called CDUMP, designed to
format the output according to bytes, not cells:

%!! does \textbf{} work inside verbatim blocks? It better should...
\begin{Code}[commandchars=\&\{\}]
: CDUMP  ( a # )
   O DO  I 16 MOD O= IF  CR  DUP I +  5 U.R  2 SPACES  THEN
   DUP I +  &poorbf{C@  4} U.R  &poorbf{LOOP} DROP ;
\end{Code}
%!! include \Fig{fig7-9 here

The code within these two definitions is identical except for the fragments
in boldface. But factoring is difficult because the fragments occur inside the
\textbf{DO LOOP}.

Here's a solution to this problem, using \textbf{DOER/MAKE}. The code that
changes has been replaced with the word .UNIT, whose behavior is vectored
by the code in DUMP and CDUMP. (Recognize that ``\textbf{1 +LOOP}'' has
the same effect as ``\textbf{LOOP}''.)

%!! does \textbf{} work inside verbatim blocks? It better should...
\begin{Code}[commandchars=\&\{\}]
DOER .UNIT ( a -- increment)  \ display byte or cell
: <DUMP>  ( a # )
   O DO  I 16 MOD O= IF  CR  DUP I +  5 U.R  2 SPACES  THEN
   DUP I + &poorbf{.UNIT}  +LOOP  DROP ;
: DUMP   ( a #)  MAKE .UNIT &poorbf{ @  6} U.R  &poorbf{2} ;AND <DUMP> ;
: CDUMP ( a #)   MAKE .UNIT &poorbf{C@  4} U.R  &poorbf{1} ;AND <DUMP> ;
\end{Code}
Notice how DUMP and CDUMP \emph{set-up} the vector, then go on to
\emph{execute} the shell (the word <DUMP>).

\item To change the state of related functions by invoking a single command. For
instance:

\begin{Code}
DOER TYPE'
DOER EMIT'
DOER SPACES'
DOER CR'
: VISIBLE     MAKE TYPE'  TYPE ;AND
              MAKE EMIT'  EMIT ;AND
              MAKE SPACES'  SPACES ;AND
              MAKE CR'  CR ;
: INVISIBLE   MAKE TYPE'  2DROP ;AND
              MAKE EMIT'  DROP ;AND
              MAKE SPACES'  DROP ;AND
              MAKE CR'  ;
\end{Code}
Here we've defined a vectorable set of output words, each name having a
``prime'' mark at the end. VISIBLE sets them to their expected functions.
INVISIBLE makes them no-ops, eating up the arguments that would normally
be passed to them. Say INVISIBLE and any words defined in terms
of these four output operators will \emph{not} produce any output.

\item To change the state for the next occurrence only, then change the state (or
reset it) again.

Suppose we're writing an adventure game. When the player first arrives at
a particular room, the game will display a detailed description. If the player
returns to the same room later, the game will show a shorter message.

We write:

\begin{Code}
DOER ANNOUNCE
: LONG MAKE ANNOUNCE
   CR ." You're in a large hall with a huge throne"
   CR  ." covered with a red velvet canopy. "
         MAKE ANNOUNCE
   CR ." You' re in the throne room. " ;
\end{Code}
The word ANNOUNCE will display either message. First we say LONG, to
initialize ANNOUNCE to the long message. Now we can test ANNOUNCE,
and find that it prints the long message. Having done that,
however, it continues to ``make'' ANNOUNCE display the short message.

If we test ANNOUNCE a second time, it prints the short message. And it
will for ever more, until we say LONG again.

In effect we're queuing behaviors. We can queue any number of behaviors,
letting each one set the next. The following example (though not terribly
practical) illustrates the point.

\begin{Code}
DOER WHERE
VARIABLE SHIRT
VARIABLE PANTS
VARIABLE DRESSER
VARIABLE CAR

: ORDER  \  specify search order
   MAKE WHERE  SHIRT   MAKE WHERE  PANTS
   MAKE WHERE  DRESSER   MAKE WHERE CAR
   MAKE WHERE  O ;

: HUNT  ( -- a|O )  \  find location containing 17
   ORDER  5 O DO  WHERE  DUP O=  OVER @  17 =  OR  IF
      LEAVE  ELSE  DROP  THEN  LOOP ;
\end{Code}
In this code we've created a list of variables, then defined an ORDER in
which they are to be searched. The word HUNT looks through each of them,
looking for the first one that contains a 17. HUNT returns either the
address of the correct variable, or a zero if none have the value.

It does this by simply executing WHERE five times. Each time, WHERE
returns a different address, as defined in ORDER, then finally zero.

We can even define a \textbf{DOER} word that toggles its own behavior endlessly:

\begin{Code}
DOER SPEECH
: ALTERNATE
   BEGIN  MAKE SPEECH ." HELLO "
   MAKE SPEECH ." GOODBYE "
   O UNTIL ;
\end{Code}
\item To implement a forward reference. A forward reference is usually needed as
a ``hook,'' that is, a word invoked in a low-level definition but reserved for
use by a component defined later in the listing.

To implement a forward reference, build the header of the word with \textbf{DOER},
before invoking its name.

\begin{Code}
DOER STILL-UNDEFINED
\end{Code}
Later in the listing, use MAKE;

\begin{Code}
MAKE STILL-UNDEFINED  ALL THAT JAZZ ;
\end{Code}
(Remember, MAKE can be used outside a colon definition.)

\item Recursion, direct or indirect.

Direct recursion occurs when a word invokes itself. A good example is the
recursive definition of greatest-common-denominator:

\begin{Code}
GCD of a, b =  a                     if b = O
               GCD of b, a mod b     if b > O
\end{Code}
This translates nicely into:

\begin{Code}
DOER GCD ( a b -- gcd)
MAKE GCD  ?DUP  IF  DUP ROT ROT  MOD  GCD  THEN ;
\end{Code}
Indirect recursion occurs when one word invokes a second word, while the
second word invokes the first. This can be done using the form:

\begin{Code}
DOER B
: A  ... B ... ;
MAKE B  ... A ... ;
\end{Code}
\item Debugging. I often define:

\begin{Code}
DOER SNAP
\end{Code}
(short for SNAPSHOT), then edit SNAP into my application at a point
where I want to see what's going on. For instance, with SNAP invoked inside
the main loop of a keystroke interpreter, I can set it up to let me watch
what's happening to a data structure as I enter keys. And I can change
what SNAP does without having to recompile the loop.
\end{enumerate}

%%! ---End-Of-List---

The situations in which it's preferable to use the tick-and-execute approach
are those in which you need control over the address of the vector,
such as when vectoring through an element in a decision table, or
attempting to save/restore the contents of the vector.

\othersidetrue\section{Summary}
\othersidefalse
In this chapter we've examined the tradeoffs between using the stack
and using variables and other data structures. Using the stack is
preferable for testing and reusability, but too many values manipulated
on the stack by a single definition hurts readability and writeability.

We also explored techniques for saving and restoring data structures,
and concluded with a study of vectored execution using \textbf{DOER}/\textbf{MAKE}.

\section{References}

%!! there's certainly some better way to do ordered lists in TeX
\begin{enumerate}
\item Michael Ham, ``Why Novices Use So Many Variables,'' \emph{FORTH
   Dimensions}, vol. 5, no. 4, November/December 1983.
\item Daniel Slater, ``A State Space Approach to Robotics,'' \emph{The Journal
   of FORTH Application and Research}, 1, 1 (September 1983), 17.
\end{enumerate}


%% Thinking Forth
%% Copyright (C) 2004 Leo Brodie
%% Initial transcription by Nils M Holm
%% Based on OCR scans by Steve Fisher
%% 
%% Chapter: Eight - Minimizing Control Structures

%!! horizontal line
%EIGHT
%!! horizontal line

\chapter{Minimizing Control~Structures}\Chapmark{8}%
\index{C!Control structure minimization|(}%

\initial Control structures aren't as important in \Forth{} as they
are in other languages. \Forth{} programmers tend to write very
complex applications in terms of short words, without much emphasis on
\forth{IF THEN} constructs.

There are several techniques for minimizing control structures.
They include:

%!! items should be preceded by bullets
\begin{itemize}
\item computing or calculating
\item hiding conditionals through re-factoring
\item using structured exits
\item vectoring
\item redesigning.
\end{itemize}
In this chapter we'll examine these techniques for simplifying and
eliminating control structures from your code.

\section{What's So Bad about Control Structures?}%
\index{C!Control structure minimization!reasons for|(}

Before we begin reeling off our list of tips, let's pause to examine why
conditionals should be avoided in the first place.

The use of conditional structures adds complexity to your code. The
more complex your code is, the harder it will be for you to read and to
maintain. The more parts a machine has, the greater are its chances of
breaking down. And the harder it is for someone to fix.

%!! horizontal line
\begin{interview}
\person{Moore}\index{M!Moore, Charles|(} tells this story:
\begin{tfquot}
I recently went back to a company we had done some work for several years
ago. They called me in because their program is now five years old, and
it's gotten very complicated. They've had programmers going in and
patching things, adding state variables and conditionals. Every statement
that I recall being a simple thing five years ago, now has gotten very
complicated.  ``If this, else if this, else if this'' \dots{} and then the
simple thing.

Reading that statement now, it's impossible for me to figure out what it's
doing and why. I'd have to remember what each variable indicated, why it
was relevant in this case, and then what was happening as a consequence of
it---or not happening.

It started innocently. They had a special case they needed to worry about.
To handle that special case, they put a conditional in one place. Then they
discovered that they also needed one here, and here. And then a few more.
Each incremental step only added a little confusion to the program. Since
they were the programmers, they were right on top of it.

The net result was disastrous. In the end they had half a dozen flags.
Test this one, reset it, set that one, and so on. As a result of this
condition, you knew you had other conditions coming up you had to look out
for. They created the logical equivalent of spaghetti code in spite of the
opportunity for a structured program.

The complexity went far beyond what they had ever intended. But they'd
committed themselves to going down this path, and they missed the simple
solution that would have made it all unnecessary---having two words
instead of one. You either say \forth{GO} or you say \forth{PRETEND}.

In most applications there are remarkably few times when you need to test
the condition. For instance in a video game, you don't really say ``If he
presses Button A, then do this; if he presses Button B, then do something
else.'' You don't go through that kind of logic.

If he presses the button, you do something. What you do is associated with
the button, not with the logic.

Conditionals aren't bad in themselves---they are an essential construct. But
a program with a lot of conditionals is clumsy and unreadable. All you can
do is question each one. Every conditional should cause you to ask, ``What
am I doing wrong?''

What you're trying to do with the conditional can be done in a different
way. The long-term consequences of the different way are preferable to the
long-term consequences of the conditional.
\end{tfquot}\index{M!Moore, Charles|)}
\end{interview}%
\index{C!Control structure minimization!reasons for|)}%
Before we introduce some detailed techniques, let's look at three
approaches to the use of conditionals in a particular example.
\Fig{fig8-1}, \Fig{fig8-2}, and \Fig{fig8-3} illustrate three versions of
a design for an automatic teller machine.\program{teller1}

\begin{figure*}[tttt]
\verbfig{fig8-1}{The structured approach.}
\begin{center}
\small\begin{BVerbatim}[commandchars=\&\{\},baselinestretch=0.85]
&underline{AUTOMATIC-TELLER}

IF card is valid DO
   IF card owner is valid DO
      IF request withdrawal DO
         IF authorization code is valid DO
            query for amount
            IF request &(&le&) current balance DO
               IF withdrawal &(&le&) available cash DO
                  vend currency
                  debit account
               ELSE
                  message
                  terminate session
            ELSE
               message
               terminate session
         ELSE
            message
            terminate session
      ELSE
         IF authorization code is valid DO
            query for amount
            accept envelope through hatch
            credit account
         ELSE
            message
            terminate session
   ELSE
      eat card
ELSE
   message
END
\end{BVerbatim}
\end{center}
\end{figure*}
The first example comes straight out of the School for Structured
Programmers. The logic of the application depends on the correct nesting
of \forth{IF} statements.

Easy to read? Tell me under what condition the user's card gets eaten. To
answer, you have to either count \forth{ELSE}s from the bottom and match
them with the same number of \forth{IF}s from the top, or use a
straightedge.

\begin{figure*}[tttt]
\verbfig{fig8-2}{Nesting conditionals within named procedures.}
\small\begin{center}
\begin{BVerbatim}[commandchars=\&\{\},baselinestretch=0.85]
&underline{AUTOMATIC-TELLER}

PROCEDURE READ-CARD
     IF  card is readable  THEN  CHECK-OWNER
          ELSE  eject card  END

PROCEDURE CHECK-OWNER
     IF  owner is valid  THEN  CHECK-CODE
          ELSE  eat card  END

PROCEDURE CHECK-CODE
     IF  code entered matches owner  THEN  TRANSACT
          ELSE message, terminate session  END

PROCEDURE TRANSACT
     IF requests withdrawal  THEN  WITHDRAW
          ELSE  DEPOSIT END

PROCEDURE WITHDRAW
     Query
     If  request &(&le&) current balance  THEN  DISBURSE  END

PROCEDURE DISBURSE
     IF disbursement &(&le&) available cash  THEN
           vend currency
           debit account
         ELSE  message  END

PROCEDURE DEPOSIT
     accept envelope
     credit account
\end{BVerbatim}
\end{center}
\end{figure*}
The second version, \Fig{fig8-2}, shows the improvement that using many
small, named procedures can have on readability. The user's card is eaten
if the owner is not valid.

But even with this improvement, the design of each word depends completely
on the \emph{sequence} in which the tests must be performed. The
supposedly ``highest'' level procedure is burdened with eliminating the
worst-case, most trivial kind of event. And each test becomes responsible
for invoking the next test.

\begin{figure*}[tttt]
\verbfig{fig8-3}{Refactoring and/or eliminating conditionals.}
\begin{center}
\small\begin{BVerbatim}[commandchars=\&\{\},baselinestretch=0.85]
&underline{AUTOMATIC-TELLER}

: RUN
     READ-CARD  CHECK-OWNER  CHECK-CODE  TRANSACT  ;

: READ-CARD
     valid code sequence NOT readable  IF  eject card  QUIT
        THEN ;

: CHECK-OWNER
     owner is NOT valid  IF  eat card  QUIT  THEN ;

: CHECK-CODE
     code entered MISmatches owner's code  IF  message  QUIT
        THEN ;

: READ-BUTTON ( -- adr-of-button's-function)
     ( device-dependent primitive) ;

: TRANSACT
     READ-BUTTON  EXECUTE ;

1 BUTTON WITHDRAW

2 BUTTON DEPOSIT

: WITHDRAW
     Query
     request &(&le&) current balance  IF  DISBURSE  THEN ;

: DISBURSE
     disbursement &(&le&) available cash  IF
            vend currency
            debit account
          ELSE  message  THEN  ;

: DEPOSIT
     accept envelope
     credit account ;
\end{BVerbatim}
\end{center}
\end{figure*}

The third version comes closest to the promise of \Forth{}. The highest
level word expresses exactly what's happening conceptually, showing only
the main path. Each of the subordinate words has its own error exit, not
cluttering the reading of the main word. One test does not have to invoke
the next test.

Also \code{TRANSACT} is designed around the fact that the user will make
requests by pressing buttons on a keypad. No conditions are necessary. One
button will initiate a withdrawal, another a deposit. This approach
readily accommodates design changes later, such as the addition of a
feature to transfer funds. (And this approach does not thereby become
dependent on hardware. Details of the interface to the keypad may be
hidden within the keypad lexicon, \code{READ-BUTTON} and \code{BUTTON}.)

Of course, \Forth{} will allow you to take any of the three approaches.
Which do you prefer?\program{teller2}

\section{How to Eliminate Control Structures}

In this section we'll study numerous techniques for simplifying or
avoiding conditionals. Most of them will produce code that is more
readable, more maintainable, and more efficient. Some of the techniques
produce code that is more efficient, but not always as readable.
Remember, therefore: Not all of the tips will be applicable in all
situations.

\subsection{Using the Dictionary}%
\index{C!Control structure minimization!dictionary|(}%
\index{D!Dictionary:!control structure minimization with|(}

\begin{tip}
Give each function its own definition.
\end{tip}
By using the \Forth{} dictionary properly, we're not actually eliminating
conditionals; we're merely factoring them out from our application code.
The \Forth{} dictionary is a giant string case statement. The match and
execute functions are hidden within the \Forth{} system.

\begin{interview}
\person{Moore}:\index{M!Moore, Charles|(}
\begin{tfquot}
In my accounting package, if you receive a check from somebody, you type
the amount, the check number, the word \forth{FROM}, and the person's
name:

\begin{Code}
200.00 127 FROM ALLIED
\end{Code}
The word \forth{FROM} takes care of that situation. If you want to bill
someone, you type the amount, the invoice number, the word \forth{BILL}
and the person's name:

\begin{Code}
1000.00 280 BILL TECHNITECH
\end{Code}
\end{tfquot}\index{M!Moore, Charles|)}
\dots{} One word for each situation. The dictionary is making the decision.\medskip
\end{interview}
This notion pervades \Forth{} itself. To add a pair of single-length
numbers we use the command \forth{+}. To add a pair of double-length
numbers we use the command \forth{D+}. A less efficient, more complex
approach would be a single command that somehow ``knows'' which type of
numbers are being added.

\Forth{} is efficient because all these words---\forth{FROM} and
\forth{BILL} and \forth{+} and \forth{D+}---can be implemented without any
need for testing and branching.

\index{D!Dumb words|(}
\begin{tip}
Use dumb words.
\end{tip}
This isn't advice for TV writers. It's another instance of using the
dictionary. A ``dumb'' word is one that is not state-dependent, but
instead, has the same behavior all the time (``referentially
transparent'').

A dumb word is unambiguous, and therefore, more trustworthy.

A few common \Forth{} words have been the source of controversy recently
over this issue. One such word is \forth{."} which prints a string.  In
its simplest form, it's allowed only inside a colon definition:

\begin{Code}
: TEST   ." THIS IS A STRING " ;
\end{Code}
Actually, this version of the word does \emph{not} print a string. It
\emph{compiles} a string, along with the address of another definition
that does the printing at run time.

This is the dumb version of the word. If you use it outside a colon
definition, it will uselessly compile the string, not at all what a
beginner might expect.

To solve this problem, the FIG model added a test inside \forth{."} that
determined whether the system was currently compiling or interpreting.  In
the first case, \forth{."} would compile the string and the address of the
primitives; in the second case it would \forth{TYPE} it.

\forth{."} became two completely different words housed together in one
definition with an \forth{IF ELSE THEN} structure. The flag that indicates
whether \Forth{} is compiling or interpreting is called \forth{STATE}.
Since the \forth{."} depends on \forth{STATE}, it is said to be
``\forth{STATE}-dependent,'' literally.

The command \emph{appeared} to behave the same inside and outside a colon
definition. This duplicity proved useful in afternoon introductions to
\Forth{}, but the serious student soon learned there's more to it than
that.

Suppose a student wants to write a new word called \forthb{B."} (for
``bright-dot-quote'') to display a string in bright characters on her
display, to be used like this:

\begin{Code}
." INSERT DISK IN "  B." LEFT "  ." DRIVE "
\end{Code}
She might expect to define \forth{B."} as

\begin{Code}
: B."   BRIGHT  ."  NORMAL ;
\end{Code}
that is, change the video mode to bright, print the string, then reset the
mode to normal.

She tries it. Immediately the illusion is destroyed; the deception is
revealed; the definition won't work.

To solve her problem, the programmer will have to study the definition of
\forth{(.")} in her own system. I'm not going to get sidetracked here with
explaining how \forth{(.")} works---my point is that smartness isn't all
it appears to be.

Incidentally, there's a different syntactical approach to our student's
problem, one that does not require having two separate words, \forth{."}
and \forth{B."} to print strings. Change the system's \forth{(.")} so that
it always sets the mode to normal after typing, even though it will
already be normal most of the time. With this syntax, the programmer need
merely precede the emphasized string with the simple word \forth{BRIGHT}.

\begin{Code}
." INSERT DISK IN "  BRIGHT ." LEFT "  ." DRIVE "
\end{Code}
The '83 Standard now specifies a dumb \forth{."} and, for those cases
where an interpretive version is wanted, the new word \forth{.(} has been
added. Happily, in this new standard we're using the dictionary to make a
decision by having two separate words.

\index{S!STATE|(}
The word \forth{'} (tick) has a similar history. It was
\forthb{STATE}-dependent in fig-\Forth{}, and is now dumb in the '83
Standard. Tick shares with dot-quote the characteristic that a programmer
might want to reuse either of these words in a higher-level definition and
have them behave in the same way they do normally.%
\index{D!Dumb words|)}

\begin{tip}
Words should not depend on \forthb{STATE} if a programmer might ever want
to invoke them from within a higher-level definition and expect them to
behave as they do interpretively.
\end{tip}
\forthb{ASCII}\index{A!ASCII} works well as a \forth{STATE}-dependent
word, and so does \forthb{MAKE}.\index{M!MAKE} (See \App{C}.)%
\index{S!STATE|)}%
\index{C!Control structure minimization!dictionary|)}%
\index{D!Dictionary:!control structure minimization with|)}

\subsection{Nesting and Combining Conditionals}%
\index{N!Nesting conditionals|(}%
\index{C!Conditionals, nesting and combining|(}%
\index{C!Control structure minimization!nesting and combining conditionals|(}
\begin{tip}
Don't test for something that has already been excluded.
\end{tip}
Take this example, please:

\begin{Code}
: PROCESS-KEY
   key  dup  LEFT-ARROW  =  IF CURSOR-LEFT  THEN
        dup  RIGHT-ARROW =  IF CURSOR-RIGHT THEN
        dup  UP-ARROW    =  IF CURSOR-UP    THEN
             DOWN-ARROW  =  IF CURSOR-DOWN  THEN ;
\end{Code}
This version is inefficient because all four tests must be made regardless
of the outcome of any of them. If the key pressed was the left-arrow key,
there's no need to check if it was some other key.

Instead, you can nest the conditionals, like this:

\begin{Code}
: PROCESS-KEY
   key  dup  LEFT-ARROW  =  IF CURSOR-LEFT  ELSE
        dup  RIGHT-ARROW =  IF CURSOR-RIGHT ELSE
        dup  UP-ARROW    =  IF CURSOR-UP    ELSE
                               CURSOR-DOWN
           THEN THEN THEN  drop ;
\end{Code}
\begin{tip}
Combine booleans of similar weight.
\end{tip}
Many instances of doubly-nested \forthb{IF }\forthb{THEN} structures can
be simplified by combining the flags with logical operators before making
the decision.
Here's a doubly-nested test:
\begin{Code}
: ?PLAY   SATURDAY? IF  WORK FINISHED? IF
     GO PARTY  THEN  THEN ;
\end{Code}
The above code uses nested \forthb{IF}s to make sure that it's both
Saturday and the chores are done before it boogies on down. Instead,
let's combine the conditions logically and make a single decision:

\begin{Code}
: ?PLAY   SATURDAY?  WORK FINISHED? and  IF
   GO PARTY  THEN ;
\end{Code}
It's simpler and more readable.

The logical ``or'' situation, when implemented with
\forthb{IF }\forthb{THEN}s, is even clumsier:

\begin{Code}
: ?RISE    PHONE RINGS? IF  UP GET  THEN
     ALARM-CLOCK RINGS?  IF UP GET THEN ;
\end{Code}
This is much more elegantly written as

\begin{Code}
: ?RISE  PHONE RINGS?  ALARM RINGS? or  IF  UP GET THEN ;
\end{Code}
One exception to this rule arises when the speed penalty for checking
some of the conditions is too great.

We might write

\begin{Code}
: ?CHOW-MEIN   BEAN-SPROUTS?  CHOW-MEIN RECIPE?  and IF
   CHOW-MEIN PREPARE  THEN ;
\end{Code}
But suppose it's going to take us a long time to hunt through our recipe
file to see if there's anything on chow mein. Obviously there's no point in
undertaking the search if we have no bean sprouts in the fridge. It would
be more efficient to write

\begin{Code}
: ?CHOW-MEIN   BEAN-SPROUTS? IF  CHOW-MEIN RECIPE? IF
   CHOW-MEIN PREPARE THEN   THEN ;
\end{Code}
We don't bother looking for the recipe if there are no sprouts.

Another exception arises if any term is probably not true. By
eliminating such a condition first, you avoid having to try the other
conditions.

\begin{tip}
When multiple conditions have dissimilar weights (in likelihood or
calculation time) nest conditionals with the term that is least likely
to be true or easiest to calculate on the outside.
\end{tip}
Trying to improve performance in this way is more difficult with the
\forth{OR} construct. For instance, in the definition

\begin{Code}
: ?RISE  PHONE RINGS?  ALARM RINGS? or  IF  UP GET THEN ;
\end{Code}
we're testing for the phone and the alarm, even though only one of them
needs to ring for us to get up. Now suppose it were much more difficult to
determine that the alarm clock was ringing. We could write

\begin{Code}
: ?RISE   PHONE RINGS? IF  UP GET  ELSE
     ALARM-CLOCK RINGS?  IF UP GET THEN THEN  ;
\end{Code}
If the first condition is true, we don't waste time evaluating the second.
We have to get up to answer the phone anyway.

The repetition of \forth{UP GET} is ugly---not nearly as readable as the
solution using \forth{OR}---but in some cases desirable.%
\index{C!Conditionals, nesting and combining|)}%
\index{C!Control structure minimization!nesting and combining conditionals|)}
\index{N!Nesting conditionals|)}%

\subsection{Choosing Control Structures}%
\index{C!Control structures:!choosing|(}
\begin{tip}
The most elegant code is that which most closely matches the problem.
Choose the control structure that most closely matches the control-flow
problem.
\end{tip}%
\index{C!Control structures:!choosing|)}

\subsubsection{Case Statements}%
\index{C!Case statements|(}%
\index{C!Control structure minimization!case statements|(}

A particular class of problem involves selecting one of several possible
paths of execution according to a numeric argument. For instance, we
want the word \forth{.SUIT} to take a number representing a suit of playing
cards, 0 through 3, and display the name of the suit. We might define this
word using nested \forthb{IF }\forthb{ELSE }\forthb{THEN}s, like this:

\begin{Code}
: .SUIT ( suit -- )
  dup  0=  IF ." HEARTS "   ELSE
  dup  1 = IF ." SPADES "   ELSE
  dup  2 = IF ." DIAMONDS " ELSE
              ." CLUBS "
  THEN THEN THEN  drop ;
\end{Code}
We can solve this problem more elegantly by using a ``case statement.''

Here's the same definition, rewritten using the ``\person{Eaker} case
statement'' format, named after Dr.\@ \person{Charles E. Eaker}, the
gentleman who proposed it \cite{eaker}.

\begin{Code}
: .SUIT ( suit -- )
  CASE
  0 OF   ." HEARTS "    ENDOF
  1 OF   ." SPADES "    ENDOF
  2 OF   ." DIAMONDS "  ENDOF
  3 OF   ." CLUBS "     ENDOF     ENDCASE ;
\end{Code}
The case statement's value lies exclusively in its readability and
writeability. There's no efficiency improvement either in object memory or
in execution speed. In fact, the case statement compiles much the same
code as the nested \forthb{IF }\forthb{THEN} statements. A case statement
is a good example of compile-time factoring.

Should all \Forth{} systems include such a case statement? That's a matter
of controversy. The problem is twofold. First, the instances in which a
case statement is actually needed are rare---rare enough to question its
value. If there are only a few cases, a nested \forthb{IF }\forthb{ELSE
}\forthb{THEN} construct will work as well, though perhaps not as
readably. If there are many cases, a decision table is more flexible.

Second, many case-like problems are not quite appropriate for the case
structure. The \person{Eaker} case statement assumes that you're testing
for equality against a number on the stack. In the instance of
\forth{.SUIT}, we have contiguous integers from zero to three. It's more
efficient to use the integer to calculate an offset and directly jump to
the right code.

In the case of our Tiny Editor, later in this chapter, we have not one, but
two, dimensions of possibilities. The case statement doesn't match that
problem either.

Personally, I consider the case statement an elegant solution to a
misguided problem: attempting an algorithmic expression of what is
more aptly described in a decision table.

A case statement ought to be part of the application when useful,
but not part of the system.%
\index{C!Case statements|)}%
\index{C!Control structure minimization!case statements|)}

\subsubsection{Looping Structures}%
\index{C!Control structure minimization!looping structures|(}%
\index{L!Loops|(}
The right looping structure can eliminate extra conditionals.
\begin{interview}
\person{Moore}:\index{M!Moore, Charles|(}
\begin{tfquot}
Many times conditionals are used to get out of loops. That particular use
can be avoided by having loops with multiple exit points.

This is a live topic, because of the multiple \forthb{WHILE} construct
which is in poly\Forth{} but hasn't percolated up to \Forth{} '83. It's a
simple way of defining multiple \forthb{WHILE}s in the same
\forthb{REPEAT}.

Also \person{Dean Sanderson}\index{S!Sanderson, Dean} [of \Forth{}, Inc.]
has invented a new construct that introduces two exit points to a
\forthb{DO }\forthb{LOOP}.\index{D!DO LOOP} Given that construction you'll
have fewer tests. Very often I leave a truth value on the stack, and if
I'm leaving a loop early, I change the truth value to remind myself that I
left the loop early. Then later I'll have an \forthb{IF} to see whether I
left the loop early, and it's just clumsy.

Once you've made a decision, you shouldn't have to make it again. With the
proper looping constructs you won't need to remember where you came from,
so more conditionals will go away.

This is not completely popular because it's rather unstructured. Or perhaps
it is elaborately structured. The value is that you get simpler programs.
And it costs nothing.
\end{tfquot}\index{M!Moore, Charles|)}
\end{interview}
Indeed, this is a live topic. As of this writing it's too early to make
any specific proposals for new loop constructs. Check your system's
documentation to see what it offers in the way of exotic looping
structures.  Or, depending on the needs of your application, consider
adding your own conditional constructs. It's not that hard in \Forth{}.

I'm not even sure whether this use of multiple exits doesn't violate the
doctrine of structured programming. In a
\forthb{BEGIN }\forthb{WHILE }\forthb{REPEAT}
loop with multiple \forthb{WHILE}s, all the exits bring you to a common
``continue'' point: the \forthb{REPEAT}. But with \person{Sanderson}'s
\index{S!Sanderson, Dean} construct, you
can exit the loop by jumping \emph{past} the end of the loop, continuing
at an \forthb{ELSE}. There are two possible ``continue'' points.

This is ``less structured,'' if we can be permitted to say that. And
yet the definition will always conclude at its semicolon and return to the
word that invoked it. In that sense it is well-structured; the module has
one entry point and one exit point.

When you want to execute special code only if you did \emph{not} leave the
loop prematurely, this approach seems the most natural structure to use.
(We'll see an example of this in a later section, ``Using Structured
Exits.'')
%
\index{C!Counts vs. terminators|(}
\index{T!Terminators vs. counts|(}
\begin{tip}
Favor counts over terminators.
\end{tip}
\Forth{} handles strings by saving the length of the string in the first
byte. This makes it easier to type, move, or do practically anything with
the string. With the address and count on the stack, the definition of
\forthb{TYPE} can be coded:

\begin{Code}
: type  ( a #)  over + swap DO  i c@ emit  LOOP ;
\end{Code}
(Although \forthb{TYPE} really ought to be written in machine code.)

This definition uses no overt conditional. \forthb{LOOP} actually hides the
conditional since each loop checks the index and returns to \forthb{DO} if it
has not yet reached the limit.

If a delimiter were used, let's say ASCII null (zero), the definition
would have to be written:

\begin{Code}
: type  ( a)  BEGIN dup c@  ?dup WHILE  emit  1+
   REPEAT  drop ;
\end{Code}
An extra test is needed on each pass of the loop. (\forthb{WHILE} is a
conditional operator.)

Optimization note: The use of \forthb{?DUP} in this solution is expensive
in terms of time because it contains an extra decision itself. A faster
definition would be:

\begin{Code}
: type  ( a)  BEGIN dup c@  dup WHILE emit 1+
    REPEAT  2drop ;
\end{Code}
The '83 Standard applied this principle to \forthb{INTERPRET}
\index{I!INTERPRET}
which now accepts a count rather than looking for a terminator.

The flip side of this coin is certain data structures in which it's
easiest to \emph{link} the structures together. Each record points to the
next (or previous) record. The last (or first) record in the chain can be
indicated with a zero in its link field.

If you have a link field, you have to fetch it anyway. You might as
well test for zero. You don't need to keep a counter of how many records
there are. If you decrement a counter to decide whether to terminate,
you're making more work for yourself. (This is the technique used to
implement \Forth{}'s dictionary as a linked list.)%
\index{C!Counts vs. terminators|)}%
\index{C!Control structure minimization!looping structures|)}%
\index{L!Loops|)}%
\index{T!Terminators vs. counts|)}%

\subsubsection{Calculating Results}%
\index{C!Control structure minimization!calculating results|(}%
\index{C!Calculations|(}%

\begin{tip}
Don't decide, calculate.
\end{tip}
Many times conditional control structures are applied mistakenly to
situations in which the difference in outcome results from a difference in
numbers. If numbers are involved, we can calculate them. (In Chapter
Four see the section called ``Calculations vs. Data Structures vs. Logic.'')
%
\index{B!Booleans, as hybrid values|(}
\begin{tip}
Use booleans as hybrid values.
\end{tip}
This is a fascinating corollary to the previous tip, ``Don't decide,
calculate.'' The idea is that booleans, which the computer represents as
numbers, can efficiently be used to effect numeric decisions. Here's one
example, found in many \Forth{} systems:

\begin{Code}
: s>d  ( n -- d)  \ sign extend s to d
     dup O<  IF -1  ELSE  0 THEN ;
\end{Code}
(The purpose of this definition is to convert a single-length number to
double-length. A double-length number is represented as two 16-bit
values on the stack, the high-order part on top. Converting a positive
integer to double-length merely means adding a zero onto the stack, to
represent its high-order part. But converting a negative integer to
double-length requires ``sign extension;'' that is, the high-order part
should be all ones.)

The above definition tests whether the single-length number is negative.
If so, it pushes a negative one onto the stack; otherwise a zero.  But
notice that the outcome is merely arithmetic; there's no change in
process. We can take advantage of this fact by using the boolean itself:

\begin{Code}
: s>d  ( n -- d)  \ sign extend s to d
     dup  O< ;
\end{Code}
This version pushes a zero or negative one onto the stack without a
moment's (in)decision.

\medbreak
(In pre-1983 systems, the definition would be:

\begin{Code}
: s>d  ( n -- d)  \ sign extend s to d
     dup  O< negate ;
\end{Code}
See \App{C}.)%
\index{B!Booleans, as hybrid values|)}\medbreak

\index{H!Hybrid values|(}%
We can do even more with ``hybrid values'':
%
\index{A!AND|(}
\begin{tip}
To effect a decision with a numeric outcome, use \forthb{AND}.
\end{tip}
In the case of a decision that produces either zero or a non-zero ``$n$,''
the traditional phrase

\begin{Code}
( ? ) IF  n  ELSE  0  THEN
\end{Code}
is equivalent to the simpler statement

\begin{Code}
( ? )  n and
\end{Code}
Again, the secret is that ``true'' is represented by $-1$ (all ones) in
'83 \Forth{} systems. \forthb{AND}ing ``$n$'' with the flag will either
produce ``$n$'' (all bits intact) or ``$0$'' (all bits cleared).

To restate with an example:

\begin{Code}
( ? )  IF  200  ELSE  0  THEN
\end{Code}
is the same as

\begin{Code}
( ? )  200 and
\end{Code}
Take a look at this example:

\begin{Code}
n  a b <  IF  45 +  THEN
\end{Code}
This phrase either adds 45 to ``$n$'' or doesn't, depending on the
relative sizes of ``$a$'' and ``$b$.'' Since ``adding 45 or not'' is the
same as ``adding 45 or adding 0,'' the difference between the two outcomes
is purely numeric.  We can rid ourselves of a decision, and simply
compute:

\begin{Code}
n  a b <  45 and  +
\end{Code}
\begin{interview}
\person{Moore}:\index{M!Moore, Charles|(}
\begin{tfquot}
The ``\forth{45 AND}'' is faster than the \forth{IF}, and certainly more
graceful. It's simpler. If you form a habit of looking for instances where
you're calculating this value from that value, then usually by doing
arithmetic on the logic you get the same result more cleanly.

I don't know what you call this. It has no terminology; it's merely doing
arithmetic with truth values. But it's perfectly valid, and someday boolean
algebra and arithmetic expressions will accommodate it.

In books you often see a lot of piece-wise linear approximations that fail to
express things clearly. For instance the expression

\begin{Code}[commandchars=\&\{\}]
x = 0 for t < 0
x = 1 for t &(&ge&) 0
\end{Code}
This would be equivalent to
\begin{Code}
t  0<  1 and
\end{Code}
as a single expression, not a piece-wise expression.
\end{tfquot}\index{M!Moore, Charles|)}
\end{interview}%%
\index{A!AND|)}

I call these flags ``hybrid values'' because they are booleans (truth
values) being applied as data (numeric values). Also, I don't know what
else to call them.%
\index{H!Hybrid values|)}

We can eliminate numeric \forth{ELSE} clauses as well (where both results
are non-zero), by factoring out the difference between the two results.
For instance,

\begin{Code}
: STEPPERS  'TESTING? @  IF 150 ELSE 151  THEN  load ;
\end{Code}
can be simplified to

\begin{Code}
: STEPPERS   150  'TESTING? @  1 and +  load ;
\end{Code}
This approach works here because conceptually we want to either load
Screen 150, or if testing, the next screen past it.%
\index{C!Calculations|)}%
\index{C!Control structure minimization!calculating results|)}

\section{A Note on Tricks}%
\index{T!Tricks|(}%
\index{C!Control structure minimization!tricks|(}

This sort of approach is often labeled a ``trick.'' In the computing
industry at large, tricks have a bad reputation.

A trick is simply taking advantage of certain properties of operation.
Tricks are used widely in engineering applications. Chimneys
eliminate smoke by taking advantage of the fact that heat rises.
Automobile tires provide traction by taking advantage of gravity.

Arithmetic Logic Units (ALUs) take advantage of the fact that
subtracting a number is the same as adding its two's complement.

These tricks allow simpler, more efficient designs. What justifies
their use is that the assumptions are certain to remain true.

The use of tricks becomes dangerous when a trick depends on something
likely to change, or when the thing it depends on is not protected by
information hiding.

Also, tricks become difficult to read when the assumptions on which
they're based aren't understood or explained. In the case of replacing
conditionals with \forth{AND}, once this technique becomes part of every
programmer's vocabulary, code can become \emph{more} readable. In the case
of a trick that is specific to a specific application, such as the order
in which data are arranged in a table, the listing must clearly document
the assumption used by the trick.%
\index{C!Control structure minimization!tricks|)}%
\index{T!Tricks|)}%

\begin{tip}
Use \forthb{MIN} and \forthb{MAX} for clipping.
\index{M!MAX}%
\index{M!MIN}%
\end{tip}
Suppose we want to decrement the contents of the variable \forth{VALUE}, but
we don't want the value to go below zero:

\begin{Code}
-1 Value +!  Value @  -1 = IF  0 Value !  THEN
\end{Code}
This is more simply written:

\begin{Code}
Value @  1-  0 max  Value !
\end{Code}
In this case the conditional is factored within the word \forthb{MAX}.

\subsection{Using Decision Tables}%
\index{D!Decision table|(}%
\index{C!Control structure minimization!decision tables|(}

\begin{tip}
Use decision tables.
\end{tip}
We introduced these in \Chap{2}. A decision table is a structure that
contains either data (a ``data table'') or addresses of functions (a
``function table'') arranged according to any number of dimensions. Each
dimension represents all the possible, mutually exclusive states of a
particular aspect of the problem. At the intersection of the ``true''
states of each dimension lies the desired element: the piece of data or
the function to be performed.

A decision table is clearly a better choice than a conditional structure
when the problem has multiple dimensions.

\subsubsection{One-Dimensional Data Table}
\index{O!One-dimensional data table|(}

Here's an example of a simple, one-dimensional data table. Our application
has a flag called \forth{'FREEWAY?} which is true when we're referring to
freeways, false when we're referring to city streets.

Let's construct the word \forth{SPEED-LIMIT}, which returns the speed
limit depending on the current state.
Using \forthb{IF }\forthb{THEN} we would write:

\begin{Code}
: SPEED-LIMIT  ( -- speed-limit)
     'FREEWAY? @  IF  55  ELSE  25  THEN ;
\end{Code}
We might eliminate the \forthb{IF }\forthb{THEN} by using a hybrid value
with \forthb{AND}:

\begin{Code}
: SPEED-LIMIT   25  'FREEWAY? @  30 and + ;
\end{Code}
But this approach doesn't match our conceptual model of the problem
and therefore isn't very readable.

Let's try a data table. This is a one-dimensional table, with only two
elements, so there's not much to it:

\begin{Code}
Create LIMITS   25 ,  55 ,
\end{Code}
The word \forth{SPEED-LIMIT?} now must apply the boolean to offset into
the data table:

\begin{Code}
: SPEED-LIMIT  ( -- speed-limit)
     LIMITS  'FREEWAY? @  2 and  +  @ ;
\end{Code}
Have we gained anything over the \forthb{IF }\forthb{THEN} approach?
Probably not, with so simple a problem.

What we have done, though, is to factor out the decision-making
process from the data itself. This becomes more cost-effective when we
have more than one set of data related to the same decision. Suppose we
also had

\begin{Code}
Create #LANES   4 ,  10 ,
\end{Code}
representing the number of lanes on a city street and on a freeway. We
can use identical code to compute the current number of lanes:

\begin{Code}
: #LANES?  ( -- #lanes)
     #LANES  'FREEWAY? @  2 and  +  @ ;
\end{Code}
Applying techniques of factoring, we simplify this to:

\begin{Code}
: ROAD  ( for-freeway for-city ) Create , ,
     DOES> ( -- data )  'FREEWAY? @  2 and  +  @ ;
55 25 ROAD SPEED-LIMIT?
10  4 ROAD #LANES?
\end{Code}
Another example of the one-dimensional data table is the ``superstring''
(\emph{Starting \Forth{}}, Chapter Ten).
\index{O!One-dimensional data table|)}

\subsubsection{Two-Dimensional Data Table}%
\index{T!Two-dimensional data table|(}%
\program{phone4}
In \Chap{2} we presented a phone-rate problem. \Fig{fig8-4} gives one
solution to the problem, using a two-dimensional data structure.

\begin{figure*}[bbbt]
\verbfig{fig8-4}{A solution to the phone rate problem.}
\begin{Screen}
\ Telephone rates                                       03/30/84
Create FULL     30 , 20 , 12 ,
Create LOWER    22 , 15 , 10 ,
Create LOWEST   12 ,  9 ,  6 ,
Variable RATE   \ points to FULL, LOWER or LOWEST
                \ depending on time of day
FULL RATE !  \ for instance
: CHARGE   ( o -- ) Create ,
   DOES>  ( -- rate )  @  RATE @ +  @ ;
0 CHARGE 1MINUTE   \ rate for first minute
2 CHARGE +MINUTES  \ rate for each additional minute
4 CHARGE /MILES    \ rate per each 100 miles
\end{Screen}

\begin{Screen}
\ Telephone rates                                       03/30/84
Variable OPERATOR?  \ 90 if operator assisted; else 0
Variable #MILES  \ hundreds of miles
: ?ASSISTANCE  ( direct-dial charge -- total charge)
   OPERATOR? @  + ;
: MILEAGE  ( -- charge )  #MILES @  /MILES * ;
: FIRST  ( -- charge )  1MINUTE  ?ASSISTANCE  MILEAGE + ;
: ADDITIONAL  ( -- charge)  +MINUTES  MILEAGE + ;
: TOTAL ( #minutes -- total charge)
   1- ADDITIONAL *  FIRST + ;
\end{Screen}
\end{figure*}

In this problem, each dimension of the data table consists of three
mutually exclusive states. Therefore a simple boolean (true/false) is
inadequate. Each dimension of this problem is implemented in a different
way.

The current rate, which depends on the time of day, is stored as an
address, representing one of the three rate-structure sub-tables. We can
say
\begin{Code}
FULL RATE !
\end{Code}
or
\begin{Code}
LOWER RATE !
\end{Code}
etc.\goodbreak

The current charge, either first minute, additional minute, or per mile,
is expressed as an offset into the table (0, 2, or 4).

An optimization note: we've implemented the two-dimensional table as a set
of three one-dimensional tables, each pointed to by \forth{RATE}. This
approach eliminates the need for a multiplication that would otherwise be
needed to implement a two-dimensional structure. The multiplication can be
prohibitively slow in certain cases.%
\index{T!Two-dimensional data table|)}%
\program{phone5}

\subsubsection{Two-Dimensional Decision Table}%
\index{T!Two-dimensional decision table|(}%
\program{editor4}

We'll hark back to our Tiny Editor example in \Chap{3} to illustrate
a two-dimensional decision table.

In \Fig{fig8-5} we're constructing a table of functions to be performed
when various keys are pressed. The effect is similar to that of a
case statement, but there are two modes, Normal Mode and Insert Mode.
Each key has a different behavior depending on the current mode.

The first screen implements the change of the modes. If we invoke

\begin{Code}
NORMAL MODE# !
\end{Code}
we'll go into Normal Mode.

\begin{Code}
INSERTING MODE# !
\end{Code}
enters Inserting Mode.

The next screen constructs the function table, called \forth{FUNCTIONS}.
The table consists of the ASCII value of a key followed by the address of
the routine to be performed when in Normal Mode, followed by the address
of the routine to be performed when in Insert Mode, when that key
is pressed. Then comes the second key, followed by the next pair of
addresses, and so on.

In the third screen, the word \forth{'FUNCTION} takes a key value,
searches through the \forth{FUNCTIONS} table for a match, then returns the
address of the cell containing the match. (We preset the variable
\forth{MATCHED} to point to the last row of the table---the functions we want
when \emph{any} character is pressed.)

The word \forth{ACTION} invokes \forth{'FUNCTION}, then adds the contents
of the variable \forth{MODE\#}. Since \forth{MODE\#} will contain either a
2 or a 4, by adding this offset we're now pointing into the table at the
address of the routine we want to perform. A simple

\begin{Code}
@ execute
\end{Code}
will perform the routine (or \forthb{@EXECUTE} if you have it).

In fig-\Forth{}, change the definition of \forth{IS} to:

\begin{Code}
: IS   [compile] '  CFA , ;
\end{Code}
\begin{figure*}[tttp]
\verbfig{fig8-5}{Implementation of the Tiny Editor.}
\setcounter{screen}{30}
\begin{Screen}
\ Tiny Editor
2 Constant NORMAL     \ offset in FUNCTIONS
4 Constant INSERTING  \        "
6 Constant /KEY       \ bytes in table for each key
Variable MODE#        \ current offset into table
NORMAL MODE# !
: INSERT-OFF   NORMAL    MODE# ! ;
: INSERT-ON    INSERTING MODE# ! ;

Variable ESCAPE?      \ t=time-to-leave-loop
: ESCAPE  true ESCAPE? ! ;





\end{Screen}
\begin{Screen}
\ Tiny Editor             function table             07/29/83
: IS   ' , ;  \   function   ( -- )    ( for '83 standard)
Create FUNCTIONS
\ keys                  normal mode        insert mode
 4 ,  ( ctrl-D)         IS DELETE          IS INSERT-OFF
 9 ,  ( ctrl-I)         IS INSERT-ON       IS INSERT-OFF
 8 ,  ( backspace)      IS BACKWARD        IS INSERT<
60 ,  ( left arrow)     IS BACKWARD        IS INSERT-OFF
62 ,  ( right arrow)    IS FORWARD         IS INSERT-OFF
27 ,  ( return)         IS ESCAPE          IS INSERT-OFF
 0 ,  ( no match)       IS OVERWRITE       IS insert
here /KEY -  Constant 'NOMATCH  \ adr of no-match key




\end{Screen}
\begin{Screen}
\ Tiny Editor cont'd                                 07/29/83
Variable MATCHED
: 'FUNCTION  ( key -- adr-of-match )  'NOMATCH  MATCHED !
   'NOMATCH FUNCTIONS DO  dup  i @ =  IF
     i MATCHED !  LEAVE  THEN  /KEY +LOOP  drop
    MATCHED @ ;
: ACTION  ( key -- )  'FUNCTION  MODE# @ +  @ execute ;
: GO   false ESCAPE? !  BEGIN  key ACTION  ESCAPE? @ UNTIL ;








\end{Screen}
\end{figure*}
In 79-Standard \Forth{}s, use:
\begin{Code}
: IS   [compile] '  , ;
\end{Code}
We've also used non-redundancy at compile time in the definition just
below the function table:
\begin{Code}
here /KEY -  Constant 'NOMATCH  \  adr of no-match key
\end{Code}
We're making a constant out of the last row in the function table. (At the
moment we invoke \forthb{HERE}, it's pointing to the next free cell after
the last table entry has been filled in. Six bytes back is the last row.)
We now have two words:
\begin{Code}
FUNCTIONS  ( adr of beginning of function table )
'NOMATCH   ( adr of "no-match" row; these are the
             routines for any key not in the table)
\end{Code}
We use these names to supply the addresses passed to \forthb{DO}:
\begin{Code}
'NOMATCH FUNCTION DO
\end{Code}
to set up a loop that runs from the first row of the table to the last. We
don't have to know how many rows lie in the table. We could even delete a
row or add a row to the table, without having to change any other piece of
code, even the code that searches through the table.

Similarly the constant \forth{/KEY} hides information about the number of
columns in the table.

Incidentally, the approach to \forth{'FUNCTION} taken in the listing is a
quick-and-dirty one; it uses a local variable to simplify stack
manipulation. A simpler solution that uses no local variable is:
\begin{Code}
: 'FUNCTION  ( key -- adr of match )
   'NOMATCH swap  'NOMATCH FUNCTIONS DO  dup
      i @ =  IF swap drop i swap  LEAVE  THEN
   /KEY +LOOP  drop ;
\end{Code}
(We'll offer still another solution later in this chapter, under ``Using
Structured Exits.'')%
\index{T!Two-dimensional decision table|)}%
\program{editor5}

\subsection{Decision Tables for Speed}

We've stated that if you can calculate a value instead of looking it up in a
table, you should do so. The exception is where the requirements for
speed justify the extra complexity of a table.

Here is an example that computes powers of two to 8-bit precision:

\begin{Code}
Create TWOS
   1 c,  2 c,  4 c,  8 c,  16 c,  32 c,
: 2**  ( n -- 2-to-the-n)
   TWOS +  c@ ;
\end{Code}
Instead of computing the answer by multiplying two times itself ``$n$''
times, the answers are all pre-computed and placed in a table. We can use
simple addition to offset into the table and get the answer.

In general, addition is much faster than multiplication.

\begin{interview}
\person{Moore}\index{M!Moore, Charles|(} provides another example:
\begin{tfquot}
If you want to compute trig functions, say for a graphics display, you don't
need much resolution. A seven-bit trig function is probably plenty. A table
look-up of 128 numbers is faster than anything else you're going to be able
to do. For low-frequency function calculations, decision tables are great.

But if you have to interpolate, you have to calculate a function anyway.
You're probably better off calculating a slightly more complicated function
and avoiding the table lookup.
\end{tfquot}\index{M!Moore, Charles|)}
\end{interview}%
\index{D!Decision table|)}%
\index{C!Control structure minimization!decision tables|)}

\subsection{Redesigning}%
\index{R!Redesigning to minimize control structures|(}%
\index{C!Control structure minimization!redesigning|(}

\begin{tip}
One change at the bottom can save ten decisions at the top.
\end{tip}
In our interview with \person{Moore}\index{M!Moore, Charles} at the
beginning of the chapter, he mentioned that much conditional testing could
have been eliminated from an application if it had been redesigned so that
there were two words instead of one: ``You either say \forth{GO} or you
say \forth{PRETEND}.''

It's easier to perform a simple, consistent algorithm while changing
the context of your environment than to choose from several algorithms
while keeping a fixed environment.

Recall from \Chap{1} our example of the word \forth{APPLES}. This was
originally defined as a variable; it was referred to many times throughout
the application by words that incremented the number of apples (when
shipments arrive), decremented the number (when apples are sold), and
checked the current number (for inventory control).

When it became necessary to handle a second type of apples, the
\emph{wrong} approach would have been to add that complexity to all
the shipment/\hy sales/\hy inventory words. The \emph{right} approach
was the one we took: to add the complexity ``at the bottom''; that is,
to \forth{APPLES} itself.

This principle can be realized in many ways. In \Chap{7} (under ``The
State Table'') we used state tables to implement the words \forth{WORKING}
and \forth{PRETENDING}, which changed the meaning of a group of variables.
Later in that chapter, we used vectored execution to define
\forth{VISIBLE} and \forth{INVISIBLE}, to change the meanings of
\forth{TYPE'}, \forth{EMIT'}, \forth{SPACES'} and \forth{CR'} and thereby
easily change all the formatting code that uses them.

\begin{tip}
Don't test for something that can't possibly happen.
\end{tip}
Many contemporary programmers are error-checking-happy.

There's no need for a function to check an argument passed by
another component in the system. The calling program should bear the
responsibility for not exceeding the limits of the called component.
%
\index{A!Algorithms|(}
\begin{tip}
Reexamine the algorithm.
\end{tip}
\begin{interview}
\person{Moore}:\index{M!Moore, Charles|(}
\begin{tfquot}
A lot of conditionals arise from fuzzy thinking about the problem. In
servo-control theory, a lot of people think that the algorithm for the
servo ought to be different when the distance is great than when it is
close. Far away, you're in slew mode; closer to the target you're in
decelerate mode; very close you're in hunt mode. You have to test how far
you are to know which algorithm to apply.

I've worked out a non-linear servo-control algorithm that will handle full
range. This approach eliminates the glitches at the transitioning points
between one mode and the other. It eliminates the logic necessary to decide
which algorithm to use. It eliminates your having to empirically determine
the transition points. And of course, you have a much simpler program with
one algorithm instead of three.

Instead of trying to get rid of conditionals, you're best to question the
underlying theory that led to the conditionals.
\end{tfquot}
\end{interview}%%
\index{A!Algorithms|)}
\begin{tip}
Avoid the need for special handling.
\end{tip}
One example we mentioned earlier in the book: if you keep the user out of
trouble you won't have to continually test whether the user has gotten
into trouble.

\begin{interview}
\person{Moore}:
\begin{tfquot}
Another good example is writing assemblers. Very often, even though an
opcode may not have a register associated with it, pretending that it has a
register---say, Register 0---might simplify the code. Doing arithmetic by
introducing bit patterns that needn't exist simplifies the solution. Just
substitute zeros and keep on doing arithmetic that you might have avoided
by testing for zero and not doing it.

It's another instance of the ``don't care.'' If you don't care, then give
it a dummy value and use it anyway.
\end{tfquot}\index{M!Moore, Charles|)}
\end{interview}
Anytime you run into a special case, try to find an algorithm for which
the special case becomes a normal case.

\begin{tip}
Use properties of the component.
\end{tip}
A well-designed component---hardware or software---will let you implement
a corresponding lexicon in a clean, efficient manner. The character
graphics set from the old Epson MX-80 printer (although now obsolete)
illustrates the point well. \Fig{fig8-6} shows the graphics characters
produced by the ASCII codes 160 to 223.

\wepsfiga{fig8-6}{The Epson MX-80 graphics character set.}

%!! include Figure 8-6 here

Each graphics character is a different combination of six tiny boxes,
either filled in or left blank. Suppose in our application we want to use
these characters to create a design. For each character, we know what we
want in each of the six positions---we must produce the appropriate
ASCII character for the printer.

A little bit of looking will tell you there's a very sensible pattern
involved. Assuming we have a six-byte table in which each byte represents
a pixel in the pattern:

%!! include inline graphic here (page 254)
\bigskip
{\sf\begin{tabular}{l|c|c|} \cline{2-3}
PIXELS~~ & ~0~ & ~1~ \\ \cline{2-3}
       & 2 & 3 \\ \cline{2-3}
       & 4 & 5 \\ \cline{2-3}
\end{tabular}}\bigskip

\noindent and assuming that each byte contains hex \code{FF} if the pixel is
``on;'' zero if it is ``off,'' then here's how little code it takes to
compute the character:

\begin{Code}
Create PIXELS  6 allot
: PIXEL  ( i -- a )  PIXELS + ;
: CHARACTER  ( -- graphics character)
   160   6 0 DO  i PIXEL c@  i 2** and  +  LOOP ;
\end{Code}
(We introduced \forth{2**} a few tips back.)

No decisions are necessary in the definition of \forth{CHARACTER}. The
graphics character is simply computed.

Note: to use the same algorithm to translate a set of six adjoining
pixels in a large grid, we can merely redefine \forth{PIXEL}. That's an example
of adding indirection backwards, and of good decomposition.

Unfortunately, external components are not always designed well.
For instance, The IBM Personal Computer uses a similar scheme for
graphics characters on its video display, but without any discernible
correspondence between the ASCII values and the pattern of pixels. The
only way to produce the ASCII value is by matching patterns in a lookup
table.

\begin{interview}
\person{Moore}:\index{M!Moore, Charles|(}
\begin{tfquot}
The 68000 assembler is another example you can break your heart over,
looking for a good way to express those op-codes with the minimal number
of operators. All the evidence suggests there is no good solution. The
people who designed the 68000 didn't have assemblers in mind. And they
could have made things a lot easier, at no cost to themselves.
\end{tfquot}\index{M!Moore, Charles|)}
\end{interview}
By using properties of a component in this way, your code becomes
dependent on those properties and thus on the component itself. This is
excusable, though, because all the dependent code is confined to a single
lexicon, which can easily be changed if necessary.%
\index{C!Control structure minimization!redesigning|)}
\index{R!Redesigning to minimize control structures|)}%

\subsection{Using Structured Exits}%
\index{S!Structured exits|(}%
\index{C!Control structure minimization!structured exits|(}

\begin{tip}
Use the structured exit.
\end{tip}
In the chapter on factoring we demonstrated the possibility of factoring
out a control structure using this technique:

\begin{Code}
: CONDITIONALLY   A B or  C and  IF  NOT r> drop  THEN ;
: ACTIVE   CONDITIONALLY   TUMBLE JUGGLE jump ;
: LAZY   CONDITIONALLY  SIT  EAT  SLEEP ;
\end{Code}
\Forth{} allows us to alter the control flow by directly manipulating the
return stack. (If in doubt, see \emph{Starting \Forth{}}, Chapter Nine.)
Indiscreet application of this trick can lead to unstructured code with nasty
side effects. But the disciplined use of the structured exit can actually
simplify code, and thereby improve readability and maintainability.

\begin{interview}
\person{Moore}:\index{M!Moore, Charles|(}
\begin{tfquot}
More and more I've come to favor \forth{R> DROP} to alter the flow of
control.  It's similar to the effect of an \forthb{ABORT"}, which has an
\forthb{IF }\forthb{THEN} built in it. But that's only one \forthb{IF
}\forthb{THEN} in the system, not at every error.

I either abort or I don't abort. If I don't abort, I continue. If I do
abort, I don't have to thread my way through the path. I short-circuit the
whole thing.

The alternative is burdening the rest of the application with checking
whether an error occurred. That's an inconvenience.
\end{tfquot}\index{M!Moore, Charles|)}
\end{interview}
The ``abort route'' circumvents the normal paths of control flow under
special conditions. \Forth{} provides this capability with the words
\forthb{ABORT"} and \forthb{QUIT}.

The ``structured exit'' extends the concept by allowing the immediate
termination of a single word, without quitting the entire application.

This technique should not be confused with the use of GOTO, which is
unstructured to the extreme. With GOTO you can go anywhere, inside or
outside the current module. With this technique, you effectively jump
directly to the final exit point of the module (the semicolon) and resume
execution of the calling word.
%
\index{E!EXIT|(}
The word \forthb{EXIT} terminates the definition in which the word
appears.  The phrase \forthb{R> DROP} terminates the definition that
called the definition in which the phrase appears; thus it has the same
effect but can be used one level down. Here are some examples of both
approaches.

If you have an \forthb{IF }\forthb{ELSE }\forthb{THEN} phrase in which no
code follows \forthb{THEN}, like this:

\begin{Code}
... HUNGRY?  IF  EAT-IT  ELSE  FREEZE-IT  THEN ;
\end{Code}
you can eliminate \forthb{ELSE} by using \forthb{EXIT}:

\begin{Code}
... HUNGRY?  IF EAT-IT EXIT  THEN  FREEZE-IT ;
\end{Code}
(If the condition is true, we eat and run; \forthb{EXIT} acts like a
semicolon.  If the condition is false, we skip to \forthb{THEN} and
\forth{FREEZE-IT}.)

The use of \forthb{EXIT} here is more efficient, saving two bytes and
extra code to perform, but it is not as readable.%
\index{E!EXIT|)}

\begin{interview}
\person{Moore}\index{M!Moore, Charles|(} comments on the value, and
danger, of this technique:
\begin{tfquot}
Especially if your conditionals are getting elaborate, it's handy to jump
out in the middle without having to match all your \forthb{THEN}s at the
end. In one application I had a word that went like this:

\begin{Code}
: TESTING
   SIMPLE  1CONDITION IF ... EXIT THEN
           2CONDITION IF ... EXIT THEN
           3CONDITION IF ... EXIT THEN ;
\end{Code}
SIMPLE handled the simple cases. SIMPLE ended up with \forthb{R> DROP}.
These other conditions were the more complex ones.

Everyone exited at the same point without having to painfully match all
the \forthb{IF}s, \forthb{ELSE}s, and \forthb{THEN}s. The final result, if
none of the conditions matched, was an error condition.

It was bad code, difficult to debug. But it reflected the nature of the
problem. There wasn't any better scheme to handle it. The \forthb{EXIT}
and \forthb{R> DROP} at least kept things manageable.
\end{tfquot}\index{M!Moore, Charles|)}
\end{interview}
Programmers sometimes also use \forthb{EXIT} to get out of a complicated
\forthb{BEGIN} loop in a graceful way. Or we might use a related technique
in the \forthb{DO }\forthb{LOOP} that we wrote for \forth{'FUNCTION} in
our Tiny Editor, earlier in this chapter. In this word, we are searching
through a series of locations looking for a match. If we find a match, we
want to return the address where we found it; if we don't find a match, we
want the address of the last row of the functions table.

We can introduce the word \forth{LEAP} (see \App{C}), which will work
like \forthb{EXIT} (it will simulate a semicolon). Now we can write:

\begin{Code}
: 'FUNCTION  ( key -- adr-of-match )
   'NOMATCH FUNCTIONS DO  dup  i @ =  IF  drop i LEAP
   THEN  /KEY +LOOP  drop  'NOMATCH ;
\end{Code}
If we find a match we \forth{LEAP}, not to \forthb{+LOOP}, but right out
of the definition, leaving \forth{I} (the address at which we found it) on
the stack. If we don't find a match, we fall through the loop and execute

\begin{Code}
drop  'NOMATCH
\end{Code}
which drops the key\# being searched for, then leaves the address of the
last row!

As we've seen, there may be times when a premature exit is appropriate,
even multiple exit points and multiple ``continue'' points.

Remember though, this use of \forthb{EXIT} and \forthb{R> DROP} is
\emph{not consistent} with structured programming in the strictest sense,
and requires great care.

For instance, you may have a value on the stack at the beginning of a
definition which is consumed at the end. A premature \forthb{EXIT} will
leave the unwanted value on the stack.

Fooling with the return stack is like playing with fire. You can get
burned. But how convenient it is to have fire.%
\index{C!Control structure minimization!structured exits|)}
\index{S!Structured exits|)}%

\subsection{Employing Good Timing}%
\index{C!Control structure minimization!timing|(}

\begin{tip}
Take the action when you know you need to, not later.
\end{tip}
Any time you set a flag, ask yourself why you're setting it. If the answer
is, ``So I'll know to do such-and-such later,'' then ask yourself if you
can do such-and-such \emph{now}. A little restructuring can greatly
simplify your design.

\begin{tip}
Don't put off till run time what you can compile today.
\end{tip}
Any time you can make a decision prior to compiling an application, do.

Suppose you had two versions of an array: one that did bounds
checking for your protection during development and one that ran faster,
though unprotected for the actual application.

Keep the two versions in different screens. When you compile your
application, load only the version you need.

By the way, if you follow this suggestion, you may go crazy editing
parentheses in and out of your load blocks to change which version gets
loaded each time. Instead, write throw-away definitions that make the
decisions for you. For instance (as already previewed in another context):

\begin{Code}
: STEPPERS   150  'TESTING? @  1 and +  load ;
\end{Code}
\begin{tip}
\forthb{DUP} a flag, don't recreate it.
\end{tip}
Sometimes you need a flag to indicate whether or not a previous piece of
code was invoked. The following definition leaves a flag which indicates
that \forth{DO-IT} was done:

\begin{Code}
: DID-I?  ( -- t=I-did)
   SHOULD-I?  IF  DO-IT  true  ELSE  false  THEN ;
\end{Code}
This can be simplified to:

\begin{Code}
: DID-I?  ( -- t=I-did)
        SHOULD-I? dup  IF  DO-IT  THEN ;
\end{Code}
\begin{tip}
Don't set a flag, set the data.
\end{tip}
If the only purpose to setting a flag is so that later some code can decide
between one number and another, you're better off saving the number
itself.

The ``colors'' example in \Chap{6}'s section called ``Factoring
Criteria'' illustrates this point.

The purpose of the word \forth{LIGHT} is to set a flag which indicates
whether we want the intensity bit to be set or not. While we could have
written\program{color4}

\begin{Code}
: LIGHT   true 'LIGHT? ! ;
\end{Code}
to set the flag, and

\begin{Code}
'LIGHT? @ IF  8 or  THEN ...
\end{Code}
to use the flag, this approach is not quite as simple as putting the
intensity bit-mask itself in the variable:

\begin{Code}
: LIGHT   8 'LIGHT? ! ;
\end{Code}
and then simply writing

\begin{Code}
'LIGHT? @  or ...
\end{Code}
to use it.

\begin{tip}
Don't set a flag, set the function. (Vector.)
\end{tip}
This tip is similar to the previous one, and lives under the same
restriction. If the only purpose to setting a flag is so that later some
code can decide between one function and another, you're better off saving
the address of the function itself.

For instance, the code for transmitting a character to a printer is
different than for slapping a character onto a video display. A poor
implementation would define:

\begin{Code}
Variable DEVICE  ( O=video | 1=printer)
: VIDEO   false DEVICE ! ;
: PRINTER   true DEVICE ! ;
: type  ( a # -- ) DEVICE @ IF
   ( ...code for printer...) ELSE
   ( ...code for video...)  THEN ;
\end{Code}
This is bad because you're deciding which function to perform every time
you type a string.

A preferable implementation would use vectored execution. For
instance:

\begin{Code}
DOER type  ( a # -- )
: VIDEO   MAKE type ( ...code for video...) ;
: PRINTER   MAKE type ( ...code for printer...) ;
\end{Code}
This is better because \forth{TYPE} doesn't have to decide which code to
use, it already knows.

(On a multi-tasked system, the printer and monitor tasks would
each have their own copies of an execution vector for \forth{TYPE}
stored in a user variable.)

The above example also illustrates the limitation of this tip. In our
second version, we have no simple way of knowing whether our current
device is the printer or the video screen. We might need to know, for
instance, to decide whether to clear the screen or issue a formfeed. Then
we're making an additional use of the state, and our rule no longer
applies.

A flag would, in fact, allow the simplest implementation of additional
state-dependent operations. In the case of \forth{TYPE}, however, we're
concerned about speed. We type strings so often, we can't afford to waste
time doing it. The best solution here might be to set the function of
\forth{TYPE} and also set a flag:

\begin{Code}
DOER type
: VIDEO   0 DEVICE !  MAKE type
     ( ...code for video...) ;
: PRINTER   1 DEVICE !  MAKE type
     ( ...code for printer...) ;
\end{Code}
Thus \forth{TYPE} already knows which code to execute, but other definitions
will refer to the flag.

Another possibility is to write a word that fetches the parameter of the
\forthb{DOER} word \forth{TYPE} (the pointer to the current code) and
compares it against the address of \forth{PRINTER}. If it's less than the
address of \forth{PRINTER}, we're using the \forth{VIDEO} routine;
otherwise we're using the \forth{PRINTER} routine.

If changing the state involves changing a small number of functions,
you can still use \forth{DOER}/\forth{MAKE}. Here are definitions of three
memory-move operators that can be shut off together.

\begin{Code}
DOER !'  ( vectorable ! )
DOER CMOVE'  ( vectorable CMOVE )
DOER FILL'  ( vectorable FILL )
: STORING   MAKE !' ! ;AND
            MAKE CMOVE'  cmove ;AND
            MAKE FILL'  fill ;
: -STORING  MAKE !'  2drop ;AND
            MAKE CMOVE'  2drop drop ;AND
            MAKE FILL'  2drop drop ;
\end{Code}
But if a large number of functions need to be vectored, a state table
would be preferable.

A corollary to this rule introduces the ``structured exit hook,'' a
\forthb{DOER} word vectored to perform a structured exit.

\begin{Code}
DOER HESITATE  ( the exit hook)
: DISSOLVE   HESITATE  FILE-DIVORCE ;
\end{Code}
(\dots{} Much later in the listing:)
\begin{Code}
: RELENT   MAKE HESITATE   SEND-FLOWERS  r> drop ;
\end{Code}
By default, \forth{HESITATE} does nothing. If we invoke \forth{DISSOLVE},
we'll end up in court. But if we \forth{RELENT} before we
\forth{DISSOLVE}, we'll send flowers, then jump clear to the semicolon,
canceling that court order before our partner ever finds out.

This approach is especially appropriate when the cancellation must be
performed by a function defined much later in the listing (decomposition
by sequential complexity). Increased complexity of the earlier code is
limited solely to defining the hook and invoking it at the right spot.%
\index{C!Control structure minimization!timing|)}

\subsection{Simplifying}%
\index{S!Simplicity|(}%
\index{C!Control structure minimization!simplification|(}

I've saved this tip for last because it exemplifies the rewards of opting
for simplicity. While other tips concern maintainability, performance,
compactness, etc., this tip relates to the sort of satisfaction that
Thoreau sought at Walden Pond.

\begin{tip}
Try to avoid altogether saving flags in memory.
\end{tip}
A flag on the stack is quite different from a flag in memory. Flags on the
stack can simply be determined (by reading the hardware, calculating, or
whatever), pushed onto the stack, then consumed by the control structure.
A short life with no complications.

But save a flag in memory and watch what happens. In addition to
having the flag itself, you now have the complexity of a location for the
flag. The location must be:

%!! items should be preceded by bullets
\begin{itemize}
\item created
\item initialized (even before anything actually changes)
\item reset (otherwise, passing a flag to a command leaves the flag in
that current state).
\end{itemize}
Because flags in memory are variables, they are not reentrant.

An example of a case in which we might reconsider the need for a flag is
one we've seen several times already. In our ``colors'' example we made
the assumption that the best syntax would be:\program{color5}

\begin{Code}
LIGHT BLUE
\end{Code}
that is, the adjective \forth{LIGHT} preceding the color. Fine. But
remember the code to implement that version? Compare it with the
simplicity of this approach:

\begin{Code}
0 Constant BLACK    1 Constant BLUE    2 Constant GREEN
3 Constant CYAN     4 Constant RED     5 Constant MAGENTA
6 Constant BROWN    7 Constant GRAY
: LIGHT   ( color -- color )  8 or ;
\end{Code}
In this version we've reversed the syntax, so that we now say

\begin{Code}
BLUE LIGHT
\end{Code}
We establish the color, then we modify the color.

We've eliminated the need for a variable, for code to fetch from the
variable and more code to reset the variable when we're done. And the
code is so simple it's impossible not to understand.

When I first wrote these commands, I took the English-like approach.
``\forth{BLUE LIGHT}'' sounded backwards, not at all acceptable. That was
before my conversations with
\person{Chuck Moore}\index{M!Moore, Charles|(}.

\begin{interview}
\person{Moore}'s philosophy is persuasive:
\begin{tfquot}
I would distinguish between reading nicely in English and reading nicely.
In other languages such as Spanish, adjectives follow nouns. We should be
independent of details like which language we're thinking in.

It depends on your intention: simplicity, or emulation of English. English
is not such a superb language that we should follow it slavishly.
\end{tfquot}\index{M!Moore, Charles|)}
\end{interview}
If I were selling my ``colors'' words in a package for graphic artists, I
would take the trouble to create the flag. But writing these words for my
own use, if I had to do it over again, I'd favor the \person{Moore}-ish
influence, and use ``\forth{BLUE LIGHT}.''%
\index{C!Control structure minimization|)}%
\index{C!Control structure minimization!simplification|)}%
\index{S!Simplicity|)}%

\section{Summary}

The use of logic and conditionals as a significant structural element in
programming leads to overly-complicated, difficult-to-maintain, and
inefficient code. In this chapter we've discussed several ways to minimize,
optimize or eliminate unnecessary conditional structures.

As a final note, \Forth{}'s downplaying of conditionals is not shared
by most contemporary languages. In fact, the Japanese are basing their
fifth-generation computer project on a language called PROLOG---for
PROgramming in LOGic---in which one programs entirely in logic. It will
be interesting to see the battle-lines forming as we ponder the question:

\begin{tfquot}
To \forthb{IF} or not to \forthb{IF}
\end{tfquot}
In this book we've covered the first six steps of the software development
cycle, exploring both the philosophical questions of designing software
and practical considerations of implementing robust, efficient, readable
software.

We have not discussed optimization, validation, debugging, documenting,
project management, \Forth{} development tools, assembler
definitions, uses and abuses of recursion, developing multiprogrammed
applications, or target compilation.

But that's another story.

\begin{references}{9}
\bibitem{eaker} \person{Charles Eaker}, ``Just in Case,'' \emph{\Forth{}
Dimensions} II/3, p. 37.
\end{references}

\section{For Further Thinking}

\ifeightyfour\begin{enumerate}
\item
You have the word \forth{CHOOSE} which takes an argument ``$n$''
and returns a random number between 0 and $n-1$. The result is always
positive or zero.  You can use \forth{CHOOSE} to produce a flag; the
phrase
\begin{Code}
2 CHOOSE
\end{Code}
produces a random flag of zero or one (false or true).

%%!!! this is not the second question, there is no question above
\item
Write a phrase to choose a number at random between 0 and 19 (inclusive)
\emph{or} between $-20$ and 0.
\end{enumerate}\else
%% replacement problem
Define the word \forth{DIRECTION}, which returns either 1,
$-1$, or 0, depending on whether the input argument is non-zero
positive, negative, or zero, respectively.
\fi

%% Thinking Forth
%% Copyright (C) 2004 Leo Brodie
%% Initial transcription by Ed Beroset
%% 
%% Chapter: Epilog
% \epilog{}  - it's not an Appendix, but not a numbered chapter, either
\part{Epilog}
\chapter{
FORTH's Effect
on Thinking}
\begin{quote}
FORTH is like the Tao: it is a Way, and is realized when followed.  Its
fragility is its strength; its simplicity is its direction {\em (Michael Ham,
winning entry in Mountain View Press's contest to describe FORTH in twenty-five 
words or less).}
\end{quote}

To help extract something of the FORTH philosophy, I conducted a poll
among several FORTH users in which I asked, ``How has FORTH affected
your thinking?  Have you found yourself applying `FORTH-like' 
principles in other areas?''

Here are some of the replies:

Mark Berstein is president of Eastgate Systems Inc\. in Cambridge, 
Massachusetts, and holds a doctorate from the department of 
chemistry at Harvard University.

\begin{quote}
I first met FORTH while working in laser chemistry.  I was trying to build a
rather complicated controller for a new laser spectrometer.  The original
plans called for a big green box full of electronics, The Interface.  nobody 
had built this particular kind of instrument before---that's why we were 
doing it---and the list of things we wanted the computer to handle changed
every couple of weeks.

After a few months, I had hundreds of pages of assembly-language 
routines, three biig circuit boards filled with ICs, and a 70-odd pin System
Bus.  Day by day, everything got more fragile and harder to fix.  The wiring 
on the circuit boards frayed, the connectors got loose, the assembler code
grew ever more tangled.

FORTH was an obvious solution to the software problem, since it provided
a decent environment in which to build and maintain a complex and rapidly-
changing program.  But the essence of good FORTH programming is the 
art of factoring procedures into useful, free-standing words.  The idea of the
FORTH word had unexpected implications for laboratory hardware design.

Instead of building a big, monolithic, all-purpose Interface, I found myself
building piles of simple little boxes which worked a lot like FORTH words:
they had a fixed set of standard inputs and standard outputs, they performed
just one function, they were designed to connect up to each other 
without much effort, and they wre simple enough that you could tell what 
a box did just by looking at its label.

\ldots The idea of ``human scale'' is, I think, today's seminal concept in 
software design.  This isn't specifically a FORTH development; the great joy of
UNIX, in its youth at least, was that you could read it (since it was written
in C), understand it (since it was small), and modify it (since it was simple).
FORTH shares these virtues, although it's designed to tackle a different 
sort of problem.

Because FORTH is small, and because FORTH gives its users control over
their machines, FORTH lets humans control their applications.  It's just
silly to expect scientists to sit in front of a lab computer playing ``twenty-
questions'' with packaged software.  FORTH, used properly, lets a scientist
instruct the computer instead of letting the computer instruct the scientist.

In the same sense that in baseball, a batter is supposed to feel the bat as an 
extension of himself, FORTH is human-scaled, and helps convince you that
the computer's achievements, and its failures, are also your own.
\end{quote}
Raymond E. Dessy is Professor of Chemistry at Virginia Polytechnic Institute
and State University, Blacksburg, Virginia.
\begin{quote}
As I attempted to understand the nature and structure of the language C, I 
found myself drawing upon the knowledge I had of the organization and
approach of FORTH.  This permitted me to understand convoluted, or hig-fog-
coefficient sections describing C.

I have found the FORTH approach is an ideal platform upon which to build
an understanding and an educational framework for other languages and
operating system concepts.
\end{quote}
Jerry Boutelle is owner of Nautilus Systems in Santa Cruz, California, 
which markets the Nautilus Cross-compiler.
\begin{quote}
FORTH has changed my thinking in many ways.  Since learning FORTH 
I've coded in other languages, including assembler, BASIC and FORTRAN.
I've found that I used the same kind of decomposition wwe do in
FORTH, in the sense of creating words and grouping them together.  For
example, in handling strings I would define subroutines analogous to
CMOVE, --TRAILING, FILL, etc.

More fundamentally, FORTH has reaffirmed my faith in simplicity.  Most 
people go out and attack problems with complicated tools.  But simpler
tools are available and more useful.

I try to simplify all the aspects of my life.  There's a quote I like from {\em Tao Te Ching}
by the Chinese philosopher Lao Tzu: ``To attain knowledge, add 
things every day; to obtain wisdom, remove things every day.''
\end{quote}

%% Thinking Forth
%% Copyright (C) 2004 Leo Brodie
%% Initial transcription by Ed Beroset
%% 
%% Chapter: Appendix A, Overview of FORTH (For Newcomers)
\appendix{}
\Chapmark{A}
\chapter{Overview
of~FORTH
(For~Newcomers)}
\section{The Dictionary}\index{D!Dictionary:!defined|(}
FORTH\index{F!FORTH!overview of|(} is expressed in words (and numbers) and is separated by spaces:
\begin{verbatim}
    HAND OPEN  ARM LOWER  HAND CLOSE  ARM RAISE 
\end{verbatim}
Such commands may be typed directly from the keyboard, or edited onto 
mass storage then ``{\bf LOAD}''ed.

All words, whether included with the system or user-defined,
exist in the ``dictionary,'' a linked list.  A ``defining word,'' is used to add
new names to the dictionary.  One defining word is {\bf :} (pronounced
``colon''), which is used to define a new word in terms of previously-
defined words.  Here is how one might define a new word called LIFT:
\begin{verbatim}
    : LIFT   HAND OPEN  ARM LOWER  HAND CLOSE  ARM RAISE ;
\end{verbatim}
The {\bf ;} terminates the definition.  The new word LIFT may now be used
instead of the long sequence of words that comprise its definition.

FORTH words can be nested like this indefinitely.  Writing a 
FORTH application consists of building increasingly powerful definitions,
such as this one, in terms of previously defined ones.

Another defining word is {\bf CODE},\index{C!CODE} which is used in place of colon to
define a command in terms of machine instructions for the native processor.
Words defined with {\bf CODE} are indistinguishable to the user from 
words defined with colon.  {\bf CODE} definitions are needed only for the most
time-critical portions of an applications, if at all.
\subsection{Data Structures}
\index{D!Data structures:!operators}
Still another defining word is {\bf CONSTANT},\index{C!CONSTANT} which
is used like this:
\begin{verbatim}
    17 CONSTANT SEVENTEEN
\end{verbatim}
The new word SEVENTEEN can now be used in place of the actual number 17.

The defining word VARIABLE\index{V!VARIABLE} creates a location
for temporary data. VARIABLE is used like this:
\begin{verbatim}
    VARIABLE BANANAS
\end{verbatim}
This reserves a location which is identified by the name BANANAS.

Fetching the contents of this location is the job of the word {\bf @}
(pronounced ``fetch'').  For instance,
\begin{verbatim}
    BANANAS @
\end{verbatim}
fetches the contents of the variable BANANAS.  Its counterpart is {\bf !}
(pronounced ``store''), which stores a value into the location, as in:
\begin{verbatim}
    100 BANANAS !
\end{verbatim}
FORTH also provides a word to increment the current value by the given
value; for instance, the phrase
\begin{verbatim}
    2 BANANAS +!
\end{verbatim}
increments the count by two, making it 102.

FORTH provides many other
data structure operators\index{D!Data structures:!operators}, but more
importantly, it provides the tools necessary for the programmer to
create any time of data structure needed for the application.
\index{D!Dictionary:!defined|)}

\subsection{The Stack}
In FORTH, variables and arrays are used for saving values that may be
required by many other routines and/or at unpredictable times.  They are {\em not}
used for the local passing of data between the definitions.  For this, FORTH
employs a much simpler mechanism: the data stack.

When you type a number, it goes on the stack.  When you invoke a 
word which has numeric input, it will take it from the stack.  Thus the
phrase
\begin{verbatim}
    17 SPACES
\end{verbatim}
will display seventeen blanks on the current output device.  ``17'' pushes
the binary value 17 onto the stack; the word {\bf SPACES} consumes it.

A constant also pushes its value onto the stack; thus the phrase:
\begin{verbatim}
    SEVENTEEN SPACES
\end{verbatim}
has the same effect.

The stack operates on a ``last-in, first-out'' (LIFO) basis.  This means
that data can be passed between words in an orderly, modular way, consistent
with the nesting of colon definitions.

For instance, a defintion called GRID might invoke the phrase 17 SPACES.  
This temporary activity on the stack will be transmparent to any
other definition that invokes GRID because the value placed on the stack
is removed before the defintion of GRID ends.  The calling definition 
might have placed some numbers of its own on the stack prior to calling 
GRID.  These will remain on the stack, unharmed, until GRID has been 
executed  and calling definition continues.   % probably should be the calling..
\subsection{Control Structures}
FORTH provides all the control structures\index{C!Control structures:!defined}
needed for structured, GOTO-less programming.

The syntax of the {\bf IF THEN} construct is as follows:
\begin{verbatim}
    ... ( flag ) IF  KNOCK  THEN  OPEN ...
\end{verbatim}
The ``flag''\index{F!Flags} is a value on the stack, consumed by IF. A non-zero value
indicates true, zero indicates false.  A true flag causes the code after {\bf IF} (in
this case, the word KNOCK) to be executed.  The word {\bf THEN} marks the
end of the conditional phrase; execution resumes with the word OPEN.  A
false flag causes the code between {\bf IF} and {\bf THEN} to {\em not} be executed.  In
either case, OPEN will be performed.

The word {\bf ELSE}\index{E!ELSE} allows an alternate phrase to be executed
in the false case. In the phrase:
\begin{verbatim}
    ( flag ) IF KNOCK  ELSE  RING  THEN  OPEN ...
\end{verbatim}
the word KNOCK will be performed if the flag is true, otherwise the word
RING will be performed.  Either way, execution will continue starting
with OPEN.

FORTH also provides for indexed loops\index{L!Loops} in the form
\begin{verbatim}
    ( limit) ( index) DO ... LOOP
\end{verbatim}
and indefinite loops in the forms:
\begin{verbatim}
    ... BEGIN  ...  ( flag) UNTIL
\end{verbatim}
and
\begin{verbatim}
    ... BEGIN  ...  ( flag) WHILE ... REPEAT ;
\end{verbatim}
\subsection{For the Whole Story}
For a complete instroduction to the FORTH command set, read {\em Starting
FORTH}, published by Prentice-Hall.
\index{F!FORTH!overview of|)}

%% Thinking Forth
%% Copyright (C) 2004 Leo Brodie
%% Initial transcription by Ed Beroset
%% 
%% Chapter: Appendix B, Defining DOER/MAKE
\Chapmark{B}
\chapter{Defining
DOER/MAKE}
\index{D!DOER/MAKE|(}
\initial If your system doesn't have \forth{DOER} and \forth{MAKE}
already defined, this appendix is meant to help you install them and,
if necessary, understand how they work.  Because by its nature this
construct is system dependent, I've included several different
implementations at the end of this appendix in the hope that one of
them will work for you.  If no, and if this section doesn't give you
enough information to get them running, you probably have an unusual
system.  Please don't ask me for help; ask your \Forth{} vendor.

Here's how it works.  \forthb{DOER} is a defining word that creates an
entry with one cell in its parameter field. That cell contains the vector
address, and is initialized to point to a no-op word called
\forthb{NOTHING}.

Children of \forthb{DOER} will execute that \forthb{DOES>} code of
\forthb{DOER}, which does only two things:  fetch the vector address and
place it on the return stack.  That's all.  \Forth{} execution then
continues with this address on the return stack, which will cause the
vectored function to be performed.  It's like saying (in '83-Standard)
\begin{Code}
' NOTHING >BODY >R <return>
\end{Code}
which executes \forth{NOTHING}.  (This trick only works with colon
definitions.)

Here's an illustration of the dictionary entry created when we enter

{\sf
\bigskip
\begin{tabular}{@{} l l@{}}
DOER JOE & \\
\hline
\vline\ JOE & \vline\ pfa of NOTHING \vline \\
\hline
\ header &\  parameter field
\end{tabular}
\bigskip
}

\noindent Now suppose we define:
\begin{Code}
: TEST   MAKE JOE  CR ;
\end{Code}
that is, we define a word that can vector \forth{JOE} to do a carriage
return.

Here's a picture of the compiled definition of \forth{TEST}:
%! This is supposed to be the cell diagram on p. 277, but I just 
%! couldn't figure out how to do it.

{\sf\bigskip\begin{tabular}{|c|c|c|c|c|c|}\hline
& adr of & & adr of & adr of & adr of \\
TEST & (MAKE) & 0 & JOE & CR & EXIT \\ \hline\noalign{\vspace{2pt}}
\multicolumn{1}{c}{header} & \multicolumn{1}{c}{} & \multicolumn{1}{c}{\boxto{adr}{MARKER}} & \multicolumn{3}{c}{} \\
\end{tabular}
\bigskip}

\noindent Let's look at the code for \forthb{MAKE}.  Since we're using
\forthb{MAKE} inside a colon definition, \forthb{STATE} will be true, and
we'll execute the phrase:
\begin{Code}
COMPILE (MAKE)  HERE MARKER !  0 ,
\end{Code}
We can see how \forthb{MAKE} has compiled the address of the run-time
routine, \forthb{(MAKE)}, followed by a zero.  (We'll explain what the
zero is for, and why we save its address in the variable
\forthb{MARKER}, later).
% should probably be "save its address in the variable..."

Now let's look at what \forthb{(MAKE)} does when we execute our new 
definition \forthb{TEST}:

\bigskip{
\begin{tabular}{l@{\hspace{4ex}}>{\parindent-2ex}p{.5\textwidth}}
\verb%R>%		& Gets an address from the return stack.
			  This address points to the cell just
			  past \forthb{(MAKE)}, where the zero is.\\ 
\verb%DUP 2+%		& Gets the address of the second cell after
			  \forthb{(MAKE)}, where the address of \forthb{JOE} is.\\
\verb%DUP 2+%		& Gets the address of the third cell after
			  \forthb{(MAKE)}, where the code we want to
			  execute begins.  The stack now has:

			  \parindent2ex (~'marker,~'joe,~'code~--{}--~)\\
\verb%SWAP @ >BODY%	& Fetches the contents of the address
			  pointing to \forth{JOE} (i.e., gets the address 
			  of \forth{JOE}) and computes \forth{JOE}'s pfa, where
			  the vector address goes.\\
\verb%!%		& Stores the address where the new code
			  begins (\forthb{CR}, etc.) into the vector
			  address of \forth{JOE}. \\
			& Now \forth{JOE} points inside the definition of
			  \forth{TEST}.  When we type \forth{JOE}, we'll do a
			  carriage return.\\
\verb%@ ?DUP IF >R THEN%& Fetches the contents of the cell
			  containing zero.  Since the cell does
			  contain zero, the \forth{IF THEN} statement is
			  not performed.\\
\end{tabular}}
\bigskip

That's the basic idea.  But what about that cell containing zero?  That's
for the use of \forthb{;AND}.  Suppose we changed \forth{TEST} to read:

\begin{Code}
: TEST   MAKE JOE  CR ;AND SPACE ;
\end{Code}
That is, when we invoke \forth{TEST} we'll vector \forth{JOE} to do a
\forthb{CR}, and we'll do a \forthb{SPACE} right now.  Here's what this
new version of \forth{TEST} will look like:
%! The diagram from page 278 is supposed to be here -- I didn't have a 
%! clue, so I just omitted it.

\begin{center}\vspace{10pt}\sf\begin{tabular}{|c|c|c|c|c|c|c|c|}\hline
& adr of & \smash{\rnode{M1}{~\large\strut}} & adr of & adr of & adr of & \rnode{M2}{adr of\large\strut} & adr of \\
TEST & (MAKE) & adr & JOE & CR & EXIT & SPACE & EXIT \\ \hline\noalign{\vspace{2pt}}
\multicolumn{1}{c}{header} & \multicolumn{1}{c}{} & \multicolumn{1}{c}{\boxto{adr}{MARKER}} & \multicolumn{5}{c}{} \\
\end{tabular}
\psset{arrows=->,arrowinset=0,arrowsize=5pt,armA=0,angleA=90,angleB=90}
\ncangle{M1}{M2}
\end{center}
Here's the definition of \forthb{;AND}:
\begin{Code}
: ;AND   COMPILE  EXIT  HERE MARKER @ ! ;   IMMEDIATE
\end{Code}
We can see that \forthb{;AND} has compiled an \forthb{EXIT}\index{E!EXIT},
just as semicolon would.

Next, recall that \forthb{MAKE} saved the address of that cell in a
variable called \forthb{MARKER}.  Now \forthb{;AND} stores \forthb{HERE}
(the location of the second string of code beginning with \forthb{SPACE})
into the cell previously containing zero.  Now \forthb{(MAKE)} has a
pointer to the place to resume execution.  The phrase
\begin{Code}
IF >R THEN
\end{Code}
will leave on the return stack the address of the code beginning with
\forthb{SPACE}.  Thus execution will skip over the code between
\forthb{MAKE} and \forthb{;AND} and continue with the remainder of the
definition up to semicolon.

The word \forthb{UNDO} ticks the name of a \forthb{DOER} word, and stores the
address of \forthb{NOTHING} into it.

One final note:  on some systems you may encounter a problem.  If you use
\forthb{MAKE} outside of a colon definition to create a forward reference,
you may not be able to find the most recently defined word.  For instance,
if you have:
\begin{Code}
: REFRAIN   DO-DAH  DO-DAH ;
MAKE SONG  CHORUS  REFRAIN ;
\end{Code}
your system might think that refrain has not been defined.  The problem is
due to the placement of \forth{SMUDGE}.  As a solution, try rearranging
the order of definitions or, if necessary, put \forth{MAKE} code inside a 
definition which you then execute:
\begin{Code}
: SETUP   MAKE SONG  CHORUS  REFRAIN ;   SETUP
\end{Code}
In Laboratory Microsystems PC/FORTH 2.0, the \forth{UNSMUDGE} on line 9
handles the problem.  This problem does not arise with the
Laxen/Perry/Harris model.

The final screen is an example of using \forthb{DOER/MAKE}.  After loading
the block, enter
\begin{Code}
RECITAL
\end{Code}
then enter 
\begin{Code}
WHY?
\end{Code}
followed by return, as many times as you like (you'll get a different
reason each time).

\vfill
\setcounter{screen}{21}
\begin{Screen}
( DOER/MAKE   Shadow screen                      LPB 12/05/83 )
NOTHING   A no-opp
DOER      Defines a word whose behavior is vectorable.
MARKER    Saves adr for optional continuation pointer.
(MAKE)    Stuffs the address of further code into the
          parameter field of a doer word.
MAKE      Used interpretively:  MAKE doer-name  forth-code ;
          or inside a definition:
             : def   MAKE doer-name  forth-code ;
          Vectors the doer-name word to the forth-code.
;AND      Allows continuation of the "making" definition
UNDO      Usage:  UNDO doer-name ; makes it safe to execute




\end{Screen}
\vfill
\begin{Screen}
\ DOER/MAKE   FORTH-83 Laxen/Perry/Harris model  LPB 12/05/83 
: NOTHING ;
: DOER   CREATE  ['] NOTHING  >BODY ,  DOES> @ >R ;
VARIABLE MARKER
: (MAKE)  R>  DUP 2+  DUP 2+  SWAP @  >BODY !
   @ ?DUP IF >R THEN ;
: MAKE   STATE @ IF ( compiling)
   COMPILE (MAKE)  HERE MARKER !  0 ,
   ELSE  HERE  [COMPILE] '  >BODY !
   [COMPILE] ]  THEN ;   IMMEDIATE
: ;AND   COMPILE EXIT  HERE MARKER @ ! ;   IMMEDIATE
: UNDO   ['] NOTHING  >BODY  [COMPILE] '  >BODY ! ;

\ The code in this screen is in the public domain.


\end{Screen}
\vfill
\begin{Screen}
( DOER/MAKE   FORTH-83 Lab. Micro PC/FORTH 2.0   LPB 12/05/83 )
: NOTHING ;
: DOER   CREATE  ['] NOTHING  >BODY ,  DOES> @ >R ;
VARIABLE MARKER
: (MAKE)  R>  DUP 2+  DUP 2+  SWAP @  >BODY !
   @ ?DUP IF >R THEN ;
: MAKE   STATE @ IF ( compiling)
   COMPILE (MAKE)  HERE MARKER !  0 ,
   ELSE  HERE  [COMPILE] '  >BODY !
   [COMPILE] ] UNSMUDGE  THEN ;   IMMEDIATE
: ;AND   COMPILE EXIT  HERE MARKER @ ! ;   IMMEDIATE
: UNDO   ['] NOTHING  >BODY  [COMPILE] '  >BODY ! ;

( The code in this screen is in the public domain.)


\end{Screen}
\vfill
\begin{Screen}
( DOER/MAKE   FIG model                          LPB 12/05/83 )
: NOTHING   ;
: DOES-PFA  ( pfa -- pfa of child of <BUILD-DOES> )   2+ ;
: DOER   <BUILDS  ' NOTHING ,  DOES> @ >R ;
0 VARIABLE MARKER
: (MAKE)  R>  DUP 2+  DUP 2+  SWAP @  2+ DOES-PFA !
   @ -DUP IF >R THEN ;
: MAKE  STATE @ IF ( compiling)
   COMPILE (MAKE)  HERE MARKER !  0 ,
   ELSE  HERE  [COMPILE] '  DOES-PFA !
   SMUDGE    [COMPILE] ] THEN ; IMMEDIATE
: ;AND   COMPILE ;S  HERE MARKER @ ! ;  IMMEDIATE
: UNDO   ' NOTHING  [COMPILE] '  DOES-PFA ! ;
;S
The code in this screen is in the public domain.

\end{Screen}
\vfill
\begin{Screen}
( DOER/MAKE   79-Standard  MVP FORTH             LPB 12/05/83 )
: NOTHING ;
: DOER   CREATE  ' NOTHING  ,  DOES> @ >R ;
VARIABLE MARKER
: (MAKE)  R>  DUP 2+  DUP 2+  SWAP @  2+ ( pfa) !
   @ ?DUP IF >R THEN ;
: MAKE   STATE @ IF ( compiling)
   COMPILE (MAKE)  HERE MARKER !  0 ,
   ELSE  HERE  [COMPILE] ' !
    [COMPILE] ]  THEN ;   IMMEDIATE
: ;AND   COMPILE EXIT  HERE MARKER @ ! ;   IMMEDIATE
: UNDO   ['] NOTHING  [COMPILE] ' ! ;


( The code in this screen is in the public domain.)

\end{Screen}
\vfill
\begin{Screen}
( TODDLER: Example of DOER/MAKE                      12/01/83 )
DOER ANSWER
: RECITAL
  CR ." Your daddy is standing on the table.  Ask him 'WHY?' "
  MAKE ANSWER  ." To change the light bulb."
  BEGIN
  MAKE ANSWER  ." Because it's burned out."
  MAKE ANSWER  ." Because it was old."
  MAKE ANSWER  ." Because we put it in there a long time ago."
  MAKE ANSWER  ." Because it was dark!"
  MAKE ANSWER  ." Because it was night time!!"
  MAKE ANSWER  ." Stop saying WHY?"
  MAKE ANSWER  ." Because it's driving me crazy."
  MAKE ANSWER  ." Just let me change this light bulb!"
  FALSE UNTIL ;
: WHY?   CR  ANSWER  QUIT ;
\end{Screen}
\index{D!DOER/MAKE|)}
\vfill

%% Thinking Forth
%% Copyright (C) 2004 Leo Brodie
%% Initial transcription by Ed Beroset
%% Translated into Spanish (es_ES) by Francisco Escobedo
%% 
%% Chapter: Appendix C, Other Utilities Described in This Book
\Chapmark{C}
\chapter{Otras~utilidades descritas en~este~libro}

\initial Este ap��ndice est�� aqu�� para ayudarle a definir algunas de
las palabras a la que se hace referencia en este libro y que pueden no
existir en su sistema. Las definiciones se hacen seg��n la norma del
\Forth{}-83.

\section{Del \Chap{4}}

Una definici��n of \forthb{ASCII}\index{A!ASCII} que funcionar�� en
Forth{}-83 es:
\begin{Code}
: ASCII  ( -- c)  \  Compilando:  c  ( -- )
\ Interpretando:   c   ( -- c)
     BL WORD 1+ C@  STATE @
     IF [COMPILE] LITERAL  THEN ; IMMEDIATE
\end{Code}

\section{Del \Chap{5}}

La palabra \forthb{\bs}\index{S!salto, ��rdenes de} puede definirse
como:
\begin{Code}
: \  ( salta el resto de la l��nea)
     >IN @  64 / 1+  64 *  >IN ! ; IMMEDIATE
\end{Code}
Si decide no usar \forthb{EXIT} para concluir una pantalla, puede
definir \forthb{\bs S} como:
\begin{Code}
: \S   1024 >IN ! ;
\end{Code}
\index{F!FH|(}%
La palabra \forthb{FH} puede definirse simplemente como:
\begin{Code}
: FH   \   ( desplazamiento -- desplazamiento-bloque)   "from here" ("desde aqu��")
    BLK @ + ;
\end{Code}
Esta descomposici��n permite usar \forth{FH} de muchas formas, p.e.:
\begin{Code}
: TEST   [ 1 FH ] LITERAL LOAD ;
\end{Code}
or
\begin{Code}
: SEE   [ 2 FH ] LITERAL LIST ;
\end{Code}
Una versi��n ligeramente m��s complicada de \forth{FH} tambi��n permite
editar o cargar una pantalla con una frase como ``\forth{14 FH
  LIST}'', relativa a la pantalla que se acaba de listar
(\forth{SCR}):
\begin{Code}
: FH   \   ( desplazamiento -- desplazamiento-bloque)   "from here" ("desde aqu��")
     BLK @  ?DUP 0= IF  SCR @  THEN  + ;
\end{Code}
\index{F!FH|)}
\forthb{BL}\index{B!Blank space (BL)} es una simple constante:
\begin{Code}
32 CONSTANT BL
\end{Code}
\forthb{TRUE}\index{T!TRUE} y \forthb{FALSE}\index{F!FALSE}
pueden definirse como:
\begin{Code}
0 CONSTANT FALSE
-1 CONSTANT TRUE
\end{Code}
(las palabras de control de \Forth{} como \forth{IF} y \forth{UNTIL}
interpretan cero como ``falso'' y cualquier valor distinto de cero
como ``verdadero''. Antes de \Forth{}-83, el convenio era indicar
``verdadero'' con el valor $1$. Desde \Forth{}-83, sin embargo,
``verdadero'' se indica con hex \code{FFFF}, que es el n��mero con
signo $-1$ (todos los bits a 1).)

\forthb{WITHIN}\index{W!WITHIN} puede definirse en alto nivel as��:
\begin{Code}
: WITHIN  ( n bajo alto+1 -- ?)
     >R  1- OVER <  SWAP R>  < AND ;
\end{Code}
o
\begin{Code}
: WITHIN ( n bajo alto+1 -- ?)
   OVER -  >R - R> U< ;
\end{Code}

\section{Del \Chap{8}}

La manera de hacer \forthb{LEAP}\index{L!LEAP} depender�� de c��mo su
sistema hace los bucles \forthb{DO} \forthb{LOOP}. Si \forthb{DO}
guarda 2 elementos en la pila de retorno (el ��ndice y el l��mite),
\forthb{LEAP} debe descartar ambos adem��s de otro elemento de la pila
de retorno para salir:
\begin{Code}
: LEAP   R> R> 2DROP  R> DROP ;
\end{Code}
Si \forthb{DO} guarda \emph{3} elemento en la pila de retorno, debe
definirse:
\begin{Code}
: LEAP   R> R> 2DROP  R> R> 2DROP ;
\end{Code}


%% Thinking Forth
%% Copyright (C) 2004 Leo Brodie
%% Initial transcription by Ed Beroset
%% 
%% Chapter: Appendix D, Answers to "Further Thinking" Problems

\Chapmark{D}
\chapter{Answers~to
\smashlt{``}Further~Thinking''
Problems}
\section{\Chap{3}}
\newcounter{exercise}
\begin{enumerate}
\item The answer depends on whether you believe that other components will
need to ``know the numeric code associated with each key.  Usually this
would \emph{not} be the case.  The simpler, more compact form is therefore
preferable.  Also in the first version, to add a new key would require a
change in two places.
\item The problem with the words \forth{RAM-ALLOT} and \forth{THERE} are
that they are \emph{time-dependent}:  we must execute them in a particular
order.  Our solution then will be to devise an interface to the RAM
allocation pointer that is not dependent on order; the way to do this is
to have a \emph{single} word which does both functions transparently.

Our word's syntax will be
\begin{Code}
: RAM-ALLOT   ( #bytes-to-allot -- starting-adr) 
    ... ;
\end{Code}
This syntax will remain the same whether we define it to allocate growing 
upward:
\begin{Code}
: RAM-ALLOT  ( #bytes-to-allot -- starting-adr)
    >RAM @  dup rot +  >RAM ! ;
\end{Code}
or to allocate growing downward:
\begin{Code}
: RAM-ALLOT  ( #bytes-to-allot -- starting-adr)
    >RAM @  swap -  dup >RAM ! ;
\end{Code}
\end{enumerate}

\section{\Chap{4}}

\ifeightyfour\begin{enumerate}
\item\fi Our solution is as follows:
\begin{Code}
\ CARDS Shuffle                              6-20-83
52 Constant #CARDS
Create DECK  #CARDS allot   \ one card per byte
: CARD ( i -- adr) DECK + ;
: INIT-DECK  #CARDS 0 DO  i  i CARD  c!  LOOP ;
INIT-DECK
: 'CSWAP  ( a1 a2 -- )  \  swap bytes at a1 and a2
   2dup c@  swap c@  rot c!  swap c! ;
: SHUFFLE   \  shuffle deck of cards
   #CARDS 0 DO  i CARD  #CARDS CHOOSE CARD  'CSWAP
   LOOP ;
\end{Code}
\ifeightyfour\end{enumerate}\fi

\section{\Chap{8}}
\ifeightyfour\begin{enumerate}
\item This will work:
\begin{Code}
20 CHOOSE  2 CHOOSE  IF negate THEN
\end{Code}
But this is simpler:
\begin{Code}
40 CHOOSE  20 -
\end{Code}
\end{enumerate}\else
\begin{Code}
: DIRECTION  ( n|-n|0 -- 1|-1|0)  dup  IF  0< 1 or  THEN ;
\end{Code}
\fi

%% Note: The question was random number between 0 and 19, and -20 and 0.
%% The first solution gives between 0 and 19, or -19 and 0.
%% The second solution doesn't give a double likelyhood for 0.
%% It's a bit tricky to get both question and answer really correctly.

%% Thinking Forth
%% Copyright (C) 2004 Leo Brodie
%% Initial transcription by Ed Beroset
%% 
%% Chapter: Appendix E, Summary of Style Conventions

\Chapmark{E}
\chapter{Summary~of Style~Conventions}

\initialb The contents of this Appendix are in the public domain. We encourage
publication without restriction, provided that you credit the source.

\section{Spacing and Indentation Guidelines}%
\index{I!Indentation}\index{S!Spacing}

\begin{tfquot}
1 space between the colon and the name

2 spaces between the name and the comment\footnote{
An often-seen alternative calls for 1 space between the name and comment
and 3 between the comment and the definition. A more liberal technique
uses 3 spaces before and after the comment. Whatever you choose, be
consistent.}

2 spaces, or a carriage return, after the comment and before the
definition\footnotemark[1]

3 spaces between the name and definition if no comment is used

3 spaces indentation on each subsequent line (or multiples of 3 for nested
indentation)

1 space between words/numbers within a phrase

2 or 3 spaces between phrases

1 space between the last word and the semicolon

1 space between semicolon and \forthb{IMMEDIATE} (if invoked)
\end{tfquot}
No blank lines between definitions, except to separate distinct groups of
definitions

\section{Stack-Comment Abbreviations}
\index{S!Stack abbreviation standards|(}
\begin{longtable}{>{\small}l>{\small}l}
{}{\setlength{\parsep}{0cm}}n&single-length signed number\\
d&double-length signed number\\
u&single-length unsigned number\\
ud&double-length unsigned number\\
t&triple-length\\
q&quadruple-length\\
c&7-bit character value\\
b&8-bit byte\\
?&boolean flag; or:\\
t=&true\\
f=&false\\
a or adr&address\\
acf&address of code field\\
apf&address of parameter field\\
`&(as prefix) address of\\
s d&(as a pair) source destination\\
lo hi&lower-limit upper-limit (inclusive)\\
\#&count\\
o&offset\\
i&index\\
m&mask\\
x&don't care (data structure notation)\\
\end{longtable}


An ``offset'' is a difference expressed in absolute units, such as bytes.

An ``index'' is a difference expressed in logical units, such as
elements or records.
\index{S!Stack abbreviation standards|)}

{\section{Input-Stream Comment Designations}
\index{I!Input-stream comment}}

\indent\begin{tabular}{ll}
{}{\setlength{\parsep}{0cm}}c&single character, blank-delimited\\
name&sequence of characters, blank delimited\\
text&sequence of characters, delimited by non-blank\\
\end{tabular}\medskip

\noindent Follow ``text'' with the actual delimiter required; e.g., text''
or text).

\section{Samples of Good Commenting Style}

Here are two sample screens to illustrate good commenting
style\index{C!Commenting style|(}.

\setcounter{screen}{126}
\begin{Screen}
\ Formatter         Data Structures -- p.2             06/06/83
 6 Constant TMARGIN \ line# where body of text begins
55 Constant BMARGIN \ line# where body of text ends

Create header 82 allot
\  { 1left-ent | 1right-cnt | 80header }
Create FOOTER 82 allot
\  { 1left-cnt | 1right-ent | 80footer }

Variable ACROSS   \ formatter's current horizontal position
Variable DOWNWARD \ formatter's current vertical position
Variable LEFT     \ current primary left margin
Variable WALL     \ current primary right margin
Variable WALL-WAS \ WALL when curr. line started being formt'd

\end{Screen}

\begin{Screen}
\ Formatter          positioning -- p.1               06/06/83
: skip  ( n)  ACROSS +! ;
: NEWLEFT  \ reset left margin
   LEFT @  PERMANENT @ +  TEMPORARY @ +  ACROSS ! ;
: \LINE  \ begin new line
   DOOR  CR'  1 DOWNWARD +!  NEWLEFT  WALL @  WALL-WAS ! ;
: AT-TOP?  ( -- t=at-top)  TMARGIN  DOWNWARD @ = ;
: >TMARGIN  \ move from crease to TMARGIN
   0 DOWNWARD !  BEGIN  \LINE  AT-TOP? UNTIL ;

\end{Screen}
\index{C!Commenting style|)}

\section{Naming Conventions}%
\index{N!Naming conventions|(}%
\begin{longtable}{ll>{\tt}l}
{}{\setlength{\parsep}{0cm}}%
\emph{Meaning}&\emph{Form}&\textrm{\emph{Example}}\\[.5ex] \endhead
\underline{Arithmetic}\\
integer 1&1name&1+\\
integer 2&2name&2*\\
takes relative input parameters&+name&+DRAW\\
takes scaled input parameters&*name&*DRAW\\[0.5ex]
\underline{Compilation}\\
start of ``high-level'' code&name:&CASE:\\
end of ``high-level'' code&;name&;CODE\\
put something into dictionary&name,&C,\\
executes at compile time&[name]&[COMPILE]\\
slightly different&name' (prime)&CR'\\
internal form or primitive&(name)&(TYPE)\\
&or <name>&<TYPE>\\
compiling word run-time part:\\
~~systems with no folding&lower-case&if\\
~~systems with folding&(NAME)&(IF)\\
defining word&:name&:COLOR\\
block-number where overlay begins&nam{\bf ING}&DISKING\\[0.5ex]
\underline{Data Structures}\\
table or array&names&EMPLOYEES\\
total number of elements&\#name&\#EMPLOYEES\\
current item number (variable)&name\#&EMPLOYEE\#\\
sets current item&( n) name&13 EMPLOYEE\\
advance to next element&+name&+EMPLOYEE\\
size of offset to item from&name+&DATE +\\
~~beginning of structure\\
size of (bytes per)&/name&/EMPLOYEE\\
~~(short for BYTES/name)\\
index pointer&>name&>IN\\
convert address of structure to&>name&>BODY\\
~~address of item\\
file index&(name)&(PEOPLE)\\
file pointer&--name&--JOB\\
initialize structure&0name&0RECORD\\[0.5ex]
\underline{Direction, Conversion}\\
backwards&name<&SLIDE<\\
forwards&name>&CMOVE>\\
from&<name&<TAPE\\
to&>name&>TAPE\\
convert to&name>name&FEET>METERS\\
downward&$\backslash$name&$\backslash$LINE\\
upward&/name&/LINE\\
open&\{name&\{FILE\\
close&\}name&\}FILE\\[0.5ex]
\underline{Logic, Control}\\
return boolean value&name?&SHORT?\\
returns reversed boolean&-name?&-SHORT?\\
address of boolean&'name?&'SHORT?\\
operates conditionally&?name&?DUP\\
&&\textrm{(maybe \forth{DUP})}\\
enable&+name&+CLOCK\\
~~or, absence of symbol&name&BLINKING\\
disable&-name&-CLOCK\\
&&-BLINKING\\[0.5ex]
\underline{Memory}\\
save value of&@name&@CURSOR\\
restore value of&!name&!CURSOR\\
store into&name!&SECONDS!\\
fetch from&name@&INDEX@\\
name of buffer&:name&:INSERT\\
address of name&'name&'S\\
address of pointer to name&'name&'TYPE\\
exchange, especially bytes&>name<&>MOVE<\\[0.5ex]
\underline{Numeric Types}\\
byte length&Cname&C@\\
2 cell size, 2's complement&Dname&D+\\
~~integer encoding\\
mixed 16 and 32-bit operator&Mname&M*\\
3 cell size&Tname&T*\\
4 cell size&Qname&Q*\\
unsigned encoding&Uname&U.\\[0.5ex]
\underline{Output, Printing}\\
print item&.name&.S\\
print numeric (name denotes type)&name.&D. , U.\\
print right justified&name.R&U.R\\[0.5ex]
\underline{Quantity}\\
``per''&/name&/SIDE\\[0.5ex]
\underline{Sequencing}\\
start&<name&<\#\\
end&name>&\#>\\[0.5ex]
\underline{Text}\\
string follows delimited by ''&name''&ABORT'' text''\\
text or string operator&``name&``COMPARE\\
~~(similar to \$ prefix in BASIC)\\
superstring array&``name''&``COLORS''\\
\end{longtable}%
\index{N!Naming conventions|)}%

\noindent\emph{How to Pronounce the Symbols}%
\index{S!Symbols, pronunciation of}\medskip

{\small\begin{tabular}{ll}
{}{\setlength{\parsep}{0cm}}
!&store\\
@&fetch\\
\#&sharp (or ``number,'' as in \#RECORDS)\\
\$&dollar\\
\%&percent\\
\^{ }&caret\\
\&&ampersand\\
{*}&star\\
(&left paren; paren\\
)&right paren; paren\\
--&dash; not\\
+&plus\\
=&equals\\
\{ \}&braces (traditionally called ``curly brackets'')\\
{[} {]}&square brackets\\
``&quote\\
'&as prefix: tick; as suffix: prime\\
\~{ }&tilde\\
\(|\)&bar\\
\(\backslash\)&backslash. (also ``under,'' ``down,'' and ``skip'')\\
/&slash. (also ``up'')\\
<&less-than\\
 &left dart\\
>&greater-than\\
 &right dart\\
?&question (some prefer ``query'')\\
,&comma\\
.&dot\\
\end{tabular}}


% Index
\cleardoublepage
\printindex
% prepare for back page
%% Alternative back page, written by Bernd Paysan
%% you'll find the original back page at version 1.3 in the CVS repository

\noindent
\textbf{Thinking \Forth{}} is a book about the philosophy of problem
solving and programming style, applied to the unique programming
language `\Forth{}.' Published first in 1984, it could be among the
timeless classics of computer books, such as \person{Fred Brook}'s
\emph{The Mythical Man-Month} and \person{Donald Knuth}'s \emph{The
Art of Computer Programming}.

Many software engineering principles discussed here have been
rediscovered in eXtreme Programming, including (re)factoring,
modularity, bottom-up and incremental design.  Here you'll find all of
those and more, described in \person{Leo Brodie}'s down-to-earth,
humorous style, with illustrations, code examples, practical real life
applications, illustrative cartoons, and interviews with \Forth{}'s
inventor, \person{Charles H. Moore} as well as other Forth thinkers.

If you program in \Forth{}, this is a must-read book.  If you don't,
the fundamental concepts are universal: \textbf{Thinking \Forth{}} is
meant for anyone interested in writing software to solve problems.
The concepts go beyond \Forth{}, but the \emph{simple beauty} of
\Forth{} throws those concepts into stark relief.

So flip open the book, and read all about the philosophy of Forth,
analysis, decomposition, problem solving, style and conventions,
factoring, handling data, and minimizing control structures.  But be
prepared: you may not be able to put it down.

\medskip
This book has been scanned, OCR'd, typeset in \LaTeX{}, and brought
back to print (and your monitor) by a collaborative effort under a Creative
Commons license.

\medskip\blackline{0pt}\medskip
%%if you want Leo's picture here, uncommend this
\ifleo
\makeatletter
\setbox\@tempboxa\hbox{\includegraphics{leobrodie}}
\rightmargin\wd\@tempboxa
\@tempdima\ht\@tempboxa
\advance\@tempdima-\ht\strutbox
\advance\rightmargin1em
\advance\textwidth-\rightmargin
\makeatother
\noindent\begin{tabular}{@{}p{\textwidth}@{\hspace{1em}}l@{}}
\parindent2.5em
\fi
\textbf{\person{Leo Brodie}} is a writer, programmer, consultant,
teacher, and world-renowned authority on \Forth. He was a technical
writer for \Forth{} Inc., and has been an independent consultant
(since 1982) for such clients as IBM, NCR, and Lockheed. He is also
the author of \emph{``Starting \Forth''} (Prentice-Hall, 1981).
\ifleo
& \makeatletter
\smash{\lower\@tempdima\hbox{\includegraphics{leobrodie}}}
\makeatother \\
\end{tabular}
\fi
\end{document}
